% ============================================================
%  Systemtheorie
%  Persönliches Skript
% ============================================================
\documentclass[
  12pt,
  a4paper,
  twoside=false,
  headings=normal,
  parskip=half
]{scrreprt}

% ============================================================
%  Preamble – Packages & Layout
%  Eingebunden in main.tex via % ============================================================
%  Preamble – Packages & Layout
%  Eingebunden in main.tex via % ============================================================
%  Preamble – Packages & Layout
%  Eingebunden in main.tex via \input{preamble}
% ============================================================

% ── Encoding & Language ──────────────────────────────────────
\usepackage[utf8]{inputenc}
\usepackage[T1]{fontenc}
\usepackage[ngerman]{babel}

% ── Fonts ────────────────────────────────────────────────────
\usepackage{lmodern}
\usepackage{microtype}

% ── Page Geometry ────────────────────────────────────────────
\usepackage[
  top=2.8cm,
  bottom=2.8cm,
  left=2.6cm,
  right=2.6cm,
  headheight=28pt
]{geometry}

% ── Colors ───────────────────────────────────────────────────
\usepackage[table]{xcolor}
\definecolor{accent}{RGB}{0, 80, 150}
\definecolor{defblue}{RGB}{220, 235, 255}
\definecolor{thmgreen}{RGB}{220, 245, 225}
\definecolor{examplegray}{RGB}{245, 245, 245}
\definecolor{remarkgold}{RGB}{255, 248, 220}
% table highlight
\definecolor{highlight}{RGB}{255, 195, 215}
% new box colors
\definecolor{satzred-bg}{RGB}{255, 235, 235}
\definecolor{satzred-fr}{RGB}{180, 30, 30}
\definecolor{beweisblue-bg}{RGB}{225, 240, 255}
\definecolor{beweisblue-fr}{RGB}{30, 90, 180}
\definecolor{gedankenpurple-bg}{RGB}{240, 232, 255}
\definecolor{gedankenpurple-fr}{RGB}{110, 50, 180}
\definecolor{exercisegreen-bg}{RGB}{225, 248, 230}
\definecolor{exercisegreen-fr}{RGB}{30, 140, 60}

% ── Header & Footer ──────────────────────────────────────────
\usepackage{fancyhdr}

\fancypagestyle{main}{%
  \fancyhf{}
  \fancyhead[L]{\small\itshape\parbox{0.45\linewidth}{\raggedright\rightmark}}
  \fancyhead[R]{\small\itshape Systemtheorie – Signale und Systeme I}
  \fancyfoot[C]{\small\thepage}
  \renewcommand{\headrulewidth}{0.4pt}
  \renewcommand{\footrulewidth}{0.4pt}
}

% Chapter start pages: use full header too
\renewcommand*{\chapterpagestyle}{main}

% plain style (used for ToC pages)
\fancypagestyle{plain}{%
  \fancyhf{}
  \fancyfoot[C]{\small\thepage}
  \renewcommand{\headrulewidth}{0pt}
  \renewcommand{\footrulewidth}{0.4pt}
}

% ── Math ─────────────────────────────────────────────────────
\usepackage{amsmath}
\usepackage{amssymb}
\usepackage{amsthm}
\usepackage{mathtools}

% ── Graphics & Floats ────────────────────────────────────────
\usepackage{graphicx}
\usepackage{float}
\usepackage{tikz}
\usetikzlibrary{shapes.geometric}
\usepackage{pgfplots}
\pgfplotsset{compat=1.15}
\usepgfplotslibrary{groupplots}

% ── Tables ───────────────────────────────────────────────────
\usepackage{booktabs}
\usepackage{array}
\usepackage{multirow}
\usepackage{makecell}

% ── Lists ────────────────────────────────────────────────────
\usepackage{enumitem}
\setlist[itemize]{itemsep=2pt, topsep=4pt}
\setlist[enumerate]{itemsep=2pt, topsep=4pt}

% ── Colored Boxes ────────────────────────────────────────────
\usepackage[most]{tcolorbox}

\tcbset{
  base style/.style={
    breakable,
    enhanced,
    fonttitle=\bfseries,
    colbacktitle=white,        % weißer Titelbereich
    colframe=black,            % schwarzer Rand
    titlerule=0.4pt,
    attach boxed title to top left={yshift=-2mm, xshift=4mm},
    boxed title style={
      colframe=black,
      colback=white,
      boxrule=0.4pt,
      arc=2pt,
    },
    left=6pt, right=6pt,
    top=6pt, bottom=6pt,
    arc=3pt,
    boxrule=0.6pt,
  }
}

% Definition (blau)
\newtcolorbox{definitionbox}[1][]{
  base style,
  colback=defblue,
  coltitle=accent,
  title={Definition\ifx&#1&\else: #1\fi},
}

% Satz (rot)
\newtcolorbox{satzbox}[1][]{
  base style,
  colback=satzred-bg,
  coltitle=satzred-fr,
  title={Satz\ifx&#1&\else: #1\fi},
}

% Beweis (blau)
\newtcolorbox{beweisbox}[1][]{
  base style,
  colback=beweisblue-bg,
  coltitle=beweisblue-fr,
  title={Beweis\ifx&#1&\else: #1\fi},
}

% Gedankenexperiment (lila)
\newtcolorbox{gedankenbox}[1][]{
  base style,
  colback=gedankenpurple-bg,
  coltitle=gedankenpurple-fr,
  title={Gedankenexperiment\ifx&#1&\else: #1\fi},
}

% Beispiel / Exercise (grün)
\newtcolorbox{exercisebox}[1][]{
  base style,
  colback=exercisegreen-bg,
  coltitle=exercisegreen-fr,
  title={Beispiel\ifx&#1&\else: #1\fi},
}

% Bemerkung (gold)
\newtcolorbox{bemerkungbox}[1][]{
  base style,
  colback=remarkgold,
  coltitle=orange!80!black,
  title={Bemerkung\ifx&#1&\else: #1\fi},
}

% ── Environments (lowercase and capitalized aliases) ──────────
\newenvironment{definition}[1][]{\begin{definitionbox}[#1]}{\end{definitionbox}}
\newenvironment{Definition}[1][]{\begin{definitionbox}[#1]}{\end{definitionbox}}
\newenvironment{satz}[1][]{\begin{satzbox}[#1]}{\end{satzbox}}
\newenvironment{Satz}[1][]{\begin{satzbox}[#1]}{\end{satzbox}}
\newenvironment{beweis}[1][]{\begin{beweisbox}[#1]}{\end{beweisbox}}
\newenvironment{Beweis}[1][]{\begin{beweisbox}[#1]}{\end{beweisbox}}
\newenvironment{gedankenexperiment}[1][]{\begin{gedankenbox}[#1]}{\end{gedankenbox}}
\newenvironment{Gedankenexperiment}[1][]{\begin{gedankenbox}[#1]}{\end{gedankenbox}}
\newenvironment{beispiel}[1][]{\begin{exercisebox}[#1]}{\end{exercisebox}}
\newenvironment{Beispiel}[1][]{\begin{exercisebox}[#1]}{\end{exercisebox}}
\newenvironment{bemerkung}[1][]{\begin{bemerkungbox}[#1]}{\end{bemerkungbox}}
\newenvironment{Bemerkung}[1][]{\begin{bemerkungbox}[#1]}{\end{bemerkungbox}}

% ── Heading Spacing ──────────────────────────────────────────
% Reduce the space below chapter, section, subsection headings
\RedeclareSectionCommand[beforeskip=0.8\baselineskip, afterskip=0.5\baselineskip]{chapter}
\RedeclareSectionCommand[afterskip=0.3\baselineskip]{section}
\RedeclareSectionCommand[afterskip=0.2\baselineskip]{subsection}

% ── Hyperlinks ───────────────────────────────────────────────
\usepackage[
  colorlinks=true,
  linkcolor=accent,
  citecolor=accent,
  urlcolor=accent,
  bookmarksnumbered=true,
  pdfauthor={Philipp Lehmann},
  pdftitle={Systemtheorie – Signale und Systeme I},
]{hyperref}

% ============================================================

% ── Encoding & Language ──────────────────────────────────────
\usepackage[utf8]{inputenc}
\usepackage[T1]{fontenc}
\usepackage[ngerman]{babel}

% ── Fonts ────────────────────────────────────────────────────
\usepackage{lmodern}
\usepackage{microtype}

% ── Page Geometry ────────────────────────────────────────────
\usepackage[
  top=2.8cm,
  bottom=2.8cm,
  left=2.6cm,
  right=2.6cm,
  headheight=28pt
]{geometry}

% ── Colors ───────────────────────────────────────────────────
\usepackage[table]{xcolor}
\definecolor{accent}{RGB}{0, 80, 150}
\definecolor{defblue}{RGB}{220, 235, 255}
\definecolor{thmgreen}{RGB}{220, 245, 225}
\definecolor{examplegray}{RGB}{245, 245, 245}
\definecolor{remarkgold}{RGB}{255, 248, 220}
% table highlight
\definecolor{highlight}{RGB}{255, 195, 215}
% new box colors
\definecolor{satzred-bg}{RGB}{255, 235, 235}
\definecolor{satzred-fr}{RGB}{180, 30, 30}
\definecolor{beweisblue-bg}{RGB}{225, 240, 255}
\definecolor{beweisblue-fr}{RGB}{30, 90, 180}
\definecolor{gedankenpurple-bg}{RGB}{240, 232, 255}
\definecolor{gedankenpurple-fr}{RGB}{110, 50, 180}
\definecolor{exercisegreen-bg}{RGB}{225, 248, 230}
\definecolor{exercisegreen-fr}{RGB}{30, 140, 60}

% ── Header & Footer ──────────────────────────────────────────
\usepackage{fancyhdr}

\fancypagestyle{main}{%
  \fancyhf{}
  \fancyhead[L]{\small\itshape\parbox{0.45\linewidth}{\raggedright\rightmark}}
  \fancyhead[R]{\small\itshape Systemtheorie – Signale und Systeme I}
  \fancyfoot[C]{\small\thepage}
  \renewcommand{\headrulewidth}{0.4pt}
  \renewcommand{\footrulewidth}{0.4pt}
}

% Chapter start pages: use full header too
\renewcommand*{\chapterpagestyle}{main}

% plain style (used for ToC pages)
\fancypagestyle{plain}{%
  \fancyhf{}
  \fancyfoot[C]{\small\thepage}
  \renewcommand{\headrulewidth}{0pt}
  \renewcommand{\footrulewidth}{0.4pt}
}

% ── Math ─────────────────────────────────────────────────────
\usepackage{amsmath}
\usepackage{amssymb}
\usepackage{amsthm}
\usepackage{mathtools}

% ── Graphics & Floats ────────────────────────────────────────
\usepackage{graphicx}
\usepackage{float}
\usepackage{tikz}
\usetikzlibrary{shapes.geometric}
\usepackage{pgfplots}
\pgfplotsset{compat=1.15}
\usepgfplotslibrary{groupplots}

% ── Tables ───────────────────────────────────────────────────
\usepackage{booktabs}
\usepackage{array}
\usepackage{multirow}
\usepackage{makecell}

% ── Lists ────────────────────────────────────────────────────
\usepackage{enumitem}
\setlist[itemize]{itemsep=2pt, topsep=4pt}
\setlist[enumerate]{itemsep=2pt, topsep=4pt}

% ── Colored Boxes ────────────────────────────────────────────
\usepackage[most]{tcolorbox}

\tcbset{
  base style/.style={
    breakable,
    enhanced,
    fonttitle=\bfseries,
    colbacktitle=white,        % weißer Titelbereich
    colframe=black,            % schwarzer Rand
    titlerule=0.4pt,
    attach boxed title to top left={yshift=-2mm, xshift=4mm},
    boxed title style={
      colframe=black,
      colback=white,
      boxrule=0.4pt,
      arc=2pt,
    },
    left=6pt, right=6pt,
    top=6pt, bottom=6pt,
    arc=3pt,
    boxrule=0.6pt,
  }
}

% Definition (blau)
\newtcolorbox{definitionbox}[1][]{
  base style,
  colback=defblue,
  coltitle=accent,
  title={Definition\ifx&#1&\else: #1\fi},
}

% Satz (rot)
\newtcolorbox{satzbox}[1][]{
  base style,
  colback=satzred-bg,
  coltitle=satzred-fr,
  title={Satz\ifx&#1&\else: #1\fi},
}

% Beweis (blau)
\newtcolorbox{beweisbox}[1][]{
  base style,
  colback=beweisblue-bg,
  coltitle=beweisblue-fr,
  title={Beweis\ifx&#1&\else: #1\fi},
}

% Gedankenexperiment (lila)
\newtcolorbox{gedankenbox}[1][]{
  base style,
  colback=gedankenpurple-bg,
  coltitle=gedankenpurple-fr,
  title={Gedankenexperiment\ifx&#1&\else: #1\fi},
}

% Beispiel / Exercise (grün)
\newtcolorbox{exercisebox}[1][]{
  base style,
  colback=exercisegreen-bg,
  coltitle=exercisegreen-fr,
  title={Beispiel\ifx&#1&\else: #1\fi},
}

% Bemerkung (gold)
\newtcolorbox{bemerkungbox}[1][]{
  base style,
  colback=remarkgold,
  coltitle=orange!80!black,
  title={Bemerkung\ifx&#1&\else: #1\fi},
}

% ── Environments (lowercase and capitalized aliases) ──────────
\newenvironment{definition}[1][]{\begin{definitionbox}[#1]}{\end{definitionbox}}
\newenvironment{Definition}[1][]{\begin{definitionbox}[#1]}{\end{definitionbox}}
\newenvironment{satz}[1][]{\begin{satzbox}[#1]}{\end{satzbox}}
\newenvironment{Satz}[1][]{\begin{satzbox}[#1]}{\end{satzbox}}
\newenvironment{beweis}[1][]{\begin{beweisbox}[#1]}{\end{beweisbox}}
\newenvironment{Beweis}[1][]{\begin{beweisbox}[#1]}{\end{beweisbox}}
\newenvironment{gedankenexperiment}[1][]{\begin{gedankenbox}[#1]}{\end{gedankenbox}}
\newenvironment{Gedankenexperiment}[1][]{\begin{gedankenbox}[#1]}{\end{gedankenbox}}
\newenvironment{beispiel}[1][]{\begin{exercisebox}[#1]}{\end{exercisebox}}
\newenvironment{Beispiel}[1][]{\begin{exercisebox}[#1]}{\end{exercisebox}}
\newenvironment{bemerkung}[1][]{\begin{bemerkungbox}[#1]}{\end{bemerkungbox}}
\newenvironment{Bemerkung}[1][]{\begin{bemerkungbox}[#1]}{\end{bemerkungbox}}

% ── Heading Spacing ──────────────────────────────────────────
% Reduce the space below chapter, section, subsection headings
\RedeclareSectionCommand[beforeskip=0.8\baselineskip, afterskip=0.5\baselineskip]{chapter}
\RedeclareSectionCommand[afterskip=0.3\baselineskip]{section}
\RedeclareSectionCommand[afterskip=0.2\baselineskip]{subsection}

% ── Hyperlinks ───────────────────────────────────────────────
\usepackage[
  colorlinks=true,
  linkcolor=accent,
  citecolor=accent,
  urlcolor=accent,
  bookmarksnumbered=true,
  pdfauthor={Philipp Lehmann},
  pdftitle={Systemtheorie – Signale und Systeme I},
]{hyperref}

% ============================================================

% ── Encoding & Language ──────────────────────────────────────
\usepackage[utf8]{inputenc}
\usepackage[T1]{fontenc}
\usepackage[ngerman]{babel}

% ── Fonts ────────────────────────────────────────────────────
\usepackage{lmodern}
\usepackage{microtype}

% ── Page Geometry ────────────────────────────────────────────
\usepackage[
  top=2.8cm,
  bottom=2.8cm,
  left=2.6cm,
  right=2.6cm,
  headheight=28pt
]{geometry}

% ── Colors ───────────────────────────────────────────────────
\usepackage[table]{xcolor}
\definecolor{accent}{RGB}{0, 80, 150}
\definecolor{defblue}{RGB}{220, 235, 255}
\definecolor{thmgreen}{RGB}{220, 245, 225}
\definecolor{examplegray}{RGB}{245, 245, 245}
\definecolor{remarkgold}{RGB}{255, 248, 220}
% table highlight
\definecolor{highlight}{RGB}{255, 195, 215}
% new box colors
\definecolor{satzred-bg}{RGB}{255, 235, 235}
\definecolor{satzred-fr}{RGB}{180, 30, 30}
\definecolor{beweisblue-bg}{RGB}{225, 240, 255}
\definecolor{beweisblue-fr}{RGB}{30, 90, 180}
\definecolor{gedankenpurple-bg}{RGB}{240, 232, 255}
\definecolor{gedankenpurple-fr}{RGB}{110, 50, 180}
\definecolor{exercisegreen-bg}{RGB}{225, 248, 230}
\definecolor{exercisegreen-fr}{RGB}{30, 140, 60}

% ── Header & Footer ──────────────────────────────────────────
\usepackage{fancyhdr}

\fancypagestyle{main}{%
  \fancyhf{}
  \fancyhead[L]{\small\itshape\parbox{0.45\linewidth}{\raggedright\rightmark}}
  \fancyhead[R]{\small\itshape Systemtheorie – Signale und Systeme I}
  \fancyfoot[C]{\small\thepage}
  \renewcommand{\headrulewidth}{0.4pt}
  \renewcommand{\footrulewidth}{0.4pt}
}

% Chapter start pages: use full header too
\renewcommand*{\chapterpagestyle}{main}

% plain style (used for ToC pages)
\fancypagestyle{plain}{%
  \fancyhf{}
  \fancyfoot[C]{\small\thepage}
  \renewcommand{\headrulewidth}{0pt}
  \renewcommand{\footrulewidth}{0.4pt}
}

% ── Math ─────────────────────────────────────────────────────
\usepackage{amsmath}
\usepackage{amssymb}
\usepackage{amsthm}
\usepackage{mathtools}

% ── Graphics & Floats ────────────────────────────────────────
\usepackage{graphicx}
\usepackage{float}
\usepackage{tikz}
\usetikzlibrary{shapes.geometric}
\usepackage{pgfplots}
\pgfplotsset{compat=1.15}
\usepgfplotslibrary{groupplots}

% ── Tables ───────────────────────────────────────────────────
\usepackage{booktabs}
\usepackage{array}
\usepackage{multirow}
\usepackage{makecell}

% ── Lists ────────────────────────────────────────────────────
\usepackage{enumitem}
\setlist[itemize]{itemsep=2pt, topsep=4pt}
\setlist[enumerate]{itemsep=2pt, topsep=4pt}

% ── Colored Boxes ────────────────────────────────────────────
\usepackage[most]{tcolorbox}

\tcbset{
  base style/.style={
    breakable,
    enhanced,
    fonttitle=\bfseries,
    colbacktitle=white,        % weißer Titelbereich
    colframe=black,            % schwarzer Rand
    titlerule=0.4pt,
    attach boxed title to top left={yshift=-2mm, xshift=4mm},
    boxed title style={
      colframe=black,
      colback=white,
      boxrule=0.4pt,
      arc=2pt,
    },
    left=6pt, right=6pt,
    top=6pt, bottom=6pt,
    arc=3pt,
    boxrule=0.6pt,
  }
}

% Definition (blau)
\newtcolorbox{definitionbox}[1][]{
  base style,
  colback=defblue,
  coltitle=accent,
  title={Definition\ifx&#1&\else: #1\fi},
}

% Satz (rot)
\newtcolorbox{satzbox}[1][]{
  base style,
  colback=satzred-bg,
  coltitle=satzred-fr,
  title={Satz\ifx&#1&\else: #1\fi},
}

% Beweis (blau)
\newtcolorbox{beweisbox}[1][]{
  base style,
  colback=beweisblue-bg,
  coltitle=beweisblue-fr,
  title={Beweis\ifx&#1&\else: #1\fi},
}

% Gedankenexperiment (lila)
\newtcolorbox{gedankenbox}[1][]{
  base style,
  colback=gedankenpurple-bg,
  coltitle=gedankenpurple-fr,
  title={Gedankenexperiment\ifx&#1&\else: #1\fi},
}

% Beispiel / Exercise (grün)
\newtcolorbox{exercisebox}[1][]{
  base style,
  colback=exercisegreen-bg,
  coltitle=exercisegreen-fr,
  title={Beispiel\ifx&#1&\else: #1\fi},
}

% Bemerkung (gold)
\newtcolorbox{bemerkungbox}[1][]{
  base style,
  colback=remarkgold,
  coltitle=orange!80!black,
  title={Bemerkung\ifx&#1&\else: #1\fi},
}

% ── Environments (lowercase and capitalized aliases) ──────────
\newenvironment{definition}[1][]{\begin{definitionbox}[#1]}{\end{definitionbox}}
\newenvironment{Definition}[1][]{\begin{definitionbox}[#1]}{\end{definitionbox}}
\newenvironment{satz}[1][]{\begin{satzbox}[#1]}{\end{satzbox}}
\newenvironment{Satz}[1][]{\begin{satzbox}[#1]}{\end{satzbox}}
\newenvironment{beweis}[1][]{\begin{beweisbox}[#1]}{\end{beweisbox}}
\newenvironment{Beweis}[1][]{\begin{beweisbox}[#1]}{\end{beweisbox}}
\newenvironment{gedankenexperiment}[1][]{\begin{gedankenbox}[#1]}{\end{gedankenbox}}
\newenvironment{Gedankenexperiment}[1][]{\begin{gedankenbox}[#1]}{\end{gedankenbox}}
\newenvironment{beispiel}[1][]{\begin{exercisebox}[#1]}{\end{exercisebox}}
\newenvironment{Beispiel}[1][]{\begin{exercisebox}[#1]}{\end{exercisebox}}
\newenvironment{bemerkung}[1][]{\begin{bemerkungbox}[#1]}{\end{bemerkungbox}}
\newenvironment{Bemerkung}[1][]{\begin{bemerkungbox}[#1]}{\end{bemerkungbox}}

% ── Heading Spacing ──────────────────────────────────────────
% Reduce the space below chapter, section, subsection headings
\RedeclareSectionCommand[beforeskip=0.8\baselineskip, afterskip=0.5\baselineskip]{chapter}
\RedeclareSectionCommand[afterskip=0.3\baselineskip]{section}
\RedeclareSectionCommand[afterskip=0.2\baselineskip]{subsection}

% ── Hyperlinks ───────────────────────────────────────────────
\usepackage[
  colorlinks=true,
  linkcolor=accent,
  citecolor=accent,
  urlcolor=accent,
  bookmarksnumbered=true,
  pdfauthor={Philipp Lehmann},
  pdftitle={Systemtheorie – Signale und Systeme I},
]{hyperref}


% ═════════════════════════════════════════════════════════════
%  DOKUMENT
% ═════════════════════════════════════════════════════════════
\begin{document}

% ── Deckblatt ────────────────────────────────────────────────
\pagenumbering{alph}% avoids page.1 duplicate with ToC/content
\begin{titlepage}
  \thispagestyle{empty}
  \centering
  \vspace*{3cm}

  {\color{accent}\rule{\linewidth}{1.5pt}}\vspace{0.5em}

  {\fontsize{28}{34}\selectfont\bfseries Systemtheorie}\\[0.4em]
  {\Large\bfseries Signale und Systeme I}

  \vspace{0.4em}
  {\color{accent}\rule{\linewidth}{1.5pt}}

  \vspace{2.5cm}
  \vfill

  \begin{tabular}{@{}rl@{}}
    \textbf{Autor:}    & Philipp Lehmann \\[0.5em]
    \textbf{Semester:} & Sommersemester 2026 \\[0.5em]
    \textbf{Stand:}    & \today
  \end{tabular}

  \vspace{2.5cm}
  {\color{accent}\rule{\linewidth}{0.5pt}}
\end{titlepage}

% ── Inhaltsverzeichnis ───────────────────────────────────────
\pagenumbering{roman}
\pagestyle{plain}
\tableofcontents
\clearpage

% ── Hauptinhalt ──────────────────────────────────────────────
\pagenumbering{arabic}
\setcounter{page}{1}
\pagestyle{main}

% ── Abstract ─────────────────────────────────────────────────
% ============================================================
%  Abstract – Modulbeschreibung
% ============================================================
\chapter*{Abstract}
\addcontentsline{toc}{chapter}{Abstract}
\markboth{Abstract}{Abstract}

\begin{tabular}{@{}ll@{}}
  \textbf{Modul:}       & Systemtheorie 1 – Signale und Systeme (141170) \\[0.3em]
  \textbf{Dozent:}      & Prof.\,Dr.-Ing.\ Rainer Martin \\[0.3em]
  \textbf{Fakultät:}    & Elektrotechnik und Informationstechnik, Ruhr-Universität Bochum \\[0.3em]
  \textbf{Turnus:}      & jedes Sommersemester \quad|\quad 4\,SWS \quad|\quad 5\,CP \\[0.3em]
  \textbf{Vorkenntnisse:} & Komplexe Zahlen, Differential- und Integralrechnung, Reihen \\
\end{tabular}

\vspace{1.2em}

Die Vorlesung \textbf{Systemtheorie 1 – Signale und Systeme} vermittelt die mathematischen
Grundlagen zur Beschreibung, Analyse und Verarbeitung von Signalen sowie zum systematischen
Umgang mit linearen und nichtlinearen Systemen. Sie bildet das Fundament für weiterführende
Veranstaltungen in der Nachrichten-, Mess- und Regelungstechnik.

\vspace{0.8em}

\section*{Modulinhalt}

Das Modul gliedert sich in drei thematische Blöcke:

\subsection*{1\quad Signale und Systeme}

Der erste und umfangreichste Teil behandelt die mathematische Beschreibung von Signalen
im Zeitbereich. Ausgangspunkt sind grundlegende Kenngrößen und Eigenschaften von Signalen
– Energie, Leistung, Periodizität und Symmetrie – sowie elementare Operationen wie
Verschiebung, Spiegelung und Skalierung. Auf dieser Basis werden Methoden der
\textbf{Signalsynthese und -analyse} eingeführt, insbesondere die Zerlegung periodischer
Signale in Fourierreihen. Ein weiterer Schwerpunkt liegt auf der
\textbf{Analog-Digital- und Digital-Analog-Umsetzung} (Abtasttheorem nach
Shannon–Nyquist, Quantisierung). Den Abschluss bildet die Klassifikation von Systemen:
lineare vs.\ nichtlineare, zeitinvariante vs.\ zeitvariante sowie kausale Systeme und
deren Beschreibung durch Impulsantwort und Faltungsintegral.

\subsection*{2\quad Einführung in die Wahrscheinlichkeitsrechnung}

Der zweite Block legt die stochastischen Grundlagen, die für die spätere Behandlung
von Rauschen und Zufallssignalen unverzichtbar sind. Behandelt werden Axiome und
Rechenregeln der Wahrscheinlichkeit, bedingte Wahrscheinlichkeit und Bayes'sches
Theorem sowie mehrstufige Zufallsexperimente. Im Mittelpunkt stehen
\textbf{diskrete und kontinuierliche Zufallsvariablen} mit ihren Verteilungsfunktionen,
Dichtefunktionen und charakteristischen Kenngrößen (Erwartungswert, Varianz, Momente).

\subsection*{3\quad Grundbegriffe der Informationstheorie}

Der dritte Block führt in die von Claude Shannon begründete Informationstheorie ein.
Zentrale Fragen sind: Was ist Information? Wie lässt sie sich messen? Wie effizient
kann eine Quelle komprimiert werden? Als mathematisches Werkzeug wird der
\textbf{Entropiebegriff} (Shannon-Entropie, bedingte Entropie, Transinformation)
eingeführt und auf praktische Anwendungen in der Quellen- und Kanalcodierung übertragen.

\vspace{0.8em}

\section*{Lernziele}

\begin{itemize}
  \item Mathematische Beschreibung von Signalen und Systemen im Zeitbereich und
        Kenntnis ihrer wesentlichen Merkmale.
  \item Sicherer Umgang mit diskreten und kontinuierlichen Zufallsvariablen sowie
        deren Kenngrössen.
  \item Verständnis der grundlegenden Konzepte der Informationstheorie und Fähigkeit
        zur Anwendung des Entropiebegriffs.
  \item Vernetzung von Signaltheorie, Stochastik und Informationstheorie als Basis
        für weiterführende Gebiete der Elektrotechnik und Informationstechnik.
\end{itemize}


% ── Prolog ───────────────────────────────────────────────────
% ============================================================
%  Prolog
% ============================================================
\chapter*{Prolog}
\addcontentsline{toc}{chapter}{Prolog}
\markboth{Prolog}{Prolog}

Dieses Skript ist meine Begleitdokument zur Vorlesung
\textbf{Systemtheorie – Signale und Systeme~I} im Sommersemester~2026.
Es soll als mathematisch strukturierte Mitschrift mit dem Ziel, die
zentralen Konzepte und Methoden der Signalverarbeitung und Systemanalyse
kompakt und nachvollziehbar festzuhalten, geschrieben werden. 
Der Inhalt soll mathematisch komplexe sProzesse verständlich erklären.

\bigskip

\begin{bemerkungbox}[Hinweis zur Nutzung]
Das Skript umfasst mehr Inhalte als die Vorlesung da Grundlagen mit erklärt 
werden.
\end{bemerkungbox}

\bigskip

\section*{Aufbau des Skripts}

Das Skript folgt dem Verlauf der Vorlesung und ist in drei Teile gegliedert:

\begin{enumerate}
  \item \textbf{Vorlesungen} – Kapitelweise Mitschrift der einzelnen
        Vorlesungseinheiten in chronologischer Reihenfolge.
  \item \textbf{Zusätzliches Wissen} – Themen, die über den direkten
        Vorlesungsinhalt hinausgehen, aber für das Verständnis relevant
        oder mathematisch interessant sind.
  \item \textbf{Formelsammlung} – Eine kompakte Zusammenstellung aller
        wichtigen Formeln und Identitäten zum schnellen Nachschlagen.
\end{enumerate}

\section*{Verwendete Boxtypen}

Zur übersichtlichen Darstellung unterschiedlicher Inhaltstypen werden
folgende farblich kodierte Boxen verwendet:

\begin{definitionbox}[Beispiel]
  Formale Definitionen mathematischer Begriffe und Konzepte.
\end{definitionbox}

\begin{satzbox}[Beispiel]
  Mathematische Sätze und Theoreme mit Anspruch auf allgemeine Gültigkeit.
\end{satzbox}

\begin{beweisbox}[Beispiel]
  Beweise zu Sätzen, sofern sie für das Verständnis relevant sind.
\end{beweisbox}

\begin{exercisebox}[Beispiel]
  Anschauliche Beispiele und Rechenaufgaben zur Illustration der Theorie.
\end{exercisebox}

\begin{bemerkungbox}[Beispiel]
  Ergänzende Hinweise, Besonderheiten und Querverweise.
\end{bemerkungbox}

\begin{gedankenbox}[Beispiel]
  Gedankenexperimente zur intuitiven Motivation von Konzepten.
\end{gedankenbox}

\vfill
\begin{flushright}
  \textit{Philipp Lehmann, Sommersemester 2026}
\end{flushright}


% ── Vorlesungen ──────────────────────────────────────────────
\chapter{Determinierte Signale}
\label{chap:determinierte-signale}


\section{Signalarten}
\label{sec:signalarten}

Je nach Zeitbasis und Wertemenge des Signals unterscheidet man in:

\begin{description}
    \item[Analoge Signale]: Zeit- und Wertekontinuierliche Signale
    \item[Zeitdiskrete Signale]: Zeitdiskret, aber Wertekontinuierliche Signale
    \item[Wertdiskrete Signale]: Wertediskret, aber Zeitkontinuierliche Signale
    \item[Diskrete Signale]: Zeit- und Wertediskrete Signale. \textit{Digitale Signale}.   
\end{description}

% ###################################################
%              Graphen der Signalarten
% ###################################################

\begin{figure}[htbp]
\centering
\begin{tikzpicture}
\pgfplotsset{
  sigstyle/.style={
    width=5.5cm, height=4.5cm,
    xmin=0, xmax=4.7,
    ymin=-0.5, ymax=0.7,
    xtick={0,1,2,3,4},
    ytick={-0.4,-0.2,0,0.2,0.4,0.6},
    xlabel={$t$/ms},
    axis x line=bottom,
    axis y line=left,
    grid=both,
    grid style={dotted, gray!50},
    tick style={color=black},
    title style={font=\small},
    label style={font=\small},
    tick label style={font=\scriptsize},
  }
}
\begin{groupplot}[
  group style={
    group size=2 by 2,
    horizontal sep=2.2cm,
    vertical sep=1.8cm,
  },
  sigstyle,
]

\nextgroupplot[title={analog}, ylabel={$x(t)$}]
\addplot[blue, thick, smooth, domain=0:4.5, samples=400]
  {0.3*cos(360*(x-2))*exp(-0.5*(x-2)^2/0.6)};

\nextgroupplot[title={zeitdiskret}, ylabel={$x(kT_A)$}]
\addplot[blue, thick, ycomb, mark=*, mark size=2pt,
  mark options={fill=blue, draw=blue},
  samples at={0,0.2,...,4.4}]
  {0.3*cos(360*(x-2))*exp(-0.5*(x-2)^2/0.6)};

\nextgroupplot[title={wertdiskret}, ylabel={$x_q(t)$}]
\addplot[blue, thick, const plot, mark=none, domain=0:4.5, samples=1000]
  {round(10*(0.3*cos(360*(x-2))*exp(-0.5*(x-2)^2/0.6)))*0.1};

\nextgroupplot[title={digital}, ylabel={$x_q(kT_A)$}]
\addplot[blue, thick, ycomb, mark=*, mark size=2pt,
  mark options={fill=blue, draw=blue},
  samples at={0,0.2,...,4.4}]
  {round(10*(0.3*cos(360*(x-2))*exp(-0.5*(x-2)^2/0.6)))*0.1};

\end{groupplot}
\end{tikzpicture}
\caption{Signaltypen: analog, zeitdiskret, wertdiskret und digital ($T_A = 0{,}2\,$ms)}
\label{fig:signaltypen}
\end{figure}


\subsection{Zeitdiskrete Signale}
\label{sec:zeitdiskrete-signale}

Das Zeitdiskrete Signal entsteht durch \textit{Abtastung} des analogen Signals.
Das Wertdiskrete Signal durch \textit{Quantisierung}.


\subsection{Komplexwertige und Reellwertige Signale}
\label{sec:komplexwertige-signale}

Signale können reellwertig $(x: \mathbb{R} \rightarrow \mathbb{R})$ oder komplexwertig 
$(x: \mathbb{R} \rightarrow \mathbb{C})$ sein.

\subsubsection{Komplexwertige Signale}

Ein komplexwertiges Signal kann über reellwertige Real- und Imaginärteilsignale oder über Amplitude
und Phasensignal formuliert werden. Wir schreiben vortan für dei Phase: $j = \sqrt{-1}$.

$$
\begin{array}{rl}
x(t) = x_R(t) + jx_I(t) = & Re\{x(t)\} + jIm\{x(t)\}\\
= & |x(t)|e^{j \phi(t)}
\end{array}
$$

Mit der Definition des zu $x(t)$ konjugierten Signals $x^*(t) = x_R(t) - jx_I(t)$ gilt:

$$
\begin{array}{rl}
x_R(t) = & \operatorname{Re}\{x(t)\} = \dfrac{1}{2}\bigl(x(t) + x^*(t)\bigr) \\[8pt]
x_I(t) = & \operatorname{Im}\{x(t)\} = \dfrac{1}{2j}\bigl(x(t) - x^*(t)\bigr) \\[8pt]
|x(t)| = & \sqrt{x_R^2(t) + x_I^2(t)} = \sqrt{x(t)\cdot x^*(t)} \\[8pt]
\phi(t) = & \arg\{x(t)\}
\end{array}
$$

\section{Elementare analoge Signale}
\label{sec:elementare-analoge-signale}

\begin{definition}[Gleichsignal]
$x(t) = A \forall t$ $T$: Zeitliches Bezugsintervall
$$
x(t) = A \forall t
$$


\begin{center}
\begin{tikzpicture}
\begin{axis}[
  width=6.5cm, height=3cm,
  xmin=-1.5, xmax=3.2, ymin=-0.3, ymax=1.5,
  axis lines=center,
  xtick=\empty, ytick={1}, yticklabels={$A$},
  xlabel={$t$},
  tick style={color=black, thin}, tick label style={font=\small},
]
\addplot[blue, very thick, domain=-1.5:3.2, samples=2] {1};
\end{axis}
\end{tikzpicture}
\end{center}

\end{definition}


\begin{definition}[Sinussignal]

$x(t) = A sin(\omega t + \phi)$

  $$
  x(t) = A sin(\omega t+ \phi) = A cos(\omega t + \phi - \frac{\pi}{2})
  $$

  \begin{itemize}
    \item $A$ die Amplitude 
    \item $\omega = 2 \pi f$ die Kreisfrequenz (gemessen in $1/s$)
    \item $f = \frac{1}{T}$ die Frequenz (Schwingungen pro Sekunde in Hertz, $Hz = 1/s$)
    \item $T$ die Periodendauer (gemessen in s)
    \item $\phi$ den Phasenwinkel (gemessen in rad)
  \end{itemize}


\begin{center}
\begin{tikzpicture}
\begin{axis}[
  width=6.5cm, height=3cm,
  xmin=-0.3, xmax=4.3, ymin=-1.5, ymax=1.5,
  axis lines=center,
  xtick={2,4}, xticklabels={$T$, $2T$},
  ytick={1,-1}, yticklabels={$A$, $-A$},
  xlabel={$t$},
  tick style={color=black, thin}, tick label style={font=\small},
]
\addplot[blue, very thick, smooth, domain=0:4, samples=300] {sin(180*x)};
\end{axis}
\end{tikzpicture}
\end{center}

\end{definition}


\begin{definition}[Einheitssprung]
$$
  x(t) = \epsilon(t) =
  \begin{cases}
    0 & t < 0 \\
    1 & t \geq 0
  \end{cases}
$$


\begin{center}
\begin{tikzpicture}
\begin{axis}[
  width=6.5cm, height=3cm,
  xmin=-1.8, xmax=3.0, ymin=-0.4, ymax=1.5,
  axis lines=center,
  xtick=\empty, ytick={1}, yticklabels={$1$},
  xlabel={$t$},
  tick style={color=black, thin}, tick label style={font=\small},
]
\addplot[blue, very thick] coordinates {(-1.8,0) (0,0)};
\addplot[blue, very thick] coordinates {(0,1) (3.0,1)};
\draw[blue, very thick] (axis cs:0,0) -- (axis cs:0,1);
\addplot[only marks, mark=*, mark options={fill=blue,  draw=blue},  mark size=2pt] coordinates {(0,1)};
\addplot[only marks, mark=*, mark options={fill=white, draw=blue},  mark size=2pt] coordinates {(0,0)};
\end{axis}
\end{tikzpicture}
\end{center}

\end{definition}


\begin{definition}[Rechteckimpuls]
Den Rechteckimpuls definieren wir zunächst allgemein für $u \in \mathbb{R}$:

$$
\operatorname{rect}(u) = \begin{cases} 0 & u < 0 \\ 1 & 0 \leq u < 1 \\ 0 & u \geq 1 \end{cases}
$$

Damit ergibt sich für ein Rechtecksignal der Breite $T$:

$$
  x(t) = \operatorname{rect}\!\left(\frac{t}{T}\right) =
  \begin{cases}
    0 & t < 0 \\
    1 & 0 \leq t < T \\
    0 & t \geq T
  \end{cases}
$$

Der Flächeninhalt des Rechtecksignals berechnet sich wie folgt:

$$
\int_{-\infty}^{\infty} x(t)\,\mathrm{d}t = \int_0^T 1\,\mathrm{d}t = \bigl[t\bigr]_0^T = T.
$$


\begin{center}
\begin{tikzpicture}
\begin{axis}[
  width=6.5cm, height=3cm,
  xmin=-0.8, xmax=2.5, ymin=-0.4, ymax=1.5,
  axis lines=center,
  xtick={1}, xticklabels={$T$},
  ytick={1}, yticklabels={$1$},
  xlabel={$t$},
  tick style={color=black, thin}, tick label style={font=\small},
]
\addplot[blue, very thick] coordinates {(-0.8,0) (0,0)};
\draw[blue, very thick] (axis cs:0,0) -- (axis cs:0,1);
\addplot[blue, very thick] coordinates {(0,1) (1,1)};
\draw[blue, very thick] (axis cs:1,1) -- (axis cs:1,0);
\addplot[blue, very thick] coordinates {(1,0) (2.5,0)};
\end{axis}
\end{tikzpicture}
\end{center}

\end{definition}


\begin{definition}[Exponentialsignal]
Das Exponentialsignal ist wie folgt definiert:

$$
  x(t) =
  \begin{cases}
    0             & t < 0 \\
    A e^{-t/\tau} & t \geq 0
  \end{cases}
$$

\begin{center}
\begin{tikzpicture}
\begin{axis}[
  width=6.5cm, height=3cm,
  xmin=-0.8, xmax=4.5, ymin=-0.2, ymax=1.4,
  axis lines=center,
  xtick={1}, xticklabels={$\tau$},
  ytick={1}, yticklabels={$A$},
  xlabel={$t$},
  tick style={color=black, thin}, tick label style={font=\small},
]
\addplot[blue, very thick] coordinates {(-0.8,0) (0,0)};
\draw[blue, very thick] (axis cs:0,0) -- (axis cs:0,1);
\addplot[blue, very thick, smooth, domain=0:4.5, samples=200] {exp(-x)};
\draw[gray, dashed, thin] (axis cs:0,0.368) -- (axis cs:1,0.368);
\draw[gray, dashed, thin] (axis cs:1,0)     -- (axis cs:1,0.368);
\end{axis}
\end{tikzpicture}
\end{center}

Die Steigung des Exponentialsignals an der Stelle $t = 0^+$, also bei rechtsseitiger Annäherung an den Punkt $t = 0$, beträgt:

$$
\frac{\mathrm{d}x(t)}{\mathrm{d}t}\bigg|_{t=0^+} = -\frac{A}{\tau}.
$$

\end{definition}


\begin{definition}[si-Funktion]
$$
  x(t) =
  \begin{cases}
    A\,\dfrac{\sin(\pi t/T)}{\pi t/T} & t \neq 0 \\[6pt]
    A                                  & t = 0
  \end{cases}
  \qquad \Leftrightarrow \qquad
  x(t) = A\,\operatorname{si}\ \left(\frac{t}{T}\right)
$$


\begin{center}
\begin{tikzpicture}
\begin{axis}[
  width=6.5cm, height=3cm,
  xmin=-3.5, xmax=3.5, ymin=-0.35, ymax=1.3,
  axis lines=center,
  xtick={-3,-2,-1,1,2,3},
  xticklabels={$-3T$,$-2T$,$-T$,$T$,$2T$,$3T$},
  ytick={1}, yticklabels={$A$},
  xlabel={$t$},
  tick style={color=black, thin}, tick label style={font=\tiny},
]
\addplot[blue, very thick, smooth, domain=-3.5:-0.05, samples=200] {sin(180*x)/(pi*x)};
\addplot[blue, very thick, smooth, domain= 0.05: 3.5, samples=200] {sin(180*x)/(pi*x)};
\addplot[only marks, mark=*, mark options={fill=blue, draw=blue}, mark size=2pt]
  coordinates {(0,1)};
\end{axis}
\end{tikzpicture}
\end{center}

\end{definition}

\subsection{Komplexe Exponentialsignale}
\label{sec:komplexe-exponentialsignale}

Im Zusammenhang mit linearen Netzwerken und Systemen im eingeschwungenen Zustand
wird häufig auch das komplexe Exponentialsignal verwendet,

\begin{align*}
  \underline{x}(t)
    &= A e^{j(\omega t + \varphi)}
     = A\bigl(\cos(\omega t + \varphi) + j\sin(\omega t + \varphi)\bigr) \\
    &= A e^{j\varphi} \cdot e^{j\omega t}
     = \underline{A}\, e^{j\omega t},
\end{align*}

wobei $\underline{A} = A e^{j\varphi}$ die \textit{komplexe Amplitude} bezeichnet.
Der Realteil und der Imaginärteil ergeben sich wie folgt,

\begin{align*}
  \operatorname{Re}\bigl\{\underline{x}(t)\bigr\} &= A\cos(\omega t + \varphi) \\
  \operatorname{Im}\bigl\{\underline{x}(t)\bigr\} &= A\sin(\omega t + \varphi).
\end{align*}

\medskip
Erzeugung des Sinus- oder Cosinussignals durch einen in der komplexen
Zahlenebene rotierenden Zeiger:

\begin{center}
\begin{tikzpicture}[font=\small, >=stealth]

%% φ = 40°, ωt₀ = φ = 40° → total phasor angle = 80°
\def\R{2.0}
\def\phideg{40}
\def\totaldeg{80}

%% ── LEFT: Phasor in the complex plane ───────────────────────
\draw[->] (-2.7,0) -- (2.7,0);
\draw[->] (0,-2.7) -- (0,3.1);
\draw[thick] (0,0) circle (\R);

%% Phasor arrow
\draw[->, very thick]
  (0,0) -- ({\R*cos(\totaldeg)},{\R*sin(\totaldeg)});
\node[above left] at ({0.58*\R*cos(\totaldeg)},{0.52*\R*sin(\totaldeg)}) {$A$};

%% φ inner arc
\draw[->] (0.55,0) arc (0:\phideg:0.55);
\node at ({0.90*cos(\phideg/2)},{0.42*sin(\phideg/2)}) {$\varphi$};

%% ωt₀+φ outer arc
\draw[->] (1.1,0) arc (0:\totaldeg:1.1);
\node[right, font=\scriptsize]
  at ({1.45*cos(\totaldeg/2)},{1.25*sin(\totaldeg/2)})
  {$\omega t_0{+}\varphi$};

%% Dashed projection lines
\draw[dashed, gray]
  ({\R*cos(\totaldeg)},{\R*sin(\totaldeg)}) -- (0,{\R*sin(\totaldeg)});
\draw[dashed, gray]
  ({\R*cos(\totaldeg)},{\R*sin(\totaldeg)}) -- ({\R*cos(\totaldeg)},0);

%% Projection labels
\node[left, font=\scriptsize]
  at (-0.1,{\R*sin(\totaldeg)}) {$A\sin(\omega t_0{+}\varphi)$};
\node[below, font=\scriptsize]
  at ({\R*cos(\totaldeg)},-0.2) {$A\cos(\omega t_0{+}\varphi)$};

%% ω rotation indicator
\draw[->, thick]
  ({2.3*cos(88)},{2.3*sin(88)}) arc(88:108:2.3);
\node[font=\scriptsize] at ({2.58*cos(99)},{2.6*sin(99)}) {$\omega$};

%% ── RIGHT: Sine wave (pgfplots) ──────────────────────────────
%% φ_rad = 0.698, sin value at ωt₀: 2·sin(80°) ≈ 1.970
%% zero crossings: π−φ ≈ 2.443, 2π−φ ≈ 5.585
\begin{axis}[
  at={(5.2cm,-2.2cm)}, anchor=south west,
  width=6.5cm, height=4.8cm,
  xmin=-0.1, xmax=6.8, ymin=-2.5, ymax=2.7,
  axis lines=center,
  xtick={0.698, 2.443, 5.585},
  xticklabels={$\omega t_0$, $\pi{-}\varphi$, $2\pi{-}\varphi$},
  ytick={2,-2}, yticklabels={$A$,$-A$},
  xlabel={$\omega t$},
  ylabel={$\operatorname{Im}\bigl\{\underline{x}(t)\bigr\}$},
  ylabel style={at={(axis description cs:0.18,1.06)}, anchor=south},
  tick style={color=black, thin},
  tick label style={font=\scriptsize},
  clip=false,
]
\addplot[very thick, smooth, domain=0:6.5, samples=400]
  {2*sin(deg(x)+40)};
\node[font=\scriptsize] at (axis cs:3.6,2.15) {$A\sin(\omega t{+}\varphi)$};

%% Dashed lines at ωt₀ = 0.698 (value 1.970)
\draw[dashed,gray,thin] (axis cs:0.698,0)     -- (axis cs:0.698,1.970);
\draw[dashed,gray,thin] (axis cs:0,1.970)     -- (axis cs:0.698,1.970);

%% φ phase indicator
\draw[<->] (axis cs:0,-0.55) -- (axis cs:0.698,-0.55)
  node[midway,below,font=\scriptsize] {$\varphi$};
\end{axis}

\end{tikzpicture}
\end{center}

\subsection{Dirac-Impuls $\delta(t)$}
\label{sec:dirac-impuls}

Der Dirac-Impuls ist ein Signal, das für die mathematische Behandlung der Abtastung analoger Signale von großer Bedeutung ist. Den Dirac-Impuls erhält man im Grenzwert $T \to 0$ aus Rechteckfunktionen der Form $x(t) = \frac{1}{T}\operatorname{rect}\!\bigl(\frac{t + T/2}{T}\bigr)$.

Sei $f:\mathbb{R} \rightarrow \mathbb{R}$ eine beliebige stetige Funktion. Der Dirac-Impuls ist
dann durch den folgenden Zusammenhang definiert:

\begin{definition}[Siebeigenschaft des Dirac-Impulses]
$$
  \int_{-\infty}^{\infty} f(\xi)\,\delta(\xi - a)\,\mathrm{d}\xi = f(a)
$$
\end{definition}

Der Dirac-Impuls liefert somit den Funktionswert $f(\xi)$ an der Stelle $\xi = a$
(\textit{Siebeigenschaft}).


Im Spezialfall $f(\xi) = 1$ folgt:
$$
  \int_{-\infty}^{\infty} \delta(\xi)\,\mathrm{d}\xi = 1\,.
$$

Ein analoges Signal $x(t)$ kann somit auch wie folgt dargestellt werden:
$$
  x(t) = \int_{-\infty}^{\infty} x(\xi)\,\delta(\xi - t)\,\mathrm{d}\xi
        = \int_{-\infty}^{\infty} x(\xi)\,\delta(t - \xi)\,\mathrm{d}\xi\,.
$$

Aufgrund der Siebeigenschaft wird der Dirac-Impuls häufig zur Beschreibung von
Abtastvorgängen eingesetzt. Dabei wird der Dirac-Impuls $\delta(t-a)$ wie ein
Signal visualisiert, wobei das Signal für $t \neq a$ Null ist.

In dem Katalog der bisher behandelten Elementarsignale nimmt der Dirac-Impuls eine Sonderstellung ein, da er nicht der Kategorie der Funktionen zuzuordnen ist. Es handelt sich hierbei um eine sogenannte \textit{Distribution}. Daher sind einerseits die Methoden der Standardanalysis nicht ohne weiteres anwendbar. Zum Beispiel ist die Frage nach der „Energie" oder der „Leistung" des Dirac-Impulses nicht sinnvoll zu beantworten. Andererseits erweitert der Dirac-Impuls den mathematischen Werkzeugkasten und ist u.\,a.\ im Zusammenhang mit der Fourier-Transformation von Bedeutung.

\begin{center}
\begin{tikzpicture}[>=stealth, font=\small]
\begin{axis}[
  width=9cm, height=5.5cm,
  xmin=-5, xmax=8,
  ymin=-0.75, ymax=1.2,
  axis lines=center,
  xtick={3.5}, xticklabels={$a$},
  ytick=\empty,
  xlabel={$t$},
  tick style={color=black, thin},
  tick label style={font=\small},
  clip=false,
]
%% f(t): left hump, dip near origin, main peak at a=3.5
\addplot[blue, very thick, smooth, domain=-5:8, samples=500]
  {0.5*exp(-0.8*(x+3)^2) - 0.55*exp(-1.2*x^2) + 0.9*exp(-0.7*(x-3.5)^2)};

%% y-axis label
\node[above right, font=\small] at (axis cs:0.1, 1.1) {$f(t)$};

%% dashed vertical reference at t=a (peak ≈ 0.9)
\draw[dashed, gray, thin] (axis cs:3.5,0) -- (axis cs:3.5,0.9);

%% f(a) annotation
\draw[->, thin] (axis cs:4.4,0.78)
  node[right, font=\small] {$f(a)$}
  -- (axis cs:3.56,0.893);
\end{axis}
\end{tikzpicture}
\end{center}

\section{Elementare Operationen}
\label{sec:elementare-operationen}

Mit Hilfe der elementaren Operationen werden aus elementaren Signalen beliebige zusammengesetzte Signale erzeugt. Die folgenden Operationen sind hierbei besonders wichtig:

\begin{itemize}
  \item \textbf{Multiplikation mit einer Konstanten:} $y(t) = a\,x(t)$.\\
  Beispiel: $x(t) = \sin(\omega t)$ und $a = 1/3$.

  \item \textbf{Verzögerung um $T$:} $y(t) = x(t - T)$.\\
  Beispiel:
  $$
    x(t) = e^{-t/\tau}\,\epsilon(t) = \begin{cases} 0 & t < 0 \\ e^{-t/\tau} & t \geq 0 \end{cases},
    \qquad T = \tau.
  $$

  \item \textbf{Addition zweier Signale:} $y(t) = x_1(t) + x_2(t)$.\\
  Beispiel: $x_1(t) = \sin(\omega t)$ und $x_2(t) = \tfrac{1}{3}\sin(3\omega t)$.

  \item \textbf{Multiplikation zweier Signale:} $y(t) = x_1(t)\,x_2(t)$.\\
  Beispiel: $x_1(t) = e^{t/\tau}\,\epsilon(t)$, $x_2(t) = \sin(\omega t)$ und $\frac{1}{f} = T = \tau$.

  \item \textbf{Zeitliche Spiegelung:} $y(t) = x(-t)$.\\
  Beispiel: $x(t) = e^{-t/\tau}\,\epsilon(t)$.

  \item \textbf{Zeitliche Skalierung:} $y(t) = x(at)$.\\
  Beispiel: $x(t) = e^{-t/\tau}\,\epsilon(t)$ und $a = 3$.
\end{itemize}

\begin{beispiel}[Übertragung digitaler Signale durch Amplitudenmodulation]
In der Kommunikationstechnik werden Informationen den Signalen durch „Modulation" aufgeprägt. Ein einfaches Beispiel hierfür ist die Amplitudenmodulation:

$$
x(t) = A(t)\cos(\omega_0 t + \varphi_0).
$$

Im einfachsten Fall werden nur zwei Amplitudenstufen verwendet, d.\,h.\ $A(t) \in \{A_0, A_1\}$. Hierbei ist $A(t)$ das Nachrichtensignal und $\cos(\omega_0 t + \varphi_0)$ das Trägersignal.

% Grafik: A(t) als binäres Rechtecksignal, x(t) als amplitudenmoduliertes Trägersignal
\end{beispiel}


\section{Elementare diskrete Signale}
\label{sec:elementare-diskrete-signale}

Auch für zeitdiskrete Signale sollen zunächst einige elementare Signale definiert werden. Das (zeit-)diskrete Signal $x$ wird als Folge von Werten $x[k]$ mit $k \in \mathbb{Z}$ definiert:

$$
x:\mathbb{Z} \rightarrow \mathbb{R}
$$

oder, als komplexwertiges Signal,

$$
\underline{x}:\mathbb{Z} \rightarrow \mathbb{C}.
$$

Wenn eine physikalische Interpretation des diskreten Signals gewünscht ist, muss die Dauer des Abtastintervalls bzw. deren Kehrwert, die Abtastrate $f_A = \frac{1}{T_A}$, bekannt sein.

\subsection*{Die diskrete Einheitsimpulsfolge}

Für den Einheitsimpuls wird das Symbol $\delta[k]$ verwendet. Er ist wie folgt definiert:

$$
\delta[k] = \begin{cases} 1 & k = 0 \\ 0 & k \neq 0. \end{cases}
$$

% Grafik: δ[k] für k von -30 bis 30, Einheitsimpuls bei k=0

\subsection*{Das diskrete Einheitssprungsignal}

Für den diskreten Einheitssprung wird das Symbol $\epsilon[k]$ verwendet. Er ist wie folgt definiert:

$$
\epsilon[k] = \begin{cases} 1 & k \geq 0 \\ 0 & k < 0. \end{cases}
$$

% Grafik: ε[k] für k von -30 bis 30, Sprung bei k=0

\subsection*{Das diskrete Rechtecksignal}

Aus zwei Einheitssprungsignalen kann ein diskretes Rechtecksignal mit $M$ von Null verschiedenen Werten gebildet werden:

$$
x[k] = \epsilon[k] - \epsilon[k - M].
$$

Die Werte dieses diskreten Signals lassen sich auch mit $x[k] = \operatorname{rect}(k/M)$, $k \in \mathbb{Z}$, angeben.

% Grafik: diskretes Rechtecksignal für M=10

Die Energie eines diskreten Rechtecksignals $x[k] = A\,(\epsilon[k] - \epsilon[k-M])$ mit der Amplitude $A$ und der Breite $M$ beträgt:

$$
E = \sum_{k=-\infty}^{\infty} x^2[k] = \sum_{k=0}^{M-1} A^2 = M A^2.
$$

\subsection*{Die diskrete Exponentialfolge}

Die Exponentialfolge ist für $a \in \mathbb{R}$ und $k \in \mathbb{Z}$ wie folgt definiert:

$$
x[k] = a^k\,\epsilon[k].
$$

Man beachte, dass der Betrag dieses Signals für $|a| > 1$ und $k \to \infty$ über alle Grenzen wächst. Für $|a| < 1$ klingen die Beträge des Signals ab.

% Grafik: Exponentialfolge für a = 0.9

\begin{bemerkung}[Geometrische Reihe]
Für das Rechnen mit Exponentialfolgen sind generell die Summenformeln

$$
a_0 \sum_{k=0}^{M-1} q^k = a_0\,\frac{1 - q^M}{1 - q}\,; \quad q \neq 1
$$

$$
a_0 \sum_{k=0}^{\infty} q^k = a_0\,\frac{1}{1 - q}\,; \quad |q| < 1
$$

hilfreich.
\end{bemerkung}

\subsection*{Das diskrete Sinussignal}

Das diskrete Sinussignal ist wie folgt definiert:

$$
x[k] = A\sin\!\left(2\pi k\,\frac{f}{f_A}\right).
$$

% Grafik: diskrete Sinusfolge für f/f_A = 1/8

\section{Elementare Operationen diskreter Signale}
\label{sec:elementare-operationen-diskret}

Die Operationen für diskrete Signale sind analog zu den Operationen kontinuierlicher Signale definiert. Es werden daher nur zwei einfache Beispiele betrachtet.

\begin{itemize}
  \item \textbf{Addition zweier Einheitsimpulsfolgen:} $y[k] = \delta[k] + \delta[k - 10]$

  % Grafik: Summe zweier Einheitsimpulse bei k=0 und k=10

  \item \textbf{Multiplikation von Exponential- und Sinusfolge:} $y[k] = a^k\,\epsilon[k]\,\sin\!\left(2\pi\dfrac{f}{f_A}\,k\right)$

  % Grafik: Produkt zweier Elementarsignale
\end{itemize}

\section{Schwebung}
\label{sec:schwebung}

\begin{beispiel}[Schwebung]
Eine Schwebung entsteht durch die Addition zweier Sinussignale $x_1(t)$ und $x_2(t)$ mit gleicher Amplitude und geringer Frequenzdifferenz $2\Delta\omega$. Zur Berechnung der Summe verwenden wir die komplexwertige Schreibweise:

$$
\begin{aligned}
\underline{y}(t) &= A e^{j(\omega_0+\Delta\omega)t} + A e^{j(\omega_0-\Delta\omega)t} \\
                 &= A e^{j\omega_0 t} \bigl( e^{j\Delta\omega t} + e^{-j\Delta\omega t} \bigr) \\
                 &= 2A \cos(\Delta\omega\, t)\, e^{j\omega_0 t},
\end{aligned}
$$

wobei $\omega_0 = 2\pi f_0$ die mittlere Kreisfrequenz angibt. Die Bildung des Real- und Imaginärteils ergibt:

$$
\operatorname{Re}\{\underline{y}(t)\} = 2A\cos(\Delta\omega\, t)\cos(\omega_0 t)
$$
$$
\operatorname{Im}\{\underline{y}(t)\} = 2A\cos(\Delta\omega\, t)\sin(\omega_0 t).
$$

In den reellwertigen Signalen ist gut zu erkennen, dass die Schwebung sich durch eine Modulation der Amplitude mit $\cos(\Delta\omega\, t)$ äußert. Die Kreisfrequenz dieser Amplitudenmodulation ist durch $\Delta\omega$ gegeben. Der Schwebungseffekt kann z.\,B.\ zum Stimmen von Saiteninstrumenten genutzt werden, da eine Amplitudenmodulation leichter wahrzunehmen ist als sehr kleine Frequenzabweichungen.

In den folgenden Beispielen verwenden wir reellwertige, zeitdiskrete Signale der Abtastrate $f_A = 1/T_A$:

$$
x_1[k] = \sin\bigl((\omega_0 - \Delta\omega)\,k T_A\bigr), \qquad
x_2[k] = \sin\bigl((\omega_0 + \Delta\omega)\,k T_A\bigr).
$$

Das Summensignal ist als Mittelwert der Teilsignale definiert:

$$
z(k T_A) = \frac{1}{2}(x_1[k] + x_2[k]) = \frac{1}{2}\operatorname{Im}\{\underline{y}(k T_A)\}.
$$

Die Mittenkreisfrequenz sei $\omega_0 = 2\pi \cdot 800\,\frac{1}{\text{s}}$. Bei einer Frequenzdifferenz $\Delta\omega = 0$ hören wir einen Ton bei 800\,Hz. Mit wachsender Frequenzdifferenz tritt der Schwebungseffekt in Erscheinung: die Amplitude wird langsam moduliert. In den $xy$-Darstellungen der Teilsignale $x_1[k]$ und $x_2[k]$ entstehen dabei die sogenannten \textbf{Lissajous-Figuren}. Während die oberen Bilder regelmäßige Lissajous-Figuren zeigen, ist in den unteren Bildern die Amplitudenmodulation deutlich zu erkennen.

% Grafik: xy-Darstellung (Lissajous) und Summensignal für Δω = 0
% Grafik: xy-Darstellung (Lissajous) und Summensignal für Δω = 0.5 ω₀
% Grafik: xy-Darstellung (Lissajous) und Summensignal für Δω = 0.1 ω₀
% Grafik: xy-Darstellung (Lissajous) und Summensignal für Δω = 0.01 ω₀
% Grafik: xy-Darstellung (Lissajous) und Summensignal für Δω = 0.001 ω₀
\end{beispiel}

\chapter{Signaleigenschaften und Kenngrößen}

\section{Symmetrie reellwertiger Signale}

\begin{definition}[gerade symmetrisch]
$$
x(t) = x(-t)\ \forall t \in \mathbb{R}\ \text{bzw.}\ x[k] = x[-k]\ \forall k \in \mathbb{Z}\
$$
\end{definition}

\begin{definition}[ungerade symmetrisch]
$$
x(t) = -x(-t)\ \forall t \in \mathbb{R}\ \text{bzw.}\ x[k] = -x[-k]\ \forall k \in \mathbb{Z}\
$$
\end{definition}

Analog dazu ergänzen sich symmetrie Eigenschaften für komplexwertige Signale.

\begin{definition}[konjugiert gerade symmetrisch]
$$
x(t) = x*(-t)\ \forall t \in \mathbb{R}\ \text{bzw.}\ x[k] = x*[-k]\ \forall k \in \mathbb{Z}\
$$
\end{definition}

Hierbei ist der Realteil von x eine gerade symmetrische Funktion und
der Imaginärteil eine ungerade symmetrische Funktion.


\begin{definition}[konjugiert ungerade symmetrisch]
$$
x(t) = -x*(-t)\ \forall t \in \mathbb{R}\ \text{bzw.}\ x[k] = -x*[-k]\ \forall k \in \mathbb{Z}\
$$
\end{definition}

Hierbei ist der Realteil von x eine ungerade symmetrische Funktion
und der Imaginärteil eine gerade symmetrische Funktion.


\section{Periodische Signale}

\begin{definition}[Periodisches analoges Signal]
Ein analoges Signal $x(t)$ heißt periodisch, falls es eine Periodendauer
$T > 0$ gibt, so dass
$$
x(t) = x(t + nT)
$$
für alle $t \in R$ und $n \in Z$ gilt.
Die kleinste Periodendauer $T$ wird Grundperiode und deren Kehrwert $f_0 = \frac{1}{T}$ wird
Grundfrequenz genannt.
\end{definition}

\begin{definition}[Periodisches zeitdiskretes Signal]
Ein zeitdiskretes Signal $x[k]$ heißt periodisch, falls es eine diskrete
Periode $N \in \mathbb{N}$ gibt, so dass
$$
x[k] = x[k + n \cdot N]
$$
für alle $k \in \mathbb{Z}$ und $n \in \mathbb{Z}$ gilt.
\end{definition}


Die Abtastung eines periodischen analogen Signals ergibt nicht notwendigerweise auch
ein periodisches zeitdiskretes Signal.


\section{Mittelwert}

\subsection{Mittelwert diskreter Signale}

\begin{definition}[Mittelwert]
\textbf{Mittelwert} $\bar{x}$ des zeitkontinuierlichen Signals $x(t)$ im Zeitintervall $[t_1, t_2]$
bzw. eines zeitdiskreten Signals $x[k]$ im Bereich $k = k_1, \ldots, k_2$:

$$
\bar{x} = \frac{1}{t_2 - t_1} \int_{t_1}^{t_2} x(t)\,\mathrm{d}t
\qquad \text{bzw.} \qquad
\bar{x} = \frac{1}{k_2 - k_1 + 1} \sum_{k=k_1}^{k_2} x[k]
$$

Ein Signal $x$ heißt mittelwertfrei, falls $\bar{x} = 0$ gilt.


Für \textbf{komplexwertige Signale} $\underline{x}$ für den Mittelwert eines zeitkontinuierlichen
Signals im Intervall $[t_1, t_2]$

$$
\bar{\underline{x}} = \frac{1}{t_2 - t_1} \int_{t_1}^{t_2} \underline{x}(t)\,\mathrm{d}t
= \overline{\operatorname{Re}\{\underline{x}(t)\}} + j\,\overline{\operatorname{Im}\{\underline{x}(t)\}}.
$$
\end{definition}

\subsection{Mittelwert periodischer Signale}

\begin{definition}[Mittelwert periodischer Signale]
Für \textbf{periodische Signale} $x(t)$ bzw. $x[k]$ wird der Mittelwert über die
Dauer einer Grundperiode $T$ bzw. $N$ berechnet:

$$
\bar{x} = \frac{1}{T} \int_{-T/2}^{T/2} x(t)\,\mathrm{d}t
\qquad \text{bzw.} \qquad
\bar{x} = \frac{1}{N} \sum_{k=0}^{N-1} x[k].
$$

bzw. für komplexwertige, periodische Signale:

$$
\bar{\underline{x}} = \frac{1}{T} \int_{-T/2}^{T/2} \underline{x}(t)\,\mathrm{d}t
\qquad \text{bzw.} \qquad
\bar{\underline{x}} = \frac{1}{N} \sum_{k=0}^{N-1} \underline{x}[k].
$$
\end{definition}

\section{Energie}

\begin{definition}[Energie und Energiesignal]
Die \textbf{Energie} $E$ des Signals $\underline{x}$ im Zeitintervall $[t_1, t_2]$ bzw. eines zeitdiskreten
Signals $\underline{x}[k]$ im Bereich $k = k_1, \ldots, k_2$:

$$
E = \int_{t_1}^{t_2} |\underline{x}(t)|^2\,\mathrm{d}t
\qquad \text{bzw.} \qquad
E = \sum_{k=k_1}^{k_2} |\underline{x}[k]|^2
$$

gegeben. Falls das Integral bzw. die Summe

$$
E = \int_{-\infty}^{\infty} |\underline{x}(t)|^2\,\mathrm{d}t < \infty
\qquad \text{bzw.} \qquad
E = \sum_{k=-\infty}^{\infty} |\underline{x}[k]|^2 < \infty
$$

existiert, wird $\underline{x}$ als \textbf{Energiesignal} bezeichnet.
\end{definition}


Ferner halten wir fest, dass für reellwertige Signale $|\underline{x}(t)|^2 = x^2(t) = x(t)x(t)$
und für komplexwertige Signale $|\underline{x}(t)|^2 = x(t)x^*(t)$ gilt, wobei $\underline{x}^*(t)$ das zu $\underline{x}(t)$
konjugiert komplexe Signal bezeichnet.


\section{Leistung}


Einige wichtige Elementarsignale (Sprungfunktion, Sinussignal) sind
keine Energiesignale. Daher wird die Leistung eingeführt.


\begin{definition}[Leistung und Leistungssignal]
\textbf{Leistung} (mittlere Energie pro Zeitintervall) des Signals $\underline{x}$:

$$
P = \lim_{T \to \infty} \frac{1}{T} \int_{-T/2}^{T/2} |\underline{x}(t)|^2\,\mathrm{d}t
\qquad \text{bzw.} \qquad
P = \lim_{K \to \infty} \frac{1}{2K+1} \sum_{k=-K}^{K} |\underline{x}[k]|^2
$$

\textbf{Leistung} des periodischen Signals $\underline{x}$:

$$
P = \frac{1}{T} \int_{t_0}^{t_0+T} |\underline{x}(t)|^2\,\mathrm{d}t
\qquad \text{bzw.} \qquad
P = \frac{1}{N} \sum_{k=k_0}^{k_0+N-1} |\underline{x}[k]|^2,
$$

wobei der Startpunkt $t_0$ bzw. $k_0$ beliebig, also auch $t_0 = 0$ bzw. $k_0 = 0$,
gewählt werden kann. Falls der Grenzwert mit $P > 0$ existiert, wird
das Signal als \textbf{Leistungssignal} bezeichnet.
\end{definition}

\chapter{Inneres Produkt (Skalarprodukt)}

\begin{definition}[Inneres Produkt]
Das innere Produkt wird für zwei \textbf{Energiesignale} $\underline{x}$ und $\underline{y}$ wie folgt definiert:

$$
\langle \underline{x}(t), \underline{y}(t) \rangle = \int_{-\infty}^{\infty} \underline{x}^*(t)\underline{y}(t)\,\mathrm{d}t
\qquad \text{bzw.} \qquad
\langle \underline{x}[k], \underline{y}[k] \rangle = \sum_{k=-\infty}^{\infty} \underline{x}^*[k]\underline{y}[k]
$$

Dabei muss das Integral bzw. die Summe einen endlichen Wert ergeben.
\end{definition}

\begin{definition}[Orthogonalität und Kollinearität]
Wenn das innere Produkt Null ist, werden die Signale $\underline{x}$ und $\underline{y}$ \textbf{orthogonal}
genannt. Zwei Signale, die bis auf einen Skalierungsfaktor identisch sind,
werden \textbf{kollinear} genannt.
\end{definition}


Für \textbf{periodische Signale} der Periode $T$ bzw. $N$ gilt entsprechend

$$
\langle \underline{x}(t), \underline{y}(t) \rangle = \frac{1}{T} \int_{t_0}^{t_0+T} \underline{x}^*(t)\underline{y}(t)\,\mathrm{d}t
\qquad \text{bzw.} \qquad
\langle \underline{x}[k], \underline{y}[k] \rangle = \frac{1}{N} \sum_{k=k_0}^{k_0+N-1} \underline{x}^*[k]\underline{y}[k].
$$


\section{Kreuzkorrelation}

\begin{definition}[Kreuzkorrelationsfunktion (KKF)]
Für zwei Energiesignale $\underline{x}$ und $\underline{y}$ wird die Ähnlichkeit (Kreuzkorrelation) wie folgt definiert:

$$
\varphi_{xy}(\tau) = \int_{-\infty}^{\infty} \underline{x}^*(t)\,\underline{y}(t+\tau)\,\mathrm{d}t
\qquad \text{bzw.} \qquad
\varphi_{xy}[\varkappa] = \sum_{k=-\infty}^{\infty} \underline{x}^*[k]\,\underline{y}[k+\varkappa].
$$

Das Integral bzw. die Summe wird für alle zeitlichen Verschiebungen $\tau$ bzw. $\varkappa$ des Signals $\underline{y}$ ausgeführt.
\end{definition}


Im Allgemeinen ist die KKF eine komplexwertige Funktion bzw. Folge. Im Fall reellwertiger Signale vereinfacht sich die Schreibweise zu

$$
\varphi_{xy}(\tau) = \int_{-\infty}^{\infty} x(t)\,y(t+\tau)\,\mathrm{d}t
\qquad \text{bzw.} \qquad
\varphi_{xy}[\varkappa] = \sum_{k=-\infty}^{\infty} x[k]\,y[k+\varkappa].
$$


\newpage
\textbf{Schema zur Kreuzkorrelation (zeitdiskret und zeitkontinuierlich)}

Das Schema gilt für beide Signalklassen. In der Tabelle steht jeweils die diskrete
Variante links und die kontinuierliche rechts.

\begin{enumerate}

  \item \textbf{Träger bestimmen.}
  Finde den Bereich, in dem jedes Signal ungleich null ist:
  \begin{center}
  \begin{tabular}{ll}
    \textit{diskret} & \textit{kontinuierlich} \\[4pt]
    $\mathrm{supp}(x) = \{k : x[k] \neq 0\}$ & $\mathrm{supp}(x) = \{t \in \mathbb{R} : x(t) \neq 0\}$ \\
    $\mathrm{supp}(y) = \{k : y[k] \neq 0\}$ & $\mathrm{supp}(y) = \{t \in \mathbb{R} : y(t) \neq 0\}$
  \end{tabular}
  \end{center}
  Im kontinuierlichen Fall ist der Träger meist ein Intervall $[a, b]$.

  \item \textbf{Verschiebebereich einschränken.}
  $\varphi_{xy}$ ist nur dort ungleich null, wo sich die Träger von $x$ und dem
  verschobenen $y(\cdot + \tau)$ überlappen.
  Der relevante Bereich ist in beiden Fällen:
  $$
    \tau\ (\text{bzw.}\ \varkappa) \;\in\;
    \bigl[\min\mathrm{supp}(y) - \max\mathrm{supp}(x),\;
          \max\mathrm{supp}(y) - \min\mathrm{supp}(x)\bigr].
  $$
  Außerhalb gilt $\varphi_{xy} = 0$.

  \item \textbf{Verschiebung skizzieren (optional).}
  Zeichne $x$ fest und schiebe $y$ um $\tau$ bzw.\ $\varkappa$.
  Markiere das Überlappungsintervall — nur dort leistet das Produkt
  $x \cdot y(\cdot + \tau)$ einen Beitrag.

  \item \textbf{Integral bzw.\ Summe auswerten.}
  Berechne $\varphi_{xy}$ für jeden relevanten Verschiebewert:
  \begin{center}
  \begin{tabular}{ll}
    \textit{diskret} & \textit{kontinuierlich} \\[4pt]
    $\displaystyle\varphi_{xy}[\varkappa] = \sum_{k} x[k]\, y[k+\varkappa]$
    &
    $\displaystyle\varphi_{xy}(\tau) = \int_{-\infty}^{\infty} x(t)\, y(t+\tau)\,\mathrm{d}t$
  \end{tabular}
  \end{center}
  Im kontinuierlichen Fall werden die Integrationsgrenzen durch das
  Überlappungsintervall aus Schritt~2 und~3 bestimmt; das Integral wird
  abschnittsweise ausgewertet, falls sich die Grenzen mit $\tau$ ändern.

  \item \textbf{Ergebnis zusammenstellen und prüfen.}
  \begin{itemize}
    \item \textit{diskret:} Die KKF-Folge hat
          $|\mathrm{supp}(x)| + |\mathrm{supp}(y)| - 1$ Elemente.
    \item \textit{kontinuierlich:} Die KKF ist eine stückweise definierte Funktion
          von $\tau$; typischerweise ein Polynom oder eine Exponentialfunktion
          je Abschnitt.
    \item In beiden Fällen gilt für reellwertige Signale:
          $\varphi_{yx}(\tau) = \varphi_{xy}(-\tau)$.
    \item Das Maximum von $|\varphi_{xy}|$ zeigt die Verschiebung,
          bei der $y$ $x$ am ähnlichsten ist.
  \end{itemize}

\end{enumerate}

\newpage
\subsection*{Rechenbeispiele}

\begin{beispiel}[Diskrete Kreuzkorrelation]

Gegeben seien die Folgen
$$
x[k] = \begin{cases} 1 & k=0 \\ 2 & k=1 \\ 1 & k=2 \\ 0 & \text{sonst} \end{cases}
\qquad \text{und} \qquad
y[k] = \begin{cases} 1 & k=0 \\ 1 & k=1 \\ 0 & \text{sonst} \end{cases}
$$

Gesucht ist $\varphi_{xy}[\varkappa] = \sum_{k=-\infty}^{\infty} x[k]\,y[k+\varkappa]$ für alle $\varkappa$.

\begin{enumerate}
  \item \textbf{$\varkappa = -2$:} $y[k+(-2)]$ ist ungleich null bei $k = 2, 3$.
  $$\varphi_{xy}[-2] = x[2]\cdot y[0] + x[3]\cdot y[1] = 1\cdot 1 + 0\cdot 1 = 1$$

  \item \textbf{$\varkappa = -1$:} $y[k-1]$ ist ungleich null bei $k = 1, 2$.
  $$\varphi_{xy}[-1] = x[1]\cdot y[0] + x[2]\cdot y[1] = 2\cdot 1 + 1\cdot 1 = 3$$

  \item \textbf{$\varkappa = 0$:} $y[k]$ ist ungleich null bei $k = 0, 1$.
  $$\varphi_{xy}[0] = x[0]\cdot y[0] + x[1]\cdot y[1] = 1\cdot 1 + 2\cdot 1 = 3$$

  \item \textbf{$\varkappa = 1$:} $y[k+1]$ ist ungleich null bei $k = -1, 0$.
  $$\varphi_{xy}[1] = x[-1]\cdot y[0] + x[0]\cdot y[1] = 0\cdot 1 + 1\cdot 1 = 1$$

  \item \textbf{Ergebnis:}
  $$\varphi_{xy}[\varkappa] = \{\ldots,\, 0,\, \underset{\varkappa=-2}{1},\, 3,\, 3,\, 1,\, 0,\, \ldots\}$$
\end{enumerate}

\end{beispiel}

\begin{beispiel}[Kontinuierliche Kreuzkorrelation zweier Rechteckimpulse]

Gegeben seien
$$
x(t) = \begin{cases} 1 & 0 \leq t \leq 1 \\ 0 & \text{sonst} \end{cases}
\qquad \text{und} \qquad
y(t) = \begin{cases} 1 & 0 \leq t \leq 2 \\ 0 & \text{sonst.} \end{cases}
$$
Gesucht ist $\varphi_{xy}(\tau) = \int_{-\infty}^{\infty} x(t)\,y(t+\tau)\,\mathrm{d}t$.

\begin{enumerate}
  \item \textbf{Träger bestimmen.}
  $$\mathrm{supp}(x) = [0,\,1], \qquad \mathrm{supp}(y) = [0,\,2].$$

  \item \textbf{Verschiebebereich einschränken.}
  $$\tau \in [\min\mathrm{supp}(y) - \max\mathrm{supp}(x),\;
              \max\mathrm{supp}(y) - \min\mathrm{supp}(x)]
         = [0-1,\; 2-0] = [-1,\,2].$$
  Außerhalb gilt $\varphi_{xy}(\tau) = 0$.

  \item \textbf{Überlappungsintervall je Abschnitt bestimmen.}
  $y(t+\tau)$ ist ungleich null für $t \in [-\tau,\, 2-\tau]$.
  Der Schnitt mit $\mathrm{supp}(x) = [0,1]$ ergibt drei Abschnitte:

  \begin{center}
  \begin{tabular}{cll}
    Abschnitt & Überlappung & Begründung \\[3pt]
    $\tau \in [-1,\,0]$ & $[-\tau,\;1]$ & $-\tau\in[0,1]$, $2{-}\tau \geq 2 > 1$ \\
    $\tau \in [\phantom{-}0,\,1]$ & $[0,\;1]$    & $-\tau\leq 0$, $2{-}\tau \geq 1$ \\
    $\tau \in [\phantom{-}1,\,2]$ & $[0,\;2{-}\tau]$ & $-\tau\leq 0$, $2{-}\tau\in[0,1]$
  \end{tabular}
  \end{center}

  \item \textbf{Integral abschnittsweise auswerten.}
  Da $x(t) = y(t) = 1$ auf dem jeweiligen Überlappungsintervall, vereinfacht sich
  das Integral zur Länge des Überlappungsintervalls:
  \begin{align*}
    \tau \in [-1,\,0]\colon\quad
      \varphi_{xy}(\tau) &= \int_{-\tau}^{1} 1\,\mathrm{d}t = 1 - (-\tau) = 1+\tau, \\
    \tau \in [\phantom{-}0,\,1]\colon\quad
      \varphi_{xy}(\tau) &= \int_{0}^{1} 1\,\mathrm{d}t = 1, \\
    \tau \in [\phantom{-}1,\,2]\colon\quad
      \varphi_{xy}(\tau) &= \int_{0}^{2-\tau} 1\,\mathrm{d}t = 2-\tau.
  \end{align*}

  \item \textbf{Ergebnis:}
  $$
  \varphi_{xy}(\tau) = \begin{cases}
    1+\tau & -1 \leq \tau \leq 0 \\
    1      & \phantom{-}0 \leq \tau \leq 1 \\
    2-\tau & \phantom{-}1 \leq \tau \leq 2 \\
    0      & \text{sonst.}
  \end{cases}
  $$
  Die KKF hat Trapezform: linearer Anstieg, Plateau, linearer Abfall.
  Das Plateau bei $\tau \in [0,1]$ entspricht dem Bereich, in dem $x$ vollständig
  im Träger von $y$ liegt.
\end{enumerate}

\end{beispiel}

\newpage
\subsection*{Grafisches Beispiel: Diskrete Kreuzkorrelation}

Die folgende Abbildung zeigt die Folgen $x[k]$ und $y[k]$ aus Beispiel~1 sowie die resultierende Kreuzkorrelationsfolge $\varphi_{xy}[\varkappa]$.

\begin{figure}[h]
\centering
\begin{tikzpicture}[scale=0.9]

  % --- x[k] (oben links) ---
  \begin{scope}[xshift=0cm, yshift=5.5cm]
    \draw[->] (-1,0) -- (4.5,0) node[right] {$k$};
    \draw[->] (0,-0.4) -- (0,2.8) node[above] {$x[k]$};
    \foreach \k/\v in {0/1, 1/2, 2/1} {
      \draw[very thick, blue] (\k,0) -- (\k,\v);
      \filldraw[blue] (\k,\v) circle (2.5pt);
      \filldraw[blue] (\k,0)  circle (2.5pt);
    }
    \foreach \k in {-0.5, 3, 4} { \filldraw[blue] (\k,0) circle (2.5pt); }
    \foreach \k/\l in {0/{$0$}, 1/{$1$}, 2/{$2$}, 3/{$3$}}
      { \node[below] at (\k,-0.15) {\small\l}; }
    \node[left] at (0,1) {\small$1$};
    \node[left] at (0,2) {\small$2$};
  \end{scope}

  % --- y[k] (oben rechts) ---
  \begin{scope}[xshift=6.5cm, yshift=5.5cm]
    \draw[->] (-1,0) -- (4.5,0) node[right] {$k$};
    \draw[->] (0,-0.4) -- (0,2.8) node[above] {$y[k]$};
    \foreach \k/\v in {0/1, 1/1} {
      \draw[very thick, red] (\k,0) -- (\k,\v);
      \filldraw[red] (\k,\v) circle (2.5pt);
      \filldraw[red] (\k,0)  circle (2.5pt);
    }
    \foreach \k in {-0.5, 2, 3, 4} { \filldraw[red] (\k,0) circle (2.5pt); }
    \foreach \k/\l in {0/{$0$}, 1/{$1$}, 2/{$2$}, 3/{$3$}}
      { \node[below] at (\k,-0.15) {\small\l}; }
    \node[left] at (0,1) {\small$1$};
  \end{scope}

  % --- phi_xy[kappa] (unten zentriert) ---
  \begin{scope}[xshift=4.75cm, yshift=0cm]
    \draw[->] (-3.5,0) -- (4,0) node[right] {$\varkappa$};
    \draw[->] (0,-0.4) -- (0,3.8) node[above] {$\varphi_{xy}[\varkappa]$};
    \foreach \k/\v in {-2/1, -1/3, 0/3, 1/1} {
      \draw[very thick, black] (\k,0) -- (\k,\v);
      \filldraw[black] (\k,\v) circle (2.5pt);
      \filldraw[black] (\k,0)  circle (2.5pt);
    }
    \foreach \k in {-3, 2, 3} { \filldraw[black] (\k,0) circle (2.5pt); }
    \foreach \k/\l in {-2/{$-2$}, -1/{$-1$}, 0/{$0$}, 1/{$1$}, 2/{$2$}}
      { \node[below] at (\k,-0.15) {\small\l}; }
    \node[left] at (0,1) {\small$1$};
    \node[left] at (0,3) {\small$3$};
  \end{scope}

\end{tikzpicture}
\caption{Kreuzkorrelation von $x[k]$ (blau) und $y[k]$ (rot) liefert $\varphi_{xy}[\varkappa]$ (schwarz).
Die KKF verschiebt das Maximum dorthin, wo die Überlappung der Folgen am größten ist.}
\end{figure}

\newpage
\section{Selbstähnlichkeit (Autokorrelation)}

\begin{definition}[Autokorrelationsfunktion (AKF)]
Die \textbf{Autokorrelationsfunktion (AKF)} ergibt sich als Spezialfall der KKF
mit $y = x$. Für ein \textbf{Energiesignal} gilt:
$$
\varphi_{xx}(\tau) = \int_{-\infty}^{\infty} x^*(t)\,x(t+\tau)\,\mathrm{d}t
\qquad \text{bzw.} \qquad
\varphi_{xx}[\varkappa] = \sum_{k=-\infty}^{\infty} x^*[k]\,x[k+\varkappa].
$$
\end{definition}

\begin{satz}[Eigenschaften der AKF]
\begin{itemize}
  \item \textbf{Maximum bei $\tau = 0$:}
  $$
    \varphi_{xx}(0) = \int_{-\infty}^{\infty} |x(t)|^2\,\mathrm{d}t
    \qquad \text{bzw.} \qquad
    \varphi_{xx}[0] = \sum_{k=-\infty}^{\infty} |x[k]|^2.
  $$
  Der Maximalwert entspricht der \textbf{Energie} von $x$.
  \item \textbf{Konjugiert gerade Symmetrie:}
  $\varphi_{xx}(\tau) = \varphi_{xx}^*(-\tau)$; für reellwertige Signale gilt
  $\varphi_{xx}(\tau) = \varphi_{xx}(-\tau)$.
\end{itemize}
\end{satz}

\begin{definition}[AKF für Leistungs- und periodische Signale]
Für \textbf{Leistungssignale} wird die AKF analog zur KKF über einen Grenzwert definiert:
$$
\varphi_{xx}^{L}(\tau) = \lim_{T \to \infty} \frac{1}{T}
  \int_{-T/2}^{T/2} x^*(t)\,x(t+\tau)\,\mathrm{d}t
\qquad \text{bzw.} \qquad
\varphi_{xx}^{L}[\varkappa] = \lim_{K \to \infty} \frac{1}{2K+1}
  \sum_{k=-K}^{K} x^*[k]\,x[k+\varkappa].
$$

Für \textbf{periodische Signale} der Periode $T$ bzw.\ $N$ wird nur über eine Periode
gemittelt:
$$
\tilde{\varphi}_{xx}(\tau) = \frac{1}{T}\int_{t_0}^{t_0+T} x^*(t)\,x(t+\tau)\,\mathrm{d}t
\qquad \text{bzw.} \qquad
\tilde{\varphi}_{xx}[\varkappa] = \frac{1}{N}\sum_{k=k_0}^{k_0+N-1} x^*[k]\,x[k+\varkappa].
$$
\end{definition}


Die AKF eines periodischen Signals ist selbst periodisch mit derselben Periode $T$
bzw.\ $N$. Maxima treten bei $\tau = n \cdot T$ mit $n \in \mathbb{Z}$ auf.


\subsection*{Allgemeines Vorgehen}

Das Schema für die AKF folgt dem der KKF mit $y = x$. Der einzige Unterschied:
da beide Signale identisch sind, vereinfacht sich Schritt~1 — es gibt nur einen
Träger $\mathrm{supp}(x)$. Der Verschiebebereich (Schritt~2) ist daher stets
$$
\tau\ (\text{bzw.}\ \varkappa) \in
\bigl[\min\mathrm{supp}(x) - \max\mathrm{supp}(x),\;
      \max\mathrm{supp}(x) - \min\mathrm{supp}(x)\bigr].
$$
Das Ergebnis ist immer eine \textbf{gerade} Funktion bzw.\ Folge.

\subsection*{Rechenbeispiel: AKF des Rechteckimpulses}

\begin{beispiel}[AKF des Rechteckimpulses]

Gegeben sei
$$
x(t) = \begin{cases} A & 0 \leq t \leq T_1 \\ 0 & \text{sonst.} \end{cases}
$$
Gesucht ist $\varphi_{xx}(\tau) = \int_{-\infty}^{\infty} x(t)\,x(t+\tau)\,\mathrm{d}t$.

\begin{enumerate}
  \item \textbf{Träger bestimmen.}
  $\mathrm{supp}(x) = [0,\,T_1]$.

  \item \textbf{Verschiebebereich einschränken.}
  $$\tau \in [0 - T_1,\; T_1 - 0] = [-T_1,\; T_1].$$

  \item \textbf{Überlappungsintervall je Abschnitt.}
  $x(t+\tau)$ ist ungleich null für $t \in [-\tau,\, T_1-\tau]$.
  Schnitt mit $[0, T_1]$:

  \begin{center}
  \begin{tabular}{cll}
    Abschnitt & Überlappung & Begründung \\[3pt]
    $\tau \in [-T_1, 0]$ & $[-\tau,\; T_1]$       & $-\tau \in [0,T_1]$,\; $T_1{-}\tau \geq T_1$ \\
    $\tau \in [\phantom{-}0, T_1]$ & $[0,\; T_1{-}\tau]$ & $-\tau \leq 0$,\; $T_1{-}\tau \in [0,T_1]$
  \end{tabular}
  \end{center}

  \item \textbf{Integral auswerten.}
  \begin{align*}
    \tau \in [-T_1,\, 0]\colon\quad
      \varphi_{xx}(\tau) &= \int_{-\tau}^{T_1} A^2\,\mathrm{d}t
        = A^2\,(T_1 + \tau), \\
    \tau \in [\phantom{-}0,\, T_1]\colon\quad
      \varphi_{xx}(\tau) &= \int_{0}^{T_1-\tau} A^2\,\mathrm{d}t
        = A^2\,(T_1 - \tau).
  \end{align*}

  \item \textbf{Ergebnis.}
  $$
  \varphi_{xx}(\tau) = \begin{cases}
    A^2\,(T_1 + \tau) & -T_1 \leq \tau \leq 0 \\
    A^2\,(T_1 - \tau) & \phantom{-}0 \leq \tau \leq T_1 \\
    0 & \text{sonst.}
  \end{cases}
  $$
  Die AKF ist ein \textbf{Dreieckimpuls} (Plateau entfällt, da $x$ und $y$ gleich lang
  sind). Das Maximum liegt bei $\tau = 0$ mit $\varphi_{xx}(0) = A^2 T_1$, was der
  Energie von $x$ entspricht.
\end{enumerate}

\end{beispiel}



\section{Orthogonalsysteme}

$M$ zueinander orthogonale Funktionen $x_i(t)$ bilden ein Orthogonalsystem. Es gilt

$$
\langle x_i(t), x_j(t) \rangle = 0 \quad \forall\, i, j \in \{1, \ldots, M\},\quad i \neq j.
$$

Gilt zusätzlich $\langle x_i(t), x_i(t) \rangle = 1 \quad \forall\, i \in \{1, \ldots, M\}$, spricht man von einem Orthonormalsystem.

Orthonormalsysteme eignen sich besonders zur Approximation komplizierter Signale durch elementare Signale.

\subsection*{Kommunikation mit orthogonalen Sequenzen}

Orthogonale Signale werden u.\,a.\ in Mobilfunksysteme verwendet um Teilnehmer $m_1, m_2, \ldots$ im Empfangssignal zu trennen.

\chapter{Signalsynthese und Signalanalyse}

\textbf{Aufgabe:} Synthese eines Signals mit Sinus- und Cosinusfunktionen.

\begin{center}
\textit{Jean Baptiste Joseph Fourier, 1768–1830}
\end{center}

\section{Fourierreihe}

Ein periodisches Signal $x(t)$ kann aus Sinusfunktionen unterschiedlicher Amplitude, Frequenz und Phase zusammengesetzt werden:

$$
x(t) = \bar{x} + \sum_{n=1}^{\infty} X[n] \cos(2\pi n f_0 t + \varphi_n)
$$

$$
= \bar{x} + \sum_{n=1}^{\infty} X_a[n] \cos(n\omega_0 t) + \sum_{n=1}^{\infty} X_b[n] \sin(n\omega_0 t)
$$

Dabei sind
\begin{itemize}
  \item $\bar{x}$ der Mittelwert (Gleichanteil),
  \item $X[n]$ die Amplitude der n-ten Cosinusfunktion,
  \item $\varphi_n$ die Phase der n-ten Cosinusfunktion,
  \item $f_0$ die Grundfrequenz, $\omega_0 = 2\pi f_0$ die Kreisgrundfrequenz,
  \item $T_0$ die Periodendauer des Signals.
\end{itemize}

\section{Fourierreihe (komplexwertige Schreibweise)}

Ein periodisches oder ein im Intervall $[t_1,\, t_1 + T_0]$ definiertes Signal der Dauer $T_0$ kann in eine Fourierreihe

$$
x(t) = \sum_{n=-\infty}^{\infty} X[n]\, e^{jn\omega_0 t}, \qquad t_1 \leq t \leq t_1 + T_0,
$$

mit $\omega_0 = 2\pi/T_0$ und

$$
X[n] = \frac{1}{T_0} \int_{t_1}^{t_1+T_0} x(t)\, e^{-jn\omega_0 t}\,\mathrm{d}t
$$

entwickelt werden. Das zeitbegrenzte Signal $x(t)$ wird dann durch die diskrete Folge $X[n]$ repräsentiert.

\subsection*{Analyse und Synthese im Vergleich}

\begin{center}
\begin{tabular}{lll}
\toprule
 & \textbf{reellwertige Schreibweise} & \textbf{komplexwertige Schreibweise} \\
\midrule
\multirow{3}{*}{Analyse}
  & $\displaystyle\bar{x} = \frac{1}{T_0}\int_{-T_0/2}^{T_0/2} x(t)\,\mathrm{d}t$
  & \multirow{3}{*}{$\displaystyle X[n] = \frac{1}{T_0}\int_{-T_0/2}^{T_0/2} x(t)\,e^{-jn\omega_0 t}\,\mathrm{d}t$} \\[10pt]
  & $\displaystyle X_a[n] = \frac{2}{T_0}\int_{-T_0/2}^{T_0/2} x(t)\cos(n\omega_0 t)\,\mathrm{d}t$ & \\[10pt]
  & $\displaystyle X_b[n] = \frac{2}{T_0}\int_{-T_0/2}^{T_0/2} x(t)\sin(n\omega_0 t)\,\mathrm{d}t$ & \\[10pt]
\midrule
\multirow{2}{*}{Synthese}
  & $\displaystyle x(t) = \bar{x} + \sum_{n=1}^{\infty} X_a[n]\cos(n\omega_0 t)$
  & \multirow{2}{*}{$\displaystyle x(t) = \sum_{n=-\infty}^{\infty} X[n]\,e^{jn\omega_0 t}$} \\[6pt]
  & $\displaystyle\phantom{x(t) = \bar{x}} + \sum_{n=1}^{\infty} X_b[n]\sin(n\omega_0 t)$ & \\[6pt]
\bottomrule
\end{tabular}
\end{center}

\subsection*{Approximation durch abgebrochene Fourierreihe}

Ferner halten wir fest, dass jede abgebrochene Fourierreihe

$$
x_K(t) = \sum_{n=-K}^{K} X[n]\,e^{jn\omega_0 t}, \qquad t_1 \leq t \leq t_1 + T_0,
$$

mit $\omega_0 = 2\pi/T_0$ die beste Approximation an $x(t)$, $t_1 \leq t \leq t_1 + T_0$, im Sinne des mittleren, quadratischen Fehlers

$$
J = \int_{-\infty}^{\infty} |x(t) - x_K(t)|^2\,\mathrm{d}t
$$

darstellt. Der Approximationsfehler ist durch

$$
J_K = \sum_{n=-\infty}^{-(K+1)} |X[n]|^2 + \sum_{n=K+1}^{\infty} |X[n]|^2
$$

gegeben.

\subsection*{Beispiel: Signalsynthese}

Die folgende Abbildung veranschaulicht die schrittweise Synthese eines Signals. Die linke Spalte zeigt die einzelnen Elementarsignale (Gleichanteil, erste und zweite Cosinusfunktion usw.), die rechte Spalte das jeweils aufaddierte rekonstruierte Signal. Mit wachsender Anzahl von Termen $K$ nähert sich $x_K(t)$ dem Zielsignal an.

\section{Bandbreite eines Signals}

Die Differenz von maximaler und minimaler Frequenz der im Signal enthaltenen Sinusfunktionen wird als \textit{Signalbandbreite $B$} des Signals bezeichnet, d.\,h.

$$
B = f_{\max} - f_{\min}
$$

Wenn das Signal einen Gleichanteil enthält, gilt $f_{\min} = 0$ und $B = f_{\max}$. Die Bandbreite ist gleich der maximalen Frequenz.

\subsection*{Begrenzung der Bandbreite}

Wird die Fourierreihe auf die ersten 60 Sinusfunktionen begrenzt, gehen höherfrequente Anteile verloren. Das Spektrum (Sinusamplituden über Frequenz/Grundfrequenz) wird bei der Grenzfrequenz abgeschnitten.

\subsection*{Approximation eines Bildsignals}

Ein Bildsignal (z.\,B.\ eine Zeile eines Röntgenbildes) lässt sich ebenfalls durch Sinusfunktionen darstellen. Bei Verwendung von 60 bzw.\ 30 Sinusfunktionen nimmt die Bildschärfe sichtbar ab:

$$
\Rightarrow \text{Je kleiner die Bandbreite, desto unschärfer ist das Bild.}
$$

\section{Diskrete Fourier-Transformation (DFT)}

\textbf{Signalsynthese mit DFT:}

$$
x[k] = \sum_{m=0}^{M-1} X[m]\, e^{j\frac{2\pi}{M}mk}, \qquad k = 0, \ldots, M-1
$$

\textbf{Signalanalyse mit DFT:}

$$
X[m] = \frac{1}{M} \sum_{k=0}^{M-1} x[k]\, e^{-j\frac{2\pi}{M}mk}, \qquad m = 0, \ldots, M-1
$$

Diese Transformation berechnet aus $M$ (möglicherweise komplexwertigen) Signalwerten $M$ in der Regel komplexwertige Spektralkoeffizienten (Fourierkoeffizienten). Sie ist eindeutig umkehrbar.

\subsection*{DFT eines Sprachlauts ($M = 164$)}

Wenn das diskrete Zeitsignal durch Abtastung mit der Abtastperiode $T_A$ aus einem analogen Signal hervorgegangen ist, gilt für die physikalische Frequenz des $m$-ten Koeffizienten mit $m < M/2$ die Beziehung

$$
f_m = \frac{m}{M} \cdot \frac{1}{T_A}.
$$

\section{Spektrogramm}

Für die Spektralanalyse eines zeitlich ausgedehnten Signals wird das Signal in kurze Segmente zerlegt und jedes Segment mit der DFT transformiert. Dieses Verfahren bezeichnet man als Kurzzeit-Fourier-Transformation (engl.\ \textit{short-time Fourier transform}, STFT).

Ein \textit{Spektrogramm} erhält man, wenn man die logarithmierten Beträge der Fourierkoeffizienten eines jeden Segments entlang der Zeitachse aneinanderreiht. Es ergibt sich dann eine \textit{Zeit-Frequenz-Darstellung} des Signals.

\chapter{Zufällige Signale}

\section{Rauschsignale}

Zufällige Signale (Rauschsignale) zeichnen sich dadurch aus, dass ihr zukünftiger Verlauf nicht exakt vorhersagbar ist. Zwei typische Rauschsignale $x_1(t)$ und $x_2(t)$ unterscheiden sich z.\,B.\ in Amplitude und spektraler Zusammensetzung.

\subsection*{Realisationen zufälliger Signale}

$\Rightarrow$ Jede Realisation eines Zufallssignals (Musterfunktion des Zufallsprozesses) hat einen anderen Amplitudenverlauf.

\section{Verteilungsdichtefunktion}

Die Verteilungsdichtefunktion (VDF) beschreibt die Wahrscheinlichkeit, mit der ein Zufallssignal bestimmte Amplitudenwerte annimmt. Sie ergibt sich als Grenzwert des normierten Histogramms bei wachsender Stichprobenzahl.

\section{Kenngrößen diskreter zufälliger Signale}

Kenngrößen zufälliger Signale lassen sich in \textbf{Lagemaße} (beschreiben den mittleren Amplitudenwert) und \textbf{Streumaße} (beschreiben die Streuung um den Mittelwert) unterteilen.

\subsection{Lagemaße}

Der \textbf{empirische Mittelwert} eines diskreten Signals mit $N$ Abtastwerten ist durch

$$
\bar{x} = \frac{1}{N} \sum_{k=0}^{N-1} x[k] = \frac{1}{N}\bigl(x[0] + \ldots + x[N-1]\bigr)
$$

gegeben.

Der \textbf{Median} $\tilde{x}$ eines diskreten Signals teilt die Werte in zwei Hälften, so dass die eine Hälfte größer oder gleich $\tilde{x}$ ist und die andere Hälfte kleiner oder gleich $\tilde{x}$ ist.

$\Rightarrow$ Der Median ist robust gegenüber Ausreißern. Oft ist er nicht eindeutig. Er kann dann innerhalb eines Intervalls frei gewählt werden.

\subsection{Streumaße}

Die \textbf{empirische Varianz} eines diskreten Signals $x[k]$ mit $N$ Abtastwerten ist durch

$$
\hat{\sigma}^2 = \frac{1}{N-1} \sum_{k=0}^{N-1} (x[k] - \hat{\bar{x}})^2
$$

definiert.

Die \textbf{empirische Standardabweichung} ist durch

$$
\hat{\sigma} = \sqrt{\hat{\sigma}^2} = \sqrt{\frac{1}{N-1} \sum_{k=0}^{N-1} (x[k] - \hat{\bar{x}})^2}
$$

gegeben.

\section{Die empirische Korrelationsfunktion}

Die Kreuzkorrelationsfunktion zweier Signale $x[k]$ und $y[k]$, wobei $x[k]$ ausserhalb des Bereich $0 \ldots N$ Null ist, ist dann durch

$$
\hat{\varphi}_{xy}[\varkappa] = \sum_{k=0}^{N} x[k]\, y[k + \varkappa]
$$

gegeben. Für die Autokorrelation erhält man entsprechend

$$
\hat{\varphi}_{xx}[\varkappa] = \sum_{k=0}^{N} x[k]\, x[k + \varkappa].
$$

Man beachte, dass in der Literatur diese Funktionen mit (unterschiedlichen) Normierungen verwendet werden.

\subsection*{Autokorrelationsfunktion eines Zufallssignals}

Das Signal $x[k]$ weist 51 von Null verschiedene Werte auf. Die AKF hat ein klares Maximum für $\varkappa = 0$ und zeigt ansonsten kaum Korrelation an.

\section{Signalentdeckung im Rauschen}

Ein Terminal sendet ein Signal $x(t)$. An der Basisstation wird das empfangene Signal durch additives Rauschen $n(t)$ gestört:

$$
y(t) = x(t) + n(t)
$$

\subsection*{Beispiel: Signalentdeckung im Rauschen}

Wir betrachten einen kurzen und einen langen Rechteckimpuls mit
$t_1 = 380T$, $T_1 = 40T$, $t_2 = 220T$, $T_2 = 360T$,

$$
x_1(t) = A\,\operatorname{rect}\!\left(\frac{t - t_1}{T_1}\right)
\qquad \text{und} \qquad
x_2(t) = \frac{A}{3}\,\operatorname{rect}\!\left(\frac{t - t_2}{T_2}\right),
$$

wobei das Bezugsintervall $T$ z.\,B.\ für $T = 1\,\text{ms}$ gewählt wird und beide Signale die Energie $E_1 = E_2 = A^2 \cdot 40T$ aufweisen.

Nun wird zu diesen Signalen Rauschen addiert:

$$
y_1(t) = x_1(t) + n_1(t) \qquad \text{und} \qquad y_2(t) = x_2(t) + n_2(t).
$$

$n_1(t)$ und $n_2(t)$ sind mittelwertfreie Rauchsignale.

\subsection*{Rechteckimpulse im Rauschen}

Ist eine Detektion durch eine einfache Schwellwertentscheidung möglich?

\subsection*{Signalentdeckung durch Korrelation}

\textbf{Berechnung der Kreuzkorrelation:}

\begin{align*}
\varphi_{y_1 x_1}(\tau)
  &= \int_{-\infty}^{\infty} y_1(t)\, x_1(t+\tau)\,\mathrm{d}t \\
  &= \int_{-\infty}^{\infty} (x_1(t) + n_1(t))\, x_1(t+\tau)\,\mathrm{d}t \\
  &= \int_{-\infty}^{\infty} x_1(t)\, x_1(t+\tau)\,\mathrm{d}t
     + \int_{-\infty}^{\infty} n_1(t)\, x_1(t+\tau)\,\mathrm{d}t \\
  &= \varphi_{x_1 x_1}(\tau) + \varphi_{n_1 x_1}(\tau)
\end{align*}

Falls $x_1(t)$ und $n_1(t)$ keine Korrelationen aufweisen, gilt

$$
\varphi_{y_1 x_1}(\tau) \approx \varphi_{x_1 x_1}(\tau)
$$

\subsection*{Rechteckimpulse im Rauschen – weiteres Beispiel}

Nun werden die folgenden Impulse betrachtet:

$$
x_1(t) = 0{,}4A\,\operatorname{rect}\!\left(\frac{t - t_1}{T_1}\right), \qquad
x_2(t) = \frac{2}{15}\,A\,\operatorname{rect}\!\left(\frac{t - t_2}{T_2}\right)
$$

\chapter{Abtastung und Quantisierung}

\section{Digitale Signalverarbeitung}

Die Verarbeitung informationstragender Signale erfolgt heute überwiegend mit Hilfe digitaler Prozessoren.

\begin{center}
\begin{tikzpicture}[>=stealth, font=\small]
  \node at (0,0) (in) {$u_1(t)$};
  \node[draw, rectangle, minimum width=2cm, minimum height=1.2cm, align=center] at (2.5,0) (ad) {A/D-\\Umsetzung};
  \node[draw, rectangle, minimum width=2cm, minimum height=1.2cm, align=center] at (5.5,0) (proc) {Prozessor};
  \node[draw, rectangle, minimum width=2cm, minimum height=1.2cm, align=center] at (8.5,0) (da) {D/A-\\Umsetzung};
  \node at (11,0) (out) {$u_2(t)$};
  \draw[->] (in) -- (ad);
  \draw[->] (ad) -- (proc);
  \draw[->] (proc) -- (da);
  \draw[->] (da) -- (out);
\end{tikzpicture}
\end{center}

\section{Analog-Digital-Umsetzung}

Die A/D-Umsetzung beinhaltet drei wesentliche Schritte:
\begin{itemize}
  \item Abtastung des zeitkontinuierlichen Signals
  \item Quantisierung der abgetasteten Signalwerte
  \item Umsetzung der quantisierten Werte in einen digitalen Code
\end{itemize}

\begin{itemize}
  \item Abtastung unter bestimmten Bedingungen ohne Informationsverlust
  \item Quantisierung in der Regel mit Informationsverlust
\end{itemize}

\section{Abtastung}

Entnahme von „Signalproben" in regelmäßigen Intervallen. Definition einer Abtastfrequenz oder Abtastrate $f_A$:

$$
f_A = \frac{\text{Zahl der Abtastwerte}}{\text{Zeiteinheit}}
$$

Typische Abtastraten:
\begin{itemize}
  \item Telefonsprache (ISDN, GSM, UMTS) \hfill 8\,kHz, 16\,kHz
  \item Audio-CD \hfill 44{,}1\,kHz
  \item DAT-Rekorder \hfill 48\,kHz
\end{itemize}

\subsection*{Abtastung des Sinussignals}

Das analoge Signal $x(t)$ wird in regelmäßigen Abständen $T_A$ abgetastet, woraus die diskrete Folge $x[k] = x(kT_A)$ entsteht.

\subsection*{Rekonstruktion des analogen Signals}

Aus den Abtastwerten lässt sich das analoge Signal näherungsweise rekonstruieren. Zwei gängige Methoden sind:
\begin{itemize}
  \item \textbf{Interpolation 0-ter Ordnung:} Jeder Abtastwert wird bis zum nächsten Abtastzeitpunkt konstant gehalten (Treppenform).
  \item \textbf{Lineare Interpolation:} Aufeinanderfolgende Abtastwerte werden durch Geraden verbunden.
\end{itemize}

\subsection*{(Ideale) Rekonstruktion mit si-Funktionen}

Die ideale Rekonstruktion des analogen Signals aus seinen Abtastwerten erfolgt durch Überlagerung von si-Funktionen, eine pro Abtastwert:

$$
\hat{x}(t) = \sum_{k=-\infty}^{\infty} x[k]\, \operatorname{si}\!\left(\frac{t - kT_A}{T_A}\right)
$$

Beispiel: $f = \dfrac{1}{T}$, $f_A = \dfrac{5}{T} = 5f$.

\subsection*{Spektrogramme}

Das Spektrogramm des Originalsignals und des nach Interpolation 0.\ Ordnung rekonstruierten Signals zeigen deutliche Unterschiede im Hochfrequenzbereich. Lineare Interpolation und Interpolation mit si-Funktionen liefern eine deutlich bessere Annäherung an das Originalsignal.

\subsection*{Ist die Rekonstruktion eindeutig?}

Die Rekonstruktion ist nur dann eindeutig, wenn die Abtastwerte das Signal vollständig beschreiben. Dabei gilt:

$$
\Rightarrow \text{Die Nullstellen der si-Funktion müssen im Raster } 1/f_A \text{ liegen!}
$$

\subsection*{Wie viele Abtastwerte sind erforderlich?}

Bei einer Abtastrate von $f_A = \dfrac{5}{2T}$ ist eine eindeutige Rekonstruktion des Signals möglich.

Bei einer Abtastrate von $f_A \approx \dfrac{4}{2T}$ (mit $\tilde{T} = 4T$) gelingt die Rekonstruktion ebenfalls noch fehlerfrei.

Offenbar geht im letzten Fall \textit{bei der Abtastung} die Information über die wahre Signalfrequenz verloren.

Die Rekonstruktion mit si-Funktionen liefert in diesem Fall ein Sinussignal bei einer tieferen Frequenz

$$
\tilde{f} = f_A - f,
$$

wobei $\tilde{f}$ die sogenannte \textit{Spiegelfrequenz} darstellt. Der oben beschriebene Effekt wird in der englischen Literatur häufig als \textit{aliasing} bezeichnet.

\section{Das Abtasttheorem}

Die in einem Signal enthaltene maximale Frequenz bestimmt die erforderliche Anzahl von Abtastwerten pro Zeiteinheit. Wenn die maximale Frequenz als Grenzfrequenz $f_g$ gegeben ist, muss also

$$
f_A > 2f_g
$$

gelten. Dieser Zusammenhang wird als \textit{Abtasttheorem} nach Nyquist oder Shannon bezeichnet. Die Abtastwerte werden dann in zeitlichen Intervallen $T_A = 1/f_A$ dem Signal entnommen und mit $x(kT_A)$, $k \in \mathbb{Z}$, bezeichnet.

\section{Quantisierung}

Die Amplitude $x(kT_A)$ eines Signalabtastwertes zum Zeitpunkt $t = kT_A$, $k \in \mathbb{Z}$, wird auf eine Menge endlicher, meist fest vorgegebener Werte, den Quantisierungsniveaus, abgebildet. Damit geht in der Regel ein Informationsverlust einher. Wird der quantisierte Wert mit $x_q(kT_A)$ bezeichnet, so kann der \textit{Quantisierungsfehler} des $k$-ten Abtastwertes mit

$$
\epsilon(kT_A) = x(kT_A) - x_q(kT_A)
$$

angegeben werden.

\section{Die Quantisierungskennlinie}

Der Zusammenhang zwischen zu quantisierenden und quantisiertem Wert kann mittels einer \textit{Quantisierungskennlinie} beschrieben werden.

Der (gleichmäßige) Quantisierer ist durch den Aussteuerbereich $[x_{\min}, x_{\max}]$ und die Stufenhöhe $\Delta x$ gekennzeichnet. Im obigen Beispiel gilt $x_{\min} = -1$, $x_{\max} = 1$ und $\Delta x = 0{,}25$.

\section{Der Quantisierungsfehler}

Der Quantisierungsfehler ist durch die Quantisierungsstufenhöhe beschränkt. Zur Berechnung der Quantisierungsfehlerleistung betrachten wir das folgende spezielle Eingangssignal (Rampe): Der Fehler $\epsilon(t) = x(t) - x_q(t)$ hat Sägezahnform mit Amplitude $\pm\tfrac{\Delta x}{2}$.

Bei einer großen Zahl von Quantisierungsstufen hat der Quantisierungsfehler rauschartigen Charakter. Beispiel einer Quantisierung mit $2^4 = 16$ Stufen:

Beispiel einer Quantisierung mit $2^{12} = 4096$ Stufen: Der Fehler $\epsilon(t)$ ist deutlich kleiner und kaum sichtbar.

\section{Dezibel}

\begin{itemize}
  \item Die Einheit Bel zeigt an, dass es sich um eine logarithmierte Verhältnisgröße handelt. Die Definition bezieht sich auf das Verhältnis zweier Leistungen $P_2$ und $P_1$:
  $$
  P = \log_{10}\!\left(\frac{P_2}{P_1}\right) \text{ [Bel]}
    = 10\log_{10}\!\left(\frac{P_2}{P_1}\right) \text{ [Dezibel]}
    = 10\,\text{lg}\!\left(\frac{P_2}{P_1}\right) \text{ [dB]}
  $$
  \item Beispiel: Die Leistung des Nutzsignals wird zur Rauschleistung ins Verhältnis gesetzt. Da das Ergebnis, der Signal-zu-Störabstand oder SNR (\textit{signal-to-noise ratio}), je nach Anzahl der Quantisierungsstufen über mehrere Größenordnungen variieren kann, ist es üblich, die Verhältnisgröße zu logarithmieren.
\end{itemize}

\subsection*{Einige dB-Werte}

\begin{center}
\begin{tabular}{cc|cc}
\toprule
\multicolumn{2}{c|}{\textbf{Leistungen}} & \multicolumn{2}{c}{\textbf{Amplituden}} \\
$\dfrac{P_2}{P_1}$ & $10\lg\!\left(\dfrac{P_2}{P_1}\right)$/dB
& $\dfrac{A_2}{A_1}$ & $20\lg\!\left(\dfrac{A_2}{A_1}\right)$/dB \\
\midrule
0{,}01      & $-20$ & 0{,}01      & $-40$ \\
0{,}1       & $-10$ & 0{,}1       & $-20$ \\
0{,}25      & $-6$  & 0{,}25      & $-12$ \\
0{,}5       & $-3$  & 0{,}5       & $-6$  \\
$1/\sqrt{2}$& $-1{,}5$ & $1/\sqrt{2}$ & $-3$ \\
1           & $0$   & 1           & $0$   \\
$\sqrt{2}$  & $1{,}5$  & $\sqrt{2}$   & $3$  \\
2           & $3$   & 2           & $6$   \\
4           & $6$   & 4           & $12$  \\
10          & $10$  & 10          & $20$  \\
100         & $20$  & 100         & $40$  \\
\bottomrule
\end{tabular}
\end{center}

\subsection*{Quantisierung eines Bildsignals}

$\Rightarrow$ Die Qualität des rekonstruierten Signals hängt wesentlich von der Zahl der Quantisierungsniveaus ab.

\section{Codierung: Natürliche binäre Zahlen}

\begin{center}
\begin{tabular}{lcc}
\toprule
\textbf{Menge} & \textbf{Dezimale Zahl} & \textbf{Natürliche binäre Zahl} \\
\midrule
       & 0 & 0    \\
$\bullet$     & 1 & 1    \\
$\bullet\bullet$    & 2 & 10   \\
$\bullet\bullet\bullet$   & 3 & 11   \\
$\bullet\bullet\bullet\bullet$  & 4 & 100  \\
$\bullet\bullet\bullet\bullet\bullet$ & 5 & 101  \\
$\bullet\bullet\bullet\bullet\bullet\bullet$ & 6 & 110  \\
$\bullet\bullet\bullet\bullet\bullet\bullet\bullet$ & 7 & 111  \\
$\bullet\bullet\bullet\bullet\bullet\bullet\bullet\bullet$ & 8 & 1000 \\
\bottomrule
\end{tabular}
\end{center}

\section{Prinzip der Codierung}

Die Quantisierungsniveaus eines Quantisierers lassen sich durch einen binären Code mit endlich vielen Stellen repräsentieren.

Bei binärer Codierung ergibt sich aus der Wortbreite $w$ die Zahl darstellbarer Quantisierungsstufen zu $2^w$.

\subsection*{Typische Wortbreiten}

\begin{center}
\begin{tabular}{lcc}
\toprule
\textbf{Signal} & \textbf{Wortbreite $w$} & \textbf{Anzahl der Quantisierungsniveaus} \\
\midrule
Sprache & 8--16 Bit      & 256 -- 65\,536         \\
Audio   & 16--24 Bit     & 65\,536 -- 16\,777\,216 \\
Video   & $3 \times 8$ Bit & $3 \times 256$         \\
\bottomrule
\end{tabular}
\end{center}

\section*{Zusammenfassung}

\begin{itemize}
  \item Informationstragende Signale entstehen auf vielfältige Weise. Es werden analoge, zeitdiskrete, wertdiskrete und digitale Signale unterschieden.
  \item Sie lassen sich häufig als Summe elementarer Signale darstellen.
  \item Die digitale Verarbeitung analoger Signale erfordert eine Abtastung, Quantisierung und Codierung.
  \item Die Abtastung ist verlustfrei, wenn das Abtasttheorem eingehalten wird.
  \item Die Quantisierung führt in der Regel zu einem Informationsverlust.
\end{itemize}
\chapter{Zufällige Signale}
\label{chap:zufaellige-signale}

\section{Rauschsignale}
\label{sec:rauschsignale}

Signale $x(t)$ oder $x[k]$, deren Verlauf für alle $t \in \mathbb{R}$ bzw.\ $k \in \mathbb{Z}$ exakt bekannt ist, werden deterministische Signale genannt. Diese Signale lassen sich in der Regel durch einen analytischen Ausdruck beschreiben. Signale, deren Amplituden zufällig variieren, werden Zufallssignale oder stochastische Signale genannt. Viele natürliche Signale (Sprache, Musik, Video) sind Zufallssignale. Zufallssignale werden nicht durch ihren Amplitudenverlauf sondern durch ein Histogramm der Amplituden und durch statistische Kenngrößen beschrieben.

Die Übertragung informationstragender Signale wird häufig durch störende Signale von zufälligem Charakter (,,Rauschen'', engl.\ ,,noise'') beeinträchtigt. Auch das Rauschen ist ein zufälliges Signal.

\begin{beispiel}[Rauschen]
Die folgende Abbildung zeigt zwei zufällige Signale.

% Grafik: Abb. 2.56 – Oben: Verrauschtes Audiosignal x_1(t); unten: Weißes, gaußsches Rauschen x_2(t) (WGN), Zeitachse 0–0,5 s
\end{beispiel}

Zufallssignale werden durch Zufallsprozesse erzeugt. Jede Realisation eines Zufallssignals (Musterfunktion des Zufallsprozesses) hat einen anderen Amplitudenverlauf.

% Grafik: Abb. 2.57 – Vier Realisationen (Musterfunktionen) x_1(t), x_2(t), x_3(t), x_4(t) eines weißen Gaußschen Rauschens

\section{Kenngrößen zufälliger Signale}
\label{sec:kenngroessen-zufaellige-signale}

Ein wichtiges Hilfsmittel zur empirischen Analyse zufälliger Signale ist das \textit{Histogramm} der Signalamplituden. Das Histogramm wird konstruiert, indem die Amplitudenwerte in $L$ Klassen (engl.\ bins) eingeteilt werden und die relative oder absolute Häufigkeit der zu einer Klasse gehörigen Daten aufgetragen wird. Die zu wählende Breite der Klassen ist von der Aufgabenstellung und der Datenmenge abhängig.

\begin{beispiel}[Histogramm und Verteilungsdichtefunktion eines Rauschsignals]
Für eines der in Abb.\ 2.57 gezeigten Rauschsignale ist das Histogramm dargestellt.

% Grafik: Abb. 2.58 – Histogramm (links) und Verteilungsdichtefunktion VDF (rechts), Amplitudenachse –5 bis 5
\end{beispiel}

Im Allgemeinen müssen die Methoden der Wahrscheinlichkeitsrechnung angewandt werden (siehe Kapitel~4). Die im Histogramm dargestellte relative Häufigkeit der Amplitudenwerte kann dann mit einer \textbf{Verteilungsdichtefunktion} (VDF) modelliert werden. Die Verteilungsdichtefunktion ist dann das der beobachteten Häufigkeitsverteilung zugrunde liegende Modell.

Wichtige Kenngrößen zufälliger Signale sind wieder der Mittelwert, die Varianz und die Leistung. Wenn diese unmittelbar aus nur einer Realisation des Zufallssignals berechnet werden, spricht man von \textbf{empirischen} Kenngrößen. Diese Kenngrößen werden häufig in ,,Lagemaße'' und ,,Streumaße'' eingeteilt.

\subsection{Lagemaße}
\label{sec:lagemasse}

Der \textbf{empirische Mittelwert} eines diskreten Signals mit $N$ Abtastwerten ist durch

$$
\hat{\bar{x}} = \frac{1}{N} \sum_{k=0}^{N-1} x[k] = \frac{1}{N}\bigl(x[0] + \ldots + x[N-1]\bigr)
$$

gegeben.

Der \textbf{Median} $\tilde{x}$ eines diskreten Signals $x[k]$ teilt die Werte in zwei Hälften, so dass die eine Hälfte größer oder gleich $\tilde{x}$ ist und die andere Hälfte kleiner oder gleich $\tilde{x}$ ist.
Der Median liegt also in der Mitte der geordneten Daten und ist robust gegenüber Ausreißern. Oft ist er durch die obige Definition nicht eindeutig bestimmt. Er kann dann innerhalb eines Intervalls frei gewählt werden.

\subsection{Streumaße}
\label{sec:streumasse}

Die \textbf{empirische Varianz} $\hat{\sigma}^2$ eines diskreten Signals $x[k]$ mit $N$ Abtastwerten ist durch

$$
\hat{\sigma}^2 = \frac{1}{N-1} \sum_{k=0}^{N-1} \bigl(x[k] - \hat{\bar{x}}\bigr)^2
$$

definiert. Die \textbf{empirische Standardabweichung} ist durch

$$
\hat{\sigma} = \sqrt{\hat{\sigma}^2} = \sqrt{\frac{1}{N-1} \sum_{k=0}^{N-1} \bigl(x[k] - \hat{\bar{x}}\bigr)^2}
$$

gegeben.

\section{Die empirische Korrelationsfunktion}
\label{sec:empirische-korrelation}

Im Fall zufälliger Signale beschränken wir uns vorerst auf diskrete und zeitlich begrenzte Signale. Dazu behandeln wir das Zufallssignal (genauer: die Musterfunktion des Zufallsprozesses) wie ein determiniertes Signal. Die Kreuzkorrelationsfunktion zweier Signale $x[k]$ und $y[k]$, wobei $x[k]$ außerhalb des Bereichs $0 \ldots N$ zu Null gesetzt wird, ist dann durch

$$
\hat{\varphi}_{xy}[\varkappa] = \sum_{k=0}^{N} x[k]\, y[k + \varkappa]
$$

gegeben. Für die Autokorrelationsfunktion erhält man entsprechend

$$
\hat{\varphi}_{xx}[\varkappa] = \sum_{k=0}^{N} x[k]\, x[k + \varkappa].
$$

Man beachte, dass in der Literatur diese Funktionen mit (unterschiedlichen) Normierungen verwendet werden. Die empirische Autokorrelationsfunktion ist für ein Rauschsignal dargestellt.

\begin{beispiel}[Autokorrelationsfunktion eines Zufallssignals]
Das Signal $x[k]$ weist 51 von Null verschiedene Werte auf. Die AKF ist symmetrisch. Sie hat ein Maximum für $\varkappa = 0$ und zeigt ansonsten kaum Korrelation an. Man erkennt, dass durch die Zufälligkeit des Signals für $\varkappa \neq 0$ nur kleine Werte entstehen. Die Selbstähnlichkeit des Signals ist gering.

% Grafik: Abb. 2.59 – Signal x[k] (oben, k von 0 bis 50) und AKF φ̂_xx[ϰ] (unten, ϰ von –50 bis 50)
\end{beispiel}

\section{Signalentdeckung im Rauschen}
\label{sec:signalentdeckung-rauschen}

\begin{beispiel}[Signalentdeckung im Rauschen]
Mit den bisher eingeführten Begriffen lassen sich bereits interessante systemtheoretische Aufgaben lösen. Mit dem nachfolgenden Beispiel zeigen wir, dass für die Entdeckung eines deterministischen Signals im Rauschen nicht etwa dessen Amplitude sondern die Energie des Signals maßgeblich ist.

Wir betrachten einen kurzen Rechteckimpuls ($t_1 = 380T$, $T_1 = 40T$)

$$
x_1(t) = A\,\operatorname{rect}\!\left(\frac{t - t_1}{T_1}\right)
$$

und einen langen Rechteckimpuls ($t_2 = 220T$, $T_2 = 360T$)

$$
x_2(t) = \frac{A}{3}\,\operatorname{rect}\!\left(\frac{t - t_2}{T_2}\right),
$$

wobei das Bezugsintervall $T$ z.\,B.\ zu $T = 1\,\text{ms}$ gewählt wird und beide Signale dieselbe Energie $E_1 = E_2 = A^2 \cdot 40T$ aufweisen. Zunächst lässt sich leicht überprüfen, dass beide Rechtecksignale die gleiche Energie aufweisen. Hierfür nutzen wir die Definition der Energie und finden

$$
E_1 = \int_{t_0}^{t_n} |x_1(t)|^2\,\mathrm{d}t = \int_{t_1}^{t_1+T_1} A^2\,\mathrm{d}t = A^2 T_1 = A^2 \cdot 40T
$$

$$
E_2 = \int_{t_0}^{t_n} |x_2(t)|^2\,\mathrm{d}t = \int_{t_2}^{t_2+T_2} \frac{A^2}{9}\,\mathrm{d}t = \frac{A^2 T_2}{9} = A^2 \cdot 40T
$$

Nun wird zu diesen Signalen mittelwertfreies Rauschen $n_1(t)$ und $n_2(t)$ addiert:
$$
y_1(t) = x_1(t) + n_1(t) \qquad \text{und} \qquad y_2(t) = x_2(t) + n_2(t).
$$

% Grafik: Abb. 2.60 – x_1(t) und y_1(t) oben; x_2(t) und y_2(t) unten, Zeitachse 0–1000 ms

\textbf{Signalentdeckung durch Schwellwertentscheidung}

Wenn die Detektion der Rechteckimpulse im verrauschten Signal mit Hilfe einer einfachen Schwellwertentscheidung durchgeführt wird, dann ist die Lage des Impulses $x_1(t)$ im Signal nur dann leicht zu erkennen, wenn die Signalamplitude deutlich größer als die Rausch\-amplituden ist. Wegen der geringeren Amplitude ist die Entdeckung des Signals $x_2(t)$ daher deutlich schwieriger.

\textbf{Signalentdeckung durch Korrelation}

Berechnen wir nun die Autokorrelationsfunktion der Signale $x_1(t)$ und $x_2(t)$ sowie die Kreuzkorrelation zwischen $y_1(t)$ und $x_1(t)$ bzw.\ $y_2(t)$ und $x_2(t)$, erkennen wir, dass die Signaldetektion auch im verrauschten Signal zuverlässig ein zuverlässiges Ergebnis liefert. Der Maximalwert der Autokorrelationsfunktion ist in beiden Fällen gleich, da er der Energie der Signale entspricht. Dies ist das Prinzip eines Korrelationsdetektors, das in der Kommunikationstechnik und der Mustererkennung vielfältig eingesetzt wird.

Die Kreuzkorrelationsfunktion soll jetzt für den Fall $\varphi_{y_1 x_1}(\tau)$ explizit ausgerechnet werden:

\begin{align*}
\varphi_{y_1 x_1}(\tau)
  &= \int_{-\infty}^{\infty} y_1(t)\, x_1(t+\tau)\,\mathrm{d}t \\
  &= \int_{-\infty}^{\infty} \bigl(x_1(t) + n_1(t)\bigr)\, x_1(t+\tau)\,\mathrm{d}t \\
  &= \int_{-\infty}^{\infty} x_1(t)\, x_1(t+\tau)\,\mathrm{d}t
     + \int_{-\infty}^{\infty} n_1(t)\, x_1(t+\tau)\,\mathrm{d}t \\
  &= \varphi_{x_1 x_1}(\tau) + \varphi_{n_1 x_1}(\tau)
\end{align*}

Falls $x_1(t)$ und $n_1(t)$ keine Korrelationen aufweisen, gilt

$$
\varphi_{y_1 x_1}(\tau) \approx \varphi_{x_1 x_1}(\tau)
$$

Die Kreuzkorrelation liefert näherungsweise die Autokorrelationsfunktion.

% Grafik: Abb. 2.61 – Autokorrelationsfunktionen φ_{x1x1} und φ_{x2x2} (ohne Rauschen) sowie φ_{y1x1} und φ_{y2x2} (mit Rauschen)

Wir wiederholen das obige Experiment nun mit Rechteckimpulsen $x_1(t)$ und $x_2(t)$, deren Amplituden auf $\tfrac{1}{3}$ des ursprünglichen Wertes reduziert sind:

$$
x_1(t) = 0{,}4A\,\operatorname{rect}\!\left(\frac{t - t_1}{T_1}\right), \qquad
x_2(t) = \frac{2}{15}\,A\,\operatorname{rect}\!\left(\frac{t - t_2}{T_2}\right)
$$

\begin{itemize}
  \item \textbf{Signalentdeckung durch Schwellwertentscheidung im Amplitudenbereich:} Wie in Abb.\ 2.62 zu sehen, scheint eine Detektion der Rechteckimpulse im Amplitudenbereich kaum noch möglich.
  \item \textbf{Signalentdeckung durch Schwellwertentscheidung im Korrelationsbereich:} In Abb.\ 2.63 ist zu erkennen, dass eine Detektion des Signals durch Kreuzkorrelation auch in diesem Fall noch möglich ist.
\end{itemize}

% Grafik: Abb. 2.62 – Zwei Rechtecksignale (reduzierte Amplitude) und die zugehörigen verrauschten Signale y_1(t) und y_2(t)
% Grafik: Abb. 2.63 – Autokorrelationsfunktionen der Signale aus Abb. 2.62
\end{beispiel}

\begin{beispiel}[Signalübertragung mit Code-division Multiple Access (CDMA)]
In Abbildung~2.64 sehen wir eine Anordnung, bestehend aus zwei Sendern $m_j$, $j \in \{1,2\}$ und einem Empfänger. Die Signale $m_j[k]$ sind die Nachrichtensignale, $y_j[k]$ die Sendesignale. $n[k]$ bezeichnet das Rauschen. Dies könnte ein stark vereinfachtes Mobilfunksystem darstellen. Zur weiteren Vereinfachung betrachten wir hier nur diskrete Signale.

% Grafik: Abb. 2.64 – Zwei mobile Endgeräte senden Signale m_1[k] und m_2[k] an einen gemeinsamen Empfänger; Rauschen n[k] addiert sich zum Empfangssignal L[k]

Alle Teilnehmer senden zur gleichen Zeit im gleichen Frequenzband. Ihre Nachrichtensignale sind jedoch mit individuellen Codes moduliert, so dass die Sendesignale unterschiedlicher Teilnehmer im Empfänger getrennt werden können (,,code division multiple access'', CDMA). Die einzelnen Codes $c_1[k]$ und $c_2[k]$ sind orthogonale, diskrete Sequenzen, sogenannte ,,Spreizcodes'' der Länge $N$. Sie sind so gestaltet, dass die Codes unterschiedlicher Sender eine verschwindend kleine Kreuzkorrelation $\varphi_{c_1 c_2}[\varkappa]$ für alle $\varkappa$ aufweisen. Dabei soll die Autokorrelation eines jeden Codes für $\varkappa = 0$ jedoch möglichst groß sein. Die Nachrichtensignale $m_1[k]$ und $m_2[k]$ werden über Abschnitte der Länge $N$ als konstant angenommen.

Für die Sendesignale erhält man somit
$$
y_1[k] = m_1[k]\,c_1[k] \quad \text{und} \quad y_2[k] = m_2[k]\,c_2[k].
$$

% Grafik: Abb. 2.65 – Nachrichtensignale m_1[k], m_2[k], Codesequenzen c_1[k], c_2[k] und Sendesignale y_1[k], y_2[k]

Die Codesequenzen $c_1[k]$ und $c_2[k]$ können z.\,B.\ mit dem Walsh-Hadamard-Verfahren erzeugt werden. Bei dem Empfänger kommen zur gleichen Zeit die Signale verschiedener Sender an; diese ergeben sich für unser einfaches Beispiel zu:

$$
L[k] = y_1[k] + y_2[k] + n[k] = m_1[k]\,c_1[k] + m_2[k]\,c_2[k] + n[k]
$$

Wir berechnen nun die Kreuzkorrelation des empfangenen Signals $L[k]$ mit dem Code $c_1[k]$ auf einem zeitlich begrenzten Segment von $N$ Werten und mit $k_0 = n \cdot N$, $n \in \mathbb{Z}$:

\begin{align*}
\varphi_{Lc_1}[\varkappa]
  &= \sum_{k=k_0}^{k_0+N-1} \bigl(y_1[k] + y_2[k] + n[k]\bigr)\,c_1[k+\varkappa] \\
  &= \sum_{k=k_0}^{k_0+N-1} m_1[k]\,c_1[k]\,c_1[k+\varkappa]
   + \sum_{k=k_0}^{k_0+N-1} m_2[k]\,c_2[k]\,c_1[k+\varkappa]
   + \sum_{k=k_0}^{k_0+N-1} n[k]\,c_1[k+\varkappa] \\
  &= m_1[k_0]\,E_{c_1,N} + m_2[k_0] \cdot 0 + \varphi_{nc_1}
\end{align*}

wobei $E_{c_1,N} > 0$ die Energie der Sequenz $c_1[k]$ für $N$ aufeinanderfolgende Werte darstellt. Die zweite Summe verschwindet, da $c_1[k]$ und $c_2[k]$ orthogonal sind. Ferner ist das Nachrichtensignal $m_1[k]$ im betrachteten Zeitraum konstant und kann deshalb mit $m_1[k_0]$ angegeben werden. Da nun das Rauschsignal $n[k]$ zufällig variiert, ist auch die dritte Summe relativ klein. Daher erhalten wir näherungsweise

$$
\varphi_{Lc_1}[\varkappa] \approx m_1[k_0] \cdot E_{c_1,N}
$$

Zum Beispiel erhalten wir für

$$
m_1[k_0] = \pm 1 \quad\Rightarrow\quad \varphi_{Lc_1}[0] = \pm E_{c_1,N}
$$

Wir können also durch eine einfache Schwellwertentscheidung die Informationen des Senders~1 aus dem Signalgemisch extrahieren. Das CDMA-Verfahren benötigt orthogonale Sequenzen, deren Konstruktion nachfolgend beschrieben wird.

\subsection*{Walsh-Hadamard-Sequenzen}

Ein einfaches Verfahren, um orthogonale Sequenzen für Spreizcodes zu erzeugen, kann wie folgt angegeben werden. Zunächst beginnt man mit Sequenzen der Länge~2,

$$
c_{11}[k] = \delta[k] + \delta[k-1]
$$
$$
c_{12}[k] = \delta[k] - \delta[k-1]
$$

deren Koeffizienten alternativ auch in einer Matrix $c_1$ zusammengefasst werden können,

$$
c_1 = \begin{pmatrix} 1 & 1 \\ 1 & -1 \end{pmatrix},
$$

wobei jede Spalte eine Sequenz $c_{11}$ bzw.\ $c_{12}$ darstellt.

Die Sequenzen $c_{11}[k]$ und $c_{12}[k]$ sind zueinander orthogonal und weisen die Energie $N$ auf:

$$
\sum_{k=0}^{1} c_{11}[k]\,c_{12}[k] = 1 \cdot 1 + 1 \cdot (-1) = 0
$$

$$
\sum_{k=0}^{1} c_{11}[k]\,c_{11}[k] = \sum_{k=0}^{1} c_{22}[k]\,c_{22}[k] = 1 \cdot 1 + 1 \cdot 1 = 2 = N.
$$

Mit der Matrix $c_1$ lassen sich nun längere Sequenzen konstruieren:

$$
c_2 = \begin{pmatrix} c_1 & c_1 \\ c_1 & -c_1 \end{pmatrix}
= \begin{pmatrix} 1 & 1 & 1 & 1 \\ 1 & -1 & 1 & -1 \\ 1 & 1 & -1 & -1 \\ 1 & -1 & -1 & 1 \end{pmatrix},
$$

wobei die vier neuen Sequenzen durch die Spalten $c_{21}$, $c_{22}$, $c_{23}$ und $c_{24}$ gegeben sind. Auch hier lässt sich die Orthogonalität der Sequenzen leicht nachweisen. Allgemein lässt sich die Matrix in einer rekursiven Form berechnen:

$$
c_{n+1} = \begin{pmatrix} c_n & c_n \\ c_n & -c_n \end{pmatrix}
$$

Wir erhalten also je nach Länge der Sequenz:
\begin{itemize}
  \item 2 orthogonale Folgen der Länge $N = 2$ oder
  \item 4 orthogonale Folgen der Länge $N = 4$ oder
  \item 8 orthogonale Folgen der Länge $N = 8$, etc.
\end{itemize}
\end{beispiel}


\chapter{Abtastung und Quantisierung}

\section{Analog-digital Umsetzung}

Die Verarbeitung wert- und zeitkontinuierlicher (,,analoger'') Signale mit einem digitalen Prozessor erfordert die Umsetzung des Signals in eine Zahlenfolge. Nach der numerischen Verarbeitung im Prozessor wird aus der resultierenden Zahlenfolge wieder ein analoges Signal rekonstruiert.

\begin{center}
\begin{tikzpicture}[>=stealth, font=\small]
  \node at (0,0) (in) {$u_1(t)$};
  \node[draw, rectangle, minimum width=2cm, minimum height=1.2cm, align=center] at (2.5,0) (ad) {A/D-\\Umsetzung};
  \node[draw, rectangle, minimum width=2cm, minimum height=1.2cm, align=center] at (5.5,0) (proc) {Prozessor};
  \node[draw, rectangle, minimum width=2cm, minimum height=1.2cm, align=center] at (8.5,0) (da) {D/A-\\Umsetzung};
  \node at (11,0) (out) {$u_2(t)$};
  \draw[->] (in) -- (ad);
  \draw[->] (ad) -- (proc);
  \draw[->] (proc) -- (da);
  \draw[->] (da) -- (out);
\end{tikzpicture}
\end{center}

Die A/D-Umsetzung beinhaltet drei wesentliche Schritte:
\begin{enumerate}[label=\alph*.]
  \item Abtastung des zeitkontinuierlichen Signals ($\to$ Abtastwerte)
  \item Quantisierung der abgetasteten Signalwerte ($\to$ quantisierte Abtastwerte)
  \item Umsetzung der quantisierten Werte in einen digitalen Code ($\to$ Binärcode)
\end{enumerate}

Während die Abtastung unter bestimmten Bedingungen ohne Informationsverlust realisiert werden kann, ist die Quantisierung in der Regel mit einem Informationsverlust verbunden.

\section{Abtastung}

Die Abtastung eines analogen Signals führt zu einer Entnahme von ,,Signalproben'' in regelmäßigen Intervallen. In digitalen Telefonsystemen (ISDN, GSM, etc.) werden dem analogen Sprachsignal z.\,B.\ 8000 Abtastwerte pro Sekunde entnommen. Dementsprechend kann man eine \textit{Abtastfrequenz} oder \textit{Abtastrate} $f_A$ definieren:

$$
f_A = \frac{\text{Zahl der Abtastwerte}}{\text{Zeiteinheit}}
$$

Die Abtastrate wird häufig in Hz angegeben. In der Sprach- und Audiosignalverarbeitung sind typische Abtastraten:
\begin{itemize}
  \item Telefonsprache (ISDN, GSM, UMTS) \hfill 8\,kHz, 16\,kHz
  \item Audio-CD \hfill 44{,}1\,kHz
  \item DAT-Rekorder \hfill 48\,kHz
\end{itemize}

Bei Bildern sind die Abtastpunkte häufig in Form eines regelmäßigen Gitters angeordnet.

\begin{beispiel}[Abtastung eines Sinussignals]
Das analoge Signal $x(t)$ wird in regelmäßigen Abständen $T_A$ abgetastet, woraus die diskrete Folge $x[k] = x(kT_A)$ entsteht.

% Grafik: Abb. 2.68 – Oben: kontinuierliches Sinussignal x(t); unten: abgetastete Werte x(kT_A) als Impulsfolge
\end{beispiel}

\subsection*{Rekonstruktion des analogen Signals}

Bevor wir untersuchen, unter welchen Bedingungen eine Abtastung ohne Informationsverlust erreicht wird, ist die Frage zu klären, wie das analoge Signal aus den Abtastwerten rekonstruiert werden kann. Hierzu betrachten wir die drei Varianten:

\begin{itemize}
  \item \textbf{Interpolation 0.\ Ordnung, Halteglied}\\
  Die einfachste Form der Interpolation besteht darin, den Amplitudenwert zum Abtastzeitpunkt $k$ für die Dauer eines Abtastintervalls $T_A$ zu halten:
  $$
  \tilde{x}(t) = x(kT_A), \qquad kT_A \leq t < (k+1)T_A
  $$
  Die Interpolation 0.\ Ordnung ergibt ein treppenförmiges Signal.

  \item \textbf{Interpolation 1.\ Ordnung, lineare Interpolation}\\
  Bei der Interpolation 1.\ Ordnung werden die Abtastwerte abschnittsweise mit Geraden verbunden:
  $$
  \tilde{x}(t) = \bigl(x((k+1)T_A) - x(kT_A)\bigr)\,\frac{t - kT_A}{T_A} + x(kT_A),
  $$
  wobei $kT_A \leq t < (k+1)T_A$ gelten muss. Durch Einsetzen kann man sich leicht überzeugen, dass diese Gerade die Punkte $(kT_A,\, x(kT_A))$ und $((k+1)T_A,\, x((k+1)T_A))$ beinhaltet.

  \item \textbf{(Ideale) Rekonstruktion mit si-Funktionen}\\
  Gesucht sind Funktionen $g_k(t)$, so dass die Bedingung
  $$
  x(t) = \sum_{k=-\infty}^{\infty} x[k]\, g_k(t)
  $$
  mit $x[k] = x(kT_A)$ erfüllt ist. Die Funktionen $g_k(t)$ können z.\,B.\ mit Hilfe der Fouriertransformation (Vorlesung Systemtheorie~II) oder der Lagrange-Polynome gefunden werden. Es handelt sich dabei um die si-Funktion,
  $$
  g_k(t) = \operatorname{si}\!\left(\frac{t - kT_A}{T_A}\right),
  $$
  wobei $T_A = \tfrac{1}{f_A}$ das Abtastintervall bezeichnet. Somit ergibt sich die ideale Rekonstruktion zu
  $$
  \tilde{x}(t) = \sum_{k=-\infty}^{\infty} x[k]\, \operatorname{si}\!\left(\frac{t - kT_A}{T_A}\right).
  $$
\end{itemize}

% Grafik: Abb. 2.69 – Signalrekonstruktion eines Sinussignals: Interpolation 0. Ordnung (oben) und Interpolation 1. Ordnung (unten)

\subsection*{Herleitung der idealen Rekonstruktion über Lagrange-Polynome}

Die si-Funktion entspricht einem Lagrange-Polynom zur Interpolation unendlich vieler, gleichmäßig verteilter Stützstellen/Abtastwerte. Dieses wird im Folgenden gezeigt. Gegeben sei die diskrete Folge $x[k] = x(kT_A)$, für $k = 0, 1, \ldots, N-1$ und die Basisfunktionen

$$
L_k(t) = \prod_{\substack{j=0,\, j\neq k}}^{N-1} \frac{t_j - t}{t_j - t_k}
$$

mit den Werten $t_k = k \cdot T_A$ und $t_j = j \cdot T_A$. Die Interpolation

$$
\tilde{x}(t) = \sum_{k=0}^{N-1} x(t_k) \cdot L_k(t)
$$

gilt auch für unendlich viele, gleichmäßig verteilte Abtastwerte. Wir definieren hierfür die Basisfunktionen:

$$
l_k(t) = \prod_{\substack{j=-\infty,\, j\neq k}}^{\infty} \frac{t_j - t}{t_j - t_k}
= \prod_{\substack{j=-\infty,\, j\neq k}}^{\infty} \frac{jT_A - t}{jT_A - kT_A}
$$

und substituieren mit $j = n + k$:

$$
l_k(t) = \prod_{\substack{n=-\infty,\, n\neq 0}}^{\infty} \frac{(n+k)T_A - t}{nT_A}.
$$

Wir zeigen nun, dass dies einer Reihenentwicklung der si-Funktion entspricht, denn es gilt

$$
\operatorname{si}(x) = \prod_{n=1}^{\infty} \left(1 - \frac{x^2}{n^2}\right)
$$

und also auch

$$
\operatorname{si}\!\left(\frac{t - kT_A}{T_A}\right) = \prod_{n=1}^{\infty} \left(1 - \left[\frac{t - kT_A}{nT_A}\right]^2\right).
$$

Mit Hilfe der dritten binomischen Formel erhalten wir:

$$
\operatorname{si}\!\left(\frac{t - kT_A}{T_A}\right) = \prod_{n=1}^{\infty} \left(1 - \frac{t - kT_A}{nT_A}\right)\left(1 + \frac{t - kT_A}{nT_A}\right).
$$

Durch geschickte Indexerweiterung kann der Term weiter vereinfacht werden:

$$
\operatorname{si}\!\left(\frac{t - kT_A}{T_A}\right)
= \prod_{\substack{n=-\infty,\, n\neq 0}}^{\infty} \left(1 - \frac{t - kT_A}{nT_A}\right)
= \prod_{\substack{n=-\infty,\, n\neq 0}}^{\infty} \frac{(n+k)T_A - t}{nT_A}.
$$

Somit erhalten wir

$$
\tilde{x}(t) = \sum_{k=-\infty}^{\infty} x[k] \cdot \operatorname{si}\!\left(\frac{t - kT_A}{T_A}\right) = x(t).
$$

% Grafik: Abb. 2.70 – Interpolation mit si-Funktionen: Abtastwerte, eine si-Funktion pro Abtastwert, Summe aller si-Funktionen

Die maximale Signalfrequenz $f_{\max}$ muss die Bedingung $2f_{\max} < f_A$ erfüllen. Dies wird im Folgenden am Beispiel der Sinusfunktion erklärt.

\section{Das Abtasttheorem}

Die vorangegangenen Experimente und weitergehende Untersuchungen zeigen: Bei sinusförmigen Vorgängen sind mehr als zwei Abtastwerte pro Signalperiode für die eindeutige Repräsentation des Signals durch die Abtastwerte erforderlich. Wird diese Bedingung verletzt, ergibt die ideale Rekonstruktion eine Sinusfunktion bei einer Spiegelfrequenz. Da man sich beliebige Signale aus Sinussignalen zusammensetzen kann, bestimmt die in einem Signal enthaltene maximale Frequenz die erforderliche Anzahl von Abtastwerten pro Zeiteinheit.

\begin{satz}[Abtasttheorem nach Nyquist und Shannon]
Wenn die maximale Frequenz als Grenzfrequenz $f_g$ gegeben ist, muss

$$
f_A > 2f_g
$$

gelten. Die Abtastwerte werden dann in gleichmäßigen, zeitlichen Intervallen $T_A = 1/f_A$ dem Signal entnommen und mit $x(kT_A)$, $k \in \mathbb{Z}$, bezeichnet.
\end{satz}

Für $f_A = 2f_g$ ist die Signalrekonstruktion nicht eindeutig, für $f_A < 2f_g$ tritt \textit{aliasing} auf. Die Rekonstruktion mit si-Funktionen liefert in diesem Fall ein Sinussignal bei einer tieferen Frequenz

$$
\tilde{f} = f_A - f,
$$

wobei $\tilde{f}$ die sogenannte \textit{Spiegelfrequenz} darstellt. Der oben beschriebene Effekt wird in der englischen Literatur häufig als \textit{aliasing} bezeichnet.

% Grafik: Abb. 2.71 – Interpolation mit si-Funktionen: fehlerhafte Rekonstruktion bei zu geringer Abtastrate
% Grafik: Abb. 2.72 – Interpolation mit si-Funktionen und f_A = 5/(2T): korrekte Rekonstruktion
% Grafik: Abb. 2.73 – Interpolation mit si-Funktionen und f_A ≈ 5/(4T): aliasing

\section{Quantisierung}

Im Zuge der Quantisierung wird die Amplitude $x(kT_A)$ eines Signalabtastwertes zum Zeitpunkt $t = kT_A$, $k \in \mathbb{Z}$, auf eine Menge endlicher, meist fest vorgegebener Werte, den Quantisierungsniveaus, abgebildet. Damit geht in der Regel ein Informationsverlust einher. Wird der quantisierte Wert mit $x_q(kT_A)$ bezeichnet, so kann der \textit{Quantisierungsfehler} des $k$-ten Abtastwertes mit

$$
e(kT_A) = x(kT_A) - x_q(kT_A)
$$

angegeben werden.

% Grafik: Abb. 2.74 – Quantisierung des Intervalls [0,1] auf zwei (links) und vier (rechts) Werte

\section{Die Quantisierungskennlinie}

Der Zusammenhang zwischen zu quantisierendem und quantisiertem Wert kann mittels einer \textit{Quantisierungskennlinie} beschrieben werden. Zwei gebräuchliche Typen sind die Midrise- und die Midtread-Kennlinie, wie in Abb.\ 2.75 dargestellt.

Der (gleichmäßige) Quantisierer ist durch den Aussteuerbereich $[x_{\min}, x_{\max}]$ und die Stufenhöhe $\Delta x$ gekennzeichnet. In den Beispielen in Abb.\ 2.75 gilt $x_{\min} = -1$, $x_{\max} = 1$ und $\Delta x = 0{,}25$. Die Quantisierungsniveaus werden häufig auf binäre Zahlen abgebildet. Mit $w$ binären Stellen lassen sich dann maximal $2^w$ Quantisierungsniveaus indizieren.

% Grafik: Abb. 2.75 – Quantisierungskennlinien: Midrise (links) und Midtread (rechts)

\section{Der Quantisierungsfehler und seine Leistung}

Der Betrag $|e(t)|$ des Quantisierungsfehlers ist durch die halbe Quantisierungsstufenhöhe $\Delta x / 2$ beschränkt. Bei einer großen Zahl von Quantisierungsstufen hat der Quantisierungsfehler rauschartigen Charakter.

Zur Berechnung der Quantisierungsfehlerleistung betrachten wir das speziell gewählte Eingangssignal (Rampe):

$$
x(t) = \frac{t}{T}, \qquad -T \leq t \leq T,
$$

das außerhalb des betrachteten Intervalls $[-T, T]$ periodisch fortgesetzt werden kann. Dieses Signal dient als Modell für ein gleichverteiltes Signal, in dem alle Amplituden $x(t)$ mit gleicher Häufigkeit auftreten. Mit der Substitution $u = t/T$ ergibt sich die Signalleistung $P_x$ zu:

$$
P_x = \frac{1}{2T} \int_{-T}^{T} \left(\frac{t}{T}\right)^2 \mathrm{d}t
= \frac{1}{2} \int_{-1}^{1} u^2\,\mathrm{d}u
= \int_0^1 u^2\,\mathrm{d}u
= \left[\frac{1}{3}u^3\right]_0^1 = \frac{1}{3}.
$$

Der Quantisierungsfehler $e(t) = x(t) - x_q(t)$ hat Sägezahnform mit einer Periodendauer von $T_e = \Delta x \cdot T$, wobei $x_q$ auf dem Intervall $[0, \Delta x\, T]$ mit $x_q = \frac{\Delta x}{2}$ beschrieben wird. Die Leistung des Quantisierungsfehlers ergibt sich somit zu:

\begin{align*}
P_e &= \frac{1}{\Delta x\, T} \int_0^{\Delta x\, T} \bigl(x(t) - x_q(t)\bigr)^2\,\mathrm{d}t
     = \frac{1}{\Delta x\, T} \int_0^{\Delta x\, T} \left(\frac{t}{T} - \frac{\Delta x}{2}\right)^2 \mathrm{d}t \\
    &= \frac{1}{\Delta x} \int_0^{\Delta x} \left(u^2 - \Delta x\, u + \frac{\Delta x^2}{4}\right) \mathrm{d}u
     = \frac{1}{\Delta x} \left[\frac{1}{3}u^3 - \frac{1}{2}u^2 \Delta x + \frac{\Delta x^2}{4}u\right]_0^{\Delta x} \\
    &= \frac{1}{\Delta x} \left[\frac{1}{3}\Delta x^3 - \frac{1}{2}\Delta x^3 + \frac{\Delta x^3}{4}\right]
     = \frac{\Delta x^2}{12},
\end{align*}

wobei wieder $u = t/T$ substituiert wurde.

Somit erhalten wir für dieses spezielle Signal $P_e = \dfrac{\Delta x^2}{12}$.

Durch weiterführende Rechnungen kann nun gezeigt werden, dass die Formel für die Leistung des Quantisierungsfehlers näherungsweise auch für andere als das gewählte, linear ansteigende Signal gilt. Das ist dann der Fall, wenn die Zahl der Quantisierungsstufen $2^w$ ,,hinreichend'' groß ist.

\section{Dezibel}

\begin{itemize}
  \item Die Einheit Bel zeigt an, dass es sich um eine logarithmierte Verhältnisgröße handelt. Die Definition bezieht sich auf das Verhältnis zweier Leistungen $P_2$ und $P_1$:
  $$
  P = \log_{10}\!\left(\frac{P_2}{P_1}\right) \text{ [Bel]}
    = 10\log_{10}\!\left(\frac{P_2}{P_1}\right) \text{ [Dezibel]}
    = 10\,\mathrm{lg}\!\left(\frac{P_2}{P_1}\right) \text{ [dB]}
  $$
  \item Beispiel: Die Leistung des Nutzsignals wird zur Rauschleistung ins Verhältnis gesetzt. Da das Ergebnis, der Signal-zu-Störabstand oder SNR (\textit{signal-to-noise ratio}), je nach Anzahl der Quantisierungsstufen über mehrere Größenordnungen variieren kann, ist es üblich, die Verhältnisgröße zu logarithmieren.
\end{itemize}

\subsection*{Einige dB-Werte}

\begin{center}
\begin{tabular}{cc|cc}
\toprule
\multicolumn{2}{c|}{\textbf{Leistungen}} & \multicolumn{2}{c}{\textbf{Amplituden}} \\
$\dfrac{P_2}{P_1}$ & $10\,\mathrm{lg}\!\left(\dfrac{P_2}{P_1}\right)$/dB
& $\dfrac{A_2}{A_1}$ & $20\,\mathrm{lg}\!\left(\dfrac{A_2}{A_1}\right)$/dB \\
\midrule
0{,}01      & $-20$ & 0{,}01      & $-40$ \\
0{,}1       & $-10$ & 0{,}1       & $-20$ \\
0{,}25      & $-6$  & 0{,}25      & $-12$ \\
0{,}5       & $-3$  & 0{,}5       & $-6$  \\
$1/\sqrt{2}$& $-1{,}5$ & $1/\sqrt{2}$ & $-3$ \\
1           & $0$   & 1           & $0$   \\
$\sqrt{2}$  & $1{,}5$  & $\sqrt{2}$   & $3$  \\
2           & $3$   & 2           & $6$   \\
4           & $6$   & 4           & $12$  \\
10          & $10$  & 10          & $20$  \\
100         & $20$  & 100         & $40$  \\
\bottomrule
\end{tabular}
\end{center}

\subsection*{Quantisierung eines Bildsignals}

$\Rightarrow$ Die Qualität des rekonstruierten Signals hängt wesentlich von der Zahl der Quantisierungsniveaus ab.

% Grafik: Abb. 2.76 – Quantisierungsfehler e(t) = x(t) - x_q(t): Rampensignal, Sägezahnfehler, Fehlerleistung

\section{Codierung: Natürliche binäre Zahlen}

\begin{center}
\begin{tabular}{lcc}
\toprule
\textbf{Menge} & \textbf{Dezimale Zahl} & \textbf{Natürliche binäre Zahl} \\
\midrule
       & 0 & 0    \\
$\bullet$     & 1 & 1    \\
$\bullet\bullet$    & 2 & 10   \\
$\bullet\bullet\bullet$   & 3 & 11   \\
$\bullet\bullet\bullet\bullet$  & 4 & 100  \\
$\bullet\bullet\bullet\bullet\bullet$ & 5 & 101  \\
$\bullet\bullet\bullet\bullet\bullet\bullet$ & 6 & 110  \\
$\bullet\bullet\bullet\bullet\bullet\bullet\bullet$ & 7 & 111  \\
$\bullet\bullet\bullet\bullet\bullet\bullet\bullet\bullet$ & 8 & 1000 \\
\bottomrule
\end{tabular}
\end{center}

\section{Prinzip der Codierung}

Die Quantisierungsniveaus eines Quantisierers lassen sich durch einen binären Code mit endlich vielen Stellen repräsentieren.

Bei binärer Codierung ergibt sich aus der Wortbreite $w$ die Zahl darstellbarer Quantisierungsstufen zu $2^w$.

\subsection*{Typische Wortbreiten}

\begin{center}
\begin{tabular}{lcc}
\toprule
\textbf{Signal} & \textbf{Wortbreite $w$} & \textbf{Anzahl der Quantisierungsniveaus} \\
\midrule
Sprache & 8--16 Bit      & 256 -- 65\,536         \\
Audio   & 16--24 Bit     & 65\,536 -- 16\,777\,216 \\
Video   & $3 \times 8$ Bit & $3 \times 256$         \\
\bottomrule
\end{tabular}
\end{center}

\section*{Zusammenfassung}

\begin{itemize}
  \item Informationstragende Signale entstehen auf vielfältige Weise. Es werden analoge, zeitdiskrete, wertdiskrete und digitale Signale unterschieden.
  \item Sie lassen sich häufig als Summe elementarer Signale darstellen.
  \item Die digitale Verarbeitung analoger Signale erfordert eine Abtastung, Quantisierung und Codierung.
  \item Die Abtastung ist verlustfrei, wenn das Abtasttheorem eingehalten wird: $f_A > 2f_g$.
  \item Die Quantisierung führt in der Regel zu einem Informationsverlust. Die Quantisierungsfehlerleistung beträgt näherungsweise $P_e = \Delta x^2 / 12$.
\end{itemize}

\chapter{Systeme und ihre Eigenschaften}
\label{chap:systeme}

Der Entwurf, die Analyse und die Berechnung informationstechnischer Methoden und Geräte erfordert eine strukturierende Betrachtungsweise. Dazu wurde der Begriff des ,,Systems'' eingeführt, der auch in vielen anderen Bereichen der technischen und physikalischen Welt eingesetzt wird.

In unserem Kontext werden Systeme im Allgemeinen wie folgt charakterisiert:

\begin{itemize}
  \item Ein System ist die Abstraktion eines Teils der physikalischen oder technischen Welt.
  \item Ein System ist von seiner Umwelt abgrenzbar. Die Wechselwirkungen der Komponenten des Systems untereinander überwiegen im Allgemeinen die Wechselwirkungen mit der Umgebung.
  \item Die Festlegung der Systemgrenzen hängt von der zu behandelnden Fragestellung ab.
  \item Das Systemverhalten wird im Allgemeinen durch Eingangs- und Ausgangssignale charakterisiert, wobei jedem Eingangssignal Ausgangssignale zugeordnet sind.
  \item Das Systemverhalten beruht auf seiner inneren Struktur, also der Komponenten und Wechselwirkungen. Wenn diese nicht bekannt oder nicht von Interesse sind, spricht man auch von einer ,,black box''.
\end{itemize}

Die Systemtheorie unterscheidet somit zwischen der inneren Struktur eines Systems und seinem Verhalten nach außen.

\section*{Überblick}

\begin{itemize}
  \item[3.1] Mathematische Systemdefinition
  \item[3.2] Systemeigenschaften
  \item[3.3] Lineare und zeitinvariante Systeme
  \item[3.4] Parametrische Systeme
  \item[3.5] Weitere Beispiele
\end{itemize}

\section{Mathematische Systemdefinition}
\label{sec:systemdefinition}

\subsection{Allgemeine Definitionen}
\label{sec:system-allgemein}

Aus den oben genannten Punkten läßt sich nun eine erste, sehr allgemeine Systemdefinition ableiten.

\begin{definition}[Mathematische Systemdefinition]
Ein System $\mathbb{S}$ mit einem Eingangssignal $x(t)$ und einem Ausgangssignal $y(t)$ wird durch eine bezüglich der Komposition abgeschlossene Menge von geordneten Signalpaaren

\begin{equation}
\mathbb{S} := \{(x(t), y(t)) \mid t_0 \leq t < t_1;\, t_0, t_1 \in \mathbb{R}\}
\label{eq:systemdefinition}
\end{equation}

definiert. Hierbei ist $x(t)$ das Eingangssignal und $y(t)$ das Ausgangssignal des Systems. Die Menge $\mathbb{S}$ heißt bezüglich der Komposition abgeschlossen, falls für alle Signalpaare $\{(x(t), y(t)) \mid t_0 \leq t < t_1\}$ und $\{(x(t), y(t)) \mid t_1 \leq t < t_2\}$ auch $\{(x(t), y(t)) \mid t_0 \leq t < t_2\}$ in $\mathbb{S}$ enthalten ist.
\end{definition}

Die Abgeschlossenheit der Menge $\mathbb{S}$ bezüglich der Komposition garantiert, dass zwei zeitlich aufeinander folgende Eingangs-Ausgangssignalpaare, die in zwei Experimenten nacheinander beobachtet wurden, prinzipiell in einem weiteren Experiment auch an einem Stück beobachtet werden können. Hiermit wird die Konsistenz der Beobachtungen an den Ein- und Ausgängen des Systems sichergestellt.

Die Definition des Systems als Menge geordneter Signal- oder Datenpaare ist sowohl auf \textbf{reale Systeme} als auch auf \textbf{Systemmodelle} anwendbar.

\subsection{Single/Multiple Input/Output Systeme}
\label{sec:siso-mimo}

Die obige Definition lässt sich sinngemäß auch auf Systeme mit mehreren Eingangs- und Ausgangssignalen erweitern.

\begin{center}
\begin{tabular}{lcc}
  & \textbf{Single Input} & \textbf{Multiple Input} \\[4pt]
  \textbf{Single Output}   & SISO & MISO \\[4pt]
  \textbf{Multiple Output} & SIMO & MIMO \\
\end{tabular}
\end{center}

MIMO Systeme haben u.\,a.\ in der drahtlosen Funkübertragung große Bedeutung erlangt. Anstatt nur eine Sende- und Empfangsantenne zu verwenden kann mit mehreren Antennen eine deutliche Erhöhung der Übertragungsraten erzielt werden. Im mehrkanaligen Fall werden die Sende- und Empfangsantennen mit dem dazwischenliegenden Funkkanal als MIMO-System modelliert.

\begin{beispiel}[Drahtlose Gigabit Netzwerke]
  % Grafik: Internet → Router → Desktop Computer (MIMO-System)
\end{beispiel}

\subsection{Systeme und innere Zustände}
\label{sec:innere-zustaende}

Die Definition der Menge $\mathbb{S}$ schließt nicht aus, dass es zu einem Eingangssignal $x(t)$ mehr als nur ein Ausgangssignal $y(t)$ gibt. Diese Mehrdeutigkeit ist Ausdruck der Abhängigkeit der Ausgangssignale von ,,inneren'' Variablen des Systems. Diese werden als Zustände $z_1(t), z_2(t), \ldots, z_K(t)$ des Systems bezeichnet. Das Systemverhalten kann dann durch einen eindeutigen funktionalen Zusammenhang

$$
y(t) = S[z_1(t), \ldots, z_K(t), x(t)]
$$

bestimmt werden. In Abhängigkeit von der Zeitvariablen (kontinuierlich oder diskret) kann dieser Zusammenhang z.\,B.\ durch eine Differentialgleichung oder durch eine Differenzengleichung beschrieben werden.

% Grafik: Blockdiagramm x(t) → [S mit z_1(t)...z_K(t)] → y(t)

\subsection{Zustandsraummodelle}
\label{sec:zustandsraum}

Eine allgemeine Form, um die Abhängigkeit von Eingangs- und Ausgangssignalen und den inneren Zustandsvariablen des Systems zu erfassen, bietet das Zustandsraummodell.

\begin{definition}[Zustandsraummodell]
Beim \textit{Zustandsraummodell} wird das Systemverhalten mithilfe einer linearen DGL erster Ordnung (oder eines Systems linearer DGLen erster Ordnung) in eine Berechnungsvorschrift überführt. Wenn hierbei im Kontext eines MIMO-Systems $K$ Zustandsvariablen, $N$-dimensionale Eingangssignale $\vec{x}(t)$ und $M$-dimensionale Ausgangssignale $\vec{y}(t)$ berücksichtigt werden sollen, ist die Formulierung mit Hilfe von Vektoren und Matrizen zweckmäßig, d.\,h.

\begin{equation}
\frac{d\vec{z}(t)}{dt} = \mathbf{A}(t)\,\vec{z}(t) + \mathbf{B}(t)\,\vec{x}(t)
\label{eq:zustandsgleichung}
\end{equation}

\begin{equation}
\vec{y}(t) = \mathbf{C}(t)\,\vec{z}(t) + \mathbf{D}(t)\,\vec{x}(t)\,.
\label{eq:ausgangsgleichung}
\end{equation}

Die Gleichung~\eqref{eq:zustandsgleichung} wird als Zustandsgleichung und die Gleichung~\eqref{eq:ausgangsgleichung} als Ausgangsgleichung bezeichnet. $\vec{z}(t)$ ist ein Spaltenvektor der Dimension $K$, der die Zustandsvariablen enthält. Weiterhin sind

\begin{itemize}
  \item $\vec{x}(t)$ \quad der Vektor der Eingangssignale der Dimension $L$
  \item $\vec{y}(t)$ \quad der Vektor der Ausgangssignale der Dimension $M$
  \item $\mathbf{A}(t)$ \quad die Systemmatrix der Dimension $K \times K$
  \item $\mathbf{B}(t)$ \quad die Eingangsmatrix der Dimension $K \times L$
  \item $\mathbf{C}(t)$ \quad die Beobachtungsmatrix der Dimension $M \times K$
  \item $\mathbf{D}(t)$ \quad die Durchgangsmatrix der Dimension $M \times L$
  \item $\vec{z}_0$ \quad der Anfangszustandsvektor der Dimension $K$ zum Zeitpunkt $t_0$.
\end{itemize}
\end{definition}

Das Systemverhalten wird durch das Zusammenwirken dieser Größen charakterisiert. Die Systembeobachtung startet bei $t_0$ mit dem Anfangszustand $\vec{z}_0$. Wenn der Anfangszustand $\vec{z}_0$ Null ist, handelt es sich um ein lineares Systemmodell. Wenn das System zeitinvariant ist, entfällt die Abhängigkeit dieser Matrizen von der Zeit. Der Ruhezustand ist durch $\vec{z}(t) = \vec{0}$ gegeben.

Für die lineare DGL erster Ordnung mit konstanten Koeffizienten lässt sich das allgemeine Zustandsraummodell wie folgt vereinfachen:

$$
\frac{dz(t)}{dt} = A\,z(t) + B\,x(t)
$$

$$
y(t) = C\,z(t) + D\,x(t)\,.
$$

Man erhält ein zeitinvariantes Modell. Eine Lösung für das Ausgangssignal findet man, indem zuerst die Lösung der homogenen Gleichung $\mathrm{d}z(t)/\mathrm{d}t = Az(t)$ und dann durch Integration die spezielle Lösung unter Berücksichtigung des Eingangssignals bestimmt wird.

Die Zustandsraumbeschreibung kann ebenso auch für zeitdiskrete Systeme entwickelt werden. Für ein System mit $K$-dimensionalen Zustandvektoren $\vec{z}[k]$, $L$-dimensionalen Eingangsvektoren $\vec{x}[k]$ und $M$-dimensionalen Ausgangsvektoren $\vec{y}[k]$ erhält man

$$
\vec{z}[k+1] = \mathbf{A}[k]\,\vec{z}[k] + \mathbf{B}[k]\,\vec{x}[k]
$$

$$
\vec{y}[k] = \mathbf{C}[k]\,\vec{z}[k] + \mathbf{D}[k]\,\vec{x}[k]\,.
$$

Im Fall eines zeitinvarianten Systems lassen sich die wiederholten Iterationen für $k > k_0$ in eine geschlossene Form überführen:

$$
\vec{z}[k] = \mathbf{A}^{k-k_0}\,\vec{z}[k_0] + \sum_{\kappa=1}^{k-k_0} \mathbf{A}^{k-k_0-\kappa}\cdot\mathbf{B}\cdot\vec{x}[k_0 + \kappa - 1]\,.
$$

Die Matrix $\Phi(k, k_0) = \mathbf{A}^{k-k_0}$ wird als Zustandsübergangsmatrix oder Transitionsmatrix bezeichnet.

\subsection{MIMO Zustandsraummodell}
\label{sec:mimo-zustandsraum}

Das MIMO Systemmodell für $t \geq t_0$:

\begin{center}
\begin{tikzpicture}[auto, node distance=1.5cm, >=latex]
  \node (x) {$\vec{x}(t)$};
  \node [draw, rectangle, right=1cm of x] (B) {$\mathbf{B}(t)$};
  \node [draw, circle, right=1cm of B] (sum1) {$+$};
  \node [draw, rectangle, right=1cm of sum1, label=above:{$\frac{d\vec{z}(t)}{dt}$}] (int) {$\int$};
  \node [draw, circle, right=1cm of int] (sum2) {$+$};
  \node [draw, rectangle, right=1cm of sum2] (C) {$\mathbf{C}(t)$};
  \node [draw, circle, right=1cm of C] (sum3) {$+$};
  \node [right=1cm of sum3] (y) {$\vec{y}(t)$};
  \node [draw, rectangle, above=1.5cm of sum1] (A) {$\mathbf{A}(t)$};
  \node [draw, rectangle, below=1.5cm of int] (z0) {$\vec{z}_0\delta(t-t_0)$};
  \node [draw, rectangle, below=2cm of sum1] (D) {$\mathbf{D}(t)$};

  \draw[->] (x) -- (B);
  \draw[->] (B) -- (sum1);
  \draw[->] (sum1) -- (int);
  \draw[->] (int) -- node[above] {$\vec{z}(t)$} (sum2);
  \draw[->] (sum2) -- (C);
  \draw[->] (C) -- (sum3);
  \draw[->] (sum3) -- (y);
  \draw[->] (z0) -- (sum2);
  \draw[->] (sum2.north) -- ++(0,0.7) -| (A.south);
  \draw[->] (A.west) -- ++(-0.5,0) |- (sum1.north);
  \draw[->] (x.south) -- ++(0,-2) -- (D.west);
  \draw[->] (D.east) -- ++(0.5,0) |- (sum3.south);
\end{tikzpicture}
\end{center}

\begin{beispiel}[Ein zeitkontinuierliches System: Das RC-Zweitor]

\begin{center}
\begin{tikzpicture}
  % Input port
  \draw (0,0) -- (0,-2);
  \draw (0,0) -- (1.5,0);
  \draw (0,-2) -- (3,-2);
  % Resistor
  \draw (1.5,0) rectangle (2.5,0.4);
  \draw (2.5,0.2) -- (3,0.2) -- (3,0);
  % Current arrow
  \draw[->] (2.8,0.5) -- node[above] {$i(t)$} (3.5,0.5);
  % Node
  \draw (3,0) -- (5,0);
  % Capacitor
  \draw (3.8,-0.5) -- (4.2,-0.5);
  \draw (3.8,-1) -- (4.2,-1);
  \draw (4,-0.5) -- (4,0);
  \draw (4,-1) -- (4,-2);
  \draw (3,-2) -- (5,-2);
  % Output port
  \draw (5,0) -- (5,-2);
  % Labels
  \node[left] at (0,-1) {$x(t)$};
  \node[right] at (5,-1) {$y(t)$};
  % Port dots
  \filldraw (0,0) circle (2pt);
  \filldraw (0,-2) circle (2pt);
  \filldraw (5,0) circle (2pt);
  \filldraw (5,-2) circle (2pt);
\end{tikzpicture}
\end{center}

Das oben dargestellte RC-Zweitor kann als System mit einer inneren Zustandsvariablen $z(t)$ modelliert werden. Das System ist für $t \geq t_0$ durch die folgende Beziehung zwischen dem Eingangssignal $x(t)$ und dem Ausgangssignal $y(t)$ charakterisiert:

$$
y(t) = z(t_0)\,e^{-\frac{t-t_0}{\tau}} + \frac{1}{\tau}\int_{t_0}^{t} e^{-\frac{t-\xi}{\tau}}\,x(\xi)\,\mathrm{d}\xi
$$

Hierbei ist $\tau$ eine Zeitkonstante und $z(t_0)$ der Zustand des Systems (Spannung am Kondensator) zum Zeitpunkt $t_0$. Offenbar klingt bei diesem System der Einfluss des Zustands zu Beginn der Beobachtung ($t = t_0$) exponentiell ab, während das Eingangssignal $x(t)$ mit fortschreitender Zeit durch Integration zunehmend das Ausgangssignal und somit den neuen Wert der Zustandsvariablen bestimmt.

Die Reaktion des Systems auf das Sprungsignal $x(t) = A\epsilon(t)$ ist mit $y(t) = A(1 - e^{-\frac{t}{\tau}})$ gegeben, wobei $\tau = RC$ eine Zeitkonstante ist.

\begin{center}
\begin{tikzpicture}
  \begin{groupplot}[
    group style={group size=2 by 1, horizontal sep=2cm},
    width=5cm, height=4cm,
    axis lines=left,
    xlabel=$t$,
    xmin=-0.5, xmax=4,
    ymin=0, ymax=1.4,
    ytick={1}, yticklabels={$A$},
    xtick=\empty,
  ]
  \nextgroupplot[title={$x(t)$}]
    \addplot[thick] coordinates {(-0.5,0)(0,0)(0,1)(4,1)};
  \nextgroupplot[title={$y(t)$}]
    \addplot[thick, domain=0:4, samples=100] {1 - exp(-x)};
    \addplot[dashed] coordinates {(0,1)(4,1)};
    \node at (axis cs:0,1) [above right] {$A$};
  \end{groupplot}
\end{tikzpicture}
\end{center}

Ferner ist die Zustandsvariable $z(t)$ am Ausgang des Systems unmittelbar beobachtbar, denn es gilt für $t = t_0$

$$
y(t_0) = z(t_0)\,e^{-\frac{t_0-t_0}{\tau}} + \frac{1}{\tau}\int_{t_0}^{t_0} e^{-\frac{t_0-\xi}{\tau}}\,x(\xi)\,\mathrm{d}\xi = z(t_0)\,.
$$

Somit können wir für beliebige $t \geq t_0$ schreiben

$$
z(t) = z(t_0)\,e^{-\frac{t-t_0}{\tau}} + \frac{1}{\tau}\int_{t_0}^{t} e^{-\frac{t-\xi}{\tau}}\,x(\xi)\,\mathrm{d}\xi
$$

$$
y(t) = z(t)\,.
$$

Dies ist eine einfache Form eines \textit{Zustandsraummodells} in Integralform.

\subsubsection*{Zustandsraummodell des RC-Zweitors}

Durch Differenzieren der Zustandsgleichung und mit Hilfe der Leibniz-Regel erhalten wir:

$$
\frac{dz(t)}{dt} = -\frac{1}{\tau}\,z(t) + \frac{1}{\tau}\,x(t)
$$

Mit $A = -1/\tau$, $B = 1/\tau$, $C = 1$ und $D = 0$ erhalten wir die Zustandsraumformulierung für das RC-System (ein System erster Ordnung):

$$
\frac{dz(t)}{dt} = A\,z(t) + B\,x(t)
$$

$$
y(t) = C\,z(t) + D\,x(t)\,.
$$

\subsubsection*{Die Impulsantwort}

Wir betrachten den Dirac-Impuls als spezielles Eingangssignal $x(t) = \delta(t)$. Für $x(t) = \delta(t)$ erhalten wir mit der Siebeigenschaft des $\delta$-Impulses (Abschnitt~\ref{sec:dirac-impuls}; implizit gilt hier $z(t_0) = 0$ für $t_0 = -\infty$):

\begin{align*}
y(t) &= \frac{1}{\tau}\int_{-\infty}^{\infty} \delta(\xi)\,e^{-\frac{t-\xi}{\tau}}\,\epsilon(t-\xi)\,\mathrm{d}\xi
      = \frac{1}{\tau}\,e^{-\frac{t}{\tau}}\,\epsilon(t) = h(t)\,.
\end{align*}

Somit ist die \textit{Impulsantwort} des Systems durch $h(t) = \epsilon(t)\,\dfrac{1}{\tau}\,e^{-\frac{t}{\tau}}$ gegeben.
\end{beispiel}

In der folgenden Fortsetzung dieses Beispiels wollen wir zeigen, wie man im Fall eines einfachen elektrischen Netzes die Beziehung zwischen Eingangs- und Ausgangssignal berechnet und damit das bereits eingeführte Ergebnis herleitet.

\begin{beispiel}[Vom elektrischen Netz zur Eingangs-Ausgangsbeziehung]

Das Ziel der folgenden Rechnung ist die Bestimmung des Ausgangssignals $y(t)$ für beliebige Eingangssignale $x(t)$ des RC-Netzes. Wir wollen damit das im vorangehenden Beispiel (RC-Zweitor) eingeführte Systemverhalten aus grundlegenden Prinzipien der elektrischen Netzberechnung ableiten.

Wir beginnen mit den Bauteilgleichungen: Mit $q(t) = Cu_C(t)$ erhalten wir für die Stromstärke (Strom = bewegte Ladung)

$$
i(t) = \frac{\mathrm{d}q(t)}{\mathrm{d}t} = C\frac{\mathrm{d}u_C(t)}{\mathrm{d}t}.
$$

Die Spannung über $R$ lässt sich leicht mit $u_R(t) = R\,i(t)$ bestimmen,

$$
u_R(t) = R\,i(t) = RC\frac{\mathrm{d}u_C(t)}{\mathrm{d}t}.
$$

Aus dem Maschenumlauf erhalten wir nun eine Differentialgleichung (DGL) 1.\ Ordnung:

$$
u_e(t) = u_R(t) + u_C(t) = RC\frac{\mathrm{d}u_C(t)}{\mathrm{d}t} + u_C(t).
$$

% Grafik: RC-System mit Maschenumlauf (Abb. 3.7)

Durch die folgenden Substitutionen bringen wir die DGL in eine vereinfachte Form:

\medskip
\begin{tabular}{lll}
  $\tau = RC$ & & Zeitkonstante \\
  $x(t) = u_e(t)$ & & Eingangssignal \\
  $y(t) = u_C(t)$ & & Ausgangssignal und gleichzeitig Zustand $z(t)$ des Systems.
\end{tabular}
\medskip

Dann folgt

$$
\frac{\mathrm{d}y(t)}{\mathrm{d}t} + \frac{1}{\tau}y(t) = \frac{1}{\tau}x(t) \qquad \Leftrightarrow \qquad \frac{\mathrm{d}y(t)}{\mathrm{d}t} = -\frac{1}{\tau}y(t) + \frac{1}{\tau}x(t).
$$

Es handelt sich um eine lineare DGL 1.\ Ordnung mit konstanten Koeffizienten. Mit Hilfe des charakteristischen Polynoms $\lambda + \tfrac{1}{\tau} = 0$ erhalten wir für $t \geq t_0$ einen Ansatz für die homogene Lösung dieser DGL,

$$
y_h(t) = a\,e^{\lambda t} = a\,e^{-\frac{t-t_0}{\tau}},
$$

wobei $a$ eine noch zu bestimmende Konstante bezeichnet. Für eine spezielle Lösung benutzen wir die Methode der Variation der Konstanten und setzen

$$
y_p(t) = G_A(t)\,e^{-\frac{t-t_0}{\tau}}; \qquad y'(t) = G'_A(t)\,e^{-\frac{t-t_0}{\tau}} - \frac{G_A(t)}{\tau}\,e^{-\frac{t-t_0}{\tau}}.
$$

Nun setzen wir die spezielle Lösung in die DGL ein,

$$
G'_A(t)\,e^{-\frac{t-t_0}{\tau}} - \frac{G_A(t)}{\tau}\,e^{-\frac{t-t_0}{\tau}} + \frac{G_A(t)}{\tau}\,e^{-\frac{t-t_0}{\tau}} = \frac{1}{\tau}\cdot x(t)
\qquad \Rightarrow\quad G'_A(t) = \frac{1}{\tau}\cdot x(t)\,e^{\frac{t-t_0}{\tau}}.
$$

Aus dieser Gleichung erhalten wir $G_A(t)$ durch Integration,

$$
G_A(t) = G_A(t_0) + \frac{1}{\tau}\int_{t_0}^{t} x(\xi)\,e^{\frac{\xi-t_0}{\tau}}\,\mathrm{d}\xi.
$$

Aus diesen Teilergebnissen bilden wir die vollständige Lösung $y(t) = y_h(t) + y_p(t)$,

$$
y(t) = a\,e^{-\frac{t-t_0}{\tau}} + G_A(t_0)\,e^{-\frac{t-t_0}{\tau}} + \frac{1}{\tau}\int_{t_0}^{t} x(\xi)\,e^{-\frac{t-\xi}{\tau}}\,\mathrm{d}\xi, \quad t \geq t_0.
$$

Durch Substitution von $z(t_0) = a + G_A(t_0)$ erhalten wir schließlich

$$
y(t) = z(t_0)\,e^{-\frac{t-t_0}{\tau}} + \underbrace{\frac{1}{\tau}\int_{t_0}^{t} x(\xi)\,e^{-\frac{t-\xi}{\tau}}\,\mathrm{d}\xi}_{\text{Faltungsintegral}}, \quad t \geq t_0.
$$

$z(t_0)$ beschreibt hierbei den inneren Zustand des Systems zum Zeitpunkt $t_0$. Man erkennt, dass sich die Lösung aus einem abklingenden Anfangszustand und einem Faltungsintegral zusammensetzt. Wenn die Kondensatorspannung (und damit der Zustand $z(t)$ und das Ausgangssignal $y(t)$) zu Beginn ($t = t_0$) Null ist, wird die Ausgangsspannung für $t > t_0$ allein durch dieses Integral beschrieben.

Um den Einfluss des Anfangszustands vollständig auszuschließen, setzen wir $z(t_0) = 0$ und $t_0 \to -\infty$ und erhalten das Ausgangssignal

$$
y(t) = \frac{1}{\tau}\int_{-\infty}^{t} x(\xi)\,e^{-\frac{t-\xi}{\tau}}\,\mathrm{d}\xi
= \int_{-\infty}^{\infty} x(\xi)\,\frac{1}{\tau}\,e^{-\frac{t-\xi}{\tau}}\,\epsilon(t-\xi)\,\mathrm{d}\xi
= \int_{-\infty}^{\infty} x(\xi)\,h(t-\xi)\,\mathrm{d}\xi.
$$

Wir erhalten somit das Ausgangssignal als Faltung (Abschnitt~\ref{sec:faltungsintegral}) des Eingangssignals mit der Funktion $h(t) = \tfrac{1}{\tau}e^{-t/\tau}\,\epsilon(t)$, die wir zuvor schon als Impulsantwort identifiziert haben. Für das Eingangssignal $x(t) = \delta(t)$ ergibt sich die kausale Impulsantwort des Systems zu

$$
y(t) = \frac{1}{\tau}\,e^{-\frac{t}{\tau}}\,\epsilon(t) = \begin{cases} 0 & t < 0 \\ \dfrac{1}{\tau}\,e^{-\frac{t}{\tau}} & t \geq 0 \end{cases}.
$$
\end{beispiel}

Die an diesem Beispiel gewonnenen Erkenntnisse lassen sich auf Systeme mit beliebig vielen inneren Zustandsvariablen und beliebig vielen Eingangs- und Ausgangssignalen verallgemeinern.

\section{Diskrete Zustandsraummodelle}
\label{sec:diskrete-zustandsraum}

Zeitvariant:

\begin{center}
\begin{tikzpicture}[auto, node distance=1.4cm, >=latex]
  \node (x) {$\vec{x}[k]$};
  \node [draw, rectangle, right=0.8cm of x] (B) {$\mathbf{B}[k]$};
  \node [draw, circle, right=0.8cm of B] (sum1) {$+$};
  \node [draw, rectangle, right=0.8cm of sum1, label=above:{$\vec{z}[k+1]$}] (T) {$T$};
  \node [draw, rectangle, right=0.8cm of T] (C) {$\mathbf{C}[k]$};
  \node [draw, circle, right=0.8cm of C] (sum2) {$+$};
  \node [right=0.8cm of sum2] (y) {$\vec{y}[k]$};
  \node [draw, rectangle, above=1.2cm of sum1] (A) {$\mathbf{A}[k]$};
  \node [draw, rectangle, below=1.2cm of B] (D) {$\mathbf{D}[k]$};

  \draw[->] (x) -- (B);
  \draw[->] (B) -- (sum1);
  \draw[->] (sum1) -- (T);
  \draw[->] (T) -- node[above] {$\vec{z}[k]$} (C);
  \draw[->] (C) -- (sum2);
  \draw[->] (sum2) -- (y);
  \draw[->] (T.north) -- ++(0,0.5) -| (A.south);
  \draw[->] (A.west) -- ++(-0.4,0) |- (sum1.north);
  \draw[->] (x.south) -- ++(0,-1.2) -- (D.west);
  \draw[->] (D.east) -- ++(3.5,0) |- (sum2.south);
\end{tikzpicture}
\end{center}

Zeitinvariant:

\begin{center}
\begin{tikzpicture}[auto, node distance=1.4cm, >=latex]
  \node (x) {$\vec{x}[k]$};
  \node [draw, rectangle, right=0.8cm of x] (B) {$\mathbf{B}$};
  \node [draw, circle, right=0.8cm of B] (sum1) {$+$};
  \node [draw, rectangle, right=0.8cm of sum1, label=above:{$\vec{z}[k+1]$}] (T) {$T$};
  \node [draw, rectangle, right=0.8cm of T] (C) {$\mathbf{C}$};
  \node [draw, circle, right=0.8cm of C] (sum2) {$+$};
  \node [right=0.8cm of sum2] (y) {$\vec{y}[k]$};
  \node [draw, rectangle, above=1.2cm of sum1] (A) {$\mathbf{A}$};
  \node [draw, rectangle, below=1.2cm of B] (D) {$\mathbf{D}$};

  \draw[->] (x) -- (B);
  \draw[->] (B) -- (sum1);
  \draw[->] (sum1) -- (T);
  \draw[->] (T) -- node[above] {$\vec{z}[k]$} (C);
  \draw[->] (C) -- (sum2);
  \draw[->] (sum2) -- (y);
  \draw[->] (T.north) -- ++(0,0.5) -| (A.south);
  \draw[->] (A.west) -- ++(-0.4,0) |- (sum1.north);
  \draw[->] (x.south) -- ++(0,-1.2) -- (D.west);
  \draw[->] (D.east) -- ++(3.5,0) |- (sum2.south);
\end{tikzpicture}
\end{center}

$T$ bezeichnet einen vektorwertigen Zustandsspeicher.

\section{Systemeigenschaften}
\label{sec:systemeigenschaften}

Ausgehend von der sehr allgemeinen Systemdefinition in Gleichung~\eqref{eq:systemdefinition} sollen Systeme nun durch spezielle Eigenschaften charakterisiert werden. Dabei werden wir sowohl den zeitkontinuierlichen als auch den zeitdiskreten Fall berücksichtigen. Letztere ist besonders für die Simulation von Systemen auf digitalen Computern wichtig. Wir beschränken uns hierbei auf den SISO Fall.

\subsection{Zeitinvarianz und Linearität kontinuierlicher Systeme}
\label{sec:zeitinvarianz-linearitaet}

\begin{definition}[Zeitinvariante und lineare Systeme (LTI)]
Ein System heißt \textbf{zeitinvariant}, wenn alle Signalpaare $(x(t), y(t)) \in \mathbb{S}$ gegenüber einer zeitlichen Verschiebung $\tau_0$ invariant sind, d.\,h.\ falls für die Relation $y(t) = \mathcal{S}[x(t)]$ eines SISO-Systems und beliebige Zeitverzögerungen $\tau_0 \in \mathbb{R}$ gilt:

$$
y(t - \tau_0) = \mathcal{S}[x(t - \tau_0)]\,.
$$

Ein System ist \textbf{linear}, falls mit $y_1(t) = \mathcal{S}[x_1(t)]$ und $y_2(t) = \mathcal{S}[x_2(t)]$ und für beliebige Faktoren $a \in \mathbb{R}$ und $b \in \mathbb{R}$ oder $a \in \mathbb{C}$ und $b \in \mathbb{C}$ gilt:

$$
a\,y_1(t) + b\,y_2(t) = \mathcal{S}[a\,x_1(t) + b\,x_2(t)]\,.
$$

Ein System, das \textbf{linear} und \textbf{zeitinvariant} ist, wird LTI-System (engl.: \textit{linear time-invariant system}) genannt.
\end{definition}

Wenn ein System linear ist, gilt auch das bekannte Superpositionsprinzip.

\subsection{Verschiebungsinvarianz und Linearität diskreter Systeme}
\label{sec:verschiebungsinvarianz}

\begin{definition}[Verschiebungsinvariante und lineare Systeme (LSI)]
Ein zeitdiskretes System $y[k] = \mathcal{S}[x[k]]$ heißt \textbf{verschiebungsinvariant}, falls für ein beliebiges $k_0 \in \mathbb{Z}$ gilt:

$$
y[k - k_0] = \mathcal{S}[x[k - k_0]]\,.
$$

Ein zeitdiskretes System $y[k] = \mathcal{S}[x[k]]$ ist \textbf{linear}, falls mit $y_1[k] = \mathcal{S}[x_1[k]]$ und $y_2[k] = \mathcal{S}[x_2[k]]$ und für beliebige Faktoren $a \in \mathbb{R}$ und $b \in \mathbb{R}$ oder $a \in \mathbb{C}$ und $b \in \mathbb{C}$ gilt:

$$
a\,y_1[k] + b\,y_2[k] = \mathcal{S}[a\,x_1[k] + b\,x_2[k]]\,.
$$

Ein System, das \textbf{linear} und \textbf{verschiebungsinvariant} ist, wird LSI-System (engl.: \textit{linear shift-invariant system}) genannt.
\end{definition}

\subsection{Weitere Systemeigenschaften}
\label{sec:weitere-systemeigenschaften}

\begin{definition}[Gedächtnislosigkeit, BIBO-Stabilität und Kausalität]
\begin{itemize}
  \item Ein zeitkontinuierliches System heißt \textbf{gedächtnislos}, falls für jedes Eingangssignal $x(t)$ und beliebiges $t_0 \in \mathbb{R}$ das Ausgangssignal $y(t = t_0)$ nur von $x(t = t_0)$ abhängt.
  \item Ein zeitdiskretes System heißt \textbf{gedächtnislos}, falls für jedes Eingangssignal $x[k]$ und beliebiges $k_0 \in \mathbb{Z}$ das Ausgangssignal $y[k = k_0]$ nur von $x[k = k_0]$ abhängt.
  \item Ein System ist \textbf{bounded input -- bounded output (BIBO) stabil}, falls jedes amplitudenbegrenzte Eingangssignal $x(t)$ bzw.\ $x[k]$ ein amplitudenbegrenztes Ausgangssignal $y(t)$ bzw.\ $y[k]$ zur Folge hat.
  \item Ein System heißt \textbf{kausal}, wenn das Ausgangssignal zur Zeit $t_0$ nur von den Eingangssignalen und Zustandsvariablen für Zeiten $t \leq t_0$ abhängt.
\end{itemize}
\end{definition}

\begin{beispiel}[Das Perzeptron -- ein nichtlineares System]
Künstliche neuronale Netze werden aus vielen gleichartigen Grundelementen, den künstlichen Neuronen, zusammengesetzt. Ein einfaches, zeitdiskretes Modell eines Neurons ist in Abb.\ 3.8 dargestellt. Dieses wird als ,,Perzeptron'' bezeichnet. Die Kombination vieler Neuronen in mehreren Schichten ergibt dann das ,,Multi-layer Perzeptron'' (MLP), eine einfache Variante tiefer neuronaler Netze (engl.\ \textit{deep neural network, DNN}).

% Grafik: Perzeptron mit Eingängen x₁[k],...,xₙ[k], Gewichten w₁,...,wₙ, Summierglied, Aktivierungsfunktion f(·) und Ausgang y[k]

Das Perzeptron weist mehrere Eingänge $x_\ell[k]$ und einen Ausgang $y[k]$ auf. Es definiert einen funktionalen Zusammenhang

$$
y[k] = f\!\left(\sum_{\ell=1}^{N} x_\ell[k]\,w_\ell + b\right),
$$

wobei $w_1, \ldots, w_N$ die Gewichte (,,weights''), $b$ eine additive Konstante (,,bias'') und $f(z)$ die Aktivierungsfunktion (,,activation'') des Perzeptrons bezeichnen. Es handelt sich um ein gedächtnisloses und kausales MISO-System.

Für die Aktivierungsfunktion $f(z)$ wird häufig eine der folgenden nichtlinearen Funktionen gewählt:

\begin{itemize}
  \item Halbwellengleichrichtung (,,rectified linear unit'', ReLU): $f(z) = \max(0, z)$
  \item Logistische Funktion: $f(z) = \dfrac{1}{1 + \exp(-z)}$
  \item Tangens Hyperbolikus: $f(z) = \tanh(z)$
\end{itemize}

Es ist leicht zu zeigen, dass das System ,,Perzeptron'' von den Eingängen bis zum Signal $z[k]$ linear ist, wenn $b = 0$ gilt. Hierfür fassen wir die Eingänge und die Gewichte zu Vektoren $\vec{x}^{\,T}[k] = (x_1[k], \ldots, x_N[k])$ bzw.\ $\vec{w}^{\,T} = (w_1, \ldots, w_N)$ zusammen und schreiben

$$
z[k] = \vec{x}^{\,T}[k]\,\vec{w} + b\,.
$$

Dabei steht $\vec{x}^{\,T}[k]$ für den zu $\vec{x}[k]$ transponierten Vektor. Offenbar gilt nun für die Linearkombination zweier Eingangsvektoren $\vec{x}_1[k]$ und $\vec{x}_2[k]$ und beliebiges $b$:

\begin{align*}
(c_1\vec{x}_1[k] + c_2\vec{x}_2[k])^T\vec{w} + b &= c_1\vec{x}_1^{\,T}[k]\vec{w} + c_2\vec{x}_2^{\,T}[k]\vec{w} + b \\
&\neq c_1\!\left(\vec{x}_1^{\,T}[k]\vec{w} + b\right) + c_2\!\left(\vec{x}_2^{\,T}[k]\vec{w} + b\right) = c_1 z_1[k] + c_2 z_2[k]\,.
\end{align*}

Dieses Teilsystem ist also nur dann linear, wenn $b = 0$ gilt. Wenn zudem eine nichtlineare Aktivierungsfunktion gewählt wird, ist das Perzeptron in jedem Fall ein nichtlineares System.
\end{beispiel}

\begin{beispiel}[Das Multi-Layer Perzeptron (MLP)]
Beispiel eines vollständig vernetzten Netzwerks mit $M$ Schichten aus jeweils $N$ Neuronen (hier: $M = 3$):

% Grafik: MLP mit Eingängen x₁[k], x₂[k], x₃[k], x₄[k], drei vollständig verbundenen Schichten und Ausgängen y₁^out[k],...,y₄^out[k]

Jeder Knoten berechnet: $y = f(\vec{x}^{\,T}\vec{w} + b)$.

Offensichtlich ist dies ein MIMO-System.

\subsubsection*{Training des Neuronalen Netzes}

\begin{itemize}
  \item Die Gewichte des neuronalen Netzes werden durch Beispieldaten trainiert.
  \item Hierfür wird eine Kostenfunktion spezifiziert, mit deren Hilfe die Abweichung zwischen dem von den Daten repräsentierten Verhalten und dem vom Netzwerk erzeugten Verhalten minimiert wird.
  \item Die Adaptionsvorschrift wird durch Ableitung der Kostenfunktion nach den einzustellenden Gewichten $\vec{w}$ und Biaswerten $b$ ermittelt (,,Backpropagation'').
  \item Netzwerke mit Millionen von Neuronen können heute (unter Verwendung von Parallelrechnern, z.\,B.\ GPUs) zielgerichtet trainiert werden.
\end{itemize}
\end{beispiel}

\section{Lineare und zeitinvariante Systeme}

\subsection{Lineare verschiebungsinvariante Systeme}

Ein LTI/LSI System wird durch seine \textit{Impulsantwort} charakterisiert:

\begin{center}
\begin{tikzpicture}[>=latex, node distance=1.2cm]
  \node (d1) {$\delta(t)$};
  \node [draw, rectangle, right=1cm of d1] (S1) {$\mathcal{S}_{\mathrm{LTI}}[\bullet]$};
  \node [right=1cm of S1] (h1) {$h(t)$};
  \draw[->] (d1) -- (S1);
  \draw[->] (S1) -- (h1);

  \node [below=0.8cm of d1] (d2) {$\delta[k]$};
  \node [draw, rectangle, right=1cm of d2] (S2) {$\mathcal{S}_{\mathrm{LSI}}[\bullet]$};
  \node [right=1cm of S2] (h2) {$h[k]$};
  \draw[->] (d2) -- (S2);
  \draw[->] (S2) -- (h2);
\end{tikzpicture}
\end{center}

Voraussetzung hierfür ist, dass das System vor der Anregung mit dem Impuls in Ruhe ist, d.\,h.\ dass alle inneren Zustände Null sind.

\subsection{Das Faltungsintegral (engl. Convolution)}

Mit der Impulsantwort $h(t) = \mathcal{S}[\delta(t)]$ eines LTI-Systems $\mathcal{S}$ und mit

$$
x(t) = \int_{-\infty}^{\infty} x(\xi)\,\delta(\xi - t)\,\mathrm{d}\xi = \int_{-\infty}^{\infty} x(\xi)\,\delta(t - \xi)\,\mathrm{d}\xi
$$

erhalten wir für ein beliebiges Eingangssignal $x(t)$

\begin{align*}
y(t) = \mathcal{S}\left[x(t)\right] &= \mathcal{S}\left[\int_{-\infty}^{\infty} x(\xi)\,\delta(t-\xi)\,\mathrm{d}\xi\right] \\
&= \int_{-\infty}^{\infty} x(\xi)\,\mathcal{S}\left[\delta(t-\xi)\right]\mathrm{d}\xi \\
&= \int_{-\infty}^{\infty} x(\xi)\,h(t-\xi)\,\mathrm{d}\xi = \int_{-\infty}^{\infty} h(\xi)\,x(t-\xi)\,\mathrm{d}\xi \\
&= x(t) * h(t)
\end{align*}

\begin{beispiel}[RC-Netz: Berechnung des Ausgangssignals]
Die Eingangs-Ausgangsbeziehung des RC-Netzes ist wie folgt gegeben:

$$
y(t) = z(t_0)\,e^{-\frac{t-t_0}{\tau}} + \frac{1}{\tau}\int_{t_0}^{t} x(\xi)\,e^{-\frac{t-\xi}{\tau}}\,\mathrm{d}\xi
$$

Für $z(t_0) = 0$ und $t_0 = -\infty$ erhält man ein Faltungsintegral:

\begin{align*}
y(t) &= \frac{1}{\tau}\int_{-\infty}^{t} x(\xi)\,e^{-\frac{t-\xi}{\tau}}\,\mathrm{d}\xi \\
     &= \int_{-\infty}^{\infty} x(\xi)\,\frac{1}{\tau}\,e^{-\frac{t-\xi}{\tau}}\,\epsilon(t-\xi)\,\mathrm{d}\xi \\
     &= \int_{-\infty}^{\infty} x(\xi)\,h(t-\xi)\,\mathrm{d}\xi
\end{align*}

$\Rightarrow$ Kausales System mit der Impulsantwort $h(t) = \epsilon(t)\,\dfrac{1}{\tau}\,e^{-\frac{t}{\tau}}$.
\end{beispiel}

\subsection{Berechnung der Faltung}

% Grafik: x(t) Rechteckimpuls [0, t₂]; h(t) Rechteckimpuls [0, t₁] und [t₁, t₂]; gespiegeltes/verschobenes x(ξ); Ausgangssignal y(t) als Trapez mit Eckpunkten bei 0, t₁, t₂

\subsection{BIBO-Stabilitätskriterium für LTI-Systeme}

$$
y(t) = \int_{-\infty}^{\infty} h(\xi)\,x(t-\xi)\,\mathrm{d}\xi
$$

Das System ist BIBO-stabil, wenn für $|x(t)| < x_{\max}$ auch $|y(t)| < \infty$ gilt:

\begin{align*}
|y(t)| &= \left|\int_{-\infty}^{\infty} h(\xi)\,x(t-\xi)\,\mathrm{d}\xi\right| \leq \int_{-\infty}^{\infty} |h(\xi)\,x(t-\xi)|\,\mathrm{d}\xi \\
       &= \int_{-\infty}^{\infty} |h(\xi)|\,|x(t-\xi)|\,\mathrm{d}\xi \leq \int_{-\infty}^{\infty} |h(\xi)|\,x_{\max}\,\mathrm{d}\xi \\
       &= x_{\max}\int_{-\infty}^{\infty} |h(\xi)|\,\mathrm{d}\xi < \infty
\end{align*}

Das System ist BIBO-stabil, wenn $\int_{-\infty}^{\infty} |h(\xi)|\,\mathrm{d}\xi < \infty$ gilt (hinreichende und notwendige Bedingung für BIBO-Stabilität).

\subsection{Die zeitdiskrete Faltung}

In gleicher Weise erhält man für ein LSI-System

\begin{align*}
y[k] = \mathcal{S}\left[x[k]\right] &= \mathcal{S}\left[\sum_{l=-\infty}^{\infty} x[l]\,\delta[k-l]\right] \\
&= \sum_{l=-\infty}^{\infty} x[l]\,\mathcal{S}\left[\delta[k-l]\right] \\
&= \sum_{l=-\infty}^{\infty} x[l]\,h[k-l] = x[k] * h[k]\,,
\end{align*}

die zeitdiskrete Faltung, so dass auch hier das Ausgangssignal $y[k]$ allein aus dem Eingangssignal $x[k]$ und der Impulsantwort $h[k]$ berechnet wird.

\subsection{Eigenschaften von LTI/LSI Systemen}

\textbf{Symmetrie} von Eingangssignal und Impulsantwort:

\begin{center}
\begin{tikzpicture}[>=latex, node distance=1cm]
  \node (x1) {$x[k]$};
  \node [draw, rectangle, right=0.8cm of x1] (h1) {$h[k]$};
  \node [right=0.8cm of h1] (y1) {$y[k]$};
  \draw[->] (x1) -- (h1); \draw[->] (h1) -- (y1);

  \node [right=1.5cm of y1] (eq) {$\equiv$};

  \node [right=1cm of eq] (x2) {$h[k]$};
  \node [draw, rectangle, right=0.8cm of x2] (h2) {$x[k]$};
  \node [right=0.8cm of h2] (y2) {$y[k]$};
  \draw[->] (x2) -- (h2); \draw[->] (h2) -- (y2);
\end{tikzpicture}
\end{center}

\textbf{Kaskadierte} Systeme sind vertauschbar:

\begin{center}
\begin{tikzpicture}[>=latex, node distance=0.8cm]
  \node (x1) {$x[k]$};
  \node [draw, rectangle, right=0.6cm of x1] (h11) {$h_1[k]$};
  \node [draw, rectangle, right=0.6cm of h11] (h12) {$h_2[k]$};
  \node [right=0.6cm of h12] (y1) {$y[k]$};
  \draw[->] (x1) -- (h11); \draw[->] (h11) -- (h12); \draw[->] (h12) -- (y1);

  \node [right=1cm of y1] (eq) {$\equiv$};

  \node [right=0.8cm of eq] (x2) {$x[k]$};
  \node [draw, rectangle, right=0.6cm of x2] (h21) {$h_2[k]$};
  \node [draw, rectangle, right=0.6cm of h21] (h22) {$h_1[k]$};
  \node [right=0.6cm of h22] (y2) {$y[k]$};
  \draw[->] (x2) -- (h21); \draw[->] (h21) -- (h22); \draw[->] (h22) -- (y2);
\end{tikzpicture}
\end{center}

\textbf{Parallele} Systeme lassen sich zusammenfassen:

\begin{center}
\begin{tikzpicture}[>=latex, node distance=0.8cm]
  \node (x) {$x[k]$};
  \node [draw, circle, right=0.6cm of x, inner sep=1pt] (sp) {$\bullet$};
  \node [draw, rectangle, above right=0.4cm and 0.6cm of sp] (h1) {$h_1[k]$};
  \node [draw, rectangle, below right=0.4cm and 0.6cm of sp] (h2) {$h_2[k]$};
  \node [draw, circle, right=2cm of sp] (sum) {$+$};
  \node [right=0.6cm of sum] (y) {$y[k]$};
  \draw[->] (x) -- (sp);
  \draw[->] (sp.north east) |- (h1); \draw[->] (h1) -| (sum);
  \draw[->] (sp.south east) |- (h2); \draw[->] (h2) -| (sum);
  \draw[->] (sum) -- (y);

  \node [right=1cm of y] (eq) {$\equiv$};

  \node [right=0.8cm of eq] (x2) {$x[k]$};
  \node [draw, rectangle, right=0.6cm of x2] (hsum) {$h_1[k]+h_2[k]$};
  \node [right=0.6cm of hsum] (y2) {$y[k]$};
  \draw[->] (x2) -- (hsum); \draw[->] (hsum) -- (y2);
\end{tikzpicture}
\end{center}

\begin{beispiel}[Schallausbreitung in halligen Räumen]
% Grafik: Raum mit Sprecher, Mikrofon y[k]; Direktschall (Freifeld), Reflexionen 1. und 2. Ordnung

\subsubsection*{Zeitdiskretes Modell der Schallausbreitung}

In einem halligen Raum gelangt der Schall nicht nur auf direktem Weg sondern auch über \textbf{Rückwürfe von den Wänden} zum Mikrofon. Auch diese verzögert eintreffenden Schallsignale werden linear überlagert. Für den Direktschall und eine Reflexion erhält man:

$$
y[k] = a_0\,x[k-k_0] + a_1\,x[k-k_1] = \sum_{l=0}^{\infty} h[l]\,x[k-l]
$$

mit

$$
h[l] = \begin{cases} a_0 & l = k_0 \\ a_1 & l = k_1 \\ 0 & \text{sonst} \end{cases}
$$

% Grafik: h[l] als zwei Impulse bei k₀ und k₁

Bei vielen und mehrfachen Reflexionen:

$$
y[k] = a_0\,x[k-k_0] + a_1\,x[k-k_1] + a_2\,x[k-k_2] + a_3\,x[k-k_3] + \ldots = \sum_{l=0}^{\infty} h[l]\,x[k-l]
$$

mit

$$
h[l] = \begin{cases} a_0 & l = k_0 \\ a_1 & l = k_1 \\ a_2 & l = k_2 \\ a_3 & l = k_3 \\ \vdots \end{cases}
$$

wobei die Überlagerung der verschiedenen Ausbreitungswege auch zu einer auslöschenden Dämpfung und damit zu negativen Koeffizienten $a_i$ führen kann.

% Grafik: h[l] mit mehreren Impulsen unterschiedlicher Vorzeichen

\subsubsection*{Raumimpulsantwort und diskrete Faltung}

\begin{itemize}
  \item Die Berücksichtigung sämtlicher Rückwürfe führt dann zur \textbf{Raumimpulsantwort} $h[l]$ mit vielen hundert von Null verschiedenen Koeffizienten.
  \item Die Impulsantwort $h[l]$ kann mit einer Impulsanregung des Raums gemessen werden.
\end{itemize}

% Grafik: Raumimpulsantwort h[l], Abtastindex l von 0 bis 2500, Amplitude −0.5 bis 0.5
\end{beispiel}

\section{Parametrische Systeme}

\subsection{Signalflussdiagramme}

\begin{center}
\begin{tabular}{lcc}
  & \textbf{Kontinuierlich} & \textbf{Diskret} \\[6pt]
  Addition: &
  \begin{tikzpicture}[>=latex, baseline=-0.5ex]
    \node (x1) {$x_1(t)$};
    \node [draw, circle, right=0.6cm of x1] (s) {$+$};
    \node [right=0.6cm of s] (y) {$y(t)$};
    \node [below=0.5cm of s] (x2) {$x_2(t)$};
    \draw[->] (x1) -- (s); \draw[->] (s) -- (y); \draw[->] (x2) -- (s);
  \end{tikzpicture}
  &
  \begin{tikzpicture}[>=latex, baseline=-0.5ex]
    \node (x1) {$x_1[k]$};
    \node [draw, circle, right=0.6cm of x1] (s) {$+$};
    \node [right=0.6cm of s] (y) {$y[k]$};
    \node [below=0.5cm of s] (x2) {$x_2[k]$};
    \draw[->] (x1) -- (s); \draw[->] (s) -- (y); \draw[->] (x2) -- (s);
  \end{tikzpicture}
  \\[18pt]
  \parbox{3cm}{Multiplikation mit\\einer Konstanten:} &
  \begin{tikzpicture}[>=latex, baseline=-0.5ex]
    \node (x) {$x(t)$};
    \node [draw, circle, right=0.6cm of x] (s) {$\otimes$};
    \node [right=0.6cm of s] (y) {$y(t)$};
    \node [below=0.5cm of s] (a) {$a$};
    \draw[->] (x) -- (s); \draw[->] (s) -- (y); \draw[->] (a) -- (s);
  \end{tikzpicture}
  &
  \begin{tikzpicture}[>=latex, baseline=-0.5ex]
    \node (x) {$x[k]$};
    \node [draw, circle, right=0.6cm of x] (s) {$\otimes$};
    \node [right=0.6cm of s] (y) {$y[k]$};
    \node [below=0.5cm of s] (a) {$a$};
    \draw[->] (x) -- (s); \draw[->] (s) -- (y); \draw[->] (a) -- (s);
  \end{tikzpicture}
  \\[18pt]
  \parbox{3cm}{Multiplikation mit\\einem Signal:} &
  \begin{tikzpicture}[>=latex, baseline=-0.5ex]
    \node (x1) {$x_1(t)$};
    \node [draw, circle, right=0.6cm of x1] (s) {$\otimes$};
    \node [right=0.6cm of s] (y) {$y(t)$};
    \node [below=0.5cm of s] (x2) {$x_2(t)$};
    \draw[->] (x1) -- (s); \draw[->] (s) -- (y); \draw[->] (x2) -- (s);
  \end{tikzpicture}
  &
  \begin{tikzpicture}[>=latex, baseline=-0.5ex]
    \node (x1) {$x_1[k]$};
    \node [draw, circle, right=0.6cm of x1] (s) {$\otimes$};
    \node [right=0.6cm of s] (y) {$y[k]$};
    \node [below=0.5cm of s] (x2) {$x_2[k]$};
    \draw[->] (x1) -- (s); \draw[->] (s) -- (y); \draw[->] (x2) -- (s);
  \end{tikzpicture}
  \\[18pt]
  Verzögerung: &
  \begin{tikzpicture}[>=latex, baseline=-0.5ex]
    \node (x) {$x(t)$};
    \node [draw, rectangle, right=0.6cm of x] (T) {$T$};
    \node [right=0.6cm of T] (y) {$y(t)$};
    \draw[->] (x) -- (T); \draw[->] (T) -- (y);
  \end{tikzpicture}
  &
  \begin{tikzpicture}[>=latex, baseline=-0.5ex]
    \node (x) {$x[k]$};
    \node [draw, rectangle, right=0.6cm of x] (T) {$T_A$};
    \node [right=0.6cm of T] (y) {$y[k]$};
    \draw[->] (x) -- (T); \draw[->] (T) -- (y);
  \end{tikzpicture}
  \\[18pt]
  Funktion: &
  \begin{tikzpicture}[>=latex, baseline=-0.5ex]
    \node (x) {$x(t)$};
    \node [draw, rectangle, right=0.6cm of x] (f) {$f(x)$};
    \node [right=0.6cm of f] (y) {$y(t)$};
    \draw[->] (x) -- (f); \draw[->] (f) -- (y);
  \end{tikzpicture}
  &
  \begin{tikzpicture}[>=latex, baseline=-0.5ex]
    \node (x) {$x[k]$};
    \node [draw, rectangle, right=0.6cm of x] (f) {$f(x)$};
    \node [right=0.6cm of f] (y) {$y[k]$};
    \draw[->] (x) -- (f); \draw[->] (f) -- (y);
  \end{tikzpicture}
\end{tabular}
\end{center}

\subsection{Parametrische LSI Systeme}

% Grafik: Vollständiges parametrisches LSI-Blockdiagramm: Vorwärtspfad x[k]→T_A→...→T_A gewichtet b₀,...,b_N; Rückkopplungspfad y[k]→T_A→...→T_A gewichtet a₁,...,a_M; alle in einer Summe

\subsection{Nichtrekursives Parametrisches LSI System (FIR)}

% Grafik: x[k] → T_A → T_A → ... → T_A, Abgriffe b₀,b₁,b₂,...,b_N, Summe → y[k]

Dieses System hat eine endliche ausgedehnte Impulsantwort und wird daher auch als ,,finite impulse-response (FIR)'' System bezeichnet.

\subsection{Rekursives Parametrisches LSI System (IIR)}

% Grafik: x[k] → b₀⊗ → (+) → y[k]; Rückkopplungspfad y[k] → T_A → T_A → ... → T_A, gewichtet a₁,a₂,a₃,...,a_M, zurück zur Summe

Dieses System hat im Allgemeinen eine unendliche ausgedehnte Impulsantwort und wird als ,,infinite impulse-response (IIR)'' System bezeichnet.

\section{Weitere Beispiele}

\begin{beispiel}[2-D Faltung in Neuronalen Netzen (CNN)]
Das Bild zeigt die grundlegende Funktion eines neuronalen Faltungsnetzes für zweidimensionale (Bild-)Daten, die durch Kreise auf der linken Seite symbolisiert sind. Die Größe des 2-D Faltungskerns (hier $M = N = 3$) ist durch das gestrichelte Quadrat angedeutet. Zur Erzeugung der Ausgangsdaten (indiziert durch $u$ und $v$) wird dann durch Verschieben des Faltungskerns eine 2-D Faltung implementiert

$$
y[u,v] = f\!\left(\sum_{m=0}^{M-1}\sum_{n=0}^{N-1} w[m,n]\,x[u-m,\,v-n] + b\right)\,,
$$

und durch einen Bias $b$ und eine nichtlineare Aktivierungsfunktion $f$ ergänzt.

% Grafik: 2-D CNN – Eingabegitter x(m,n), gestricheltes 3×3-Faltungskernelement w(m,n), Summierglied (+), Aktivierungsfunktion f → Ausgabe y(u,v)
\end{beispiel}

\begin{beispiel}[Signalübertragung mit PCM und DPCM]
Für jeden Abtastwert wird, entsprechend seiner Amplituden und der daraus resultierenden Quantisierungsstufe, ein digitaler Code übertragen. Dieses Verfahren wird deshalb \textit{pulse code modulation} (PCM) genannt.

Um die Abhängigkeit aufeinander folgender Abtastwerte für die effiziente Quantisierung zu nutzen, werden häufig differentielle Verfahren eingesetzt. Im einfachsten Fall wird vom Sender die Differenz $d(kT_A) = x(kT_A) - x((k-1)T_A)$ quantisiert, codiert und übertragen. Der Empfänger muss dann durch eine laufende Summation die Abtastwerte des Signals $x(kT_A)$ aus den Differenzwerten rekonstruieren. Mit diesem Verfahren, das als \textit{differential pulse code modulation} (DPCM) bezeichnet wird, kann Abtastrate gegen Wortbreite abgetauscht werden.

\subsubsection*{DPCM Sender}

\begin{center}
\begin{tikzpicture}[>=latex, node distance=1cm]
  \node (xt) {$x(t)$};
  \node [draw, rectangle, right=0.8cm of xt] (abt) {Abtastung};
  \node [draw, circle, right=0.8cm of abt] (sum) {$+$};
  \node [right=0.8cm of sum] (dk) {$d[k]$};
  \node [draw, rectangle, below=1cm of sum, text width=1.4cm, align=center] (zsp) {Zwischen\-speicher};
  \node [below=0.5cm of xt] (TA) {$T_A$};

  \draw[->] (xt) -- (abt);
  \draw[->] (abt) -- node[above] {$x[k]$} (sum);
  \draw[->] (sum) -- (dk);
  \draw[->] (TA) -- (abt);
  \draw[->] (sum.south) -- (zsp.north);
  \draw[->] (zsp.west) -| node[left, near end] {$z_s[k]$} (sum.south west);
\end{tikzpicture}
\end{center}

Berechnung der zeitdiskreten Ausgangssignale für $k \geq 0$ mit Differenzengleichungen. Für den Sender erhält man:

$$
d[k] = x[k] - z_s[k] = \begin{cases} x[0] - z_s[0] & k = 0 \\ x[k] - x[k-1] & k \geq 1 \end{cases}
$$

\subsubsection*{DPCM Empfänger}

\begin{center}
\begin{tikzpicture}[>=latex, node distance=1cm]
  \node (dk) {$d[k]$};
  \node [draw, circle, right=0.8cm of dk] (sum) {$+$};
  \node [right=0.8cm of sum] (xhk) {$\widehat{x}[k]$};
  \node [draw, rectangle, right=1cm of xhk] (rek) {Rekonstruktion};
  \node [right=0.8cm of rek] (xht) {$\widehat{x}(t)$};
  \node [draw, rectangle, below=1cm of sum, text width=1.4cm, align=center] (zsp) {Zwischen\-speicher};
  \node [below=0.5cm of rek] (TA) {$T_A$};

  \draw[->] (dk) -- (sum);
  \draw[->] (sum) -- (xhk);
  \draw[->] (xhk) -- (rek);
  \draw[->] (rek) -- (xht);
  \draw[->] (TA) -- (rek);
  \draw[->] (xhk.south) |- (zsp.east);
  \draw[->] (zsp.west) -| node[left, near end] {$z_E[k]$} (sum.south west);
\end{tikzpicture}
\end{center}

Für den Empfänger erhält man:

$$
\widehat{x}[k] = d[k] + z_E[k] = z_E[0] + \sum_{n=0}^{k} d[n]
$$

Für $z_S[0] = z_E[0]$ gilt $\widehat{x}[k] = x[k]\ \forall\ k \geq 0$. Sender und Empfänger enthalten Speicher für einen Abtastwert bzw.\ die laufende Summe. Diese stellen die inneren Zustände dar.

\subsubsection*{Differenzsignale}

% Grafik: Zwei Plots bei f_A = 16 kHz; oben x(kT_A) (Originalsignal), unten d(kT_A) (Differenzsignal), t/s von 0 bis 6

% Grafik: Zwei Plots bei f_A = 64 kHz; oben x(kT_A), unten d(kT_A), t/s von 0 bis 6

\subsubsection*{Original und Differenzsignal}

% Grafik (4 Plots): Vergleich f_A = 16 kHz und f_A = 64 kHz, jeweils x(kT_A) und d(kT_A) im Ausschnitt kT_A von 0.156562 bis ~0.1578

\subsubsection*{Differenzsignal und rekonstruiertes Signal}

% Grafik (4 Plots): f_A = 16 kHz und f_A = 64 kHz – oben d(kT_A), unten x̂(kT_A), Ausschnitt kT_A ≈ 0.156562–0.1578

\subsubsection*{Quantisiertes (2 Bit) und rekonstruiertes Signal}

% Grafik (6 Plots): f_A = 16 kHz und f_A = 64 kHz – oben Q_q(d(kT_A)) (2-Bit-quantisiertes Differenzsignal), Mitte x̂(kT_A) (rekonstruiert), unten e²(kT_A) (quadratischer Fehler, Skala ×10⁻³)
\end{beispiel}

\section*{Zusammenfassung}

\begin{itemize}
  \item Systeme bilden Eingangssignale auf Ausgangssignale ab.
  \item Hierbei spielen die inneren Zustandvariablen des Systems eine große Rolle.
  \item Wichtige Systemeigenschaften sind die Linearität, die Zeitinvarianz oder Verschiebungsinvarianz und die Stabilität.
  \item LTI bzw.\ LSI Systeme werden durch ihre Impulsantwort beschrieben.
  \item Zeitinvariante und zeitvariante Systeme können mit einem Zustandsraummodell beschrieben und berechnet werden.
\end{itemize}



\chapter{Einführung in die Wahrscheinlichkeitsrechnung}

\section{Einführung und Definition}

Experimente, deren Ergebnisse (uns) ungewiss erscheinen, werden Zufallsexperimente genannt und deren Ergebnisse Elementarereignisse $\xi$.

Die Menge aller Elementarereignisse wird mit $\Omega$ bezeichnet.

In der Wahrscheinlichkeitsrechnung werden den Ergebnissen von Experimenten Wahrscheinlichkeiten zugeordnet, so dass bei mehrfacher Ausführung des Zufallsexperiments aus den Wahrscheinlichkeiten auf die Häufigkeiten der Ergebnisse geschlossen werden kann.

Damit wird es möglich, aus einer Vielzahl von Beobachtungen auch auf das zukünftige Verhalten von Systemen zu schließen.

\subsection*{Analyse von Videodaten}

% Grafik: Videobild (Baumkronen), Schaltfläche „play"

\subsection*{Analyse von Audiodaten}

Signalform des Wortes „one", gesprochen von 4 verschiedenen Sprechern.

% Grafik: 4 Zeitverläufe x(t) für speaker 1–4, t/s von 0 bis 1, Amplitude −0{,}2 bis 0{,}2

$\to$ Um die Information „one" sicher zu erkennen, muss ein automatischer Spracherkenner auf die statistische Vielfalt der Audiodaten trainiert werden.

\subsection{Mengen und Ereignisse}

\begin{itemize}
  \item Gegeben sei eine Menge $\Omega = \{\xi_1, \xi_2, \ldots, \xi_K\}$ von Elementen $\xi_i$, die man \textit{Elementarereignisse} nennt.
  \item Ein Ereignis $A$ ist dann eine Teilmenge von $\Omega$, d.\,h.\ $A \subseteq \Omega$.
  \item Wir sagen, dass ein Ereignis $A$ eintritt, wenn das Ergebnis eines Experiments einem Element der Menge $A$ entspricht.
  \item Die Menge $\Omega$ enthält alle in einem Experiment möglichen Elementarereignisse. $\Omega$ wird deshalb als das sichere (stets eintretende) Ereignis bezeichnet.
  \item $\phi = \{\}$ bezeichnet die leere Menge, das nie eintretende (unmögliche) Ereignis.
  \item Bei einer Beschränkung auf endlich viele Elementarereignisse kann dann jede Teilmenge ein Ereignis sein, dann $A = \{\xi_i\} \subseteq \Omega$.
\end{itemize}

\subsection{Endliche und unendliche Ereignismengen}

Die Menge der Elementarereignisse $\Omega$ kann endlich viele Elemente, oder (abzählbar oder überabzählbar) unendlich viele Elemente aufweisen. Beispiele:

\begin{enumerate}
  \item Das Zufallsexperiment besteht aus dem Werfen einer Münze, $\Omega = \{\text{Kopf}, \text{Zahl}\}$.
  \item Das Zufallsexperiment besteht aus dem Werfen eines Würfels, $\Omega = \{1,2,3,4,5,6\}$.
  \item Das Zufallsexperiment besteht in der Detektion einer Zahl von Röntgenquanten in einem Röntgengerät, $\Omega = \mathbb{N}$.
  \item Das Zufallsexperiment besteht in der Messung einer Spannung $U$, die durch Messfehler (Rauschen) gestört ist, $\Omega = \mathbb{R}$.
\end{enumerate}

\subsection{Ereignisse}

Ausgehend von der Grundmenge der Elementarereignisse können nun eine Vielzahl von speziellen, an eine Fragestellung angepassten Ereignisse formuliert werden. Diese werden durch Teilmengen $A \subseteq \Omega$ repräsentiert.

Beispielsweise kann ein Ereignis in dem Würfeln einer geraden Zahl bestehen. Es gilt dann:

$$
A_g = \{2, 4, 6\} \subset \{1, 2, 3, 4, 5, 6\} = \Omega
$$

\subsection{Die Vereinigungsmenge}

Die Vereinigungsmenge zweier Mengen $A$ und $B$ enthält alle Elemente dieser Mengen. Elemente, die sowohl in $A$ als auch in $B$ enthalten sind, werden nur einmal in der Vereinigungsmenge aufgenommen.

\begin{itemize}
  \item Notation: $A + B$ oder $A \cup B$
\end{itemize}

Für drei Mengen $A$, $B$ und $C$ sowie die leere Menge $\phi$ gelten die folgenden Gesetzmäßigkeiten:

\begin{itemize}
  \item $A + B = B + A$ \hfill (Kommutativgesetz)
  \item $(A + B) + C = A + (B + C)$ \hfill (Assoziativgesetz)
  \item $A + A = A$ \hfill (Idempotenz)
  \item $A + \phi = A$ \hfill (Identität)
  \item Wenn $B \subseteq A$, dann gilt $A + B = A$.
\end{itemize}

% Grafik: Venn-Diagramm A ∪ B

\subsection{Die Schnittmenge}

Die Schnittmenge zweier Mengen $A$ und $B$ enthält nur die Elemente dieser Mengen, die sowohl in $A$ als auch in $B$ enthalten sind.

\begin{itemize}
  \item Notation: $AB$ oder $A \cap B$
\end{itemize}

Für drei Mengen $A$, $B$ und $C$ sowie die leere Menge gelten die folgenden Gesetzmäßigkeiten:

\begin{itemize}
  \item $AB = BA$ \hfill (Kommutativgesetz)
  \item $(AB)C = A(BC)$ \hfill (Assoziativgesetz)
  \item $(A + B)C = AC + BC$ \hfill (Distributivgesetz)
  \item Wenn die Mengen keine gemeinsamen Elemente aufweisen, gilt: $AB = \phi$, also auch $A\phi = \phi$.
  \item Wenn $A \subseteq B$, dann gilt $AB = A$.
\end{itemize}

% Grafik: Venn-Diagramm A ∩ B

\subsection{Die Differenzmenge}

Die Differenzmenge $A - B$ von zwei Mengen $A$ und $B$ erhält man, in dem man die in $A$ enthaltenen Elemente von $B$ aus $A$ entfernt.

Beispiel: $A = \{1,2\}$; $B = \{2,3\}$

$$
A - B = \{1\}, \qquad B - A = \{3\}
$$

Für die Differenzbildung gilt weder das Kommutativgesetz noch das Assoziativgesetz.

% Grafik: Venn-Diagramm A − B (linke Halbmonde von A schattiert)

\subsection{Das Mengenkomplement}

Das Komplement einer Menge $A$, die Teilmenge einer Grundmenge $\Omega$ ist, wird mit

$$
\overline{A} = \Omega - A
$$

berechnet.

Beispiel: $\Omega = \{1,2,3\}$; $A = \{2,3\}$

$$
\overline{A} = \Omega - A = \{1\}
$$

% Grafik: Venn-Diagramm mit Ā schattiert (Komplement von A in Ω)

\subsection{Mengenalgebra}

\begin{itemize}
  \item Unter einem Mengensystem versteht man eine Menge, deren sämtliche Elemente selbst Mengen sind.
  \item Ein Mengensystem $\mathcal{F}$ heißt Mengenkörper, wenn die Vereinigungsmenge, die Schnittmenge und die Differenz $A - B$ von zwei Mengen $A$ und $B$ des Systems wieder dem Mengensystem angehören.
  \item Jeder nicht leere Mengenkörper enthält die leere Menge $\phi$.
  \item Ein Mengensystem $\mathcal{F}$ von Teilmengen von $\Omega$ heißt $\sigma$-Algebra, wenn Folgendes gilt:
    \begin{enumerate}
      \item $\Omega \in \mathcal{F}$
      \item $A \in \mathcal{F} \Rightarrow \overline{A} \in \mathcal{F}$
      \item $A_1, A_2, \ldots \in \mathcal{F} \Rightarrow \bigcup_{i=1}^{\infty} A_i \in \mathcal{F}$
    \end{enumerate}
\end{itemize}

\begin{beispiel}[Endliche Mengen]
Es wird eine Quelle betrachtet, die einen Strom binärer Zeichen liefert. Wenn die Ausgabe eines einzelnen Symbols als Ereignis angesehen wird, ist die Menge der Elementarereignisse durch $\Omega = \{\xi_1, \xi_2\} = \{0, 1\}$ gegeben.

% Grafik: Nachrichtenquelle → ... 0 0 1 0 1 0 1 1 ...

Aus dieser Menge von Elementarereignissen $\Omega = \{0, 1\}$ lässt sich der Mengenkörper $\mathcal{F}$ mit den folgenden Elementen ableiten:

$$
A_1 = \{\xi_1\} = \overline{A_2}, \qquad A_2 = \{\xi_2\} = \overline{A_1}
$$
$$
A_3 = \{\xi_1, \xi_2\} = A_1 \cup A_2 = \Omega, \qquad A_4 = A_1 \cap A_2 = \phi
$$
\end{beispiel}

\subsection{Überabzählbare Mengen}

Wenn die Menge der Elementarereignisse einem Intervall $[0,1] \in \mathbb{R}$ entspricht, erfordert die Definition eines Mengenkörpers zusätzliche Überlegungen. Man verwendet dann eine Borel-Algebra:

Die Borel-Algebra $\mathcal{B}$ ist die kleinste $\sigma$-Algebra auf $[0,1]$, die alle Intervalle $(a, b)$ mit $0 \leq a < b \leq 1$ enthält.

Im Fall kontinuierlicher Ereignisräume (z.\,B.\ $\Omega = \mathbb{R}$) muss also mit Intervallen gearbeitet werden. Den Elementarereignissen sind Intervalle zugeordnet.

\subsection{Axiomatische Definition der Wahrscheinlichkeit}

Die Wahrscheinlichkeit eines Ereignisses $A_i$ aus einem Mengenkörper $\mathcal{F} = \{A_1, \ldots, A_i, \ldots, A_N\}$, wobei $A_i$ Teilmengen von $\Omega$ bezeichnen und in $\mathcal{F}$ enthalten sind, ist eine reelle Zahl $P(A_i)$ mit den folgenden Eigenschaften:

\begin{description}
  \item[Axiom 1:] Die Wahrscheinlichkeit $P(A_i)$ ist nicht negativ, d.\,h.\ $P(A_i) \geq 0$.
  \item[Axiom 2:] Die Wahrscheinlichkeit $P(\Omega)$ des sicheren Ereignisses ist Eins, d.\,h.\ $P(\Omega) = 1$.
  \item[Axiom 3:] Die Wahrscheinlichkeit, dass eines von zwei sich gegenseitig ausschließenden Ereignissen $A_i$ und $A_j$, $A_i \cap A_j = A_i \cdot A_j = \phi$ auftritt, ist $P(A_i + A_j) = P(A_i) + P(A_j)$.
\end{description}

\subsection{Wahrscheinlichkeitsraum}

Mit der Zuordnung von Wahrscheinlichkeiten zu Mengen wird ein Wahrscheinlichkeitsmaß $\mathbb{P}$ definiert.

Das Triple $(\Omega, \mathcal{F}, \mathbb{P})$ wird Wahrscheinlichkeitsraum genannt. Dabei ist $\Omega \neq \phi$ eine Menge, $\mathcal{F}$ eine $\sigma$-Algebra über $\Omega$ und $\mathbb{P}$ ein Wahrscheinlichkeitsmaß auf $\Omega$.

\subsection{Konstruktion des Wahrscheinlichkeitsraums}

Im Fall von Experimenten mit endlich vielen Elementarereignissen werden der Mengenkörper $\mathcal{F}$ und die zugehörigen Wahrscheinlichkeiten $P$ wie folgt konstruiert. Man definiert

\begin{itemize}
  \item eine beliebige Menge aus sich ausschließender Elementarereignisse $\Omega = \{\xi_1, \xi_2, \ldots, \xi_N\}$ und
  \item eine beliebige Menge $\{P(\xi_1), \ldots, P(\xi_N)\}$ von nicht-negativen Zahlen mit der Summe $P(\xi_1) + \cdots + P(\xi_N) = 1$. Diese bezeichnen die Wahrscheinlichkeiten der elementaren Ereignisse $\xi_i$.
\end{itemize}

$\mathcal{F}$ wird dann aus der Gesamtheit aller Teilmengen von $\Omega$ gebildet, wobei man für ein Ereignis $A_i = \{\xi_{i_1}, \ldots, \xi_{i_k}\}$ die Wahrscheinlichkeit zu $P(A_i) = P(\xi_{i_1}) + \cdots + P(\xi_{i_k})$ setzt.

\subsection{Laplace-Experimente}

Treten alle Ergebnisse eines Experiments gleichwahrscheinlich auf, so wird das Experiment Laplace-Experiment genannt. Es gilt dann

$$
P(\xi) = \frac{1}{|\Omega|},
$$

wobei $|\Omega|$ die Mächtigkeit (Kardinalität, Anzahl der Elemente einer Menge) der Menge $\Omega$ bezeichnet.

\begin{beispiel}[Idealer Würfel]
Die Seiten des Würfels zeigen sechs verschiedene Symbole. Diese definieren die möglichen Elementarereignisse, wobei das nach einem Wurf oben liegende Symbol als Ergebnis eines Würfs angesehen wird. Die Elementarereignisse schließen sich gegenseitig aus. Der \textit{ideale} Würfel ist dadurch definiert, dass die sechs möglichen Elementarereignisse gleichwahrscheinlich auftreten.

\begin{center}
\begin{tabular}{lcccccc}
  \hline
  Elementarereignis & $\xi_1$ & $\xi_2$ & $\xi_3$ & $\xi_4$ & $\xi_5$ & $\xi_6$ \\
  Symbol            & 1 & 2 & 3 & 4 & 5 & 6 \\
  $P(\xi_i)$        & $\frac{1}{6}$ & $\frac{1}{6}$ & $\frac{1}{6}$ & $\frac{1}{6}$ & $\frac{1}{6}$ & $\frac{1}{6}$ \\
  \hline
\end{tabular}
\end{center}

Die Menge $\Omega = \{1,2,3,4,5,6\}$ repräsentiert das sichere Ereignis. Es gilt

$$
P(\xi_1) + P(\xi_2) + P(\xi_3) + P(\xi_4) + P(\xi_5) + P(\xi_6) = 1
$$
\end{beispiel}

\subsection{Abzählbare unendliche Mengen}

Das Konstruktionsprinzip ist auch auf abzählbar unendliche Mengen, zum Beispiel für $|\Omega| = |\mathbb{N}|$, anwendbar. Auch hier muss sichergestellt werden, dass $\sum_{i=1}^{\infty} P(\xi_i) = 1$ gilt. Dies ist zum Beispiel für $P(\xi_i) = 2^{-i}$ gegeben, denn

$$
\sum_{i=1}^{\infty} 2^{-i} = \frac{1}{1 - 0.5} - 1 = 1
$$

Für die Borel-Algebra $\mathcal{B}$ wird ein Wahrscheinlichkeitsmaß durch die Länge der Intervalle $(a, b)$, $0 \leq a \leq b \leq 1$, also durch $b - a$ definiert.

\begin{beispiel}[Ampelphase]
Sie fahren mit dem Auto durch eine Kurve und erblicken plötzlich eine Ampel.

% Grafik: Ampelphasen – Zeitleiste mit grün, rot, grün, t/min

Wie groß ist die Wahrscheinlichkeit, dass die Ampel in weniger als 2 Sekunden die Farbe wechselt (Gelb-Phase wird vernachlässigt)?
\end{beispiel}

\subsection{Kartesische Produktmengen}

\begin{itemize}
  \item Gegeben sind zwei Mengen $\Omega_1 = \{\xi_1, \ldots, \xi_L\}$ und $\Omega_N = \{\eta_1, \ldots, \eta_N\}$.
  \item $(\xi_i, \eta_j)$ mit $\xi_i \in \Omega_1$ und $\eta_j \in \Omega_N$ bezeichnet ein geordnetes Paar von Elementarereignissen.
  \item Unter dem kartesischen Produkt $\Omega_1 \times \Omega_N$ der Mengen $\Omega_1$ und $\Omega_N$ versteht man die Menge aller geordneten Paare $(\xi_i, \eta_j)$ mit $\xi_i \in \Omega_1$ und $\eta_j \in \Omega_N$.
\end{itemize}

Anmerkungen:

\begin{itemize}
  \item Das kartesische Produkt erzeugt eine Menge geordneter Paare, wobei jeweils ein Element aus $\Omega_1$ mit einem Element aus $\Omega_N$ zu einem Paar kombiniert wird.
  \item Kartesische Produktmengen können aus beliebigen Ereignismengen konstruiert werden.
  \item Für endliche Mengen $\Omega_1, \ldots, \Omega_L$ gilt $|\Omega_1 \times \cdots \times \Omega_L| = |\Omega_1| \cdot \ldots \cdot |\Omega_L| = \prod_{i=1}^{L} |\Omega_i|$.
\end{itemize}

\begin{beispiel}[Kartesische Produktmengen: Zweimaliger Würfelwurf]
Für den idealen Würfel können wir die Menge der Elementarereignisse mit $\Omega_w = \{\xi_1, \xi_2, \xi_3, \xi_4, \xi_5, \xi_6\} = \{1,2,3,4,5,6\}$ angeben. Zur Beschreibung eines zweimaligen Wurfes wird das kartesische Produkt $\Omega_w \times \Omega_w$ gebildet. Beim zweimaligen Werfen eines idealen Würfels sind dann die folgenden Elementarereignisse möglich:

$$
\Omega = \Omega_w \times \Omega_w = \{(1,1),(1,2),(1,3),\ldots\}
$$

bzw. in zweidimensionaler Anordnung

\begin{center}
\begin{tabular}{c|cccccc}
  $\xi_i \backslash \xi_j$ & 1 & 2 & 3 & 4 & 5 & 6 \\
  \hline
  1 & (1,1) & (1,2) & (1,3) & (1,4) & (1,5) & (1,6) \\
  2 & (2,1) & (2,2) & (2,3) & (2,4) & (2,5) & (2,6) \\
  3 & (3,1) & (3,2) & (3,3) & (3,4) & (3,5) & (3,6) \\
  4 & (4,1) & (4,2) & (4,3) & (4,4) & (4,5) & (4,6) \\
  5 & (5,1) & (5,2) & (5,3) & (5,4) & (5,5) & (5,6) \\
  6 & (6,1) & (6,2) & (6,3) & (6,4) & (6,5) & (6,6) \\
\end{tabular}
\end{center}
\end{beispiel}

\subsection{Mehrfache, unabhängige Versuche}

Nun werden mehrere Experimente, die unabhängig voneinander unter gleichen Bedingungen wiederholt werden, gemeinsam betrachtet.

Im allgemeinen Fall sind die Elementarereignisse des $L$-fach wiederholten ($L$-stufigen) Experiments durch die $L$-Tupel der kartesischen Produktmenge

$$
\Omega = \Omega_1 \times \Omega_2 \times \cdots \times \Omega_L
$$

repräsentiert.

Wenn wir die Versuche als \textit{unabhängig} bezeichnen, ist im Fall $L = 2$

$$
P\left((\xi_i, \xi_j)\right) = P(\xi_i) P(\xi_j) \qquad \forall\, i, j
$$

die einzige sinnvolle Zuordnung.

\begin{beispiel}[Mehrfaches Werfen einer Münze]
\begin{enumerate}
  \item Wie groß ist die Wahrscheinlichkeit, dass beim dritten Wurf einer fairen Münze „Kopf" erstmalig erscheint?
  \item Wie groß ist die Wahrscheinlichkeit, dass auch nach $l$ Würfen kein „Kopf" aufgetreten ist?
  \item Wie groß ist die Wahrscheinlichkeit, dass nach $l$ Würfen mindestens einmal „Kopf" aufgetreten ist?
\end{enumerate}
\end{beispiel}

\subsection{Kombinatorik}

Der Wahrscheinlichkeitsrechnung liegen häufig kombinatorische Aufgaben zugrunde. Daher ist es sinnvoll, sich zunächst einige kombinatorische Formeln anzueigenen:

\begin{itemize}
  \item Die Anzahl der Permutationen von $n$ Elementen ohne Wiederholung (Aneinanderreihung) ist durch $n!$ gegeben.
  \item Die Anzahl der Variationen von $s$ geordneten, voneinander verschiedenen Elementen einer $n$-elementigen Menge ist durch $\dfrac{n!}{(n-s)!}$ gegeben, wenn keine Wiederholung vorliegt.
  \item Die Anzahl aller voneinander verschiedenen Kombinationen einer $s$-elementigen Teilmenge einer $n$-elementigen Menge ohne Wiederholung ist durch $\dfrac{n!}{s!(n-s)!} = \binom{n}{s}$ gegeben.
\end{itemize}

\subsection{Stichproben}

Die Vielfalt der Ergebnisse eines mehrstufigen Experiments wird unter anderem durch die beiden folgenden Prinzipien bestimmt:

\begin{itemize}
  \item Auswahl der Ergebnisse
  \item Anordnung der Ergebnisse
\end{itemize}

Im Folgenden wird dies anhand des $s$-fachen Ziehens aus einer Menge mit $n$ nummerierten Kugeln behandelt. In der Statistik wird dies auch als Stichprobe bezeichnet.

\subsection{Geordnete Stichproben}

Es wird eine von $n$ Kugeln gezogen und nach dem Notieren der Nummer zurückgelegt. Das Experiment wird $s$-fach wiederholt. Da in jeder Stufe $n$ Ergebnisse möglich sind, enthält die Menge aller Ergebnisse $n^s$ Elemente.

Werden die gezogenen Kugeln nicht zurückgelegt, erhält man

$$
n \cdot (n-1) \cdot \ldots \cdot (n-s+1) = \frac{n!}{(n-s)!}
$$

Variationen.

Daraus folgt auch, dass $n$ Elemente auf $n!$ verschiedene Arten angeordnet werden können. Es existieren $n!$ Permutationen.

\subsection{Ungeordnete Stichproben}

Wenn die Reihenfolge, mit der die Kugeln gezogen werden, keine Rolle spielt, wird die Menge aller möglichen Ergebnisse kleiner. Für das $s$-stufige Experiment erhalten wir dann

$$
\binom{n}{s} = \frac{n!}{s!(n-s)!}
$$

Kombinationen, wenn die gezogenen Kugeln nicht zurückgelegt werden. Werden die gezogenen Kugeln zurückgelegt, erhält man

$$
\binom{s + n - 1}{s}
$$

mögliche Ergebnisse.

\begin{beispiel}[Binäres Codewort]
Eine Quelle sendet ein binäres Codewort mit 16 Bit, wobei das Ereignis $\xi_1 = 0$ mit der Wahrscheinlichkeit $P(0) = 0{,}6$ im Codewort auftritt. Die Bits des Codeworts sind statistisch unabhängig.

Erläutern Sie im Zusammenhang mit diesem Beispiel die Anwendung der Axiome der Wahrscheinlichkeitsrechnung, der Definition der statistischen Unabhängigkeit und der kombinatorischen Theoreme.

Wie groß ist die Wahrscheinlichkeit, dass genau sechs Bits mit $\xi_2 = 1$ im Codewort enthalten sind?
\end{beispiel}

\subsection{Verbundwahrscheinlichkeiten}

Generell wird die Wahrscheinlichkeit, dass zwei Ereignisse $A_i \in \mathcal{A}$ und $B_j \in \mathcal{B}$ gemeinsam auftreten, als Verbundwahrscheinlichkeit bezeichnet und mit $P(A_i, B_j) = P(B_j, A_i)$ angegeben. Im Fall, dass diese Ereignisse statistisch unabhängig sind (und nur dann), wird die Verbundwahrscheinlichkeit zu

$$
P(A_i, B_j) = P(A_i) \cdot P(B_j)
$$

berechnet.

\subsection{Partition von $\Omega$}

Ein Mengensystem $\mathcal{B} = \{B_1, B_2, \ldots\}$ von Teilmengen einer Menge $\Omega$ heißt Partition von $\Omega$, falls

\begin{enumerate}
  \item $B_i B_j = \phi$ \quad (paarweise disjunkt)
  \item $\Omega = B_1 + B_2 + \ldots$ \quad gilt.
\end{enumerate}

% Grafik: Tabelle mit ξ₁–ξ₃₂, zeilenweise angeordnet als Partition von Ω

\subsection{Randwahrscheinlichkeiten}

Die Wahrscheinlichkeiten für das Auftreten der einzelnen sich gegenseitig ausschließenden Ereignisse $A_i \in \mathcal{A}$ oder $B_j \in \mathcal{B}$ kann wieder aus den Verbundwahrscheinlichkeiten berechnet werden, indem man die Wahrscheinlichkeiten über die Ereignisse einer Partition summiert. Man erhält

$$
P(A_i) = \sum_{j=1}^{N} P(A_i, B_j) = P(A_i, \mathcal{B})
$$

$$
P(B_j) = \sum_{i=1}^{N} P(A_i, B_j) = P(B_j, \mathcal{A})
$$

Wenn die Wahrscheinlichkeiten der einzelnen Ereignisse aus den Verbundwahrscheinlichkeiten $P(A_i, B_j)$ gewonnen werden, werden die Wahrscheinlichkeiten $P(A_i)$ und $P(B_j)$ als Randwahrscheinlichkeiten bezeichnet.

\begin{beispiel}[Randwahrscheinlichkeiten]
\begin{center}
\begin{tabular}{c|cc}
  & $B_1$ & $B_2$ \\
  \hline
  $A_1$ & $P(A_1, B_1)$ & $P(A_1, B_2)$ \\
  $A_2$ & $P(A_2, B_1)$ & $P(A_2, B_2)$ \\
\end{tabular}
\end{center}
\end{beispiel}

\section{Bedingte Wahrscheinlichkeiten und der Satz von Bayes}

% Grafik: Videobild (Baumkronen), Schaltfläche „play"

Die Bildpunkte eines Videosignals weisen ein hohes Maß an statistischer Abhängigkeit auf.

\subsection{Bedingte Wahrscheinlichkeiten}

Die Wahrscheinlichkeit $P(A_i|B_j)$, dass ein Ereignis $A_i \in \mathcal{A}$ eintritt, wenn das Ergebnis $B_j \in \mathcal{B}$ eines anderen Experiments mit $P(B_j) > 0$ vorausgesetzt wird, wird als bedingte Wahrscheinlichkeit bezeichnet und ist zu

$$
P(A_i|B_j) = \frac{P(A_i, B_j)}{P(B_j)}
$$

definiert. Es gilt $0 \leq P(A_i|B_j) \leq 1$.

Wenn die Ereignisse $A_i$ und $B_j$ statistisch unabhängig sind gilt

$$
P(A_i|B_j) = \frac{P(A_i, B_j)}{P(B_j)} = \frac{P(A_i)P(B_j)}{P(B_j)} = P(A_i).
$$

In diesem Fall bleibt bei der Kenntnis des Ereignisses $B_j$ ohne Auswirkung auf die Wahrscheinlichkeit von $A_i$.

\begin{beispiel}[Bedingte Wahrscheinlichkeiten]
Wir betrachten die Erzeugung der Symbole $A_1$ und $A_2$ durch die Quelle $Q_A$ und der Symbole $B_1$ und $B_2$ durch die Quelle $Q_B$, also $\Omega_A = \{A_1, A_2\}$ und $\Omega_B = \{B_1, B_2\}$.

Außerdem ist bekannt, dass die Wahrscheinlichkeiten für das gemeinsame Auftreten der Symbole durch

$$
P(A_1, B_1) = 0{,}3 \qquad P(A_2, B_1) = 0{,}2
$$
$$
P(A_1, B_2) = 0{,}2 \qquad P(A_2, B_2) = 0{,}3
$$

gegeben sind. Die Randwahrscheinlichkeiten für das Auftreten eines Symbols ergeben sich somit zu $P(A_1) = P(A_2) = P(B_1) = P(B_2) = 0{,}5$.

\begin{itemize}
  \item Sind die Symbole der Quellen $Q_A$ und $Q_B$ statistisch abhängig?
  \item Berechnen Sie die bedingte Wahrscheinlichkeit $P(A_1|B_1)$.
\end{itemize}
\end{beispiel}

\begin{beispiel}[Gameshow]
Hinter einer der Türen verbirgt sich ein Auto, hinter den anderen jeweils eine Ziege. Wenn Sie die richtige Tür raten, gehört das Auto Ihnen. Welche Tür wählen Sie?

% Grafik: 3 Türen – eine mit Ziege, eine mit Auto, eine mit „?"

Der Showmaster öffnet nun eine der beiden nicht ausgewählten Türen, hinter der sich eine Ziege verbirgt.

Frage: Ist es sinnvoll, sich nun für die bisher nicht gewählte Tür zu entscheiden?
\end{beispiel}

\subsection{Totale Wahrscheinlichkeit}

Für einen Wahrscheinlichkeitsraum $(\Omega, \mathcal{F}, \mathbb{P})$ mit der Partition $B_1, \ldots, B_N \in \mathcal{F}$ und $P(B_i) > 0$ für jedes $A_k \in \mathcal{F}$

$$
P(A_k) = \sum_{i=1}^{N} P(A_k|B_i)\,P(B_i)
$$

Das Theorem von der totalen Wahrscheinlichkeit erlaubt die Berechnung einer unbekannten Wahrscheinlichkeit $P(A_k)$ durch Entwicklung in eine Summe von oft einfacher zu ermittelnden bedingten Wahrscheinlichkeiten $P(A_k|B_i)$ und den \textit{a priori} Wahrscheinlichkeiten $P(B_i)$.

\subsection{Der Satz von Bayes}

Sei $(\Omega, \mathcal{F}, \mathbb{P})$ ein Wahrscheinlichkeitsraum und $B_1, \ldots, B_N \in \mathcal{F}$ eine Partition von $\Omega$ mit $P(B_i) > 0$ für alle $i = 1, \ldots, N$. Dann gilt für jedes $A_k \in \mathcal{F}$ mit $P(A_k) > 0$ und $B_j$

$$
P(B_j|A_k) = \frac{P(A_k|B_j)P(B_j)}{\displaystyle\sum_{i=1}^{N} P(A_k|B_i)P(B_i)} = \frac{P(A_k|B_j)P(B_j)}{\displaystyle\sum_{i=1}^{N} P(A_k, B_i)} = \frac{P(A_k|B_j)P(B_j)}{P(A_k)}
$$

\begin{beispiel}[Symmetrischer Binärkanal]
Eine Datenquelle erzeugt die Symbole $\xi_1 = 0$ und $\xi_2 = 1$ mit den Wahrscheinlichkeiten $P(\xi_1) = P_1$ und $P(\xi_2) = P_2 = 1 - P_1$. Bei der Übertragung dieser Symbole treten Fehler für die beiden Fälle $0 \to 1$ und $1 \to 0$ mit der gleichen Wahrscheinlichkeit $P_e$ auf.

\begin{center}
\begin{tikzpicture}[>=latex, node distance=1.5cm]
  \node (xi1) {$\xi_1 = 0$};
  \node [below=1.2cm of xi1] (xi2) {$\xi_2 = 1$};
  \node [right=4cm of xi1] (eta1) {$0 = \eta_1$};
  \node [right=4cm of xi2] (eta2) {$1 = \eta_2$};
  \draw[->] (xi1) -- node[above] {$1 - P_e$} (eta1);
  \draw[->] (xi1) -- node[above, sloped, near end] {$P_e$} (eta2);
  \draw[->] (xi2) -- node[below, sloped, near end] {$P_e$} (eta1);
  \draw[->] (xi2) -- node[below] {$1 - P_e$} (eta2);
\end{tikzpicture}
\end{center}

Auf der Grundlage dieses Modells und mit Hilfe des Satzes von Bayes lässt sich nun auf die gesendeten Bits schließen, z.\,B.

$$
P(\xi_1|\eta_1) = \frac{P(\eta_1|\xi_1) \cdot P(\xi_1)}{\displaystyle\sum_{i=1}^{2} P(\eta_1|\xi_i)P(\xi_i)} = \frac{(1-P_e)P_1}{(1-P_e)P_1 + P_e P_2}
$$

mit

$$
P(\eta_1|\xi_1) = 1 - P_e, \quad P(\eta_1|\xi_2) = P_e, \quad P(\eta_2|\xi_1) = P_e, \quad P(\eta_2|\xi_2) = 1 - P_e
$$

Anhand dieses Beispiels interpretieren wir

\begin{itemize}
  \item $P(\xi_i)$ als \textit{a priori} Wahrscheinlichkeiten (Vorabwissen)
  \item $P(\eta_i|\xi_j)$ als Übergangswahrscheinlichkeiten
  \item $P(\xi_j|\eta_i)$ als \textit{a posteriori} Wahrscheinlichkeiten
\end{itemize}
\end{beispiel}

\begin{beispiel}[Internetwerbung]
Ein Verlag möchte mehr über seine Leser erfahren, um diese gezielt mit Werbung zu Büchern des Genres Belletristik (B), Ratgeber (R) oder Sachbuch (S) zu versorgen. So kann dann gezielt Werbung in die Browser potentieller Käufer eingeblendet werden.

Es soll in diesem Beispiel ermittelt werden, mit welcher Wahrscheinlichkeit eine der drei Personen $x_1$, $x_2$ oder $x_3$ ein Buch einer der drei Genres gekauft hat.

Wir bezeichnen die Wahrscheinlichkeit, dass Person $x_i$ ein Buch im Internet kauft mit $P(K_i)$ und die Wahrscheinlichkeit, dass kein Buch gekauft wird mit $P(\overline{K_i})$.

Aus alten Facebook-Daten ist nun bekannt, dass

\begin{itemize}
  \item die Person $x_1$ in der Schulzeit viele Bücher gelesen hat,
  \item die Person $x_2$ schon einmal in der Nähe eines Bücherregals gesehen wurde (Foto auf dem Facebook-Server),
  \item die Person $x_3$ gute Noten in den Schulfächern Deutsch und Englisch hatte.
\end{itemize}

Auf der Grundlage dieser Daten wird das generell zu vermutende Interesse am Kauf eines Buchs $K_i$ der Person $x_i$ als \textit{a priori} Information mit den Wahrscheinlichkeiten $P(K_i)$ modelliert:

$$
P(K_1) = 0{,}6 \qquad P(K_2) = 0{,}15 \qquad P(K_3) = 0{,}25
$$

Aus den Daten der Internetkäufe ist ferner bekannt, mit welchen relativen Häufigkeiten (als Wahrscheinlichkeiten zu modellieren) Bücher der Kategorie Belletristik (B), Ratgeber (R) oder Sachbuch (S) gekauft wurden, wenn von der Person $x_i$ ein Buch gekauft wurde. Dieses Wissen wird durch die bedingten Wahrscheinlichkeiten $P(\text{Genre} \mid \text{Person})$ modelliert:

\begin{center}
\begin{tabular}{lll}
  $P(B|K_1) = 0{,}4$ & $P(R|K_1) = 0{,}1$ & $P(S|K_1) = 0{,}5$ \\
  $P(B|K_2) = 0{,}1$ & $P(R|K_2) = 0{,}8$ & $P(S|K_2) = 0{,}1$ \\
  $P(B|K_3) = 0{,}5$ & $P(R|K_3) = 0{,}2$ & $P(S|K_3) = 0{,}3$ \\
\end{tabular}
\end{center}

In Kombination mit den \textit{a priori} Informationen, z.\,B.\ gilt $P(B, K) = P(B|K_1) \cdot P(K_1)$, ergeben sich nun die Verbundwahrscheinlichkeiten $P(\text{Genre}, \text{Person})$ wie folgt:

\begin{center}
\begin{tabular}{lll}
  $P(B, K_1) = 0{,}2400$ & $P(R, K_1) = 0{,}0600$ & $P(S, K_1) = 0{,}3000$ \\
  $P(B, K_2) = 0{,}0150$ & $P(R, K_2) = 0{,}1200$ & $P(S, K_2) = 0{,}0150$ \\
  $P(B, K_3) = 0{,}1250$ & $P(R, K_3) = 0{,}0500$ & $P(S, K_3) = 0{,}0750$ \\
\end{tabular}
\end{center}

Auf der Grundlage dieser Wahrscheinlichkeiten, ist es naheliegend, z.\,B.\ die Person $x_1$ vorrangig mit Buchwerbung zu den Genres Belletristik und Sachbuch zu versorgen.

Weiterhin lassen sich die Wahrscheinlichkeiten für den Kauf eines Buches eines Genres als Randwahrscheinlichkeit berechnen:

$$
P(B) = 0{,}3800 \qquad P(R) = 0{,}23 \qquad P(S) = 0{,}39
$$

Schließlich kann auf die Wahrscheinlichkeiten geschlossen werden, dass der Kauf eines Buches einer der drei Genres durch eine bestimmte Person erfolgte, also die \textit{a posteriori} Wahrscheinlichkeiten $P(\text{Person} \mid \text{Genre})$:

\begin{center}
\begin{tabular}{lll}
  $P(K_1|B) = 0{,}6316$ & $P(K_1|R) = 0{,}2609$ & $P(K_1|S) = 0{,}7692$ \\
  $P(K_2|B) = 0{,}0395$ & $P(K_2|R) = 0{,}5217$ & $P(K_2|S) = 0{,}0385$ \\
  $P(K_3|B) = 0{,}3289$ & $P(K_3|R) = 0{,}2174$ & $P(K_3|S) = 0{,}1923$ \\
\end{tabular}
\end{center}

Wurde z.\,B.\ ein Ratgeber verkauft, ist der Kauf durch Person $x_2$ sehr viel wahrscheinlicher als durch die anderen Personen.

Die grundlegende Idee des \textbf{maschinellen Lernens} beruht nun darauf, die so ermittelte \textit{a posteriori} Wahrscheinlichkeiten im nächsten Rechenschritt als \textit{a priori} Information zu verwenden. Damit werden die alten Vermutungen über das Verhalten von Personen (oder Systemen) durch aktualisierte, auf Beobachtungen beruhende Vermutungen ersetzt.

In unserem Beispiel ergeben sich mit den aktualisierten \textit{a priori} Wahrscheinlichkeiten als einer erneuten Berechnung der \textit{a posteriori} Wahrscheinlichkeiten:

\begin{center}
\begin{tabular}{lll}
  $P(K_1|B) = 0{,}6000$ & $P(K_1|R) = 0{,}0536$ & $P(K_1|S) = 0{,}8621$ \\
  $P(K_2|B) = 0{,}0094$ & $P(K_2|R) = 0{,}8571$ & $P(K_2|S) = 0{,}0086$ \\
  $P(K_3|B) = 0{,}3906$ & $P(K_3|R) = 0{,}0893$ & $P(K_3|S) = 0{,}1293$ \\
\end{tabular}
\end{center}

wobei hier jetzt die Zusammenhänge zwischen Genre und Person noch deutlicher sichtbar sind.
\end{beispiel}

\section{Zufallsvariablen, Verteilungen, Momente}

Zufallsvariablen bilden Ereignisse (also Mengen) auf Zahlen ab.

$$
X : \Omega \to \mathbb{R}
$$

Die Zufallsvariable (ZV) $X$ definiert somit eine Partition von $\Omega$. Jedem Element $A_i$ der Partition wird eine Zahl $x_i$ zugeordnet.

Für die Zufallsvariable $X$ kann dann auch eine Wahrscheinlichkeitsverteilung $P(A_i) = p_X(x_i)$ bestimmt werden. Zufallsvariablen können diskret oder kontinuierlich sein.

\subsection{Wahrscheinlichkeitsverteilung}

Wenn eine Menge von $N$ sich gegenseitig ausschließenden Ereignissen $A_i$ die Menge $\Omega$ vollständig abdeckt und $\sum_{i=1}^{N} P(A_i) = 1$ gilt, dann beschreiben die Zahlen $P(A_i) = p_X(x_i) > 0$ eine \textit{Wahrscheinlichkeitsverteilung}.

Die kumulative Wahrscheinlichkeitsverteilung wird mit

$$
F_X(x) = \sum_{x_i \leq x} p_X(x_i)
$$

angegeben. Es gilt:

$$
F_X(-\infty) = \lim_{x \to -\infty} F_X(x) = 0 \qquad F_X(\infty) = \lim_{x \to \infty} F_X(x) = 1
$$

$F_X$ ist monoton steigend.

\begin{beispiel}[Idealer Würfel]
Im Fall des idealen Würfels ergibt sich eine Gleichverteilung.

% Grafik: Stabdiagramm p_X(xᵢ) über Symbol xᵢ (1–6), alle Werte bei 1/6 ≈ 0{,}167
\end{beispiel}

\subsection{Relative Häufigkeiten}

Die im theoretischen Modell des idealen Würfels angenommenen Wahrscheinlichkeiten lassen sich im Experiment mit einem realen Würfel näherungsweise verifizieren.

\begin{beispiel}[Würfelwurf mit $M = 50$]
Wir werfen einen Würfel $M = 50$ mal und registrieren die Ergebnisse:

\begin{center}
\begin{tabular}{lcccccc}
  \hline
  Elementarereignis              & $x_1$ & $x_2$ & $x_3$ & $x_4$ & $x_5$ & $x_6$ \\
  Zahlensymbol                   & 1 & 2 & 3 & 4 & 5 & 6 \\
  beob.\ Häufigkeit $M_i$        & 9 & 8 & 10 & 7 & 8 & 8 \\
  relative Häufigkeit $M_i/M$    & $\frac{9}{50}$ & $\frac{8}{50}$ & $\frac{10}{50}$ & $\frac{7}{50}$ & $\frac{8}{50}$ & $\frac{8}{50}$ \\
  \hline
\end{tabular}
\end{center}
\end{beispiel}

Wenn das Experiment „unendlich" oft unter gleichen Bedingungen wiederholt wird, nähern sich die relativen Häufigkeiten einer festen Zahl, z.\,B.\ im Fall eines Würfels den Wahrscheinlichkeiten $\frac{1}{6}$.

\subsection{Häufigkeit und Histogramm}

Berechnung relativer Häufigkeiten:

\begin{itemize}
  \item Zahl aller Experimente: $M$
  \item Anzahl der Ereignisse $x_i$: $M_i$
  \item Relative Häufigkeit des Ereignisses $x_i$: $\dfrac{M_i}{M}$
\end{itemize}

% Grafik: Histogramm – links absolut (Häufigkeit Mᵢ über Symbol 1–6), rechts relativ (rel. Häufigkeit Mᵢ/M über Symbol 1–6)

\subsection{Das Gesetz der großen Zahlen}

\begin{bemerkung}
Aus dem obigen, einmal durchgeführten Experiment des fünfzigfachen Würfelns mit den dargestellten relativen Häufigkeiten darf man nicht den Schluss ziehen, dass der Würfel \textit{generell} die Zahl drei häufiger produziert als die anderen Zahlen. Die beobachteten Unterschiede sind der Zufälligkeit des Würfelvorgangs geschuldet. Es ist ein zentrales Anliegen der Wahrscheinlichkeitsrechnung, die Unsicherheiten, die sich in der Ausführung von Experimenten ergeben, zu quantifizieren, so dass gesicherte Schlussfolgerungen möglich werden. Das nachfolgende Gesetz der großen Zahlen betrachtet hierzu die wiederholte, sehr häufige Ausführung von Experimenten.
\end{bemerkung}

Ein Experiment werde $M$-mal unabhängig durchgeführt und es sei $\xi_i \in \{0,1\}$ die Indikatorvariable für ein Ereignis $A$ im $i$-ten Versuch. Es gilt $\xi_i = 1$, falls $A$ im $i$-ten Versuch eintritt und sonst $\xi_i = 0$. Für $\varepsilon > 0$ gilt dann das (schwache) Gesetz der großen Zahlen:

$$
\lim_{M \to \infty} P\!\left(\left\{\left|\frac{1}{M}\sum_{i=1}^{M} \xi_i - P(A)\right| \leq \varepsilon\right\}\right) = 1.
$$

\subsection{Erwartungswert}

Allgemein erhält man, wenn eine Zufallsvariable $X$ die Werte $x_1, x_2, \ldots$ annehmen kann, den statistischen Mittelwert $\mathrm{E}\{X\}$, der häufig auch als \textit{Erwartungswert} bezeichnet wird, aus

$$
\mu_X = \mathrm{E}\{X\} = \sum_i x_i\, p_X(x_i)
$$

Weiterhin erhält man den quadratischen Mittelwert

$$
\mathrm{E}\{X^2\} = \sum_i x_i^2\, p_X(x_i)
$$

und die Varianz

$$
\mathrm{var}\{X\} = \mathrm{E}\{(X - \mu_X)^2\} = \sum_i (x_i - \mu_X)^2\, p_X(x_i).
$$

\begin{beispiel}[Erwartungswert: Bitgruppen]
Eine Quelle sendet Bitgruppen $A_i$ unterschiedlicher Länge $l_i \in \{1,2\}$ mit den Wahrscheinlichkeiten $P(A_i)$ aus.

\begin{center}
\begin{tabular}{c|c|c}
  \hline
  Ereignis & gesendete Bitgruppe & Wahrscheinlichkeit $P(A_i)$ \\
  \hline
  1 & 0  & 0{,}1 \\
  2 & 1  & 0{,}1 \\
  3 & 00 & 0{,}4 \\
  4 & 11 & 0{,}4 \\
  \hline
\end{tabular}
\end{center}

Wie viele Bits werden pro Bitgruppe im Mittel gesendet?
\end{beispiel}

\subsection{Binomialverteilung}

Es wird ein Experiment mit zwei möglichen Ereignissen $A_1$ und $A_2$ und den zugehörigen Wahrscheinlichkeiten $P(A_1)$ und $P(A_2) = 1 - P(A_1)$ betrachtet. Dieses Experiment wird $N$-fach wiederholt, wobei die Versuche als statistisch unabhängig anzusehen sind. Die Wahrscheinlichkeit, dass das Ereignis $A_1$ dabei $k$ mal auftritt, beträgt

$$
P(k \text{ mal } A_1) = \binom{N}{k} P(A_1)^k \cdot (1 - P(A_1))^{N-k}
$$

Wir definieren somit eine Zufallsvariable $K$ mit der Wahrscheinlichkeitsverteilung $p_K(k) = P(k)$ und $k \in \{1, \ldots, N\}$.

% Grafik: Stabdiagramme für N=10, P(A₁)=0{,}5 (links) und N=10, P(A₁)=0{,}1 (rechts)

\subsection{Poisson-Verteilung}

Die Poisson-Verteilung beschreibt die Zahl $K$ der Ereignisse innerhalb eines Beobachtungsintervalls, wenn diese Ereignisse unabhängig auftreten und die mittlere Auftrittsrate $\lambda$ bekannt ist. Es gilt:

$$
p_X(k) = e^{-\lambda} \frac{\lambda^k}{k!}
$$

$$
\mathrm{E}\{X\} = \sum_{k=0}^{\infty} k \frac{\lambda^k}{k!} e^{-\lambda} = \lambda
$$

$$
\mathrm{E}\{X^2\} = \sum_{k=0}^{\infty} k^2 \frac{\lambda^k}{k!} e^{-\lambda} = \lambda^2 + \lambda
$$

$$
\mathrm{var}\{X\} = \mathrm{E}\{X^2\} - \mathrm{E}\{X\}^2 = \lambda
$$

% Grafik: Stabdiagramme für λ=1 (links) und λ=10 (rechts), jeweils p_K(k) und F_K(k) über k von 0 bis 20

\subsection{Übergang zu Verteilungsdichten}

\begin{itemize}
  \item Bisher wurde einem diskreten Wert $x_i$, bzw.\ der einelementigen Menge $\{A_i\}$ eine positive Wahrscheinlichkeit zugeordnet.
  \item Für kontinuierliche Zufallsvariablen besitzen nur Intervalle positive Wahrscheinlichkeiten, der einzelne Punkt hat die Wahrscheinlichkeit 0.
  \item Die Wahrscheinlichkeit, dass die Variable $X$ im Intervall $[a,b]$ liegt, lässt sich dann über den Flächeninhalt einer Dichtefunktion angeben:
\end{itemize}

$$
P(X \in (a,b)) = \int_a^b p_X(u)\,\mathrm{d}u,
$$

wobei $p_X(u)$ die Wahrscheinlichkeitsdichte bezeichnet.

\subsection{Eigenschaften der Wahrscheinlichkeitsdichte}

Die WS-Dichte $p_X(u)$ soll den folgenden Eigenschaften genügen:

\begin{itemize}
  \item $p_X(u) \geq 0$ für alle $u \in \mathbb{R}$
  \item $p_X(u)$ ist stückweise stetig
  \item $\displaystyle\int_{-\infty}^{\infty} p_X(u)\,\mathrm{d}u = 1$
\end{itemize}

Durch diese Festlegungen wird ein Wahrscheinlichkeitsmaß definiert.

\subsection{Erwartungswerte kontinuierlicher ZV}

Der Erwartungswert einer kontinuierlichen Zufallsvariablen $X$ ist mit

$$
\mathrm{E}\{X\} = \int_{-\infty}^{\infty} x\, p_X(x)\,\mathrm{d}x
$$

gegeben, wobei $\displaystyle\int_{-\infty}^{\infty} |x|\, p_X(x)\,\mathrm{d}x < \infty$ gelten muss.

Für eine stückweise stetige Funktion $g\colon \mathbb{R} \to \mathbb{R}$ einer Zufallsvariablen $X$ mit $\displaystyle\int_{-\infty}^{\infty} |g(x)|\, p_X(x)\,\mathrm{d}x < \infty$ gilt:

$$
\mathrm{E}\{g(X)\} = \int_{-\infty}^{\infty} g(x)\, p_X(x)\,\mathrm{d}x
$$

\subsection{Momente einer Wahrscheinlichkeitsdichte}

\begin{itemize}
  \item Mittelwert: $\mathrm{E}\{X\} = \displaystyle\int_{-\infty}^{\infty} x\, p_X(x)\,\mathrm{d}x$ \hfill (erstes Moment)
  \item Leistung: $\mathrm{E}\{X^2\} = \displaystyle\int_{-\infty}^{\infty} x^2\, p_X(x)\,\mathrm{d}x$ \hfill (zweites Moment)
\end{itemize}

Die Leistung der um den Mittelpunkt bereinigten Zufallsvariablen wird Varianz genannt.

$$
\mathrm{var}\{X\} = \mathrm{E}\{(X - \mathrm{E}\{X\})^2\} = \int_{-\infty}^{\infty} (x - \mathrm{E}\{X\})^2\, p_X(x)\,\mathrm{d}x = \sigma_X^2
$$

Es gilt $\mathrm{E}\{(X - \mathrm{E}\{X\})^2\} = \mathrm{E}\{X^2\} - \mathrm{E}\{X\}^2$.

Standardabweichung: $\sigma_X = \mathrm{std}\{X\} = \sqrt{\mathrm{var}\{X\}}$

\subsection{Die Verteilungsfunktion}

Die kumulative Verteilungsfunktion $F_X(x)$ einer Zufallsvariablen ist eine Funktion, die durch die Bedingung

$$
F_X(x) = P(X \leq x)
$$

definiert ist. Sie gibt an, mit welcher Wahrscheinlichkeit die ZV $X$ kleiner oder gleich einem Wert $x$ ist. $F_X(x)$ ist eine monoton steigende Funktion mit $0 \leq F_X(x) \leq 1$. Offenbar lässt sich diese Funktion aus der Verteilungsdichte $p_X(x)$ durch Integration berechnen, so dass wir $F_X(x)$ und $p_X(x)$ wie folgt schreiben können:

$$
\int_{-\infty}^{x} p_X(u)\,\mathrm{d}u = F_X(x) \quad \Leftrightarrow \quad p_X(x) = \frac{\mathrm{d}}{\mathrm{d}x} F_X(x).
$$

\subsection{Die Gleichverteilung auf dem Intervall $[a,b]$}

Für Zahlen $a, b \in \mathbb{R}$ und $b > a$ ist die Gleichverteilung einer ZV $X$ durch

$$
p_X(x) = \begin{cases} \dfrac{1}{b - a} & , \quad x \in [a, b] \\ 0 & , \quad \text{sonst} \end{cases}
$$

definiert.

% Grafik: Rechteckige WS-Dichte auf [a,b] mit Höhe 1/(b−a)

% Grafik: p_X(x) (blau) und F_X(x) (grün gestrichelt) für a=−1, b=1 (links) und a=−1, b=2 (rechts), x von −2 bis 2

Für eine gleichverteilte Zufallsvariable $X$ gilt:

$$
\mathrm{E}\{X\} = \frac{a + b}{2} \qquad \mathrm{E}\{X^2\} = \frac{a^2 + ab + b^2}{3} \qquad \mathrm{var}\{X\} = \frac{(a-b)^2}{12}
$$

\subsection{Die Normalverteilung}

Die Normalverteilung zu den Parametern $\mu \in \mathbb{R}$ und $\sigma > 0$ ist durch die Dichtefunktion

$$
p_X(x) = \mathcal{N}(x;\, \mu,\, \sigma^2) = \frac{1}{\sqrt{2\pi\sigma^2}}\, e^{-\frac{(x-\mu)^2}{2\sigma^2}}
$$

definiert. Ferner gilt für eine normalverteilte ZV $X$:

\begin{itemize}
  \item Mittelwert: $\mathrm{E}\{X\} = \mu$
  \item Varianz: $\mathrm{var}\{X\} = \sigma^2$
  \item Leistung: $P = \sigma^2 + \mu^2$
\end{itemize}

% Grafik: p_X(x) (blau) und F_X(x) (grün gestrichelt) für μ=0, σ=1 (links) und μ=1, σ=0{,}5 (rechts), x von −2 bis 2

Die Verteilungsfunktion $F_X(x)$ einer normalverteilten Zufallsvariablen $X$ mit der Dichtefunktion $p_X(x) = \dfrac{1}{\sqrt{2\pi}} e^{-\frac{x^2}{2}}$ wird häufig mit Hilfe der Gaußschen Fehlerfunktion

$$
\mathrm{erf}(x) = \frac{2}{\sqrt{\pi}} \int_0^x e^{-t^2}\,\mathrm{d}t = \sqrt{\frac{2}{\pi}} \int_0^{x\sqrt{2}} e^{-\frac{t^2}{2}}\,\mathrm{d}t
$$

oder der komplementären Fehlerfunktion

$$
\mathrm{erfc}(x) = \frac{2}{\sqrt{\pi}} \int_x^{\infty} e^{-t^2}\,\mathrm{d}t = 1 - \mathrm{erf}(x)
$$

angegeben. Somit gilt mit $\mathrm{erf}(-x) = -\mathrm{erf}(x)$

$$
F_X(x) = \frac{1}{\sqrt{2\pi}} \int_{-\infty}^{x} e^{-\frac{t^2}{2}}\,\mathrm{d}t = \frac{1}{2}\!\left(1 + \mathrm{erf}\!\left(\frac{x}{\sqrt{2}}\right)\right) = \frac{1}{2}\,\mathrm{erfc}\!\left(\frac{-x}{\sqrt{2}}\right).
$$

\subsection{Die Exponentialverteilung}

Die Exponentialverteilung zum Parameter $\lambda > 0$ ist zu

$$
p_X(x) = \begin{cases} \lambda e^{-\lambda x} & , \quad x \geq 0 \\ 0 & , \quad \text{sonst} \end{cases}
$$

definiert. Ferner gilt für eine exponentialverteilte ZV $X$:

$$
\mathrm{E}\{X\} = \frac{1}{\lambda} \qquad \mathrm{E}\{X^2\} = \frac{2}{\lambda^2} \qquad \mathrm{var}\{X\} = \frac{1}{\lambda^2}
$$

% Grafik: p_X(x) (blau) und F_X(x) (grün gestrichelt) für λ=1 (links) und λ=1{,}5 (rechts), x von −2 bis 2

\subsection{Die Laplaceverteilung}

Die zweiseitige Laplaceverteilung zu den Parametern $\mu$ und $\sigma$ ist zu

$$
p_X(x) = \frac{1}{2\sigma}\, e^{-\frac{|x-\mu|}{\sigma}}
$$

definiert. Ferner gilt für eine laplaceverteilte ZV $X$:

$$
\mathrm{E}\{X\} = \mu \qquad \mathrm{E}\{X^2\} = 2\sigma^2 + \mu^2 \qquad \mathrm{var}\{X\} = 2\sigma^2
$$

% Grafik: p_X(x) (blau) und F_X(x) (grün gestrichelt) für μ=0, σ=1 (links) und μ=1, σ=0{,}5 (rechts), x von −2 bis 2

\begin{beispiel}[Quantisierung eines Audiosignals]
Die Abtastwerte eines Audiosignals seien Laplace-verteilt mit $\mu = 0$ und $\sigma = 1$. Die Abtastwerte werden nun mit der folgenden Quantisierungskennlinie quantisiert.

% Grafik: Midrise-Quantisierungskennlinie x_q über x, x und x_q von −1 bis 1

Wie groß ist die Wahrscheinlichkeit, dass der Quantisierer übersteuert ($|x| > 1$) wird?
\end{beispiel}

\subsection{Transformation einer Zufallsvariablen}

% Grafik: Transformationsprinzip – Funktion y = g(x) (nichtlinear, ansteigend), darunter p_X(x) (Gleichverteilung, 4 Blöcke à 25\%), links p_Y(y) (transformierte Verteilung)

% Grafik: Infinitesimale Betrachtung – y = g(x), p_X(x)|dx| = p_Y(y)|dy|; p_Y(y)|dy| und p_X(x)|dx| als schattierte Streifen eingezeichnet

\subsection{Streng monotone Funktion einer Zufallsvariablen}

Falls $Y = g(X)$ gilt und $g(\cdot)$ eine streng monotone Funktion ist, so dass die Umkehrfunktion $X = g^{-1}(Y)$ existiert, dann gilt für die Lösung $x = g^{-1}(y)$

$$
p_Y(y) = p_X(x)\frac{|\mathrm{d}x|}{|\mathrm{d}y|} = p_X(g^{-1}(y))\left|\frac{\mathrm{d}g^{-1}(y)}{\mathrm{d}y}\right| = \frac{p_X(g^{-1}(y))}{\left|\dfrac{\mathrm{d}g(x)}{\mathrm{d}x}\right|} = \frac{p_X(g^{-1}(y))}{|g'(g^{-1}(y))|}
$$

wobei für das Differenzieren der Umkehrfunktion $\dfrac{\mathrm{d}g^{-1}(y)}{\mathrm{d}y} = \left(\dfrac{\mathrm{d}g(x)}{\mathrm{d}x}\right)^{-1}$ gilt.

\begin{beispiel}[Mehrfaches Werfen einer Münze]
Wie oft muss eine faire Münze im Mittel geworfen werden, bis zum ersten Mal „Kopf" erscheint?

Lösungsansatz: Es wird der Erwartungswert

$$
\mathrm{E}\{I\} = \sum_{i=1}^{\infty} i\, 2^{-i}
$$

einer Zufallsvariable $I$, die die Werte $i \in \mathbb{N}$ annehmen kann, gesucht.
\end{beispiel}

\subsection{Transformation für nichtmonotone Funktionen}

Gegeben sei eine Funktion $Y = g(X)$ einer Zufallsvariablen $X$. Wenn die Funktion $g(\cdot)$ nicht streng monoton ist, hat die Gleichung $Y = g(x)$ mehrfache Lösungen $x_1 = g^{-1}(y), \ldots, x_n = g^{-1}(y)$. Dann berechnet sich die Wahrscheinlichkeitsdichtefunktion $p_Y(y)$ zu

$$
p_Y(y) = \frac{p_X(x_1)}{|g'(x_1)|} + \ldots + \frac{p_X(x_n)}{|g'(x_n)|}
$$

wobei die $x_i$ die Lösungen der Gleichung $y = g(x_i)$ und $g'(x)$ die Ableitung der Funktion $g(x)$ bezeichnen.

\begin{beispiel}[Beispiele zur Transformation von Zufallsvariablen]
\textbf{a)} $Y = aX + b = g(X)$

Diese Gleichung hat genau eine Lösung

$$
x_1 = \frac{y - b}{a}.
$$

Ferner gilt $g'(x) = a$.

Also gilt für die WS-Dichte $p_Y(y)$:

$$
p_Y(y) = \frac{p_X\!\left(\dfrac{y-b}{a}\right)}{|a|}
$$

\textbf{b)} $Y = X^2 = g(X)$

Diese Gleichung hat für $y \geq 0$ zwei Lösungen

$$
x_1 = +\sqrt{y} \quad \text{und} \quad x_2 = -\sqrt{y}.
$$

Für $y < 0$ existiert keine reellwertige Lösung.

Ferner gilt $g'(x) = 2x$ und somit $|g'(x_1)| = |g'(x_2)| = 2\sqrt{y}$.

Also gilt für die WS-Dichte $p_Y(y)$:

$$
p_Y(y) = \begin{cases}
\dfrac{p_X(\sqrt{y})}{2\sqrt{y}} + \dfrac{p_X(-\sqrt{y})}{2\sqrt{y}} & , \quad y \geq 0 \\[6pt]
0 & , \quad y < 0
\end{cases}
= \left(\frac{p_X(\sqrt{y})}{2\sqrt{y}} + \frac{p_X(-\sqrt{y})}{2\sqrt{y}}\right) \cdot \varepsilon(y)
$$
\end{beispiel}

\begin{beispiel}[Beispiel: Normalverteilung]
$$
p_X(x) = \frac{1}{\sqrt{2\pi\sigma^2}}\,e^{-\frac{(x-\mu)^2}{2\sigma^2}} \quad \text{mit } \mu = 0
$$

$Y = X^2$ mit den Lösungen $x_1 = \sqrt{y}$ und $x_2 = -\sqrt{y}$

$$
p_Y(y) = \frac{1}{2\sqrt{y}} \cdot \frac{1}{\sqrt{2\pi\sigma^2}} \left(e^{-\frac{(\sqrt{y})^2}{2\sigma^2}} + e^{-\frac{(-\sqrt{y})^2}{2\sigma^2}}\right) \cdot \varepsilon(y)
= \frac{1}{\sqrt{2\pi\sigma^2 y}}\,e^{-\frac{y}{2\sigma^2}} \cdot \varepsilon(y)
$$
\end{beispiel}


\section{Zwei Zufallsvariablen}

Wenn zwei Zufallsvariablen $X$ und $Y$ gegeben sind, ist die Verbundverteilungsfunktion durch

$$
F_{XY}(x,y) = P(X \leq x, Y \leq y)
$$

definiert. Es gilt

$$
F_{XY}(-\infty, y) = 0, \quad F_{XY}(x, -\infty) = 0 \quad \text{und} \quad F_{XY}(\infty, \infty) = 1.
$$

Die Verbundverteilungsdichtefunktion ist dann durch

$$
p_{XY}(x,y) = \frac{\partial^2 F_{XY}(x,y)}{\partial x \, \partial y}
$$

gegeben. Somit gilt auch

$$
F_{XY}(x,y) = \int_{-\infty}^{x} \int_{-\infty}^{y} p_{XY}(\xi, \mu) \, \mathrm{d}\xi \, \mathrm{d}\mu.
$$

\subsection{Statistische Unabhängigkeit}

Für zwei statistisch unabhängige Zufallsvariablen $X$ und $Y$ gilt:

$$
F_{XY}(x,y) = F_X(x) \cdot F_Y(y)
$$

und

$$
p_{XY}(x,y) = p_X(x) \cdot p_Y(y).
$$

\subsection{Randverteilungen}

Die univariate Verteilungsdichtefunktion der Zufallsvariablen $X$ und $Y$ lässt sich mit

$$
p_X(x) = \int_{-\infty}^{\infty} p_{XY}(x,y)\,\mathrm{d}y \quad \text{und} \quad p_Y(y) = \int_{-\infty}^{\infty} p_{XY}(x,y)\,\mathrm{d}x
$$

bestimmen. Diese werden Randverteilungen genannt. Für die Verteilungsfunktion ergibt sich

$$
F_X(x) = F_{XY}(x, \infty) \quad \text{und} \quad F_Y(y) = F_{XY}(\infty, y).
$$

\begin{beispiel}[Bivariate Normalverteilung]
Zwei Zufallsvariablen $X$ und $Y$ sind gemeinsam normalverteilt, wenn ihre Verbundverteilungsdichtefunktion durch

$$
p_{XY}(x,y) = \frac{1}{2\pi\sigma_X\sigma_Y\sqrt{1-r_{XY}^2}} \exp\!\left(-\frac{1}{2(1-r_{XY}^2)}\left[\frac{(x-\mu_X)^2}{\sigma_X^2} - 2r_{XY}\frac{(x-\mu_X)(y-\mu_Y)}{\sigma_X\sigma_Y} + \frac{(y-\mu_Y)^2}{\sigma_Y^2}\right]\right)
$$

gegeben ist. Hierbei sind $\mu_X$ und $\mu_Y$ die Mittelwerte, $\sigma_X^2$ und $\sigma_Y^2$ die Varianz und $r_{XY}$ der Korrelationskoeffizient

$$
r_{XY} = \frac{\mathrm{E}\{(X-\mu_X)(Y-\mu_Y)\}}{\sigma_X\sigma_Y}
$$

der Zufallsvariablen $X$ und $Y$. Es gilt $|r_{XY}| \leq 1$.

% Grafik: 3D-Oberfläche und Konturplot von p_XY(x,y) für μ_X=0, μ_Y=0, σ_X²=1, σ_Y²=1, r_XY=0

% Grafik: 3D-Oberfläche und Konturplot von p_XY(x,y) für μ_X=1, μ_Y=1, σ_X²=1, σ_Y²=1, r_XY=0

% Grafik: 3D-Oberfläche und Konturplot von p_XY(x,y) für μ_X=0, μ_Y=0, σ_X²=1, σ_Y²=0.5, r_XY=0

% Grafik: 3D-Oberfläche und Konturplot von p_XY(x,y) für μ_X=0, μ_Y=0, σ_X²=1, σ_Y²=1, r_XY=0.5
\end{beispiel}

\subsection{Momente der bivariaten Verteilungsdichte}

Die Momente zweier Zufallsvariablen $X$ und $Y$ mit der VDF $p_{XY}(x,y)$ lassen sich mit dem folgenden Theorem ermitteln:

$$
\mathrm{E}\{g(X,Y)\} = \int_{-\infty}^{\infty} \int_{-\infty}^{\infty} g(x,y)\,p_{XY}(x,y)\,\mathrm{d}x\,\mathrm{d}y.
$$

Hierbei ist $g(X,Y)$ eine beliebige Funktion der Zufallsvariablen $X$ und $Y$.

Speziell gilt im Fall $Z = g(X,Y) = aX + bY$

$$
\mathrm{E}\{aX + bY\} = a\,\mathrm{E}\{X\} + b\,\mathrm{E}\{Y\}.
$$

Die Erwartungswertbildung ist eine lineare Operation.

\subsection{Linearität des Erwartungswertes}

Für zwei Zufallsvariablen $X$ und $Y$ gilt:

$$
\mathrm{E}\{aX + bY\} = a\,\mathrm{E}\{X\} + b\,\mathrm{E}\{Y\} \quad \text{für alle } a, b \in \mathbb{R}
$$

\subsection{Kovarianz, Korrelation und Orthogonalität}

Die Kovarianz von zwei Zufallsvariablen $X$ und $Y$ ist zu

$$
C = \mathrm{E}\{(X - \mu_X)(Y - \mu_Y)\}
$$

definiert. Es gilt $C = \mathrm{E}\{XY\} - \mu_X\mu_Y$.

Der Korrelationskoeffizient ist zu

$$
r_{XY} = \frac{\mathrm{E}\{(X - \mu_X)(Y - \mu_Y)\}}{\sigma_X\sigma_Y}
$$

definiert. Es gilt $|r_{XY}| \leq 1$.

Zwei Zufallsvariablen werden orthogonal genannt, wenn

$$
\mathrm{E}\{XY\} = \int_{-\infty}^{\infty} \int_{-\infty}^{\infty} x\,y\,p_{XY}(x,y)\,\mathrm{d}x\,\mathrm{d}y = 0
$$

gilt.

\subsection{Eine Funktion von zwei Zufallsvariablen}

Es sei eine Funktion von zwei Zufallsvariablen $X$ und $Y$ gegeben:

$$
Z = g(X, Y)
$$

Die Verteilungsfunktion $F_Z(z)$ ist dann durch

$$
F_Z(z) = P(Z \leq z) = P(g(X,Y) \leq z) = \iint_{D_z} p_{XY}(x,y)\,\mathrm{d}x\,\mathrm{d}y
$$

gegeben, wobei $D_z$ das nicht unbedingt zusammenhängende Gebiet in der $xy$-Ebene bezeichnet, für das $g(x, y) \leq z$ gilt.

\begin{beispiel}
Für $Z = X + Y$ gilt:

% Grafik: Gebiet z ≤ x+y in der xy-Ebene (Halbebene unterhalb der Linie y = z-x)

$$
F_Z(z) = \int_{-\infty}^{\infty} \int_{-\infty}^{z-x} p_{XY}(x,y)\,\mathrm{d}y\,\mathrm{d}x
$$

$$
p_Z(z) = \int_{-\infty}^{\infty} p_{XY}(x, z-x)\,\mathrm{d}x
$$

und im Fall unabhängiger Zufallsvariablen $X$ und $Y$

$$
p_Z(z) = \int_{-\infty}^{\infty} p_X(x)\,p_Y(z-x)\,\mathrm{d}x
$$
\end{beispiel}

\subsection{Zwei Funktionen von zwei Zufallsvariablen}

Gegeben sind zwei Zufallsvariablen $X$ und $Y$ und zwei Funktionen

$$
Z = g(X, Y) \quad \text{und} \quad W = h(X, Y).
$$

Die Verbundverteilungsdichtefunktion der Zufallsvariablen $Z$ und $W$ ist dann durch

$$
p_{ZW}(z,w) = \frac{p_{XY}(x_1, y_1)}{|\det(J(x_1, y_1))|} + \ldots + \frac{p_{XY}(x_n, y_n)}{|\det(J(x_n, y_n))|}
$$

gegeben, wobei die $x_i$ die Lösungen des Gleichungssystems $g(x_i, y_i) = z$, $h(x_i, y_i) = w$ sind und $J$ die Jakobi-Matrix

$$
J(x,y) = \begin{pmatrix} \dfrac{\partial g}{\partial x} & \dfrac{\partial g}{\partial y} \\[6pt] \dfrac{\partial h}{\partial x} & \dfrac{\partial h}{\partial y} \end{pmatrix}
$$

bezeichnet. Hierbei ist $\det(J(x_i, y_i))$ die Determinante dieser Matrix.

\subsection{Statistische Mittelwerte und Zeitmittelwerte}

% Grafik: Drei diskrete Realisierungen X_1[k], X_2[k], X_3[k] eines Zufallsprozesses als Impulsfolgen

\subsection{Ausblick: Zufallsprozesse}

Zufällige diskrete Signale lassen sich nun als Folge von Zufallsvariablen $X[k]$ definieren. Dabei kann für jede Zufallsvariable eine Verteilungsdichtefunktion und entsprechende statistische Mittelwerte und Momente angegeben werden. Zudem wird die statistische Abhängigkeit der Zufallsvariablen mit Verbundverteilungsdichtefunktionen beschrieben. Für die Modellbildung sind besonders

\begin{itemize}
  \item stationäre Zufallsprozesse (alle $X[k]$ weisen die selbe Verteilungsdichtefunktion auf) und
  \item ergodische Zufallsprozesse (die statistischen Mittelwerte entsprechen den durch Mittelung über eine Folge $X[k]$ gewonnenen empirischen Mittelwerten).
\end{itemize}



\chapter{Grundbegriffe der Informationstheorie}
\label{chap:informationstheorie}

Als Anwendung der Wahrscheinlichkeitsrechnung (Kapitel~\ref{chap:wahrscheinlichkeit}) betrachten wir in diesem Kapitel grundlegende Konzepte der Informationstheorie.

\section{Fragestellungen der Informationstheorie}
\label{sec:informationstheorie-fragen}

Die Informationstheorie ist die Grundlage für die Verarbeitung, Übertragung und Speicherung von Text-, Audio- und Videoinformationen. Sie ist ein Kerngebiet der Informations- und Kommunikationstechnik und stützt sich auf eine bahnbrechende Arbeit von C.E. Shannon (1948)\footnote{C.E. Shannon: A Mathematical Theory of Communication, Bell Systems Technical Journal, Vol. 27, S. 379-423, Juli 1948}. U.a. gibt sie Antworten auf die folgenden Fragen:

\begin{itemize}
  \item Wie kann man eine Informationsmenge messen?
  \item Wie werden Signale optimal übertragen?
  \item Wie kann man die Kapazität eines Übertragungskanals berechnen?
  \item Welche Informationsrate ist mindestens erforderlich um ein Signal verzerrungsfrei zu übertragen?
  \item Welchen Einfluss hat das Rauschen und wie lässt er sich reduzieren?
\end{itemize}

Die mathematische Formulierung dieser Frage und ihrer Antworten hat der Informationstheorie breite Anwendungsgebiete eröffnet.

In der Informationstechnik werden Informationen durch Zeichen (oder Symbole) dargestellt. Diese Symbole sind in der Regel einer endlichen Menge zulässiger Zeichen, dem \textit{Alphabet}, entnommen.

Beispiele für die Kommunikation mittels diskreter Symbole (Zeichen) sind

\begin{itemize}
  \item in der Nachrichtentechnik: Bits oder Bitgruppen
  \item im Strassenverkehr: Schilder (visuelle Symbole)
  \item in der sprachlichen Kommunikation: Phoneme (Lautzeichen)
\end{itemize}

Unter einem Code verstehen wir die Zuordnung von Symbolen eines ersten Alphabets zu den Symbolen eines zweiten Alphabets. In den dargestellten Anwendungen geschieht dies zum Zweck einer effizienten und/oder fehlergesicherten Übertragung. Der Aspekt der Verschlüsselung der Informationen spielt hier zunächst keine Rolle. Der Schutz der Informationen gegen z. B. unberechtigtes Lesen wird in der Kryptologie behandelt.

Weiterhin spielen die folgenden Begriffe der Informationstheorie eine wichtige Rolle:

\begin{itemize}
  \item \textbf{Redundante Daten:}

    Daten, die im Prinzip vollständig aus bereits vorhandenen Daten berechnet werden können.

  \item \textbf{Irrelevante Daten:}

    Daten, die zur Verarbeitung oder Interpretation der vorhandenen Daten nichts beitragen.

  \item \textbf{Wechselseitige Informationen:}

    Informationen, die sowohl in Datenmenge $X$ als auch in Datenmenge $Y$ enthalten sind.
\end{itemize}

Unter \textit{Quellencodierung} verstehen wir dementsprechend die (effiziente) Darstellung der Zeichen einer Quelle durch einen Code. Häufig handelt es sich dabei um einen Binärcode, der mit den Symbolen „0" und „1" notiert werden kann.

Zum Beispiel kann mit dem MP3-Verfahren ein digitalisiertes Musiksignal „komprimiert", d.h. in eine Folge von Codesymbolen umgerechnet werden, die weniger Speicherplatz benötigt, als die ursprüngliche digitale Darstellung. Dies wird durch das Entfernen redundanter und irrelevanter Daten erreicht.

Ähnliches gilt für die Übertragung von Audio- oder Videodaten über DSL-Verbindungen oder einen Mobilfunkkanal. Auch hier werden die Quellensignale komprimiert. Allerdings muss das quellencodierte Signal zusätzlich auch gegen Übertragungsfehler gesichert werden. Dies erfolgt durch das Hinzufügen redundanter Bits. Man spricht dabei von \textit{Kanalcodierung}.

Abbildung 5.1 erläutert das der Informationsübertragung zugrunde liegende Modell: Die von der Informationsquelle erzeugten Nachrichten werden in einen Sender einem physikalischen Signal aufgeprägt. Dieses Signal wird bei der Übertragung durch Störungen beeinflusst und im Empfänger in Nachrichtensymbole umgesetzt. Nur wenn eine wirkungsvolle Kanalcodierung eingesetzt wird, sind die empfangenen Symbole gleich den gesendeten Symbolen. Die Informationssenke interpretiert schließlich die empfangenen Informationen.

% Grafik: Abbildung 5.1: Blockdiagramm: Informationsquelle → Sender → Übertragungskanal (Störungen/Rauschen von oben) → Empfänger → Informationssenke

In Abhängigkeit von der Fragestellung, sind unterschiedliche Aspekte der Kommunikation von Bedeutung:

\begin{itemize}
  \item[\textbf{A.}] Mit welchen Fehlern ist die Übertragung von Symbolen (Zeichen) behaftet (das \textit{technische} Kommunikationsproblem)?
  \item[\textbf{B.}] Welche Bedeutung (Inhalt) wird von den Symbolen (Zeichen) vermittelt (das \textit{semantische} Kommunikationsproblem)?
  \item[\textbf{C.}] Welche Art von Handlungen werden durch die empfangenen Symbole (und ihrer Bedeutung) ausgelöst?
\end{itemize}

Wir stellen fest: Information (im technischen Sinn) muss von Bedeutung unterschieden werden!

\begin{beispiel}[Fernbedienung eines Fernsehers]
Die drei Ebenen der Kommunikation sollen anhand der Infrarotfernbedienung eines Fernsehers verdeutlicht werden.

\textbf{Ebene A:} Nach dem Drücken eines Programmwahlknopfes wird ein optisches Signal (im infraroten Spektrum) zum Fernseher übertragen. Diese Übertragung kann durch ungünstige Bedingungen (z. B. keine Sichtverbindung zwischen Sender und Empfänger) gestört sein.

\textbf{Ebene B:} Den verschiedenen Knöpfen der Fernbedienung sind verschiedene Fernsehprogramme zugeordnet.

\textbf{Ebene C:} Bei korrekter Bedienung der Fernbedienung, korrektem Empfang und korrekter Decodierung des Infrarotsignals schaltet der Fernseher auf einen anderen Kanal um.
\end{beispiel}

Die klassische Informationstheorie behandelt in erster Linie die Frage der Ebene A. Allerdings rücken mit der Entwicklung „intelligenter" Systeme zunehmend auch die Ebenen B und C in den Fokus.

\subsection{Definition des Informationsmaßes}
\label{sec:informationsmass}

\begin{itemize}
  \item Information im technischen Sinne beschreibt die \textbf{Wahlmöglichkeiten} der Informationsquelle bei der Auswahl eines möglichen Nachrichtensymbols.

  \item Aufgrund der Wahlmöglichkeiten des Senders entsteht im Empfänger eine \textbf{Unsicherheit} bezüglich des möglicherweise gesendeten Symbols. Durch den Empfang des vom Sender ausgewählten Symbols wird diese Unsicherheit reduziert.

  \item Im einfachsten Fall besteht die Auswahl einer Nachrichtenquelle aus \textbf{zwei} möglichen Symbolen (z. B. zwei Quantisierungsstufen, „0"/„1" oder „Ja" / „Nein"). Wenn beide Symbole mit gleicher Wahrscheinlichkeit auftreten (also größte Unsicherheit bezüglich der Auswahl einer Nachricht herrscht), ist die Informationsmenge, die mit der Auswahl einer dieser Symbole verbunden ist, genau \textbf{1 bit}.
\end{itemize}

Wichtige Unterscheidung:

1 Bit $\mathrel{\widehat{=}}$ Binary Digit $\rightarrow$ beschreibt die Ziffer einer binären Zahl

1 bit $\rightarrow$ ist die SI-Einheit der Informationsmenge

Der Fall einer größeren Auswahl wird zunächst an dem Beispiel in Abbildung 5.2 erläutert: Wie viele binäre Entscheidungen muss man treffen, um ein ausgewähltes Symbol $m$ aus einer Menge $\{1,2,3,4,5,6,7,8\}$ von $N = 8$ gleichwahrscheinlichen Symbolen zu ermitteln? Wir lösen diese Aufgabe mit dem in Abbildung 5.2 dargestellten Entscheidungsbaum.

% Grafik: Abbildung 5.2: Binärer Entscheidungsbaum für N=8 gleichwahrscheinliche Symbole; 3 Ebenen, Blätter m=1,...,8; an jeder Verzweigung ja/nein-Frage

In obigem Beispiel sind $3 = \log_2(8)$ ja/nein-Fragen mit den entsprechenden binären Entscheidungen erforderlich! Wenn man die Antwort auf jede dieser Fragen mit einem Bit notiert wird, entsteht ein einfacher \textit{binärer Code} zur Darstellung der ausgewählten Symbole.

\begin{definition}[Entscheidungsgehalt]
Allgemein wird die Informationsmenge bei Auswahl aus $N$ gleichwahrscheinlichen Symbolen zu
$$
H_0 = \log_2(N) \quad \text{[bit]} \qquad \text{oder} \qquad H_0 = \ln(N) \; \text{[nats]}
$$
berechnet. $H_0$ wird als \textit{Entscheidungsgehalt} der Nachrichtenquelle bezeichnet.
\end{definition}

Zum Beispiel ergibt sich bei einer Auswahl aus $N = 16$ gleichwahrscheinlichen Symbolen eine Informationsmenge von $H_0 = \log_2(N) = \log_2(2^4) = 4$ bit.

Die Wahl des Logarithmus als Informationsmaß hat einige Vorteile, die wie folgt begründet sind:

\begin{itemize}
  \item Der Entscheidungsgehalt $H_0$ wächst streng monoton mit der Größe des Alphabets $N$.
  \item Der Entscheidungsgehalt mehrfacher Entscheidungen ist additiv.
\end{itemize}

Die letztere Aussage lässt sich am Beispiel des Entscheidungsbaums in Abbildung 5.2 verdeutlichen, indem wir den Entscheidungsbaum von „unten nach oben" lesen. Auf der untersten Ebene sind binäre 4 Entscheidungen erforderlich. Dies entspricht einer Informationsmenge von 4 bit. Der darüberliegende „Baumstumpf" weist 4 Enden auf, was einem Entscheidungsgehalt von $\log_2(4) = 2$ bit entspricht. Da jedes dieser Enden nur in einem Viertel aller Fälle erreicht wird, beträgt die Informationsmenge insgesamt $1/4 \cdot 4\,\text{bit} + 2\,\text{bit} = 3\,\text{bit}$.

Wenn die Symbole nun mit unterschiedlicher Wahrscheinlichkeit auftreten, lässt sich das in Abbildung 5.2 dargestellte „Fragespiel" \textbf{im Mittel} deutlich abkürzen. Wenn \textit{a priori} bekannt ist, dass z. B. das Symbol „5" in 60\,\% aller Fälle von einem Sender ausgewählt wird, ist es sinnvoll zunächst zu fragen, ob das Symbol „5" ausgewählt wurde und erst danach alle anderen Möglichkeiten zu überprüfen. In 60\,\% aller Fälle ist man dann schon nach nur einer ja/nein-Frage am Ziel. Wenn dieses Experiment nun „unendlich oft" wiederholt wird, ist nach obigen Überlegungen die \textit{mittlere Informationsmenge} kleiner als $\log_2(N)$, denn es werden im Mittel weniger ja/nein-Fragen benötigt. Offensichtlich hängt also die Informationsmenge von den Wahrscheinlichkeiten des Auftretens der Symbole ab. Die \textit{a priori} Kenntnis der Auftrittswahrscheinlichkeiten verringert im Empfänger die Unsicherheit bezüglich der gesendeten Symbole.

Ähnliche Effekte treten auf, wenn man aufeinanderfolgende Symbole, z. B. die Buchstaben eines Wortes oder die Worte eines Textes betrachtet. Z.B. ist die Wahrscheinlichkeit, dass auf das Wort „grüner" das Wort „Schokoladenpudding" folgt außerordentlich gering. Allerdings ist sie auch nicht gleich Null, denn die Wortfolge „grüner Schokoladenpudding" kommt tatsächlich in Texten deutscher Sprache vor (allerdings nur in Diesem). Die Kenntnis dieser Statistik kann bei der Vermeidung von Übertragungsfehlern helfen. Ebenso lassen sich fehlende oder vertauschte Buchstaben in Worten ergänzen. Abbildung 5.3 und Abbildung 5.4 geben hierzu Beispiele. Die Redundanz der Daten ist in diesen Fällen hilfreich.

% Grafik: Abbildung 5.3: „I.F.R.A.I.N.T.C.N.K" — Jeder zweite Buchstabe fehlt. Welches Wort ist hier gemeint? (INFORMATIONSTECHNIK)

% Grafik: Abbildung 5.4: Artikel aus FAZ vom 24.09.03 — Buchstaben in Wörtern vertauscht, Text bleibt lesbar (Redundanz in der Schriftsprache)

Wir halten fest: Eine Nachrichtenquelle weist Redundanz auf,

\begin{itemize}
  \item wenn die Symbole der Quelle nicht gleichwahrscheinlich sind oder/und
  \item wenn aufeinanderfolgende Symbole voneinander abhängen.
\end{itemize}

In der Regel gibt es also signifikante Abhängigkeiten zwischen aufeinander folgenden Symbolen. Diese Abhängigkeiten vermindern die tatsächlich übertragene Informationsmenge, denn sie vermindern die Wahlfreiheit bei der Symbolauswahl, bzw. die Unsicherheit im Empfänger bezüglich der gesendeten Symbole.

Die Statistik der Symbole spielt bei der Ermittlung des Informationsgehaltes und der Redundanz offensichtlich eine große Rolle! Es ist also sehr zweckmäßig, den Informationsbegriff mit Hilfe von Wahrscheinlichkeiten zu definieren.

\section{Entropie und Redundanz}
\label{sec:entropie-redundanz}

Offensichtlich trägt ein Nachrichtensymbol um so mehr Information, je seltener es auftritt, d.h. je weniger wahrscheinlich es von der Quelle ausgewählt wird.

\begin{definition}[Informationsgehalt]
Der \textbf{Informationsgehalt} des $i$-ten Symbols $a_i$ wird daher zu

$$
I(a_i) = -\log_2(P(a_i)) = \log_2\!\left(\frac{1}{P(a_i)}\right) = -\mathrm{ld}(P(a_i))
$$

festgelegt, wobei $P(a_i)$ die Auftrittswahrscheinlichkeit des Symbols $a_i$ angibt. Wegen $0 \leq P(a_i) \leq 1$ ist dieses Informationsmaß nicht negativ.
\end{definition}

\begin{definition}[Entropie]
Der \textbf{mittlere Informationsgehalt} einer Nachrichtenquelle $A$, die $N$ verschiedene Symbole $a_i$ mit den Wahrscheinlichkeiten $P(a_i)$ aussenden kann, ist der statistische Mittelwert (Erwartungswert, Abschnitt~\ref{sec:erwartungswert}) des Informationsgehalts und wird dann zu

$$
H(A) = \sum_{i=1}^{N} P(a_i)\, I(a_i) = -\sum_{i=1}^{N} P(a_i) \log_2(P(a_i))
$$

berechnet. $H$ wird als \textit{Entropie} bezeichnet.
\end{definition}

\textbf{Anmerkung zu dieser Definition der Entropie:}

\begin{itemize}
  \item Für $N$ \textit{gleichwahrscheinliche} Ereignisse mit $\sum_{i=1}^{N} P(a_i) = 1$ gilt
$$
H(A) = \sum_{i=1}^{N} P(a_i)\, I(a_i) = -\sum_{i=1}^{N} \frac{1}{N} \log_2\!\left(\frac{1}{N}\right) = \log_2(N) = H_0.
$$
Die Entropie $H(A)$ ist gleich dem Entscheidungsgehalt $H_0$.

  \item $H$ nimmt zu wenn $N$ größer wird.

  \item Es besteht eine enge Beziehung zwischen dem hier eingeführten Informationsmaß und dem Entropiebegriff in der Thermodynamik. In der Thermodynamik beschreibt die Entropie den Grad an „Unordnung" in einem System. Auch hier wird die Entropie bei einer Gleichverteilung, z. B. von Gasmolekülen im Raum, maximal. Die informationstheoretische Entropie kann in diesem Zusammenhang als die Information verstanden werden, die in der mikroskopischen Anordnung in Gasmoleküle enthalten ist und einem makroskopischen Beobachter entgeht. Durch Zufuhr von Information kann im Prinzip die Unordnung eines Systems verringert werden. Die zugeführte Information ist somit negativer (thermodynamischer) Entropie gleichzusetzen. Wenn der 2. Hauptsatz der Thermodynamik feststellt, dass die Entropie in einem geschlossenen System nicht abnimmt, gilt dies ebenso auch für die Information. (siehe auch: H.D. Zeh: Entropie, Fischer (Tb.), Frankfurt, 2005 und J. Monod: Zufall und Notwendigkeit, Pieper 1983, DTV 1975).
\end{itemize}

\subsection{Berechnung der Entropie für $N = 2$}
\label{sec:entropie-binaer}

Eine Quelle sendet zwei Symbole $a_1$ und $a_2$, von denen \textit{a priori} bekannt ist, dass sie mit den Wahrscheinlichkeiten $P(a_1)$ und $P(a_2) = 1 - P(a_1)$ auftreten.

Der mittlere Informationsgehalt (Entropie) dieser Nachrichten berechnet sich dann zu

$$
H = -\sum_{i=1}^{2} P(a_i)\log_2(P(a_i))
= -P(a_1)\log_2(P(a_1)) - (1-P(a_1))\log_2(1-P(a_1))
$$

$$
= -\frac{1}{\log_e(2)}\left[P(a_1)\log_e(P(a_1)) + (1-P(a_1))\log_e(1-P(a_1))\right].
$$

Diese Funktion ist in Abbildung 5.5 dargestellt.

\begin{beweis}[Maximum der binären Entropie]
Das Maximum dieser Funktion finden wir durch Differenzieren nach $P(a_1)$ und Nullsetzen.

$$
\frac{dH}{dP(a_1)} = \frac{1}{\log_e(2)}\left[-\frac{P(a_1)}{P(a_1)} - \log_e(P(a_1)) + \frac{1-P(a_1)}{1-P(a_1)} + \log_e(1-P(a_1))\right] = 0
$$

Diese Funktion nimmt für

$$
\log_e\!\left(\frac{1-P(a_1)}{P(a_1)}\right) = 0 \;\Leftrightarrow\; P(a_1) = P(a_2) = 0.5
$$

ein Maximum an, denn $\dfrac{d^2H}{dP(a_1)^2} < 0$. Dabei haben wir die Beziehungen $\dfrac{d}{\mathrm{d}x}\log_e(x) = \dfrac{1}{x}$ und $\log_2(x) = \dfrac{\log_e(x)}{\log_e(2)}$ verwendet. Wenn die Quelle immer eines der Symbole auswählt, also $P(a_1) = 1$ oder $P(a_2) = 1$ gilt, dann ist die Entropie Null.
\end{beweis}

% Grafik: Abbildung 5.5: Entropie H [bit] als Funktion von P(a_1) für N=2; glockenförmige Kurve, Maximum H=1 bei P(a_1)=0.5

\subsection{Relative Entropie und Quellenredundanz}
\label{sec:relative-entropie}

Aus dem vorangegangenen Beispiel lässt sich folgern:

\begin{itemize}
  \item Der mittlere Informationsgehalt ist maximal, wenn die Nachrichten \textit{a priori} gleichwahrscheinlich ausgewählt werden (maximale Wahlfreiheit)!

  \item Offensichtlich kann der mittlere Informationsgehalt pro Symbol aber auch geringer als $H_0 = \log_2(N)$ sein. In den Nachrichtensymbolen ist dann weniger als $H_0 = \log_2(N)$ bit an Information enthalten!

  \item Für die \textit{relative Entropie}
$$
\frac{H}{H_0} = \frac{H}{\log_2(N)}
$$
gilt $0 \leq H/H_0 \leq 1$.

  \item Die Differenz
$$
r_q = 1 - \frac{H}{H_0} = \frac{H_0 - H}{H_0}
$$
wird als Quellenredundanz bezeichnet und wird oft in Prozent angegeben.
\end{itemize}

\begin{beispiel}[Relative Entropie einer Nachrichtenquelle]
Gegeben sei die binäre Quelle in Abbildung 5.6, deren Zeichen 0 und 1 mit den Wahrscheinlichkeiten $P(0) = 0.8$ und $P(1) = 0.2$ auftreten. Aufeinanderfolgende Zeichen seien statistisch unabhängig. Die folgenden informationstheoretischen Maße lassen sich nun wie folgt berechnen:

\begin{itemize}
  \item Der Entscheidungsgehalt der Quelle ist $H_0 = \log_2(2) = 1$ bit/Symbol.

  \item Die Entropie berechnet sich zu
$$
H = \sum_{i=1}^{N} P(a_i)\,I(a_i) = -0.8\log_2(0.8) - 0.2\log_2(0.2) \approx 0.7219
$$

  \item Die relative Entropie ist demnach $\dfrac{H}{H_0} = \dfrac{H}{1} \approx 0.7219$

  \item Die Quellenredundanz berechnet sich zu
$$
r_q = 1 - \frac{H}{H_0} \approx 1 - 0.7219 = 0.2781 \equiv 27.81\%
$$
\end{itemize}

% Grafik: Abbildung 5.6: Nachrichtenquelle erzeugt binäre Folge …0 0 1 0 1 0 1 1…
\end{beispiel}

\section{Kanalkapazität}
\label{sec:kanalkapazitaet}

Ein \textit{diskreter} Kanal ist ein Übertragungskanal, auf dem in einem vorgegebenen Zeitintervall $T$ nur eine endliche Anzahl $N$ verschiedener Signalformen übertragen werden kann.

\textbf{Beispiel:} Vier mögliche Signalformen pro Zeitintervall $T$.

% Grafik: Treppensignal u(t) mit 4 Stufen (−2, −1, 1, 2) über Zeitintervalle T, 2T, …

\begin{definition}[Kanalkapazität]
Die Kapazität eines rauschfreien, diskreten Kanals ist durch

$$
C = \lim_{\tau \to \infty} \frac{\log_2(N(\tau))}{\tau}
$$

gegeben, wobei $N(\tau)$ die Anzahl zulässiger Signalformen der Dauer $\tau$ bezeichnet. Die Kanalkapazität hat die Dimension Informationsmenge/Zeit und kann z. B. in bit/s angegeben werden.
\end{definition}

\section{Verlustlose Quellencodierung}
\label{sec:quellencodierung}

Die Informationsübertragung über rauschfreie, diskrete Kanäle ist durch das folgende wichtige Theorem gegeben:

\begin{satz}[Quellencodierungstheorem (Shannon)]
„Let a source have entropy $H$ (bits per symbol) and a channel have a capacity $C$ (bits per second). Then it is possible to encode the output of the source in such a way as to transmit at the average rate $\frac{C}{H} - \epsilon$ symbols per second over the channel where $\epsilon$ is arbitrarily small. It is not possible to transmit at an average rate greater than $\frac{C}{H}$."

\hfill C.E. Shannon, „A Mathematical Theory of Communication", 1948, Theorem 9
\end{satz}

Der mittlere Informationsgehalt und nicht der Entscheidungsgehalt der Quelle bestimmt die Zahl der im Mittel übertragbaren Symbole. Das Theorem sagt allerdings nicht, \textbf{wie} die Symbole codiert werden müssen, um diese Grenze zu erreichen. Hierfür wurden nachfolgend verschiedene Codierverfahren entwickelt, die auch unter dem Begriff „Entropiecodierung" zusammengefasst werden.

Unter Codierung versteht man die Zuordnung von „Ersatz"-Symbolen zu den ursprünglichen Symbolen einer Nachricht. In allen hier betrachteten Fällen werden diese Ersatzsymbole als Folge \textit{binärer} Ziffern konstruiert. Die Zuordnung der binären Ziffern zu den Symbolen der Nachricht wird dann hinsichtlich eines festzulegenden Kriteriums optimiert. In der Quellencodierung wird eine Minimierung der in dem Code befindlichen Redundanz angestrebt, wobei sich die Coderedundanz zu

$$
r_c = \frac{\mathrm{E}\{L\} - H}{H}
$$

berechnet und $\mathrm{E}\{L\}$ die mittlere Länge der Codeworte angibt.

Der mittlere Informationsgehalt einer Menge von Nachrichtensymbolen ist dann maximal (und gleich dem Entscheidungsgehalt), wenn diese Symbole mit gleicher Wahrscheinlichkeit auftreten. Die binären Ziffern „0" und „1" der Codeworte werden deshalb den Nachrichtensymbolen so zugeordnet, dass sie mit möglichst gleicher Wahrscheinlichkeit auftreten. Diese Vorgehensweise resultiert in Codes mit variabler Codewortlänge, wie sie schon lange zum Beispiel in der Form des Morse-Codes bekannt sind. Dabei werden häufig auftretenden Nachrichten kurze Codeworte und selten auftretenden Symbolen lange Codeworte zugeordnet. Die Codewortlänge wird dadurch im statistischen Mittel minimiert.

\subsection{Codierung nach Shannon und Fano}
\label{sec:shannon-fano}

Ein erstes Codierverfahren wurde von C.E. Shannon\footnote{C.E. Shannon: A Mathematical Theory of Communication, Bell Systems Technical Journal, Vol. 27, S. 379-423, Juli 1948} und R.M. Fano\footnote{R.M. Fano: The Transmission of Information, Technical Report No. 65, MIT, Cambridge, 1949.} vorgeschlagen. Das Codierschema kann wie folgt zusammengefasst werden:

\begin{enumerate}
  \item Ordnen der Originalsymbole nach fallenden Auftrittswahrscheinlichkeiten;
  \item Aufteilen der Originalsymbole in zwei Gruppen (möglichst) gleicher Summenwahrscheinlichkeiten;
  \item Zuordnung eines Bits, so dass die erste Gruppe mit „0" und die zweite Gruppe mit „1" adressiert wird;
  \item Wiederholung der Schritte 2. und 3. mit den im vorangegangenen Codierschritt entstandenen Gruppen.
\end{enumerate}

Der Code ist eindeutig decodierbar, da kein Codewort linker Teil eines anderen Codewortes ist („Präfix"-Freiheit).

Der erste Schritt des Sortierens nach fallenden Wahrscheinlichkeiten ist für dieses Verfahren essentiell. Eine kleine mittlere Codewortlänge wird nur dann erreicht, wenn Symbole mit ähnlich großen Wahrscheinlichkeiten in eine Gruppe sortiert werden. Ansonsten erhalten selten auftretende Symbole zu kurze und häufig auftretende Symbole zu lange Codeworte.

\begin{beispiel}[Codierung nach Fano]
Gegeben ist eine Quelle, die vier verschiedene Symbole $a_1$, $a_2$, $a_3$ und $a_4$ mit den Auftrittswahrscheinlichkeiten $P(a_1) = 0.1$, $P(a_2) = 0.4$, $P(a_3) = 0.2$ und $P(a_4) = 0.3$ aussendet. Der Code wird in der nachfolgenden Tabelle konstruiert. Daraus ergibt sich ein Entscheidungsbaum, der in Abb. 5.7 dargestellt ist.

\begin{center}
\begin{tabular}{c|c|ccc|c}
Symbol & $P(a_i)$ & \multicolumn{3}{c|}{Code} & Codewortlänge $l_i$ / bit \\
\hline
$a_2$ & 0.4 & 0 & $\times$ & $\times$ & 1 \\
$a_4$ & 0.3 & 1 & 0 & $\times$ & 2 \\
$a_3$ & 0.2 & 1 & 1 & 0       & 3 \\
$a_1$ & 0.1 & 1 & 1 & 1       & 3 \\
\end{tabular}
\end{center}

Weiterhin lassen sich die folgenden Größen berechnen:

Entropie der Quelle: $H = -\sum_{i=1}^{N} P(a_i)\log_2(P(a_i)) \approx 1.8464$ bit/Symbol

Relative Entropie der Quelle: $\dfrac{H}{H_0} = \dfrac{-\sum_{i=1}^{N} P(a_i)\log_2(P(a_i))}{\log_2(N)} \approx 0.9232$

Redundanz der Quelle: $r_q = 1 - \dfrac{H}{H_0} = 0.0768 \mathrel{\widehat{=}} 7.68\,\%$

Mittlere zu erwartende Codewortlänge bei Fano-Codierung:
$\mathrm{E}\{L\} = \sum_{i=1}^{N} l_i P(a_i) = 0.4 \cdot 1 + 0.3 \cdot 2 + 0.2 \cdot 3 + 0.1 \cdot 3 = 1.9$ bit/Symbol

Mit Hilfe des Codes und des Entscheidungsbaums wird nun ein Bitstrom wie folgt decodiert:
$1\;0\;1\;1\;1\;0\;1\;0\;0\;1\;1\;0\;0 \;\rightarrow\; a_4\,a_1\,a_2\,a_4\,a_2\,a_3\,a_2$.

Somit ergibt sich eine Coderedundanz: $r_c = \dfrac{\mathrm{E}\{L\} - H}{H} = \dfrac{1.9 - 1.8464}{1.8464} = 0.029$.

% Grafik: Abbildung 5.7: Entscheidungsbaum für den Fano-Code aus Beispiel 5.3
\end{beispiel}

Würde man hingegen die Symbole $a_2$ und $a_1$ in die erste Gruppe und $a_3$ und $a_4$ in die zweite Gruppe sortieren, dann erhält man einen schlechteren Code mit einer mittleren Codewortlänge von 2 Bit/Symbol.

\subsection{Codierung nach Huffman}
\label{sec:huffman}

Das Codierschema nach Huffman kann wie folgt zusammengefasst werden:

\begin{enumerate}
  \item Ordnen der Originalsymbole nach fallenden Auftrittswahrscheinlichkeiten;
  \item Zusammenfassen der Symbole mit den niedrigsten Wahrscheinlichkeiten zu einem neuen Ersatzsymbol;
  \item Zuordnung eines Bits zu den beiden Originalsymbolen des Ersatzsymbols, so dass das erste Originalsymbol mit „0" und das zweite Originalsymbol mit „1" (oder umgekehrt) adressiert wird;
  \item Ordnen der verbliebenen Originalsymbole und des neuen Ersatzsymbols nach fallenden Auftrittswahrscheinlichkeiten, wobei die Ersatzsymbole maximal aufrücken;
  \item Wiederholung der Schritte 2., 3. und 4.
\end{enumerate}

\begin{beispiel}[Codierung nach Huffman]
Ein Huffman Code kann für die Quelle aus dem vorangehenden Fano-Beispiel wie folgt entwickelt werden (Ersatzsymbole rücken maximal auf). Durch Rückwärtslesen ergibt sich die folgende Codetabelle:

\begin{center}
\begin{tabular}{c|c|c}
Symbol & Codewort & $l_i$ \\
\hline
$a_1$ & 001 & 3 \\
$a_2$ & 1   & 1 \\
$a_3$ & 000 & 3 \\
$a_4$ & 01  & 2 \\
\end{tabular}
\end{center}

% Grafik: Abbildung 5.8: Entscheidungsbaum für den Huffman-Code aus Beispiel 5.4 (Schritte: a3+a1→I(0.3); I+a4→II(0.6); II+a2→III(1.0))
\end{beispiel}

Aus dem Konstruktionsverfahren folgt, dass der Huffman-Code die mittlere Codewortlänge minimiert!

\begin{itemize}
  \item Man erkennt, dass den Knoten in jedem Schritt die minimalen Auftrittswahrscheinlichkeiten zugeordnet werden.
  \item In jedem Knoten wird zudem ein Bit vergeben. Somit werden die Bits den insgesamt wachsenden, aber in jedem Schritt minimal möglichen Wahrscheinlichkeiten zugeordnet.
  \item Somit wird der Erwartungswert $\mathrm{E}\{L\}$ minimal.
\end{itemize}

Da es häufig mehrere Huffman-Codes gleicher minimaler Codewortlänge gibt, ist es sinnvoll, mit einem weiteren Kriterium den besten Code auszuwählen. Ein häufig gewähltes Kriterium ist die minimale Varianz der Codewortlängen. Die Optimierung dieses Kriteriums wird durch das „maximale Aufrücken" der Kombinationssymbole in dem oben dargestellten Codierschema erreicht.

\subsection{Lauflängencodierung}
\label{sec:laufllaengencodierung}

Eine Form der redundanzreduzierenden Codierung, die oft im Zusammenhang mit der Entropiecodierung eingesetzt wird, ist die sogenannte Lauflängencodierung. Diese ist besonders leicht anwendbar, wenn der zu codierende Datenstrom aus binären Symbolen besteht. In diesem Fall handelt es sich um eine verlustlose Codierung. Bei der Lauflängencodierung wird die Anzahl aufeinanderfolgender gleicher Symbole erfasst. Codierung dieser „Lauflängen" kann dann z.B. mittels Huffman-Codes erfolgen, wobei die Codetabellen anhand der Häufigkeitsverteilungen der Lauflängen aus typischen Bitströmen ermittelt wurden. Der Wechsel von einer Reihe von Null-Bits auf eine Reihe von Eins-Bits muss hierbei nicht signalisiert werden, da sich diese automatisch aus dem Wechsel der beiden Bitmuster ergibt.

Dieses Prinzip wurde z.B. bei der inzwischen weitgehend verdrängten Faksimile-Übertragung („FAX") eingesetzt. Hierbei wird eine Bildvorlage zeilenweise in schwarz-weiß abgetastet. Die Lauflänge entspricht dann der Anzahl aufeinanderfolgender schwarzer oder weißer Bildpunkte. Bei schwarzer Schrift auf weißem Untergrund sind die weißen Lauflängen sehr viel länger als die schwarzen Lauflängen. Für weiße und schwarze Lauflängen werden daher unterschiedliche Huffman-Codes verwendet. Da auf eine weiße Lauflänge stets eine schwarze Lauflänge folgt (und umgekehrt), muss der Beginn einer neuen Lauflänge nicht explizit angezeigt werden. Allerdings muss das Ende einer Zeile mit einem speziellen Symbol signalisiert werden.

% Grafik: Abbildung 5.9: Haussilhouette, zeilenweise abgetastet; weiße und schwarze Lauflängen entstehen durch horizontales Abtasten der Bildvorlage

\subsection{Verlustlose und verlustbehaftete Audiocodierung}
\label{sec:audiocodierung}

Die verlustlose Quellencodierung spielt in vielen technischen Anwendungen eine große Rolle, z.B. bei der Speicherung von Audio Dateien. Die Huffman-Codierung ergibt in diesem Fall einen binären Code, der die zu speichernde Datenmenge im Mittel minimiert. Bei fehlerfreier Speicherung kann die Symbolfolge der Quelle ohne Verluste im Empfänger aus dem gespeicherten Code rekonstruiert werden. Diese Verfahren werden daher als \textit{verlustlose Quellencodierung} bezeichnet. Demgegenüber werden bei vielen Verfahren der Signalübertragung, z. B. bei der Mobilfunktelefonie oder MP3-Audiocodierung, irrelevante Daten im Zuge der Codierung entfernt. Die Symbolfolge der Quelle kann dann im Empfänger nicht exakt reproduziert werden. Jedoch ist das Empfangssignal hinsichtlich der vom Menschen wahrgenommenen Qualität nicht vom Original zu unterscheiden oder weist nur geringfügige und mithin akzeptable Verzerrungen auf. Man spricht dann von \textit{verlustbehafteter Quellencodierung}. Im Unterschied zur verlustlosen Quellencodierung ist bei einer verlustbehafteten Codierung ein mehrfaches Umcodieren des Signals mit Einbußen an Qualität verbunden und muss deshalb vermieden werden.

\textbf{FLAC Audiocodierung}

Das FLAC-Verfahren („free lossless audio codec") ist ein frei verfügbares Verfahren für die Codierung von Audiosignalen im PCM-Format, das im Mittel eine Reduktion der Datenmenge um ca. 50\,\% erreicht. Seit 2024 ist es als RFC 9639 („request for comment") als „Internet Engineering Task Force (IETF)" Standard etabliert und besonders auch für die paketierte Übertragung von Audiosignalen im Internet geeignet. Dabei wird das Audiosignal in Rahmen zerlegt, die neben den eigentlichen Audiodaten auch einen Kopf („header") mit Informationen über die Abtastrate und die Wortbreite enthalten. Jedes Signalsegment wird dann durch ein parametrisches Prädiktionsmodell approximiert und der Fehler („Residualsignal") zwischen dem parametrischen Modell und dem tatsächlichen Eingangssignal berechnet. Dieses Residualsignal wird dann einer Entropiecodierung (hier: Rice Code) unterzogen, die im Unterschied zur Huffman-Codierung keine Konstruktion und Übertragung von Codetabellen erfordert. Die Codierung und Decodierung benötigt nur einen geringen Rechenaufwand.

\textbf{MP3 und AAC Audiocodierung}

In der \textit{verlustbehafteten} Audiocodierung macht man sich den Umstand zunutze, dass zu jedem Zeitpunkt eines Audiosignals nur Teile des verfügbaren Frequenzspektrums eine nennenswerte Rolle spielen. Zum Beispiel befinden sich die leistungsstärksten Anteile eines Melodieinstruments in der Regel bei relativ tiefen Frequenzen während der Frequenzbereich nahe der halben Abtastrate („Nyquist-Frequenz") kaum Leistung gemessen werden kann. Aufgrund einer Eigenschaften des Gehörs, derzufolge leistungsstarke Anteile die leistungsschwachen Anteile (besonders bei hohen Frequenzen) verdecken, können dann die leistungsschwachen Frequenzanteile vor der Codierung entfernt werden, d.h. faktisch zu Null gesetzt, oder weniger genau codiert werden. Hierfür ist zunächst eine Frequenzanalyse des Signals erforderlich, so dass mit Hilfe eines „psychoakustischen Modells" eine Entscheidung über die zu codierenden Frequenzanteile getroffen werden kann. Im „MP3" Standard (ISO/IEC 11172-3, Layer III) werden bei tiefen Frequenzen die Frequenzwerte paarweise zusammengefasst und dann Huffman-codiert. Hierfür stellt der Standard Codiertabellen zur Verfügung, deren Statistik auf typischen Audiosignalen ermittelt wurde. Eine Folge des oben genannten Prinzips ist, dass der Decoder das originale Audiosignal nicht mehr bitgenau aus dem Bitstrom rekonstruieren kann. Ähnliche Prinzipien werden auch in neueren Standards eingesetzt, z.B. dem ISO/IEC 13818-7:2006 MPEG-2 Advanced Audio Coder (AAC), der auch Mehrkanalton unterstützt. Bei einem Stereosignal, sind Bitraten ab ca. 192 kbit/s mit der Qualität einer 16-Bit PCM Codierung vergleichbar („CD-Qualität"). Diese Standards werden kontinuierlich weiterentwickelt, z.B. in der Form des MPEG-4 High-efficiency Advanced Audio Coding AAC+ (ISO/IEC 14496-3:2005).

% Grafik: Abbildung 5.10: MP3-Encoder (oben): PCM audio samples → mapping → quantizer and coding (psychoacoustic model) → frame packing → encoded bitstream; MP3-Decoder (unten): encoded bitstream → frame unpacking → reconstruction → inverse mapping → PCM audio samples

\subsection{Kullbach-Leibler Divergenz und Kreuzentropie}
\label{sec:kullback-leibler}

Die Entropie

$$
H = \sum_i P(a_i) \log_2\!\left(\frac{1}{P(a_i)}\right)
$$

gibt an, wie viele Bits für die Übertragung der Symbole $a_i$ mindestens aufgewandt werden müssen. Auf der Grundlage der Auftrittswahrscheinlichkeiten $P(a_i)$ kann ein optimaler Code entwickelt werden. Treten die Symbole $a_i$ nun tatsächlich mit anderen Auftrittswahrscheinlichkeiten $\widehat{P}(a_i)$ auf, müssen prinzipiell mehr Bits für eine Codierung und Übertragung aufgewandt werden. Diese Abweichung wird durch die \textit{Kullbach-Leibler Divergenz} quantifiziert.

Die pro Symbol $a_i$ mindestens aufzuwendenden Bits sind durch den Informationsgehalt

$$
I(a_i) = \log_2\!\left(\frac{1}{P(a_i)}\right)
$$

gegeben. Tatsächlich werden die Symbole nun aber mit der Verteilung $\widehat{P}(a_i)$ erzeugt, so dass den Symbolen der Informationsgehalt

$$
\widehat{I}(a_i) = \log_2\!\left(\frac{1}{\widehat{P}(a_i)}\right)
$$

zugrunde liegt. Daraus ergibt sich im statistischen Mittel eine Differenz

$$
D_{\mathrm{KL}}(P(a_i) \| \widehat{P}(a_i)) = \sum_i P(a_i) \log_2\!\left(\frac{1}{\widehat{P}(a_i)}\right) - \sum_i P(a_i) \log_2\!\left(\frac{1}{P(a_i)}\right)
= \sum_i P(a_i) \log_2\!\left(\frac{P(a_i)}{\widehat{P}(a_i)}\right) \quad \text{[bit].}
$$

Diese Abweichung wird Kullbach-Leibler Divergenz genannt. Sie ist ein Maß für die Ähnlichkeit der Wahrscheinlichkeitsverteilungen $P(a_i)$ und $\widehat{P}(a_i)$. Es gilt

$$
0 \leq D_{\mathrm{KL}}(P(a_i) \| \widehat{P}(a_i)) \leq \sum_i \frac{P^2(a_i)}{\widehat{P}(a_i)} - 1 \quad \text{für } P(a_i),\, \widehat{P}(a_i) > 0.
$$

Weiterhin gilt

$$
D_{\mathrm{KL}}(P(a_i) \| \widehat{P}(a_i)) = H(P, \widehat{P}) - H(P)
$$

und

$$
D_{\mathrm{KL}}(\widehat{P}(a_i) \| P(a_i)) = H(\widehat{P}, P) - H(\widehat{P}).
$$

Der Term $H(\widehat{P}, P) = -\sum_i \widehat{P}(a_i) \log_2 P(a_i)$ wird auch als Kreuzentropie bezeichnet. Die Kreuzentropie wird häufig als Kriterium in der Adaption neuronaler Netze (mit der „Backpropagation"-Methode) eingesetzt. Hiermit kann die Abweichung der Verteilung $P(a_i)$ gegebener Daten zu der Verteilung $\widehat{P}(a_i)$ der von dem neuronalen Netz erzeugten Daten durch ein zielgerichtetes Training des Netzwerks minimiert werden.

\section{Zusammenfassung}
\label{sec:informationstheorie-zusammenfassung}

\begin{itemize}
  \item Digitale Daten werden häufig mit Hilfe binärer Zeichen codiert.
  \item Die Zahl der zu übertragenden binären Zeichen oder die Datenrate lässt per se keinen Rückschluss auf den Informationsgehalt dieser Daten zu.
  \item Der mittlere Informationsgehalt der Zeichen wird aus den Auftrittswahrscheinlichkeiten der Zeichen berechnet.
  \item Information entsteht erst durch Wahlfreiheit des Senders und die damit verbundene Unsicherheit bezüglich der gesendeten Daten beim Empfänger.
  \item Das Ziel der Quellencodierung ist die Verminderung der Datenmenge und mithin auch der Redundanz der Quellendaten.
  \item Für die über einen gestörten Kanal übertragbare Information ist der mittlere Informationsgehalt der Quelle in Relation zur Übertragungskapazität des Kanals von Bedeutung, nicht aber die Zahl der binären Zeichen.
\end{itemize}


\chapter{Übungen}

\section{Übung 1 – Analoge Signale}

Dieses Übungsblatt behandelt die grundlegenden Eigenschaften analoger Signale:
den Signalverlauf, die Unterscheidung in Energie- und Leistungssignale
(Abschnitte~\ref{sec:norm-energie-leistung} und~\ref{sec:leistung}), die
Symmetrieeigenschaften (Abschnitt~\ref{sec:symmetrie-signale}), Periodizität
(Abschnitt~\ref{sec:periodische-signale}), den Mittelwert
(Abschnitt~\ref{sec:mittelwert}) sowie komplexwertige Signale
(Abschnitt~\ref{sec:komplexwertige-signale}). Für die Rechteck- und
Sprungfunktion wird durchgehend die Skript-Definition aus
Abschnitt~\ref{sec:elementare-analoge-signale} verwendet:
$\operatorname{rect}(u) = 1$ für $0 \leq u < 1$ (sonst $0$) und
$\epsilon(t) = 1$ für $t \geq 0$ (sonst $0$).

% ================================================================
\subsection*{Vorrechenaufgabe 1 – Signalanalyse}
% ================================================================

Für die folgenden Signale ist der Signalverlauf zu skizzieren und zu prüfen, ob
ein Energie- oder ein Leistungssignal vorliegt. Zur Klassifikation dienen
Energie und Leistung (Abschnitte~\ref{sec:norm-energie-leistung}
und~\ref{sec:leistung}):

$$
E_x = \int_{-\infty}^{\infty} |x(t)|^2\,\mathrm{d}t,
\qquad
P_x = \lim_{T_0 \to \infty} \frac{1}{2T_0} \int_{-T_0}^{T_0} |x(t)|^2\,\mathrm{d}t.
$$

Ein \emph{Energiesignal} besitzt $0 < E_x < \infty$ (und $P_x = 0$), ein
\emph{Leistungssignal} $0 < P_x < \infty$ (und $E_x = \infty$).

% ---- a) --------------------------------------------------------
\subsubsection*{a) Rechteckimpuls}

\begin{beispiel}[Verschobener Rechteckimpuls]

\textbf{Gegeben:} $x(t) = \operatorname{rect}\!\left(\dfrac{t + T}{2T}\right)$
mit $T > 0$ (Abschnitt~\ref{sec:elementare-analoge-signale}).

\textbf{Gesucht:} Signalverlauf; Energie- oder Leistungssignal?

\textbf{Formel:} Rechteck-Definition sowie Energie
(Abschnitt~\ref{sec:norm-energie-leistung}).

\textbf{Rechnung:}

Das Argument liegt im Trägerbereich, wenn
$0 \leq \dfrac{t+T}{2T} < 1$, also $-T \leq t < T$. Dort ist $x(t) = 1$, sonst
$0$. Es handelt sich um einen Rechteckimpuls der Höhe $1$ und Breite $2T$,
zentriert um $t = 0$.

$$
E_x = \int_{-T}^{T} 1^2\,\mathrm{d}t = 2T < \infty.
$$

\begin{center}
\begin{tikzpicture}
\begin{axis}[
  width=10cm, height=4.5cm,
  xmin=-2.2, xmax=2.2, ymin=-0.3, ymax=1.5,
  axis lines=center,
  xlabel={$t$}, ylabel={$x(t)$},
  xtick={-1,1}, xticklabels={$-T$,$T$},
  ytick={1}, yticklabels={$1$},
  every axis plot/.append style={very thick, blue},
]
\addplot coordinates {(-2.2,0) (-1,0) (-1,1) (1,1) (1,0) (2.2,0)};
\end{axis}
\end{tikzpicture}
\end{center}

\textbf{Lösung:}

$$
\boxed{\;E_x = 2T < \infty \;\Rightarrow\; \text{Energiesignal.}\;}
$$

\end{beispiel}

% ---- b) --------------------------------------------------------
\subsubsection*{b) Abklingende Exponentialfunktion}

\begin{beispiel}[Einseitig abklingende Exponentialfunktion]

\textbf{Gegeben:} $x(t) = 2\,e^{-t/T}\,\epsilon(t)$ mit $T > 0$.

\textbf{Gesucht:} Signalverlauf; Energie- oder Leistungssignal?

\textbf{Formel:} Energie (Abschnitt~\ref{sec:norm-energie-leistung}).

\textbf{Rechnung:}

Für $t < 0$ ist $x(t) = 0$, für $t \geq 0$ fällt $x(t)$ von $x(0) = 2$
monoton gegen $0$ ab.

$$
E_x = \int_0^{\infty} \left(2\,e^{-t/T}\right)^2 \mathrm{d}t
    = 4 \int_0^{\infty} e^{-2t/T}\,\mathrm{d}t
    = 4 \cdot \frac{T}{2} = 2T < \infty.
$$

\begin{center}
\begin{tikzpicture}
\begin{axis}[
  width=10cm, height=4.5cm,
  xmin=-1.5, xmax=6, ymin=-0.3, ymax=2.4,
  axis lines=center,
  xlabel={$t/T$}, ylabel={$x(t)$},
  ytick={2}, yticklabels={$2$},
  samples=100,
  every axis plot/.append style={very thick, blue},
]
\addplot coordinates {(-1.5,0) (0,0)};
\addplot[domain=0:6] {2*exp(-x)};
\end{axis}
\end{tikzpicture}
\end{center}

\textbf{Lösung:}

$$
\boxed{\;E_x = 2T < \infty \;\Rightarrow\; \text{Energiesignal.}\;}
$$

\end{beispiel}

% ---- c) --------------------------------------------------------
\subsubsection*{c) Aufladende Exponentialfunktion}

\begin{beispiel}[Einseitig aufladende Exponentialfunktion]

\textbf{Gegeben:} $x(t) = 4\left(1 - e^{-2t/T}\right)\epsilon(t)$ mit $T > 0$.

\textbf{Gesucht:} Signalverlauf; Energie- oder Leistungssignal?

\textbf{Formel:} Energie und Leistung (Abschnitte~\ref{sec:norm-energie-leistung}
und~\ref{sec:leistung}).

\textbf{Rechnung:}

Für $t \geq 0$ steigt $x(t)$ von $0$ monoton gegen den Grenzwert $4$. Da $x(t)$
für $t \to \infty$ konstant $4$ bleibt, ist die Energie unbeschränkt
($E_x = \infty$). Die Leistung ist dagegen endlich:

$$
P_x = \lim_{T_0 \to \infty} \frac{1}{2T_0} \int_0^{T_0}
      16\left(1 - e^{-2t/T}\right)^2 \mathrm{d}t
    = \lim_{T_0 \to \infty} \frac{16\,T_0 - 12\,T}{2T_0}
    = \frac{16}{2} = 8.
$$

\begin{center}
\begin{tikzpicture}
\begin{axis}[
  width=10cm, height=4.5cm,
  xmin=-1.5, xmax=6, ymin=-0.5, ymax=5,
  axis lines=center,
  xlabel={$t/T$}, ylabel={$x(t)$},
  ytick={4}, yticklabels={$4$},
  samples=100,
  every axis plot/.append style={very thick, blue},
]
\addplot coordinates {(-1.5,0) (0,0)};
\addplot[domain=0:6] {4*(1-exp(-2*x))};
\draw[dashed, gray] (axis cs:-1.5,4) -- (axis cs:6,4);
\end{axis}
\end{tikzpicture}
\end{center}

\textbf{Lösung:}

$$
\boxed{\;E_x = \infty,\quad P_x = 8 \;\Rightarrow\; \text{Leistungssignal.}\;}
$$

\end{beispiel}

% ---- d) --------------------------------------------------------
\subsubsection*{d) Sinusförmiges Signal}

\begin{beispiel}[Sinusförmiges Signal mit Symmetrie, Mittelwert und Leistung]

\textbf{Gegeben:}
$x(t) = 2\sin\!\left(\dfrac{2\pi}{1\,\mathrm{s}}\,t\right)
= 2\sin\!\left(2\pi\dfrac{t}{T}\right) = 2\sin(\omega t) = 2\sin(2\pi f t)$.

\textbf{Gesucht:} Signalverlauf, Kenngrößen $T,\, f,\, \omega$,
Symmetrieeigenschaft, Mittelwert und Leistung.

\textbf{Formel:} Periodendauer, Symmetrie, Mittelwert und Leistung
(Abschnitte~\ref{sec:periodische-signale}, \ref{sec:symmetrie-signale},
\ref{sec:mittelwert-periodisch} und~\ref{sec:leistung}).

\textbf{Rechnung:}

Koeffizientenvergleich liefert $\omega = \dfrac{2\pi}{1\,\mathrm{s}}$, also

$$
\omega = 2\pi\,\mathrm{s}^{-1}, \qquad
f = \frac{\omega}{2\pi} = 1\,\mathrm{Hz}, \qquad
T = \frac{1}{f} = 1\,\mathrm{s}.
$$

Wegen $\sin(-\omega t) = -\sin(\omega t)$ ist $x(t)$ \emph{ungerade}. Der
Mittelwert über eine Periode verschwindet, und die Leistung eines
Sinussignals der Amplitude $A$ beträgt $A^2/2$:

$$
\bar{x} = \frac{1}{T}\int_0^T 2\sin(\omega t)\,\mathrm{d}t = 0,
\qquad
P_x = \frac{1}{T}\int_0^T 4\sin^2(\omega t)\,\mathrm{d}t
    = \frac{A^2}{2} = \frac{2^2}{2} = 2.
$$

\begin{center}
\begin{tikzpicture}
\begin{axis}[
  width=11cm, height=4.8cm,
  xmin=-0.3, xmax=2.3, ymin=-2.6, ymax=2.6,
  axis lines=center,
  xlabel={$t/T$}, ylabel={$x(t)$},
  ytick={-2,2}, yticklabels={$-2$,$2$},
  samples=200, domain=-0.3:2.3,
  every axis plot/.append style={very thick, blue},
]
\addplot {2*sin(360*x)};
\end{axis}
\end{tikzpicture}
\end{center}

\textbf{Lösung:}

$$
\boxed{\;T = 1\,\mathrm{s},\;\; f = 1\,\mathrm{Hz},\;\;
\omega = 2\pi\,\mathrm{s}^{-1};\quad
\text{ungerade},\;\; \bar{x} = 0,\;\; P_x = 2 \;\Rightarrow\;
\text{Leistungssignal.}\;}
$$

\end{beispiel}

% ---- e) --------------------------------------------------------
\subsubsection*{e) Komplexwertiges Exponentialsignal}

\begin{beispiel}[Zusammenhang zwischen Betrag, Phase und Real- bzw. Imaginärteil]

\textbf{Gegeben:} Komplexe Zahl $x = A\cdot e^{j\varphi}$ sowie das
komplexwertige Signal $\underline{x}(t) = e^{j\pi t}$
(Abschnitt~\ref{sec:komplexe-exponentialsignale}).

\textbf{Gesucht:} Grafischer und rechnerischer Zusammenhang zwischen
$\sin\varphi$, $\cos\varphi$ und Real- bzw.\ Imaginärteil; anschließend
$\operatorname{Re}\{\underline{x}(t)\}$, $\operatorname{Im}\{\underline{x}(t)\}$,
$T$, $f$, Mittelwert und Leistung von $\underline{x}(t)$.

\textbf{Formel} (Eulersche Formel, Abschnitt~\ref{sec:komplexe-exponentialsignale}):

$$
A\, e^{j\varphi} = A\cos\varphi + j\,A\sin\varphi,
\qquad
\operatorname{Re}\{x\} = A\cos\varphi, \quad
\operatorname{Im}\{x\} = A\sin\varphi.
$$

\textbf{Rechnung:}

Grafisch entspricht $x = A\,e^{j\varphi}$ einem Zeiger der Länge $A$ unter dem
Winkel $\varphi$ in der komplexen Ebene. Die Projektion auf die reelle Achse
liefert $A\cos\varphi$, die Projektion auf die imaginäre Achse $A\sin\varphi$:

\begin{center}
\begin{tikzpicture}[scale=2.2]
  \draw[->] (-1.3,0) -- (1.4,0) node[right] {$\operatorname{Re}$};
  \draw[->] (0,-0.4) -- (0,1.3) node[above] {$\operatorname{Im}$};
  \draw[gray, dashed] (0,0) circle (1);
  \def\ang{55}
  \draw[very thick, blue, ->] (0,0) -- (\ang:1) node[midway, above left] {$A$};
  \draw[dashed] (\ang:1) -- ({cos(\ang)},0);
  \draw[dashed] (\ang:1) -- (0,{sin(\ang)});
  \draw[red, thick] (0,0) -- ({cos(\ang)},0)
        node[midway, below] {$A\cos\varphi$};
  \draw[red, thick] (0,{sin(\ang)}) -- (\ang:1)
        node[midway, above] {$A\sin\varphi$};
  \draw (0.3,0) arc (0:\ang:0.3);
  \node at (28:0.42) {$\varphi$};
\end{tikzpicture}
\end{center}

Für $\underline{x}(t) = e^{j\pi t}$ gilt $A = 1$ und $\varphi(t) = \pi t$, also

$$
\operatorname{Re}\{\underline{x}(t)\} = \cos(\pi t), \qquad
\operatorname{Im}\{\underline{x}(t)\} = \sin(\pi t).
$$

Aus $\omega = \pi\,\mathrm{s}^{-1}$ folgt

$$
f = \frac{\omega}{2\pi} = \frac{1}{2}\,\mathrm{Hz}, \qquad
T = \frac{1}{f} = 2\,\mathrm{s}.
$$

Wegen $|\underline{x}(t)| = 1$ und der vollständigen Umdrehung pro Periode:

$$
\bar{\underline{x}} = \frac{1}{T}\int_0^T e^{j\pi t}\,\mathrm{d}t = 0,
\qquad
P_{\underline{x}} = \frac{1}{T}\int_0^T |e^{j\pi t}|^2\,\mathrm{d}t = 1.
$$

\begin{center}
\begin{tikzpicture}
\begin{axis}[
  width=11cm, height=4.8cm,
  xmin=-1.2, xmax=3.2, ymin=-1.5, ymax=1.5,
  axis lines=center,
  xlabel={$t$}, ylabel={},
  ytick={-1,1}, yticklabels={$-1$,$1$},
  samples=200, domain=-1.2:3.2,
  legend pos=north east, legend style={font=\small},
  every axis plot/.append style={very thick},
]
\addplot[blue] {cos(180*x)}; \addlegendentry{$\operatorname{Re}\{\underline{x}(t)\}$}
\addplot[red, dashed] {sin(180*x)}; \addlegendentry{$\operatorname{Im}\{\underline{x}(t)\}$}
\end{axis}
\end{tikzpicture}
\end{center}

\textbf{Lösung:}

$$
\boxed{\;\operatorname{Re}\{\underline{x}\} = \cos(\pi t),\;\;
\operatorname{Im}\{\underline{x}\} = \sin(\pi t);\;\;
T = 2\,\mathrm{s},\;\; f = \tfrac12\,\mathrm{Hz},\;\;
\bar{\underline{x}} = 0,\;\; P_{\underline{x}} = 1.\;}
$$

$\underline{x}(t)$ ist ein Leistungssignal mit konstantem Betrag $1$.

\end{beispiel}

% ================================================================
\subsection*{Aufgabe 1.1 – Einheitssprung}
% ================================================================

Die folgenden Signale sind unter Beachtung aller charakteristischen Werte als
Funktion der Zeit zu zeichnen; abschließend erfolgt die Klassifikation in
Energie- oder Leistungssignale.

\subsubsection*{a) $x(t) = 2 - 2\epsilon(t)$}

\begin{beispiel}[Absteigende Sprungfunktion]

\textbf{Gegeben:} $x(t) = 2 - 2\epsilon(t)$
(Abschnitt~\ref{sec:elementare-analoge-signale}).

\textbf{Gesucht:} Signalverlauf; Energie- oder Leistungssignal?

\textbf{Formel:} Sprung-Definition sowie Leistung
(Abschnitt~\ref{sec:leistung}).

\textbf{Rechnung:} Für $t < 0$ ist $\epsilon(t) = 0$, also $x = 2$; für
$t \geq 0$ ist $x = 2 - 2 = 0$. Da $x(t) = 2$ für alle $t < 0$ gilt, ist
$E_x = \infty$; die Leistung ist

$$
P_x = \lim_{T_0 \to \infty} \frac{1}{2T_0}\int_{-T_0}^{0} 2^2\,\mathrm{d}t
    = \lim_{T_0 \to \infty} \frac{4\,T_0}{2T_0} = 2.
$$

\begin{center}
\begin{tikzpicture}
\begin{axis}[
  width=10cm, height=4cm,
  xmin=-3, xmax=3, ymin=-0.5, ymax=2.6,
  axis lines=center, xlabel={$t$}, ylabel={$x(t)$},
  ytick={2}, yticklabels={$2$},
  every axis plot/.append style={very thick, blue},
]
\addplot coordinates {(-3,2) (0,2) (0,0) (3,0)};
\end{axis}
\end{tikzpicture}
\end{center}

\textbf{Lösung:} $\boxed{\;P_x = 2 \Rightarrow \text{Leistungssignal.}\;}$

\end{beispiel}

\subsubsection*{b) $x(t) = \epsilon(t) - \epsilon(t-1\,\mathrm{ms}) + \epsilon(t-2\,\mathrm{ms}) - \epsilon(t-3\,\mathrm{ms})$}

\begin{beispiel}[Zwei Rechteckimpulse aus Sprungfunktionen]

\textbf{Gegeben:} $x(t) = \epsilon(t) - \epsilon(t-1\,\mathrm{ms})
+ \epsilon(t-2\,\mathrm{ms}) - \epsilon(t-3\,\mathrm{ms})$.

\textbf{Gesucht:} Signalverlauf; Energie- oder Leistungssignal?

\textbf{Rechnung:} Abschnittsweise Auswertung ergibt zwei Rechteckimpulse der
Höhe $1$: auf $[0,1\,\mathrm{ms})$ und auf $[2\,\mathrm{ms},3\,\mathrm{ms})$.
Endliche Dauer, also Energiesignal mit

$$
E_x = 1^2\cdot 1\,\mathrm{ms} + 1^2\cdot 1\,\mathrm{ms}
    = 2\,\mathrm{ms}.
$$

\begin{center}
\begin{tikzpicture}
\begin{axis}[
  width=10cm, height=4cm,
  xmin=-1, xmax=4.5, ymin=-0.4, ymax=1.5,
  axis lines=center, xlabel={$t/\mathrm{ms}$}, ylabel={$x(t)$},
  xtick={1,2,3}, ytick={1}, yticklabels={$1$},
  every axis plot/.append style={very thick, blue},
]
\addplot coordinates {(-1,0) (0,0) (0,1) (1,1) (1,0) (2,0) (2,1) (3,1) (3,0) (4.5,0)};
\end{axis}
\end{tikzpicture}
\end{center}

\textbf{Lösung:} $\boxed{\;E_x = 2\,\mathrm{ms} \Rightarrow \text{Energiesignal.}\;}$

\end{beispiel}

\subsubsection*{c) $x(t) = \operatorname{rect}\!\left(\dfrac{t-2\,\mathrm{ms}}{4\,\mathrm{ms}}\right)$}

\begin{beispiel}[Rechteckimpuls über die Rechteckfunktion]

\textbf{Gegeben:} $x(t) = \operatorname{rect}\!\left(\dfrac{t-2\,\mathrm{ms}}{4\,\mathrm{ms}}\right)$.

\textbf{Gesucht:} Signalverlauf; Energie- oder Leistungssignal?

\textbf{Rechnung:} Trägerbedingung $0 \leq \dfrac{t-2\,\mathrm{ms}}{4\,\mathrm{ms}} < 1$
liefert $2\,\mathrm{ms} \leq t < 6\,\mathrm{ms}$; dort $x = 1$. Rechteck der
Höhe $1$ und Breite $4\,\mathrm{ms}$.

$$
E_x = 1^2 \cdot 4\,\mathrm{ms} = 4\,\mathrm{ms}.
$$

\begin{center}
\begin{tikzpicture}
\begin{axis}[
  width=10cm, height=4cm,
  xmin=-1, xmax=8, ymin=-0.4, ymax=1.5,
  axis lines=center, xlabel={$t/\mathrm{ms}$}, ylabel={$x(t)$},
  xtick={2,6}, ytick={1}, yticklabels={$1$},
  every axis plot/.append style={very thick, blue},
]
\addplot coordinates {(-1,0) (2,0) (2,1) (6,1) (6,0) (8,0)};
\end{axis}
\end{tikzpicture}
\end{center}

\textbf{Lösung:} $\boxed{\;E_x = 4\,\mathrm{ms} \Rightarrow \text{Energiesignal.}\;}$

\end{beispiel}

\subsubsection*{d) $x(t) = \pi\left[\epsilon(t-T) - \epsilon(t-3T)\right]$}

\begin{beispiel}[Skalierter Rechteckimpuls aus Sprungfunktionen]

\textbf{Gegeben:} $x(t) = \pi\left[\epsilon(t-T) - \epsilon(t-3T)\right]$,
$T > 0$.

\textbf{Gesucht:} Signalverlauf; Energie- oder Leistungssignal?

\textbf{Rechnung:} Rechteckimpuls der Höhe $\pi$ auf $[T, 3T)$, Breite $2T$.

$$
E_x = \pi^2 \cdot 2T = 2\pi^2 T.
$$

\begin{center}
\begin{tikzpicture}
\begin{axis}[
  width=10cm, height=4cm,
  xmin=-0.5, xmax=5, ymin=-0.5, ymax=4,
  axis lines=center, xlabel={$t/T$}, ylabel={$x(t)$},
  xtick={1,3}, ytick={3.1416}, yticklabels={$\pi$},
  every axis plot/.append style={very thick, blue},
]
\addplot coordinates {(-0.5,0) (1,0) (1,3.1416) (3,3.1416) (3,0) (5,0)};
\end{axis}
\end{tikzpicture}
\end{center}

\textbf{Lösung:} $\boxed{\;E_x = 2\pi^2 T \Rightarrow \text{Energiesignal.}\;}$

\end{beispiel}

\begin{bemerkung}[Klassifikation der Signale aus Aufgabe 1.1]
Die Signale b), c) und d) besitzen endliche Dauer und sind daher
\emph{Energiesignale}. Signal a) reicht für $t \to -\infty$ auf konstantem
Niveau $2$ und ist somit ein \emph{Leistungssignal}.
\end{bemerkung}

% ================================================================
\subsection*{Aufgabe 1.2 – Darstellung einer Rechteckfunktion}
% ================================================================

Die im Bild dargestellte, stückweise konstante Funktion ist als Summe von
Sprung- bzw.\ Rechteckfunktionen darzustellen. Ausgewertet ergibt sich

$$
x(t) = \begin{cases}
2 & , \quad 0 \leq t < T \\
-1 & , \quad T \leq t < 2T \\
1 & , \quad 2T \leq t < 3T \\
2 & , \quad 3T \leq t < 4T \\
0 & , \quad \text{sonst.}
\end{cases}
$$

\begin{center}
\begin{tikzpicture}
\begin{axis}[
  width=12cm, height=5cm,
  xmin=-0.5, xmax=5, ymin=-2.4, ymax=2.6,
  axis lines=center, xlabel={$t/T$}, ylabel={$x(t)$},
  xtick={1,2,3,4}, ytick={-2,-1,1,2},
  every axis plot/.append style={very thick, blue},
]
\addplot coordinates {(-0.5,0) (0,0) (0,2) (1,2) (1,-1) (2,-1) (2,1) (3,1) (3,2) (4,2) (4,0) (5,0)};
\end{axis}
\end{tikzpicture}
\end{center}

\subsubsection*{a) Darstellung als Summe von Sprungfunktionen}

\begin{beispiel}[Zerlegung in Sprungfunktionen]

\textbf{Gegeben:} Die obige Treppenfunktion.

\textbf{Gesucht:} Darstellung als Summe gewichteter, verschobener
Sprungfunktionen.

\textbf{Formel} (Sprungfunktion, Abschnitt~\ref{sec:elementare-analoge-signale}):
An jeder Sprungstelle $t_i$ wird ein Term $(\Delta x_i)\,\epsilon(t - t_i)$
addiert, wobei $\Delta x_i$ die Sprunghöhe ist.

\textbf{Rechnung:} Die Sprunghöhen betragen

$$
\begin{aligned}
&t=0:\; 0 \to 2 \;(+2), \quad
 t=T:\; 2 \to -1 \;(-3), \quad
 t=2T:\; -1 \to 1 \;(+2), \\
&t=3T:\; 1 \to 2 \;(+1), \quad
 t=4T:\; 2 \to 0 \;(-2).
\end{aligned}
$$

Probe: $2 - 3 + 2 + 1 - 2 = 0$ \checkmark (das Signal endet bei $0$).

\textbf{Lösung:}

$$
\boxed{\;x(t) = 2\,\epsilon(t) - 3\,\epsilon(t-T) + 2\,\epsilon(t-2T)
+ \epsilon(t-3T) - 2\,\epsilon(t-4T).\;}
$$

\end{beispiel}

\subsubsection*{b) Darstellung als Summe von Rechteckfunktionen}

\begin{beispiel}[Zerlegung in Rechteckfunktionen]

\textbf{Gegeben:} Dieselbe Treppenfunktion.

\textbf{Gesucht:} Darstellung als Summe gewichteter Rechteckfunktionen.

\textbf{Formel:} Jeder konstante Abschnitt $[kT,(k+1)T)$ der Höhe $h_k$ wird
durch $h_k \operatorname{rect}\!\left(\dfrac{t-kT}{T}\right)$ dargestellt.

\textbf{Rechnung:} Die vier nichtverschwindenden Stufen besitzen die Höhen
$2,\, -1,\, 1,\, 2$.

\textbf{Lösung:}

$$
\boxed{\;x(t) = 2\operatorname{rect}\!\left(\frac{t}{T}\right)
- \operatorname{rect}\!\left(\frac{t-T}{T}\right)
+ \operatorname{rect}\!\left(\frac{t-2T}{T}\right)
+ 2\operatorname{rect}\!\left(\frac{t-3T}{T}\right).\;}
$$

\end{beispiel}

% ================================================================
\subsection*{Aufgabe 1.3 – Exponentialfunktion}
% ================================================================

Die Funktionen sind für $T = 1\,\mathrm{ms}$ als Funktion von $t$ und $t/T$ zu
zeichnen; abschließend erfolgt die Klassifikation.

\subsubsection*{a) $x(t) = -e^{-t/2T}\,\epsilon(t)$}

\begin{beispiel}[Negative abklingende Exponentialfunktion]

\textbf{Gegeben:} $x(t) = -e^{-t/2T}\,\epsilon(t)$, $T = 1\,\mathrm{ms}$.

\textbf{Gesucht:} Signalverlauf; Energie- oder Leistungssignal?

\textbf{Rechnung:} Für $t \geq 0$ steigt $x(t)$ von $x(0) = -1$ monoton gegen
$0$ (von unten); Zeitkonstante $2T$. Die Energie ist endlich:

$$
E_x = \int_0^{\infty} \left(-e^{-t/2T}\right)^2 \mathrm{d}t
    = \int_0^{\infty} e^{-t/T}\,\mathrm{d}t = T = 1\,\mathrm{ms}.
$$

\begin{center}
\begin{tikzpicture}
\begin{axis}[
  width=10cm, height=4.5cm,
  xmin=-1.5, xmax=6, ymin=-1.3, ymax=0.4,
  axis lines=center, xlabel={$t/T$}, ylabel={$x(t)$},
  ytick={-1}, yticklabels={$-1$},
  samples=100,
  every axis plot/.append style={very thick, blue},
]
\addplot coordinates {(-1.5,0) (0,0)};
\addplot[domain=0:6] {-exp(-x/2)};
\end{axis}
\end{tikzpicture}
\end{center}

\textbf{Lösung:} $\boxed{\;E_x = T \Rightarrow \text{Energiesignal.}\;}$

\end{beispiel}

\subsubsection*{b) $x(t) = 4\left(1 - \exp\!\left(-\frac{t-2T}{T}\right)\right)\epsilon(t-2T)$}

\begin{beispiel}[Verzögert aufladende Exponentialfunktion]

\textbf{Gegeben:} $x(t) = 4\left(1 - \exp\!\left(-\frac{t-2T}{T}\right)\right)
\epsilon(t-2T)$, $T = 1\,\mathrm{ms}$.

\textbf{Gesucht:} Signalverlauf; Energie- oder Leistungssignal?

\textbf{Rechnung:} Für $t \geq 2T$ steigt $x(t)$ von $0$ monoton gegen den
Grenzwert $4$ (Zeitkonstante $T$). Da $x(t) \to 4$ konstant bleibt, ist
$E_x = \infty$ und

$$
P_x = \lim_{T_0 \to \infty} \frac{1}{2T_0}\int_{2T}^{T_0}
      16\left(1 - e^{-(t-2T)/T}\right)^2 \mathrm{d}t = \frac{16}{2} = 8.
$$

\begin{center}
\begin{tikzpicture}
\begin{axis}[
  width=10cm, height=4.5cm,
  xmin=-0.5, xmax=8, ymin=-0.5, ymax=5,
  axis lines=center, xlabel={$t/T$}, ylabel={$x(t)$},
  xtick={2}, ytick={4}, yticklabels={$4$},
  samples=100,
  every axis plot/.append style={very thick, blue},
]
\addplot coordinates {(-0.5,0) (2,0)};
\addplot[domain=2:8] {4*(1-exp(-(x-2)))};
\draw[dashed, gray] (axis cs:-0.5,4) -- (axis cs:8,4);
\end{axis}
\end{tikzpicture}
\end{center}

\textbf{Lösung:} $\boxed{\;P_x = 8 \Rightarrow \text{Leistungssignal.}\;}$

\end{beispiel}

% ================================================================
\subsection*{Aufgabe 1.4 – Harmonische Zeitfunktionen}
% ================================================================

Für jedes Signal sind Periodendauer $T$, Frequenz $f = 1/T$, Kreisfrequenz
$\omega = 2\pi/T$, Mittelwert, Leistung und Symmetrieeigenschaft zu bestimmen.
Für die Leistung periodischer Signale gilt
$P_x = \frac{1}{T}\int_0^T |x(t)|^2\,\mathrm{d}t$
(Abschnitt~\ref{sec:leistung}); die Leistung einer harmonischen Schwingung der
Amplitude $A$ beträgt $A^2/2$.

\subsubsection*{a) $x(t) = -\sin\!\left(\frac{2\pi}{4\,\mathrm{s}}t + \frac{\pi}{2}\right)$}

\begin{beispiel}[Phasenverschobene Sinusschwingung]

\textbf{Gegeben:} $x(t) = -\sin\!\left(\dfrac{2\pi}{4\,\mathrm{s}}t + \dfrac{\pi}{2}\right)$.

\textbf{Gesucht:} $T,\, f,\, \omega$, Mittelwert, Leistung, Symmetrie.

\textbf{Rechnung:} Mit $\sin(\theta + \tfrac{\pi}{2}) = \cos\theta$ folgt
$x(t) = -\cos\!\left(\dfrac{\pi}{2}\,t\right)$. Damit

$$
\omega = \frac{\pi}{2}\,\mathrm{s}^{-1}, \qquad
T = \frac{2\pi}{\omega} = 4\,\mathrm{s}, \qquad
f = \frac{1}{T} = \frac{1}{4}\,\mathrm{Hz}.
$$

Da $-\cos$ gerade ist, ist $x(t)$ \emph{gerade}. Mittelwert $\bar{x} = 0$,
Leistung $P_x = \dfrac{1^2}{2} = \dfrac{1}{2}$.

\begin{center}
\begin{tikzpicture}
\begin{axis}[
  width=11cm, height=4.5cm,
  xmin=-0.3, xmax=8.3, ymin=-1.4, ymax=1.4,
  axis lines=center, xlabel={$t/\mathrm{s}$}, ylabel={$x(t)$},
  ytick={-1,1}, samples=200, domain=-0.3:8.3,
  every axis plot/.append style={very thick, blue},
]
\addplot {-cos(90*x)};
\end{axis}
\end{tikzpicture}
\end{center}

\textbf{Lösung:}
$\boxed{\;T = 4\,\mathrm{s},\; f = \tfrac14\,\mathrm{Hz},\;
\omega = \tfrac{\pi}{2}\,\mathrm{s}^{-1};\; \text{gerade},\;
\bar{x} = 0,\; P_x = \tfrac12.\;}$

\end{beispiel}

\subsubsection*{b) $x(t) = 1 + \cos\!\left(\frac{\pi t}{1\,\mathrm{s}} + \pi\right)$}

\begin{beispiel}[Kosinusschwingung mit Gleichanteil]

\textbf{Gegeben:} $x(t) = 1 + \cos\!\left(\dfrac{\pi t}{1\,\mathrm{s}} + \pi\right)$.

\textbf{Gesucht:} $T,\, f,\, \omega$, Mittelwert, Leistung, Symmetrie.

\textbf{Rechnung:} Mit $\cos(\theta + \pi) = -\cos\theta$ folgt
$x(t) = 1 - \cos(\pi t)$. Damit

$$
\omega = \pi\,\mathrm{s}^{-1}, \qquad T = 2\,\mathrm{s}, \qquad
f = \tfrac12\,\mathrm{Hz}.
$$

$x(t)$ ist \emph{gerade}. Der Gleichanteil liefert den Mittelwert
$\bar{x} = 1$. Die Leistung setzt sich aus Gleich- und Wechselanteil zusammen:

$$
P_x = \frac{1}{T}\int_0^T (1 - \cos\pi t)^2\,\mathrm{d}t
    = 1^2 + \frac{1^2}{2} = \frac{3}{2}.
$$

\begin{center}
\begin{tikzpicture}
\begin{axis}[
  width=11cm, height=4.5cm,
  xmin=-0.3, xmax=4.3, ymin=-0.4, ymax=2.4,
  axis lines=center, xlabel={$t/\mathrm{s}$}, ylabel={$x(t)$},
  ytick={1,2}, samples=200, domain=-0.3:4.3,
  every axis plot/.append style={very thick, blue},
]
\addplot {1-cos(180*x)};
\end{axis}
\end{tikzpicture}
\end{center}

\textbf{Lösung:}
$\boxed{\;T = 2\,\mathrm{s},\; f = \tfrac12\,\mathrm{Hz},\;
\omega = \pi\,\mathrm{s}^{-1};\; \text{gerade},\;
\bar{x} = 1,\; P_x = \tfrac32.\;}$

\end{beispiel}

\subsubsection*{c) $x(t) = \cos\!\left(\frac{100\pi}{1\,\mathrm{s}}t\right)$}

\begin{beispiel}[Hochfrequente Kosinusschwingung]

\textbf{Gegeben:} $x(t) = \cos\!\left(\dfrac{100\pi}{1\,\mathrm{s}}\,t\right)$.

\textbf{Gesucht:} $T,\, f,\, \omega$, Mittelwert, Leistung, Symmetrie.

\textbf{Rechnung:}

$$
\omega = 100\pi\,\mathrm{s}^{-1}, \qquad
T = \frac{2\pi}{\omega} = \frac{1}{50}\,\mathrm{s} = 20\,\mathrm{ms}, \qquad
f = 50\,\mathrm{Hz}.
$$

$x(t)$ ist \emph{gerade}, $\bar{x} = 0$, $P_x = \dfrac{1^2}{2} = \dfrac{1}{2}$.

\begin{center}
\begin{tikzpicture}
\begin{axis}[
  width=11cm, height=4.5cm,
  xmin=-0.2, xmax=2.2, ymin=-1.4, ymax=1.4,
  axis lines=center, xlabel={$t/T$}, ylabel={$x(t)$},
  ytick={-1,1}, samples=200, domain=-0.2:2.2,
  every axis plot/.append style={very thick, blue},
]
\addplot {cos(360*x)};
\end{axis}
\end{tikzpicture}
\end{center}

\textbf{Lösung:}
$\boxed{\;T = 20\,\mathrm{ms},\; f = 50\,\mathrm{Hz},\;
\omega = 100\pi\,\mathrm{s}^{-1};\; \text{gerade},\;
\bar{x} = 0,\; P_x = \tfrac12.\;}$

\end{beispiel}

% ================================================================
\subsection*{Aufgabe 1.5 – Zusammenfassung zweier Harmonischer}
% ================================================================

\begin{beispiel}[Amplitude und Phase der Summe zweier Harmonischer]

\textbf{Gegeben:} $x(t) = \sin(\omega t) + \cos(\omega t)$.

\textbf{Gesucht:} Amplitude $A$ und Phase $\varphi$ so, dass
$x(t) = A\sin(\omega t + \varphi)$.

\textbf{Formel} (Additionstheorem):

$$
A\sin(\omega t + \varphi) = A\cos\varphi\,\sin(\omega t)
+ A\sin\varphi\,\cos(\omega t).
$$

\textbf{Rechnung:} Koeffizientenvergleich mit
$\sin(\omega t) + \cos(\omega t)$:

$$
A\cos\varphi = 1, \qquad A\sin\varphi = 1.
$$

Quadrieren und Addieren liefert $A^2(\cos^2\varphi + \sin^2\varphi) = 2$, also
$A = \sqrt{2}$. Der Quotient ergibt $\tan\varphi = 1$, und da beide
Komponenten positiv sind, liegt $\varphi$ im ersten Quadranten:
$\varphi = \dfrac{\pi}{4}$.

\textbf{Lösung:}

$$
\boxed{\;x(t) = \sqrt{2}\,\sin\!\left(\omega t + \frac{\pi}{4}\right).\;}
$$

\end{beispiel}

% ================================================================
\subsection*{Aufgabe 1.6 – Signaleigenschaften}
% ================================================================

Zu bestimmen sind Periodizität, Symmetrie, der Mittelwert im Zeitintervall mit
$x(t) \neq 0$ sowie die Klassifikation als Energie- oder Leistungssignal (mit
Energie- bzw.\ Leistungswert).

\subsubsection*{a) $x(t) = \sin(\omega t)$, $\omega = 2\pi\frac{1}{T}$}

\begin{beispiel}[Reine Sinusschwingung]

\textbf{Gegeben:} $x(t) = \sin(\omega t)$ mit $\omega = 2\pi/T$.

\textbf{Gesucht:} Periodizität, Symmetrie, Mittelwert, Klassifikation.

\textbf{Rechnung:} $x(t)$ ist \emph{periodisch} mit Periode $T$ und
\emph{ungerade}. Über eine Periode ist $\bar{x} = 0$. Als periodisches Signal
mit $E_x = \infty$ liegt ein Leistungssignal vor:

$$
P_x = \frac{1}{T}\int_0^T \sin^2(\omega t)\,\mathrm{d}t = \frac{1}{2}.
$$

\textbf{Lösung:}
$\boxed{\;\text{periodisch, ungerade},\; \bar{x} = 0,\;
P_x = \tfrac12 \Rightarrow \text{Leistungssignal.}\;}$

\end{beispiel}

\subsubsection*{b) $x(t) = \operatorname{rect}\!\left(\frac{t}{2T}\right)\cdot\frac{t}{2T}$}

\begin{beispiel}[Einseitige Rampe]

\textbf{Gegeben:} $x(t) = \operatorname{rect}\!\left(\dfrac{t}{2T}\right)
\cdot \dfrac{t}{2T}$.

\textbf{Gesucht:} Periodizität, Symmetrie, Mittelwert, Klassifikation.

\textbf{Rechnung:} Die Rechteckfunktion ist $1$ für $0 \leq t < 2T$; dort ist
$x(t) = \dfrac{t}{2T}$, also eine Rampe von $0$ auf $1$. Sonst $0$. Das Signal
ist \emph{nicht periodisch} und weist \emph{keine} Symmetrie auf. Mittelwert
über den Träger $[0, 2T)$:

$$
\bar{x} = \frac{1}{2T}\int_0^{2T} \frac{t}{2T}\,\mathrm{d}t
        = \frac{1}{2T}\cdot\frac{1}{2T}\cdot\frac{(2T)^2}{2} = \frac{1}{2}.
$$

Endliche Dauer, also Energiesignal:

$$
E_x = \int_0^{2T} \left(\frac{t}{2T}\right)^2 \mathrm{d}t
    = \frac{1}{4T^2}\cdot\frac{(2T)^3}{3} = \frac{2T}{3}.
$$

\begin{center}
\begin{tikzpicture}
\begin{axis}[
  width=10cm, height=4cm,
  xmin=-1, xmax=3, ymin=-0.3, ymax=1.3,
  axis lines=center, xlabel={$t/T$}, ylabel={$x(t)$},
  xtick={2}, ytick={1}, yticklabels={$1$},
  every axis plot/.append style={very thick, blue},
]
\addplot coordinates {(-1,0) (0,0) (2,1) (2,0) (3,0)};
\end{axis}
\end{tikzpicture}
\end{center}

\textbf{Lösung:}
$\boxed{\;\text{nicht periodisch, keine Symmetrie},\; \bar{x} = \tfrac12,\;
E_x = \tfrac{2T}{3} \Rightarrow \text{Energiesignal.}\;}$

\end{beispiel}

\subsubsection*{c) $x(t) = \left[\epsilon(t+T) - \epsilon(t-T)\right]\frac{|t|}{T}$}

\begin{beispiel}[Symmetrisches Dreiecksignal]

\textbf{Gegeben:} $x(t) = \left[\epsilon(t+T) - \epsilon(t-T)\right]\dfrac{|t|}{T}$.

\textbf{Gesucht:} Periodizität, Symmetrie, Mittelwert, Klassifikation.

\textbf{Rechnung:} Die Klammer ist $1$ für $-T \leq t < T$; dort ist
$x(t) = \dfrac{|t|}{T}$ (V-Form, $0$ bei $t = 0$, $1$ bei $t = \pm T$). Das
Signal ist \emph{nicht periodisch}, aber \emph{gerade} (wegen $|t|$).
Mittelwert über $[-T, T)$:

$$
\bar{x} = \frac{1}{2T}\int_{-T}^{T} \frac{|t|}{T}\,\mathrm{d}t
        = \frac{1}{2T}\cdot\frac{2}{T}\cdot\frac{T^2}{2} = \frac{1}{2}.
$$

Energiesignal mit

$$
E_x = \int_{-T}^{T} \frac{t^2}{T^2}\,\mathrm{d}t
    = \frac{2}{T^2}\cdot\frac{T^3}{3} = \frac{2T}{3}.
$$

\begin{center}
\begin{tikzpicture}
\begin{axis}[
  width=10cm, height=4cm,
  xmin=-2.2, xmax=2.2, ymin=-0.3, ymax=1.3,
  axis lines=center, xlabel={$t/T$}, ylabel={$x(t)$},
  xtick={-1,1}, ytick={1}, yticklabels={$1$},
  every axis plot/.append style={very thick, blue},
]
\addplot coordinates {(-2.2,0) (-1,0) (-1,1) (0,0) (1,1) (1,0) (2.2,0)};
\end{axis}
\end{tikzpicture}
\end{center}

\textbf{Lösung:}
$\boxed{\;\text{nicht periodisch, gerade},\; \bar{x} = \tfrac12,\;
E_x = \tfrac{2T}{3} \Rightarrow \text{Energiesignal.}\;}$

\end{beispiel}

\subsubsection*{d) $x(t) = \operatorname{si}\!\left(\frac{t}{T}\right)$}

\begin{beispiel}[Spaltfunktion (si-Funktion)]

\textbf{Gegeben:} $x(t) = \operatorname{si}\!\left(\dfrac{t}{T}\right)
= \dfrac{\sin(t/T)}{t/T}$.

\textbf{Gesucht:} Periodizität, Symmetrie, Mittelwert, Klassifikation.

\textbf{Rechnung:} Die si-Funktion ist \emph{nicht periodisch} (nur ihre
Nullstellen liegen periodisch) und \emph{gerade}. Sie klingt für
$|t| \to \infty$ ab; ihr zeitlicher Mittelwert verschwindet:
$\bar{x} = 0$. Mit der Substitution $u = t/T$ und
$\int_{-\infty}^{\infty}\frac{\sin^2 u}{u^2}\,\mathrm{d}u = \pi$ folgt die
endliche Energie:

$$
E_x = \int_{-\infty}^{\infty} \operatorname{si}^2\!\left(\frac{t}{T}\right)
      \mathrm{d}t
    = T\int_{-\infty}^{\infty} \frac{\sin^2 u}{u^2}\,\mathrm{d}u
    = \pi T.
$$

\begin{center}
\begin{tikzpicture}
\begin{axis}[
  width=11cm, height=4.5cm,
  xmin=-16, xmax=16, ymin=-0.4, ymax=1.2,
  axis lines=center, xlabel={$t/T$}, ylabel={$x(t)$},
  ytick={1}, yticklabels={$1$},
  samples=300, domain=-16:16,
  every axis plot/.append style={very thick, blue},
]
\addplot {sin(deg(x))/x};
\end{axis}
\end{tikzpicture}
\end{center}

\textbf{Lösung:}
$\boxed{\;\text{nicht periodisch, gerade},\; \bar{x} = 0,\;
E_x = \pi T \Rightarrow \text{Energiesignal.}\;}$

\end{beispiel}

% ================================================================
\subsection*{Aufgabe 1.7 – Komplexwertige Signale}
% ================================================================

Für die folgenden komplexwertigen Signale
(Abschnitt~\ref{sec:komplexwertige-signale}) sind Real- und Imaginärteil,
Periodendauer $T$, Frequenz $f$, Mittelwert und Leistung zu bestimmen. Für ein
komplexwertiges periodisches Signal gilt

$$
\bar{\underline{x}} = \frac{1}{T}\int_0^T \underline{x}(t)\,\mathrm{d}t,
\qquad
P_{\underline{x}} = \frac{1}{T}\int_0^T |\underline{x}(t)|^2\,\mathrm{d}t.
$$

Wegen $|\underline{x}|^2 = \operatorname{Re}\{\underline{x}\}^2
+ \operatorname{Im}\{\underline{x}\}^2$ und
$\underline{x} = \operatorname{Re}\{\underline{x}\}
+ j\operatorname{Im}\{\underline{x}\}$ gilt für die Mittelwerte und Leistungen
stets

$$
\bar{\underline{x}} = \overline{\operatorname{Re}\{\underline{x}\}}
+ j\,\overline{\operatorname{Im}\{\underline{x}\}},
\qquad
P_{\underline{x}} = P_{\operatorname{Re}\{\underline{x}\}}
+ P_{\operatorname{Im}\{\underline{x}\}}.
$$

\subsubsection*{a) $\underline{x}(t) = \sqrt{2}\,e^{j\left(\frac{2\pi t}{2\,\mathrm{ms}} - \frac{\pi}{2}\right)}$}

\begin{beispiel}[Komplexes Exponentialsignal mit konstantem Betrag]

\textbf{Gegeben:}
$\underline{x}(t) = \sqrt{2}\,e^{j\left(\frac{2\pi t}{2\,\mathrm{ms}}
- \frac{\pi}{2}\right)}$.

\textbf{Gesucht:} $\operatorname{Re}$, $\operatorname{Im}$, $T$, $f$,
Mittelwert, Leistung.

\textbf{Rechnung:} Aus $\omega = \dfrac{2\pi}{2\,\mathrm{ms}}
= \dfrac{\pi}{1\,\mathrm{ms}}$ folgt $T = 2\,\mathrm{ms}$, $f = 500\,\mathrm{Hz}$.
Mit $\cos(\omega t - \tfrac{\pi}{2}) = \sin(\omega t)$ und
$\sin(\omega t - \tfrac{\pi}{2}) = -\cos(\omega t)$:

$$
\operatorname{Re}\{\underline{x}(t)\} = \sqrt{2}\,\sin(\omega t),
\qquad
\operatorname{Im}\{\underline{x}(t)\} = -\sqrt{2}\,\cos(\omega t).
$$

Der Betrag $|\underline{x}| = \sqrt{2}$ ist konstant, das Signal rotiert einmal
pro Periode:

$$
\bar{\underline{x}} = 0, \qquad
P_{\underline{x}} = |\underline{x}|^2 = 2.
$$

Kontrolle über Real- und Imaginärteil: beide sind Sinusgrößen der Amplitude
$\sqrt{2}$, also
$P_{\operatorname{Re}} = P_{\operatorname{Im}} = \dfrac{(\sqrt{2})^2}{2} = 1$
und $\overline{\operatorname{Re}} = \overline{\operatorname{Im}} = 0$; damit
$P_{\underline{x}} = 1 + 1 = 2$ \checkmark.

\textbf{Lösung:}

$$
\boxed{\;\operatorname{Re} = \sqrt{2}\sin(\omega t),\;
\operatorname{Im} = -\sqrt{2}\cos(\omega t);\;
T = 2\,\mathrm{ms},\; f = 500\,\mathrm{Hz},\;
\bar{\underline{x}} = 0,\; P_{\underline{x}} = 2.\;}
$$

\end{beispiel}

\subsubsection*{b) $\underline{x}(t) = j + 2e^{j\left(\frac{2\pi t}{4\,\mathrm{ms}} + \frac{\pi}{4}\right)}$}

\begin{beispiel}[Komplexes Exponentialsignal mit Gleichanteil]

\textbf{Gegeben:}
$\underline{x}(t) = j + 2e^{j\left(\frac{2\pi t}{4\,\mathrm{ms}}
+ \frac{\pi}{4}\right)}$.

\textbf{Gesucht:} $\operatorname{Re}$, $\operatorname{Im}$, $T$, $f$,
Mittelwert, Leistung.

\textbf{Rechnung:} Aus $\omega = \dfrac{2\pi}{4\,\mathrm{ms}}
= \dfrac{\pi}{2\,\mathrm{ms}}$ folgt $T = 4\,\mathrm{ms}$, $f = 250\,\mathrm{Hz}$.
Mit $\theta = \omega t + \tfrac{\pi}{4}$:

$$
\operatorname{Re}\{\underline{x}(t)\} = 2\cos\theta,
\qquad
\operatorname{Im}\{\underline{x}(t)\} = 1 + 2\sin\theta.
$$

Der rotierende Term mittelt sich zu $0$, der konstante Anteil $j$ bleibt:

$$
\bar{\underline{x}} = j.
$$

Für die Leistung mit $|\underline{x}|^2 = (2\cos\theta)^2
+ (1 + 2\sin\theta)^2 = 5 + 4\sin\theta$:

$$
P_{\underline{x}} = \frac{1}{T}\int_0^T (5 + 4\sin\theta)\,\mathrm{d}t = 5
= \underbrace{|j|^2}_{1} + \underbrace{|2|^2}_{4}.
$$

Kontrolle: $P_{\operatorname{Re}} = \overline{4\cos^2\theta} = 2$,
$P_{\operatorname{Im}} = \overline{(1 + 2\sin\theta)^2} = 1 + 2 = 3$; Summe
$= 5$ \checkmark. Ferner $\overline{\operatorname{Re}} = 0$,
$\overline{\operatorname{Im}} = 1$.

\textbf{Lösung:}

$$
\boxed{\;\operatorname{Re} = 2\cos\theta,\;
\operatorname{Im} = 1 + 2\sin\theta,\;\theta = \omega t + \tfrac{\pi}{4};\;
T = 4\,\mathrm{ms},\; f = 250\,\mathrm{Hz},\;
\bar{\underline{x}} = j,\; P_{\underline{x}} = 5.\;}
$$

\end{beispiel}

\subsubsection*{c) $\underline{x}(t) = 5 - j + 3\exp\!\left(\frac{j2\pi t}{10\,\mathrm{ms}}\right)$}

\begin{beispiel}[Komplexes Exponentialsignal mit komplexem Gleichanteil]

\textbf{Gegeben:}
$\underline{x}(t) = 5 - j + 3\exp\!\left(\dfrac{j2\pi t}{10\,\mathrm{ms}}\right)$.

\textbf{Gesucht:} $\operatorname{Re}$, $\operatorname{Im}$, $T$, $f$,
Mittelwert, Leistung.

\textbf{Rechnung:} Aus $\omega = \dfrac{2\pi}{10\,\mathrm{ms}}$ folgt
$T = 10\,\mathrm{ms}$, $f = 100\,\mathrm{Hz}$. Mit
$3e^{j\omega t} = 3\cos\omega t + j\,3\sin\omega t$:

$$
\operatorname{Re}\{\underline{x}(t)\} = 5 + 3\cos\omega t,
\qquad
\operatorname{Im}\{\underline{x}(t)\} = -1 + 3\sin\omega t.
$$

Der rotierende Term mittelt sich zu $0$:

$$
\bar{\underline{x}} = 5 - j.
$$

Leistung (Gleichanteil plus rotierender Term):

$$
P_{\underline{x}} = |5 - j|^2 + |3|^2 = (25 + 1) + 9 = 35.
$$

Kontrolle: $P_{\operatorname{Re}} = \overline{(5 + 3\cos\omega t)^2}
= 25 + \tfrac{9}{2} = 29{,}5$ und
$P_{\operatorname{Im}} = \overline{(-1 + 3\sin\omega t)^2}
= 1 + \tfrac{9}{2} = 5{,}5$; Summe $= 35$ \checkmark. Ferner
$\overline{\operatorname{Re}} = 5$, $\overline{\operatorname{Im}} = -1$.

\textbf{Lösung:}

$$
\boxed{\;\operatorname{Re} = 5 + 3\cos\omega t,\;
\operatorname{Im} = -1 + 3\sin\omega t;\;
T = 10\,\mathrm{ms},\; f = 100\,\mathrm{Hz},\;
\bar{\underline{x}} = 5 - j,\; P_{\underline{x}} = 35.\;}
$$

Der Mittelwert des komplexen Signals ist stets die Summe der Mittelwerte von
Real- und Imaginärteil (letzterer mit $j$ gewichtet), die Gesamtleistung die
Summe der Leistungen von Real- und Imaginärteil.

\end{beispiel}

\section{Übung 2 – Diskrete Signale, Korrelation und Orthogonalität}

Dieses Aufgabenblatt behandelt zeitdiskrete Elementarsignale, ihre
Klassifikation als Energie- oder Leistungssignale sowie die Kreuz- und
Autokorrelation zeit\-diskreter und zeitkontinuierlicher Signale. Grundlage
sind die Abschnitte zu den elementaren diskreten Signalen
(\ref{sec:elementare-diskrete-signale}), zur Kreuzkorrelation
(\ref{sec:kreuzkorrelation}), zur Autokorrelation
(\ref{sec:autokorrelation}) und zu den Orthogonalsystemen
(\ref{sec:orthogonalsysteme}).

\begin{bemerkung}[Verwendete \texttt{rect}-Konvention]
Das Skript definiert den Rechteckimpuls asymmetrisch (Abschnitt~\ref{sec:elementare-analoge-signale}):
$$
\operatorname{rect}(u) = \begin{cases} 1 & , \quad 0 \leq u < 1 \\ 0 & , \quad \text{sonst.} \end{cases}
$$
Diese Konvention wird auf diesem Blatt durchgängig verwendet.
\end{bemerkung}

% ================================================================
\subsection*{Vorrechenaufgabe 2 – Zeitdiskrete Signale}
% ================================================================

Gegeben sind die Signale
$$
x_1[k] = \delta(k-1) + 2\,\delta(k-2) + \delta(k-3)
\qquad\text{und}\qquad
x_2[k] = \operatorname{rect}\!\left(\frac{k-2}{3}\right).
$$

% ---- a) --------------------------------------------------------
\subsubsection*{a) Skizze von $x_1[k]$ und $x_2[k]$}

\begin{beispiel}[Wertetabelle und Stem-Plot der Signale]

\textbf{Gegeben:} $x_1[k]=\delta(k-1)+2\delta(k-2)+\delta(k-3)$ und
$x_2[k]=\operatorname{rect}\!\big(\tfrac{k-2}{3}\big)$.

\textbf{Gesucht:} Darstellung beider Signale als Funktion des Index $k$.

\textbf{Formel} (diskrete Elementarsignale, Abschnitt~\ref{sec:elementare-diskrete-signale}):
Der Einheitsimpuls $\delta[k-k_0]$ ist genau bei $k=k_0$ gleich $1$; der
Rechteckimpuls $\operatorname{rect}\!\big(\tfrac{k-2}{3}\big)$ ist gleich $1$,
wenn $0 \leq \tfrac{k-2}{3} < 1$, also für $2 \leq k < 5$.

\textbf{Rechnung:}
$$
x_1[k]:\quad x_1[1]=1,\; x_1[2]=2,\; x_1[3]=1,\; \text{sonst }0.
$$
$$
x_2[k]:\quad x_2[2]=x_2[3]=x_2[4]=1,\; \text{sonst }0.
$$

\textbf{Lösung:}

\begin{center}
\begin{tikzpicture}
\begin{axis}[
  width=7.5cm, height=4.2cm, name=plot1,
  axis lines=middle, xlabel={$k$}, ylabel={$x_1[k]$},
  xmin=-1, xmax=5, ymin=0, ymax=2.4,
  xtick={0,1,2,3,4}, ytick={1,2},
  ycomb, mark=*, mark options={fill=accent}, color=accent, very thick,
]
\addplot coordinates {(1,1) (2,2) (3,1)};
\end{axis}
\begin{axis}[
  at={(plot1.right of east)}, xshift=1.2cm, anchor=west,
  width=7.5cm, height=4.2cm,
  axis lines=middle, xlabel={$k$}, ylabel={$x_2[k]$},
  xmin=-1, xmax=6, ymin=0, ymax=1.6,
  xtick={0,1,2,3,4,5}, ytick={1},
  ycomb, mark=*, mark options={fill=exercisegreen-fr}, color=exercisegreen-fr, very thick,
]
\addplot coordinates {(2,1) (3,1) (4,1)};
\end{axis}
\end{tikzpicture}
\end{center}

\end{beispiel}

% ---- b) --------------------------------------------------------
\subsubsection*{b) Kreuzkorrelationsfunktion $\varphi_{x_1 x_2}[\varkappa]$}

\begin{beispiel}[Diskrete Kreuzkorrelation zweier endlicher Folgen]

\textbf{Gegeben:} $x_1[k]$ und $x_2[k]$ aus Teil~a).

\textbf{Gesucht:} $\varphi_{x_1 x_2}[\varkappa]$ und ihr Verlauf.

\textbf{Formel} (Kreuzkorrelation reeller Energiesignale, Abschnitt~\ref{sec:kreuzkorrelation}):
$$
\varphi_{x_1 x_2}[\varkappa] = \sum_{k=-\infty}^{\infty} x_1[k]\,x_2[k+\varkappa]
$$

\textbf{Rechnung:}
Der Summand ist nur dort ungleich null, wo $x_1[k]\neq 0$ ($k\in\{1,2,3\}$)
\emph{und} $x_2[k+\varkappa]\neq 0$ ($2\leq k+\varkappa\leq 4$). Man wertet
$\varphi_{x_1 x_2}$ für jede relevante Verschiebung aus:
$$
\begin{aligned}
\varkappa=-1:&\quad x_1[3]\,x_2[2] = 1\cdot 1 = 1,\\
\varkappa=0:&\quad x_1[2]\,x_2[2] + x_1[3]\,x_2[3] = 2\cdot 1 + 1\cdot 1 = 3,\\
\varkappa=1:&\quad x_1[1]\,x_2[2] + x_1[2]\,x_2[3] + x_1[3]\,x_2[4] = 1 + 2 + 1 = 4,\\
\varkappa=2:&\quad x_1[1]\,x_2[3] + x_1[2]\,x_2[4] = 1 + 2 = 3,\\
\varkappa=3:&\quad x_1[1]\,x_2[4] = 1\cdot 1 = 1.
\end{aligned}
$$
Kontrolle: $\sum_\varkappa \varphi_{x_1 x_2}[\varkappa] = 1+3+4+3+1 = 12
= \big(\textstyle\sum_k x_1[k]\big)\big(\sum_k x_2[k]\big) = 4\cdot 3$. \checkmark

\textbf{Lösung:}
$$
\boxed{\;\varphi_{x_1 x_2}[\varkappa] = \{\,\underset{\varkappa=-1}{1},\;3,\;\underset{\varkappa=1}{4},\;3,\;1\,\},\qquad
\text{sonst }0.\;}
$$

\begin{center}
\begin{tikzpicture}
\begin{axis}[
  width=9cm, height=4.5cm,
  axis lines=middle, xlabel={$\varkappa$}, ylabel={$\varphi_{x_1 x_2}[\varkappa]$},
  xmin=-3, xmax=5, ymin=0, ymax=4.6,
  xtick={-2,-1,0,1,2,3,4}, ytick={1,2,3,4},
  ycomb, mark=*, mark options={fill=black}, very thick,
]
\addplot coordinates {(-1,1) (0,3) (1,4) (2,3) (3,1)};
\end{axis}
\end{tikzpicture}
\end{center}

Das Maximum liegt bei $\varkappa=1$: Um dieses Maß muss $x_2$ nach links
verschoben werden, damit sein Träger $\{2,3,4\}$ mit dem (gewichteten)
Schwerpunkt von $x_1$ zur Deckung kommt.

\end{beispiel}

% ================================================================
\subsection*{Aufgabe 2.1 – Zeitdiskrete Signale zeichnen}
% ================================================================

% ---- a) --------------------------------------------------------
\subsubsection*{a) $x[k] = 2 - 2\,\epsilon[k]$}

\begin{beispiel}[Konstante minus Sprungfunktion]

\textbf{Gegeben:} $x[k] = 2 - 2\,\epsilon[k]$ mit dem diskreten Einheitssprung
$\epsilon[k]$ ($\epsilon[k]=1$ für $k\geq 0$, sonst $0$).

\textbf{Gesucht:} Verlauf von $x[k]$.

\textbf{Formel} (Fallunterscheidung des Einheitssprungs, Abschnitt~\ref{sec:elementare-diskrete-signale}).

\textbf{Rechnung:}
$$
x[k] = \begin{cases} 2 - 0 = 2 & , \quad k < 0 \\ 2 - 2 = 0 & , \quad k \geq 0. \end{cases}
$$

\textbf{Lösung:}
\begin{center}
\begin{tikzpicture}
\begin{axis}[
  width=11cm, height=4cm, axis lines=middle, xlabel={$k$}, ylabel={$x[k]$},
  xmin=-5, xmax=5, ymin=0, ymax=2.4, xtick={-4,-3,-2,-1,0,1,2,3,4}, ytick={1,2},
  ycomb, mark=*, mark options={fill=accent}, color=accent, very thick,
]
\addplot coordinates {(-4,2) (-3,2) (-2,2) (-1,2)};
\addplot[only marks, mark=*, black] coordinates {(0,0) (1,0) (2,0) (3,0) (4,0)};
\end{axis}
\end{tikzpicture}
\end{center}
$x[k]$ ist eine „fallende`` Sprungfunktion: konstant $2$ für $k<0$ und $0$ für
$k\geq 0$.
\end{beispiel}

% ---- b) --------------------------------------------------------
\subsubsection*{b) $x[k] = \epsilon[k] - \epsilon[k-4] + \epsilon[k-8] - \epsilon[k-12]$}

\begin{beispiel}[Überlagerung verschobener Sprungfunktionen]

\textbf{Gegeben:} $x[k] = \epsilon[k] - \epsilon[k-4] + \epsilon[k-8] - \epsilon[k-12]$.

\textbf{Gesucht:} Verlauf von $x[k]$.

\textbf{Formel:} Jeder Term $\pm\epsilon[k-k_0]$ schaltet bei $k=k_0$ um $\pm 1$ um.

\textbf{Rechnung:}
$$
x[k] = \begin{cases}
0 & , \quad k < 0 \\
1 & , \quad 0 \leq k \leq 3 \\
1-1 = 0 & , \quad 4 \leq k \leq 7 \\
1-1+1 = 1 & , \quad 8 \leq k \leq 11 \\
1-1+1-1 = 0 & , \quad k \geq 12.
\end{cases}
$$

\textbf{Lösung:} Zwei Rechteckblöcke der Höhe $1$ auf $[0,3]$ und $[8,11]$.
\begin{center}
\begin{tikzpicture}
\begin{axis}[
  width=13cm, height=4cm, axis lines=middle, xlabel={$k$}, ylabel={$x[k]$},
  xmin=-2, xmax=14, ymin=0, ymax=1.5, xtick={0,2,4,6,8,10,12}, ytick={1},
  ycomb, mark=*, mark options={fill=accent}, color=accent, very thick,
]
\addplot coordinates {(0,1) (1,1) (2,1) (3,1) (8,1) (9,1) (10,1) (11,1)};
\end{axis}
\end{tikzpicture}
\end{center}
\end{beispiel}

% ---- c) --------------------------------------------------------
\subsubsection*{c) $x[k] = a^{-k}\,\epsilon[-k]$ mit $a=\tfrac12$}

\begin{beispiel}[Linksseitiges Exponentialsignal]

\textbf{Gegeben:} $x[k] = a^{-k}\,\epsilon[-k]$ mit $a=\tfrac12$.

\textbf{Gesucht:} Verlauf von $x[k]$.

\textbf{Rechnung:} Es ist $\epsilon[-k]=1$ für $k\leq 0$ und
$a^{-k} = \left(\tfrac12\right)^{-k} = 2^{k}$. Somit
$$
x[k] = \begin{cases} 2^{k} & , \quad k \leq 0 \\ 0 & , \quad k > 0. \end{cases}
\qquad
x[0]=1,\; x[-1]=\tfrac12,\; x[-2]=\tfrac14,\; \ldots
$$

\textbf{Lösung:} Das Signal ist linksseitig und klingt für $k\to -\infty$
geometrisch ab.
\begin{center}
\begin{tikzpicture}
\begin{axis}[
  width=11cm, height=4.2cm, axis lines=middle, xlabel={$k$}, ylabel={$x[k]$},
  xmin=-6, xmax=3, ymin=0, ymax=1.2, xtick={-5,-4,-3,-2,-1,0,1,2}, ytick={0.5,1},
  ycomb, mark=*, mark options={fill=accent}, color=accent, very thick,
]
\addplot coordinates {(-5,0.03125) (-4,0.0625) (-3,0.125) (-2,0.25) (-1,0.5) (0,1)};
\end{axis}
\end{tikzpicture}
\end{center}
\end{beispiel}

% ================================================================
\subsection*{Aufgabe 2.2 – Sprung und Einheitsimpuls}
% ================================================================

Aus den Bildern liest man die beiden Folgen ab:
$$
x_1[k] = (-1)^k \;\; \text{für } 0\leq k\leq 7,
\qquad
x_2[k] = \begin{cases} +1 & , \; 0\leq k\leq 4 \\ -1 & , \; 5\leq k\leq 9 \\ +1 & , \; 10\leq k\leq 14 \\ 0 & , \; \text{sonst.} \end{cases}
$$

% ---- x1 --------------------------------------------------------
\subsubsection*{$x_1[k]$ -- alternierende Folge}

\begin{beispiel}[Darstellung der alternierenden Folge]

\textbf{Gegeben:} $x_1[k]=1,-1,1,-1,1,-1,1,-1$ für $k=0,\ldots,7$.

\textbf{Gesucht:} Darstellung i) mit Einheitsimpulsen, ii) mit Sprungfunktionen.

\textbf{Formel:} $\delta[k-n]$ „setzt`` einen Einzelwert; $\epsilon[k-n]$ schaltet
eine Stufe.

\textbf{Rechnung / Lösung:}

\emph{i) Einheitsimpulse.}
$$
x_1[k] = \sum_{n=0}^{7} (-1)^n\,\delta[k-n]
= \delta[k] - \delta[k-1] + \delta[k-2] - \cdots - \delta[k-7].
$$

\emph{ii) Sprungfunktionen.} Jeder Vorzeichenwechsel entspricht einer Stufe der
Höhe $\pm 2$; die letzte Stufe $+1$ beendet die Folge:
$$
x_1[k] = \epsilon[k] + 2\sum_{n=1}^{7} (-1)^{n}\,\epsilon[k-n] + \epsilon[k-8].
$$
Kontrolle bei $k=7$: $1 - 2 + 2 - 2 + 2 - 2 + 2 - 2 = -1$; bei $k\geq 8$ addiert
der Term $+\epsilon[k-8]$ die $+1$, sodass $x_1[k]=0$. \checkmark

\begin{center}
\begin{tikzpicture}
\begin{axis}[
  width=12cm, height=4.5cm, axis lines=middle, xlabel={$k$}, ylabel={$x_1[k]$},
  xmin=-1, xmax=9, ymin=-1.4, ymax=1.4, xtick={0,1,2,3,4,5,6,7}, ytick={-1,1},
  ycomb, mark=*, mark options={fill=accent}, color=accent, very thick,
]
\addplot coordinates {(0,1) (1,-1) (2,1) (3,-1) (4,1) (5,-1) (6,1) (7,-1)};
\end{axis}
\end{tikzpicture}
\end{center}
\end{beispiel}

% ---- x2 --------------------------------------------------------
\subsubsection*{$x_2[k]$ -- Blockfolge}

\begin{beispiel}[Darstellung der Blockfolge]

\textbf{Gegeben:} $x_2[k]=+1$ auf $[0,4]$, $-1$ auf $[5,9]$, $+1$ auf $[10,14]$.

\textbf{Gesucht:} Darstellung i) mit Einheitsimpulsen, ii) mit Sprungfunktionen.

\textbf{Rechnung / Lösung:}

\emph{i) Einheitsimpulse.}
$$
x_2[k] = \sum_{n=0}^{4}\delta[k-n] \;-\; \sum_{n=5}^{9}\delta[k-n] \;+\; \sum_{n=10}^{14}\delta[k-n].
$$

\emph{ii) Sprungfunktionen.} Die Blockgrenzen liegen bei $k=0,5,10,15$:
$$
x_2[k] = \epsilon[k] - 2\,\epsilon[k-5] + 2\,\epsilon[k-10] - \epsilon[k-15].
$$
Kontrolle: $0\!\leq\! k\!\leq\!4\!:1$;\; $5\!\leq\! k\!\leq\!9\!:1-2=-1$;\;
$10\!\leq\! k\!\leq\!14\!:1-2+2=1$;\; $k\geq 15\!:1-2+2-1=0$. \checkmark

\begin{center}
\begin{tikzpicture}
\begin{axis}[
  width=14cm, height=4.5cm, axis lines=middle, xlabel={$k$}, ylabel={$x_2[k]$},
  xmin=-1, xmax=16, ymin=-1.4, ymax=1.4,
  xtick={0,2,4,6,8,10,12,14}, ytick={-1,1},
  ycomb, mark=*, mark options={fill=accent}, color=accent, very thick,
]
\addplot coordinates {(0,1) (1,1) (2,1) (3,1) (4,1) (5,-1) (6,-1) (7,-1) (8,-1) (9,-1) (10,1) (11,1) (12,1) (13,1) (14,1)};
\end{axis}
\end{tikzpicture}
\end{center}
\end{beispiel}

% ================================================================
\subsection*{Aufgabe 2.3 – Signaleigenschaften}
% ================================================================

Für zeitdiskrete Signale gelten (Abschnitt~\ref{sec:norm-energie-leistung} und~\ref{sec:leistung}):
$$
E = \sum_{k=-\infty}^{\infty} |x[k]|^2,
\qquad
P = \lim_{N\to\infty} \frac{1}{2N+1}\sum_{k=-N}^{N} |x[k]|^2.
$$
Ein \emph{Energiesignal} hat $0<E<\infty$ (und $P=0$), ein \emph{Leistungssignal}
$0<P<\infty$ (und $E\to\infty$). Der Bereichsmittelwert der nichtperiodischen
Signale wird über $k\in[-4,5]$ (also $10$ Werte) gebildet:
$\bar{x} = \tfrac{1}{10}\sum_{k=-4}^{5} x[k]$.

% ---- a) --------------------------------------------------------
\subsubsection*{a) $x[k] = e^{ak}\,\epsilon[k]$, $a\in\mathbb{R}$}

\begin{beispiel}[Kausales Exponentialsignal]

\textbf{Gegeben:} $x[k]=e^{ak}\epsilon[k]$.

\textbf{Gesucht:} Energie/Leistung, Klassifikation und Bereichsmittelwert.

\textbf{Formel} (geometrische Summe, vgl.\ Hinweis):
$\sum_{k=0}^{\infty} q^{k} = \tfrac{1}{1-q}$ für $|q|<1$.

\textbf{Rechnung:} Mit $|x[k]|^2 = e^{2ak}$ für $k\geq 0$:
$$
E = \sum_{k=0}^{\infty} e^{2ak} = \sum_{k=0}^{\infty}\big(e^{2a}\big)^{k}
= \frac{1}{1-e^{2a}} \quad\text{für } a<0.
$$
Für $a=0$ ist $x[k]=\epsilon[k]$ mit $P=\lim_{N\to\infty}\tfrac{N+1}{2N+1}=\tfrac12$;
für $a>0$ wächst $x[k]$ unbeschränkt ($E\to\infty$, $P\to\infty$).

Bereichsmittelwert ($a\neq 0$):
$$
\bar{x} = \frac{1}{10}\sum_{k=0}^{5} e^{ak}
= \frac{1}{10}\cdot\frac{1-e^{6a}}{1-e^{a}}, \qquad
\bar{x}\big|_{a=0} = \frac{6}{10}=0{,}6.
$$

\textbf{Lösung:}
$$
\boxed{\;a<0:\ \text{Energiesignal},\ E=\tfrac{1}{1-e^{2a}};\quad
a=0:\ \text{Leistungssignal},\ P=\tfrac12;\quad
a>0:\ \text{weder noch.}\;}
$$
\end{beispiel}

% ---- b) --------------------------------------------------------
\subsubsection*{b) $x_1[k]=\delta[k]$ und $x_2[k]=A\,\delta[k-2]$}

\begin{beispiel}[Einzelimpulse]

\textbf{Gegeben:} $x_1[k]=\delta[k]$, $x_2[k]=A\,\delta[k-2]$.

\textbf{Gesucht:} Energie/Leistung und Bereichsmittelwert.

\textbf{Rechnung:}
$$
E_{x_1} = \sum_k |\delta[k]|^2 = 1,\qquad
E_{x_2} = \sum_k |A\,\delta[k-2]|^2 = A^2.
$$
Beide besitzen endliche Energie, also $P=0$. Bereichsmittelwerte:
$$
\bar{x}_1 = \frac{1}{10}\cdot 1 = 0{,}1,\qquad
\bar{x}_2 = \frac{1}{10}\cdot A = \frac{A}{10}.
$$

\textbf{Lösung:}
$$
\boxed{\;\text{Beides Energiesignale: } E_{x_1}=1,\; E_{x_2}=A^2,\; P=0.\;}
$$
\end{beispiel}

% ---- c) --------------------------------------------------------
\subsubsection*{c) $x[k] = \epsilon[1-k]$}

\begin{beispiel}[Rückwärtsgerichtete Sprungfunktion]

\textbf{Gegeben:} $x[k]=\epsilon[1-k]$.

\textbf{Gesucht:} Energie/Leistung und Bereichsmittelwert.

\textbf{Rechnung:} $\epsilon[1-k]=1$ für $1-k\geq 0$, d.h.\ $k\leq 1$; sonst $0$.
Das Signal ist für unendlich viele $k$ ungleich null $\Rightarrow E\to\infty$.
$$
P = \lim_{N\to\infty}\frac{1}{2N+1}\underbrace{(N+2)}_{k=-N,\ldots,1}
= \frac{1}{2}.
$$
Bereichsmittelwert (nichtzahl $k\leq 1$ im Fenster: $k=-4,\ldots,1$, also $6$ Werte):
$$
\bar{x} = \frac{1}{10}\cdot 6 = 0{,}6.
$$

\textbf{Lösung:}
$$
\boxed{\;\text{Leistungssignal},\ P=\tfrac12,\ \bar{x}=0{,}6.\;}
$$
\end{beispiel}

% ---- d) --------------------------------------------------------
\subsubsection*{d) $x[k] = \operatorname{Im}\big\{\sqrt{2}\,e^{j\pi/4}\,e^{j\Omega k}\big\}$}

\begin{beispiel}[Reelles Sinussignal]

\textbf{Gegeben:} $\Omega = \tfrac{2\pi f}{f_A}$ mit $f=1\,\text{kHz}$, $f_A=8\,\text{kHz}$,
also $\Omega=\tfrac{\pi}{4}$.

\textbf{Gesucht:} Energie/Leistung, Klassifikation, Mittelwert.

\textbf{Rechnung:} Mit $\sqrt{2}\,e^{j\pi/4}=1+j$ folgt
$$
x[k] = \operatorname{Im}\big\{(1+j)\,e^{j\pi k/4}\big\}
= \sqrt{2}\,\sin\!\Big(\frac{\pi(k+1)}{4}\Big).
$$
Das Signal ist mit $N=8$ periodisch ($\Omega N = 2\pi$), also ein
Leistungssignal. Die Leistung eines Sinus der Amplitude $\hat{x}=\sqrt2$ ist
$P=\hat{x}^2/2 = 1$; der Mittelwert über eine Periode verschwindet.

\textbf{Lösung:}
$$
\boxed{\;\text{Leistungssignal},\ P=1,\ \bar{x}=0,\ \text{periodisch mit } N=8.\;}
$$
\end{beispiel}

% ---- e) --------------------------------------------------------
\subsubsection*{e) $\underline{x}[k] = \sqrt{2}\,e^{j\pi/4}\,e^{j\Omega k}$}

\begin{beispiel}[Komplexes Exponentialsignal]

\textbf{Gegeben:} $\Omega=\tfrac{2\pi f}{f_A}=\tfrac{\pi}{4}$ (wie d).

\textbf{Gesucht:} Energie/Leistung, Klassifikation, Mittelwert.

\textbf{Rechnung:} Der Betrag ist konstant:
$$
|\underline{x}[k]|^2 = \big|\sqrt{2}\,e^{j\pi/4}\big|^2\cdot\big|e^{j\Omega k}\big|^2 = 2.
$$
Damit $E\to\infty$ und
$$
P = \lim_{N\to\infty}\frac{1}{2N+1}\sum_{k=-N}^{N} 2 = 2.
$$
Über eine Periode ($N=8$) mittelt sich der rotierende Zeiger zu null:
$\bar{\underline{x}} = \sqrt2\,e^{j\pi/4}\cdot\tfrac1N\sum_{k=0}^{7} e^{j\pi k/4}=0$.

\textbf{Lösung:}
$$
\boxed{\;\text{Leistungssignal},\ P=2,\ \bar{\underline{x}}=0,\ \text{periodisch mit } N=8.\;}
$$
\end{beispiel}

% ================================================================
\subsection*{Aufgabe 2.4 – Auto- und Kreuzkorrelation}
% ================================================================

Für Energiesignale gilt (Abschnitte~\ref{sec:kreuzkorrelation},~\ref{sec:autokorrelation})
$\varphi_{xy}(\tau)=\int_{-\infty}^{\infty} x(t)\,y(t+\tau)\,\mathrm{d}t$.

% ---- a) --------------------------------------------------------
\subsubsection*{a) $x_1(t)=\operatorname{rect}(t/T)$, $x_2(t)=\operatorname{rect}(t/2T)$}

\begin{beispiel}[Korrelation zweier Rechtecke unterschiedlicher Breite]

\textbf{Gegeben:} $x_1(t)=1$ auf $[0,T)$, $x_2(t)=1$ auf $[0,2T)$.

\textbf{Gesucht:} $\varphi_{x_1 x_1}(\tau)$ und $\varphi_{x_1 x_2}(\tau)$.

\textbf{Rechnung:} Die AKF eines Rechtecks der Breite $T$ ist ein Dreieck; die KKF
(Breiten $T$ und $2T$) ein Trapez mit Plateau der Breite $T$:
$$
\varphi_{x_1 x_1}(\tau) = \int x_1(t)\,x_1(t+\tau)\,\mathrm{d}t = T-|\tau|,\quad |\tau|\leq T,
$$
$$
\varphi_{x_1 x_2}(\tau) = \begin{cases}
T+\tau & , \; -T\leq \tau \leq 0 \\
T & , \; 0\leq \tau \leq T \\
2T-\tau & , \; T\leq \tau \leq 2T \\
0 & , \; \text{sonst.}
\end{cases}
$$

\textbf{Lösung:}
$$
\boxed{\;\varphi_{x_1 x_1}(\tau)=T-|\tau|\ (|\tau|\leq T);\quad
\varphi_{x_1 x_2}(\tau)\ \text{Trapez, Plateau } T \text{ auf } [0,T].\;}
$$
\end{beispiel}

% ---- b) --------------------------------------------------------
\subsubsection*{b) $x_1(t)=\cos(2\pi t)$, $x_2(t)=\sin(2\pi t)$}

\begin{beispiel}[Korrelation periodischer Signale]

\textbf{Gegeben:} $x_1(t)=\cos(2\pi t)$, $x_2(t)=\sin(2\pi t)$; Periode $T_0=1$.

\textbf{Gesucht:} $\tilde{\varphi}_{x_1 x_1}(\tau)$ und $\tilde{\varphi}_{x_1 x_2}(\tau)$.

\textbf{Formel} (periodische Korrelation, Abschnitt~\ref{sec:kreuzkorrelation}):
$\tilde{\varphi}_{xy}(\tau)=\tfrac{1}{T_0}\int_0^{T_0} x(t)\,y(t+\tau)\,\mathrm{d}t$.

\textbf{Rechnung:} Mit den Additionstheoremen
$\cos a\cos b=\tfrac12[\cos(a-b)+\cos(a+b)]$ bzw.\
$\sin a\cos b=\tfrac12[\sin(a-b)+\sin(a+b)]$ und Wegmittelung der doppelten
Frequenz:
$$
\tilde{\varphi}_{x_1 x_1}(\tau) = \tfrac12\cos(2\pi\tau),\qquad
\tilde{\varphi}_{x_1 x_2}(\tau) = \tfrac12\sin(2\pi\tau).
$$

\textbf{Lösung:}
$$
\boxed{\;\tilde{\varphi}_{x_1 x_1}(\tau)=\tfrac12\cos(2\pi\tau),\quad
\tilde{\varphi}_{x_1 x_2}(\tau)=\tfrac12\sin(2\pi\tau).\;}
$$
Wegen $\tilde{\varphi}_{x_1 x_2}(0)=0$ sind $\cos$ und $\sin$ orthogonal.
\end{beispiel}

% ---- c) --------------------------------------------------------
\subsubsection*{c) $x_1(t)=\operatorname{rect}\!\big(\tfrac{t-T}{T}\big)$, $x_2(t)=\operatorname{rect}\!\big(\tfrac{t+T}{2T}\big)$}

\begin{beispiel}[Korrelation verschobener Rechtecke]

\textbf{Gegeben:} $x_1(t)=1$ auf $[T,2T)$, $x_2(t)=1$ auf $[-T,T)$.

\textbf{Gesucht:} $\varphi_{x_1 x_1}(\tau)$ und $\varphi_{x_1 x_2}(\tau)$.

\textbf{Rechnung:} Die AKF ist verschiebungsinvariant, also identisch zu a):
$\varphi_{x_1 x_1}(\tau)=T-|\tau|$ für $|\tau|\leq T$. Die KKF ergibt sich als
Überlappung von $[T,2T)$ mit dem um $\tau$ verschobenen Träger $[-T-\tau, T-\tau)$
von $x_2$:
$$
\varphi_{x_1 x_2}(\tau) = \begin{cases}
3T+\tau & , \; -3T\leq \tau \leq -2T \\
T & , \; -2T\leq \tau \leq -T \\
-\tau & , \; -T\leq \tau \leq 0 \\
0 & , \; \text{sonst.}
\end{cases}
$$

\textbf{Lösung:}
$$
\boxed{\;\varphi_{x_1 x_1}(\tau)=T-|\tau|\ (|\tau|\leq T);\quad
\varphi_{x_1 x_2}\ \text{Trapez, Plateau } T \text{ auf } [-2T,-T].\;}
$$
Das Maximum bei negativem $\tau$ spiegelt die Zeitverschiebung der beiden
Rechtecke gegeneinander wider.
\end{beispiel}

% ================================================================
\subsection*{Aufgabe 2.5 – Kreuzkorrelation}
% ================================================================

\begin{beispiel}[KKF eines Treppensignals mit einem Exponential-Rechteck]

\textbf{Gegeben:}
$$
z(t) = \begin{cases} A & , \; 2T\leq t < 4T \\ 2A & , \; 4T\leq t < 6T \\ 0 & , \; \text{sonst,} \end{cases}
\qquad
w(t) = A\,\operatorname{rect}\!\Big(\frac{t}{2T}\Big)\,e^{-t/T}
= A\,e^{-t/T}\ \text{auf } [0,2T).
$$

\textbf{Gesucht:} $\varphi_{wz}(\tau) = \int_{-\infty}^{\infty} w(t)\,z(t+\tau)\,\mathrm{d}t$.

\textbf{Formel} (KKF, Abschnitt~\ref{sec:kreuzkorrelation}) mit
$\int e^{-t/T}\,\mathrm{d}t = -T\,e^{-t/T}$.

\textbf{Rechnung:} $w(t)$ ist nur auf $[0,2T)$ ungleich null. Das um $\tau$ nach
links verschobene $z(t+\tau)$ besitzt die Sprungstellen bei $t=2T-\tau$,
$4T-\tau$, $6T-\tau$. Fallweise Integration liefert:

\emph{Fall $0\leq\tau\leq 2T$} (nur der $A$-Teil überlappt, von $2T-\tau$ bis $2T$):
$$
\varphi_{wz}(\tau) = A^2\!\!\int_{2T-\tau}^{2T}\! e^{-t/T}\,\mathrm{d}t
= A^2 T\,e^{-2}\big(e^{\tau/T}-1\big).
$$

\emph{Fall $2T\leq\tau\leq 4T$} ($A$-Teil auf $[0,4T-\tau)$, $2A$-Teil auf $[4T-\tau,2T)$):
$$
\varphi_{wz}(\tau) = A^2 T\Big(1 - 2e^{-2} + e^{(\tau-4T)/T}\Big).
$$

\emph{Fall $4T\leq\tau\leq 6T$} (nur der $2A$-Teil auf $[0,6T-\tau)$):
$$
\varphi_{wz}(\tau) = 2A^2 T\Big(1 - e^{(\tau-6T)/T}\Big).
$$
Nahtstellen-Kontrolle: bei $\tau=2T$ liefern beide oberen Fälle
$A^2T(1-e^{-2})$, bei $\tau=4T$ beide unteren $2A^2T(1-e^{-2})$; an $\tau=0$ und
$\tau=6T$ ist $\varphi_{wz}=0$. \checkmark

\textbf{Lösung:}
$$
\boxed{\;\varphi_{wz}(\tau)=\begin{cases}
0 & , \; \tau\leq 0 \\
A^2 T\,e^{-2}\big(e^{\tau/T}-1\big) & , \; 0\leq\tau\leq 2T \\
A^2 T\big(1-2e^{-2}+e^{(\tau-4T)/T}\big) & , \; 2T\leq\tau\leq 4T \\
2A^2 T\big(1-e^{(\tau-6T)/T}\big) & , \; 4T\leq\tau\leq 6T \\
0 & , \; \tau\geq 6T.
\end{cases}\;}
$$
Das Maximum $\varphi_{wz}(4T)=2A^2T(1-e^{-2})\approx 1{,}73\,A^2T$ tritt auf, wenn
der starke $2A$-Block von $z$ genau das (bei $t=0$ große) Fenster von $w$
überdeckt.
\end{beispiel}

% ================================================================
\subsection*{Aufgabe 2.6 – Energie und Korrelation}
% ================================================================

% ---- a) --------------------------------------------------------
\subsubsection*{a) KKF von Gaußimpuls und Deltafunktion}

\begin{beispiel}[Kreuzkorrelation mit einer verschobenen Deltafunktion]

\textbf{Gegeben:} $x_1(t)=e^{-\pi t^2}$, $x_2(t)=\delta(t-t_0)$.

\textbf{Gesucht:} $\varphi_{x_1 x_2}(\tau)$.

\textbf{Formel} (KKF und Ausblendeigenschaft der Deltafunktion, Abschnitt~\ref{sec:dirac-impuls}):
$$
\varphi_{x_1 x_2}(\tau) = \int_{-\infty}^{\infty} x_1(t)\,\delta(t+\tau-t_0)\,\mathrm{d}t.
$$

\textbf{Rechnung:} Die Deltafunktion blendet $t=t_0-\tau$ aus:
$$
\varphi_{x_1 x_2}(\tau) = e^{-\pi(t_0-\tau)^2} = e^{-\pi(\tau-t_0)^2}.
$$

\textbf{Lösung:}
$$
\boxed{\;\varphi_{x_1 x_2}(\tau) = e^{-\pi(\tau-t_0)^2}.\;}
$$
Die Korrelation mit einer Deltafunktion verschiebt den Gaußimpuls lediglich nach
$\tau=t_0$ (Abtast-/Verschiebeeigenschaft).
\end{beispiel}

% ---- b) --------------------------------------------------------
\subsubsection*{b) Energie und AKF des Gaußimpulses}

\begin{beispiel}[Energie und Autokorrelation des Gaußimpulses]

\textbf{Gegeben:} $x(t)=e^{-\pi t^2}$.

\textbf{Gesucht:} Energie $E$ und AKF $\varphi_{xx}(\tau)$.

\textbf{Formel} (Hinweis-Integral): $\int_0^{\infty} e^{-a^2 x^2}\,\mathrm{d}x=\tfrac{\sqrt{\pi}}{2a}$,
somit $\int_{-\infty}^{\infty} e^{-a^2 x^2}\,\mathrm{d}x=\tfrac{\sqrt{\pi}}{a}$.

\textbf{Rechnung:} Mit $a^2=2\pi$, d.h.\ $a=\sqrt{2\pi}$:
$$
E = \int_{-\infty}^{\infty} e^{-2\pi t^2}\,\mathrm{d}t
= \frac{\sqrt{\pi}}{\sqrt{2\pi}} = \frac{1}{\sqrt{2}}.
$$
Für die AKF wird der Exponent quadratisch ergänzt:
$$
-\pi t^2 - \pi(t+\tau)^2 = -2\pi\Big(t+\tfrac{\tau}{2}\Big)^2 - \tfrac{\pi}{2}\tau^2,
$$
$$
\varphi_{xx}(\tau) = e^{-\pi\tau^2/2}\!\int_{-\infty}^{\infty}\! e^{-2\pi(t+\tau/2)^2}\,\mathrm{d}t
= \frac{1}{\sqrt{2}}\,e^{-\pi\tau^2/2}.
$$

\textbf{Lösung:}
$$
\boxed{\;E = \frac{1}{\sqrt{2}}\approx 0{,}707,\qquad
\varphi_{xx}(\tau) = \frac{1}{\sqrt{2}}\,e^{-\pi\tau^2/2}.\;}
$$
Es gilt die Probe $\varphi_{xx}(0)=E$; die AKF eines Gaußimpulses ist wieder ein
(breiterer) Gaußimpuls.
\end{beispiel}

% ================================================================
\subsection*{Aufgabe 2.7 – Autokorrelation und Orthogonalität diskreter Signale}
% ================================================================

Aus den Bildern liest man (auf $k=0$ ausgerichtet) ab:
$$
\begin{aligned}
x_1[k] &= [\,1,1,1,1,1,1,1,1\,], & k&=0,\ldots,7, \\
x_2[k] &= [\,1,2,3,4,3,2,1\,], & k&=0,\ldots,6, \\
x_3[k] &= [\,1,1,1,1,-1,-1,-1,-1\,], & k&=0,\ldots,7, \\
x_4[k] &= [\,0{,}5,\,0{,}5,\,-0{,}5,\,-0{,}5,\,-0{,}5,\,-0{,}5,\,0{,}5,\,0{,}5\,], & k&=0,\ldots,7.
\end{aligned}
$$

% ---- AKFs ------------------------------------------------------
\subsubsection*{Autokorrelationsfunktionen}

\begin{beispiel}[AKF der vier Signale]

\textbf{Gegeben:} $x_1,\ldots,x_4$ wie oben.

\textbf{Gesucht:} $\varphi_{x_i x_i}[\varkappa]=\sum_k x_i[k]\,x_i[k+\varkappa]$ (Abschnitt~\ref{sec:autokorrelation}).

\textbf{Rechnung / Lösung:} (Die AKF ist gerade, es genügt $\varkappa\geq 0$.)

\emph{$x_1$ (Rechteck, $8$ Einsen):} Dreieck-AKF
$$
\varphi_{x_1 x_1}[\varkappa] = 8 - |\varkappa|,\quad |\varkappa|\leq 7.
$$

\emph{$x_2$ (Dreieck):}
$$
\varphi_{x_2 x_2}[\varkappa]\big|_{\varkappa=0,\ldots,6} = 44,\,40,\,31,\,20,\,10,\,4,\,1.
$$

\emph{$x_3$ (Sprung $+/-$):}
$$
\varphi_{x_3 x_3}[\varkappa]\big|_{\varkappa=0,\ldots,7} = 8,\,5,\,2,\,-1,\,-4,\,-3,\,-2,\,-1.
$$

\emph{$x_4$:}
$$
\varphi_{x_4 x_4}[\varkappa]\big|_{\varkappa=0,\ldots,7} = 2,\;0{,}75,\;-0{,}5,\;-0{,}75,\;-1,\;-0{,}25,\;0{,}5,\;0{,}25.
$$

\begin{center}
\begin{tikzpicture}
\begin{axis}[
  width=7.5cm, height=4cm, name=akf1,
  axis lines=middle, xlabel={$\varkappa$}, ylabel={$\varphi_{x_1 x_1}$},
  xmin=-8, xmax=8, ymin=0, ymax=9, xtick={-6,-4,-2,0,2,4,6}, ytick={4,8},
  ycomb, mark=*, mark size=1.2pt, mark options={fill=accent}, color=accent, thick,
]
\addplot coordinates {(-7,1)(-6,2)(-5,3)(-4,4)(-3,5)(-2,6)(-1,7)(0,8)(1,7)(2,6)(3,5)(4,4)(5,3)(6,2)(7,1)};
\end{axis}
\begin{axis}[
  at={(akf1.right of east)}, xshift=1.3cm, anchor=west,
  width=7.5cm, height=4cm,
  axis lines=middle, xlabel={$\varkappa$}, ylabel={$\varphi_{x_3 x_3}$},
  xmin=-8, xmax=8, ymin=-5, ymax=9, xtick={-6,-4,-2,0,2,4,6}, ytick={-4,4,8},
  ycomb, mark=*, mark size=1.2pt, mark options={fill=satzred-fr}, color=satzred-fr, thick,
]
\addplot coordinates {(-7,-1)(-6,-2)(-5,-3)(-4,-4)(-3,-1)(-2,2)(-1,5)(0,8)(1,5)(2,2)(3,-1)(4,-4)(5,-3)(6,-2)(7,-1)};
\end{axis}
\end{tikzpicture}
\end{center}

Alle AKF besitzen ihr Maximum bei $\varkappa=0$ (dort gleich der Signalenergie:
$8,\,44,\,8,\,2$).
\end{beispiel}

% ---- Orthogonalität --------------------------------------------
\subsubsection*{Orthogonalität}

\begin{beispiel}[Welche Signale sind orthogonal?]

\textbf{Gegeben:} $x_1,\ldots,x_4$.

\textbf{Gesucht:} Alle orthogonalen Paare.

\textbf{Formel} (Orthogonalität, Abschnitt~\ref{sec:orthogonalsysteme}):
$$
x_i \perp x_j \quad\Longleftrightarrow\quad
\varphi_{x_i x_j}[0] = \sum_k x_i[k]\,x_j[k] = 0.
$$

\textbf{Rechnung:} (Ausrichtung an $k=0$)
$$
\begin{aligned}
\langle x_1,x_3\rangle &= 4\cdot 1 + 4\cdot(-1) = 0, &
\langle x_1,x_4\rangle &= 2\cdot 0{,}5 + 4\cdot(-0{,}5) + 2\cdot 0{,}5 = 0, \\
\langle x_3,x_4\rangle &= 0{,}5+0{,}5-0{,}5-0{,}5+0{,}5+0{,}5-0{,}5-0{,}5 = 0, &
\langle x_1,x_2\rangle &= 1+2+3+4+3+2+1 = 16, \\
\langle x_2,x_3\rangle &= 1+2+3+4-3-2-1 = 4, &
\langle x_2,x_4\rangle &= -4.
\end{aligned}
$$

\textbf{Lösung:}
$$
\boxed{\;x_1\perp x_3,\quad x_1\perp x_4,\quad x_3\perp x_4
\;\Rightarrow\; \{x_1,x_3,x_4\}\ \text{sind paarweise orthogonal.}\;}
$$
$x_2$ ist zu keinem der anderen Signale orthogonal. Anschaulich hat $x_1$
(Gleichanteil) mit den vorzeichensymmetrischen Signalen $x_3$ und $x_4$ keinen
Überlapp, während der triangelförmige, rein positive Verlauf von $x_2$ stets
einen von null verschiedenen Beitrag liefert.
\end{beispiel}

% ================================================================
\subsection*{Aufgabe 2.8 – Log-Swept Chirp Signal}
% ================================================================

% ---- a) --------------------------------------------------------
\subsubsection*{a) Instantane Frequenz}

\begin{beispiel}[Herleitung der instantanen Frequenz]

\textbf{Gegeben:} Instantane Phase
$$
\varphi(t) = \frac{2\pi f_1 T}{\ln(f_2/f_1)}\left[\exp\!\Big(\frac{\ln(f_2/f_1)\,t}{T}\Big) - 1\right].
$$

\textbf{Gesucht:} Instantane (Kreis-)Frequenz $\omega(t)=\varphi'(t)$ bzw.\ $f(t)$.

\textbf{Rechnung:} Ableiten nach $t$ (innere Ableitung $\tfrac{\ln(f_2/f_1)}{T}$):
$$
\omega(t) = \varphi'(t)
= \frac{2\pi f_1 T}{\ln(f_2/f_1)}\cdot\frac{\ln(f_2/f_1)}{T}\,\exp\!\Big(\frac{\ln(f_2/f_1)\,t}{T}\Big)
= 2\pi f_1 \left(\frac{f_2}{f_1}\right)^{t/T}.
$$

\textbf{Lösung:}
$$
\boxed{\;f(t) = \frac{\omega(t)}{2\pi} = f_1\left(\frac{f_2}{f_1}\right)^{t/T}.\;}
$$
Es gilt $f(0)=f_1$ und $f(T)=f_2$; die Frequenz durchläuft den Bereich
$[f_1,f_2]$ \emph{exponentiell} (logarithmischer Sweep).
\end{beispiel}

% ---- b) --------------------------------------------------------
\subsubsection*{b) Erzeugung des Chirp-Signals}

\begin{beispiel}[Log-Sweep von 5\,Hz auf 15\,Hz]

\textbf{Gegeben / Gesucht:} Diskretes Chirp-Signal $x(t_k)=A\sin(\varphi(t_k))$
für $t_k = 0,\,\tfrac{1}{f_A},\,\ldots,\,T$; grafische Darstellung.

\textbf{Rechnung (Python):}
\begin{lstlisting}[style=python]
import numpy as np
import matplotlib.pyplot as plt

f1, f2, fA, T, A = 5.0, 15.0, 150.0, 2.0, 1.0
t = np.arange(0, T, 1.0/fA)          # diskrete Zeitachse t_k

# instantane Phase (log-swept chirp)
L = np.log(f2/f1)
phi = (2*np.pi*f1*T/L) * (np.exp(L*t/T) - 1.0)
x = A*np.sin(phi)

plt.figure(figsize=(9,3))
plt.plot(t, x)
plt.xlabel("t / s"); plt.ylabel("x(t)")
plt.title("Logarithmischer Frequenz-Sweep 5 Hz -> 15 Hz")
plt.tight_layout(); plt.show()
\end{lstlisting}

\textbf{Lösung:} Man erhält ein Sinussignal, dessen Schwingungen zum Ende hin
sichtbar dichter werden (die Momentanfrequenz steigt von $5$ auf $15$\,Hz an).
\end{beispiel}

% ---- c) --------------------------------------------------------
\subsubsection*{c) Audio-Chirp und Autokorrelation}

\begin{beispiel}[Audio-Chirp von 500\,Hz auf 2\,kHz]

\textbf{Gegeben / Gesucht:} Moduliertes Signal, seine Autokorrelation sowie
Speicherung als wav-Datei.

\textbf{Rechnung (Python):}
\begin{lstlisting}[style=python]
import numpy as np
import matplotlib.pyplot as plt
from scipy.io import wavfile

f1, f2, fA, T = 500.0, 2000.0, 8000.0, 2.0
t = np.arange(0, T, 1.0/fA)
L = np.log(f2/f1)
phi = (2*np.pi*f1*T/L) * (np.exp(L*t/T) - 1.0)
x = np.sin(phi)

# Autokorrelation (volle, symmetrische Faltung mit gespiegeltem Signal)
akf = np.correlate(x, x, mode="full")
lags = np.arange(-len(x)+1, len(x))

plt.plot(lags, akf); plt.xlabel("Verschiebung"); plt.ylabel("phi_xx")
plt.title("Autokorrelation des Log-Chirps"); plt.show()

# als 16-bit wav speichern
wavfile.write("chirp.wav", int(fA), (x*32767).astype(np.int16))
\end{lstlisting}

\textbf{Lösung:} Die Autokorrelation besitzt einen \emph{scharfen, schmalen Peak}
bei der Verschiebung $\varkappa=0$ und ist außerhalb nahezu null. Ein Chirp
ist damit -- trotz seiner zeitlichen Ausdehnung -- \emph{selbst\-un\-ähnlich} bei
Verschiebung; genau diese Eigenschaft nutzt man in der Messtechnik zur
Impulskompression (die AKF wirkt wie eine kurze Anregung mit hoher Energie). Beim
Anhören steigt der Ton hörbar von $500$\,Hz auf $2$\,kHz an.
\end{beispiel}

% ================================================================
\subsection*{Aufgabe 2.9 – Perfekte Sequenzen}
% ================================================================

\begin{beispiel}[Zyklische Korrelation einer perfekten Sequenz]

\textbf{Gegeben:} Die Folge (Länge $N=13$, $k=0,\ldots,12$)
$$
x[k] = [\,0,0,1,0,1,1,1,-1,-1,0,1,-1,1\,],
$$
sechsfach periodisch fortgesetzt.

\textbf{Gesucht:} Die von null verschiedenen Werte der (zyklischen)
Kreuz\-/Autokorrelation
$\varphi_{xx}(m)=\sum_{k=0}^{12} x[k]\,x[(k+m)_{\bmod 13}]$; Begriff „perfekte
Sequenz``.

\textbf{Rechnung (Python):}
\begin{lstlisting}[style=python]
import numpy as np
import matplotlib.pyplot as plt

x = np.array([0,0,1,0,1,1,1,-1,-1,0,1,-1,1])
N = len(x)                      # 13

phi = np.zeros(N)
for m in range(N):              # 'for'-Schleife ueber die Verschiebung
    s = 0
    for k in range(N):
        s += x[k]*x[(k+m) % N]  # zyklischer Index
    phi[m] = s

print(phi)                      # -> [9 0 0 0 0 0 0 0 0 0 0 0 0]
plt.stem(np.arange(N), phi)
plt.xlabel("m"); plt.ylabel("phi_xx(m)"); plt.show()
\end{lstlisting}

\textbf{Lösung:}
$$
\boxed{\;\varphi_{xx}(0) = \sum_k x[k]^2 = 9,\qquad \varphi_{xx}(m) = 0\ \text{für } m=1,\ldots,12.\;}
$$
Die zyklische Autokorrelation ist also ein reiner Diracstoß:
$\varphi_{xx}(m)=9\,\delta[m]$ (modulo $13$). Genau deshalb heißt die Folge
\emph{perfekte Sequenz} -- sie ist zu jeder von null verschiedenen zyklischen
Verschiebung ihrer selbst \emph{orthogonal}. Solche Sequenzen sind ideal für die
Kanalvermessung: Die Korrelation mit dem Empfangssignal liefert unmittelbar
einen sauberen Impuls (keine Nebenmaxima).
\end{beispiel}

\section{Übung 3 – Zufallssignale}

In dieser Übung werden empirische Kenngrößen (Lage- und Streumaße) sowie die
empirische Korrelation zeitdiskreter Zufallssignale bestimmt. Alle Signalwerte
sind den Stem-Plots bzw.\ dem Histogramm des Aufgabenblattes entnommen. Es
gelten durchgängig die Definitionen aus Abschnitt~\ref{sec:kenngroessen-zufaellige-signale}
(empirischer Mittelwert und Median: Abschnitt~\ref{sec:lagemasse}; empirische
Varianz und Standardabweichung: Abschnitt~\ref{sec:streumasse}; empirische
Kreuzkorrelation: Abschnitt~\ref{sec:empirische-korrelation}).

% ================================================================
\subsection*{Vorrechenaufgabe 3 – Statistische Kenngrößen von Signalen}
% ================================================================

Aus dem Stem-Plot liest man für $k = 0, \ldots, 16$ (also $N = 17$ Abtastwerte)
die folgenden Amplituden ab:

\begin{center}
\renewcommand{\arraystretch}{1.2}
\begin{tabular}{c|ccccccccccccccccc}
$k$    & 0 & 1 & 2 & 3 & 4 & 5 & 6 & 7 & 8 & 9 & 10 & 11 & 12 & 13 & 14 & 15 & 16 \\
\hline
$x[k]$ & $-1$ & $1$ & $0$ & $2$ & $\tfrac12$ & $-1$ & $2$ & $1$ & $1$ & $2$
       & $0$ & $-2$ & $1$ & $2$ & $1$ & $-1$ & $\tfrac12$ \\
\end{tabular}
\end{center}

\begin{center}
\begin{tikzpicture}
\begin{axis}[
  width=13cm, height=6cm,
  xmin=-0.6, xmax=16.6, ymin=-2.6, ymax=2.6,
  axis lines=middle,
  xlabel={$k$}, ylabel={$x[k]$},
  xtick={0,1,2,3,4,5,6,7,8,9,10,11,12,13,14,15,16},
  ytick={-2,-1,1,2},
  tick label style={font=\small},
  ymajorgrids=false,
]
\addplot+[ycomb, blue, thick, mark=*, mark options={blue}] coordinates {
  (0,-1) (1,1) (2,0) (3,2) (4,0.5) (5,-1) (6,2) (7,1) (8,1) (9,2)
  (10,0) (11,-2) (12,1) (13,2) (14,1) (15,-1) (16,0.5)
};
\end{axis}
\end{tikzpicture}
\end{center}

% ---- a) --------------------------------------------------------
\subsubsection*{a) Histogramm der Amplituden}

\begin{beispiel}[Histogramm der Signalamplituden]

\textbf{Gegeben:} Die $N = 17$ Amplitudenwerte des Signals $x[k]$.

\textbf{Gesucht:} Das Histogramm $f[x]$ (absolute Häufigkeit der auftretenden
Amplitudenwerte).

\textbf{Formel} (Histogramm als absolute Häufigkeit der Klassen,
Abschnitt~\ref{sec:kenngroessen-zufaellige-signale}):

$$
f[x] = \#\{\, k : x[k] = x \,\}
$$

\textbf{Rechnung:} Auszählen der auftretenden Amplitudenwerte:

\begin{center}
\renewcommand{\arraystretch}{1.2}
\begin{tabular}{c|cccccc|c}
$x$    & $-2$ & $-1$ & $0$ & $\tfrac12$ & $1$ & $2$ & $\sum$ \\
\hline
$f[x]$ & $1$  & $3$  & $2$ & $2$        & $5$ & $4$ & $17$ \\
\end{tabular}
\end{center}

Probe: $1 + 3 + 2 + 2 + 5 + 4 = 17 = N$. \checkmark

\begin{center}
\begin{tikzpicture}
\begin{axis}[
  width=11cm, height=6cm,
  ybar, bar width=18pt,
  xmin=-2.8, xmax=2.8, ymin=0, ymax=6,
  axis lines=left,
  xlabel={$x$}, ylabel={$f[x]$ (Häufigkeit)},
  xtick={-2,-1,0,0.5,1,2},
  xticklabels={$-2$,$-1$,$0$,$\tfrac12$,$1$,$2$},
  ytick={0,1,2,3,4,5},
  tick label style={font=\small},
  nodes near coords, nodes near coords style={font=\footnotesize},
]
\addplot[fill=blue!30, draw=blue!70] coordinates {
  (-2,1) (-1,3) (0,2) (0.5,2) (1,5) (2,4)
};
\end{axis}
\end{tikzpicture}
\end{center}

\textbf{Lösung:} Die Verteilung ist leicht linkslastig (Häufung bei $x = 1$ und
$x = 2$), mit einzelnen kleinen Ausreißern bei $x = -2$.

\end{beispiel}

% ---- b) --------------------------------------------------------
\subsubsection*{b) Empirischer Mittelwert}

\begin{beispiel}[Empirischer Mittelwert]

\textbf{Gegeben:} $N = 17$ Amplitudenwerte $x[k]$.

\textbf{Gesucht:} Der empirische Mittelwert $\hat{\bar{x}}$.

\textbf{Formel} (empirischer Mittelwert, Abschnitt~\ref{sec:lagemasse}):

$$
\hat{\bar{x}} = \frac{1}{N} \sum_{k=0}^{N-1} x[k]
$$

\textbf{Rechnung:} Am bequemsten über die Histogramm-Häufigkeiten:

$$
\sum_{k=0}^{16} x[k]
= (-2)\cdot 1 + (-1)\cdot 3 + 0\cdot 2 + \tfrac12\cdot 2 + 1\cdot 5 + 2\cdot 4
= -2 - 3 + 0 + 1 + 5 + 8 = 9
$$

$$
\hat{\bar{x}} = \frac{9}{17}
$$

\textbf{Lösung:}

$$
\boxed{\;\hat{\bar{x}} = \frac{9}{17} \approx 0{,}529\;}
$$

\end{beispiel}

% ---- c) --------------------------------------------------------
\subsubsection*{c) Median}

\begin{beispiel}[Median]

\textbf{Gegeben:} $N = 17$ Amplitudenwerte $x[k]$.

\textbf{Gesucht:} Der Median $\tilde{x}$.

\textbf{Formel} (Median als Mitte der geordneten Daten,
Abschnitt~\ref{sec:lagemasse}): Bei ungeradem $N$ ist

$$
\tilde{x} = x_{\left(\frac{N+1}{2}\right)}^{\text{sort}}
$$

das mittlere Element der aufsteigend sortierten Wertefolge.

\textbf{Rechnung:} Sortieren der $17$ Werte:

$$
-2,\; -1,\, -1,\, -1,\; 0,\, 0,\; \tfrac12,\, \tfrac12,\;
\underbrace{1}_{9.\ \text{Wert}},\, 1,\, 1,\, 1,\, 1,\; 2,\, 2,\, 2,\, 2
$$

Bei $N = 17$ ist der $\tfrac{N+1}{2} = 9$-te Wert der Median.

\textbf{Lösung:}

$$
\boxed{\;\tilde{x} = 1\;}
$$

\end{beispiel}

% ---- d) --------------------------------------------------------
\subsubsection*{d) Empirische Varianz}

\begin{beispiel}[Empirische Varianz]

\textbf{Gegeben:} $N = 17$, $\hat{\bar{x}} = \tfrac{9}{17}$.

\textbf{Gesucht:} Die empirische Varianz $\hat{\sigma}^2$.

\textbf{Formel} (empirische Varianz, Abschnitt~\ref{sec:streumasse}):

$$
\hat{\sigma}^2 = \frac{1}{N-1} \sum_{k=0}^{N-1} \bigl(x[k] - \hat{\bar{x}}\bigr)^2
= \frac{1}{N-1}\left( \sum_{k=0}^{N-1} x[k]^2 - N\,\hat{\bar{x}}^{\,2} \right)
$$

\textbf{Rechnung:} Zunächst die Summe der Quadrate (über die Häufigkeiten):

$$
\sum_{k=0}^{16} x[k]^2
= (-2)^2\cdot 1 + (-1)^2\cdot 3 + 0^2\cdot 2 + \left(\tfrac12\right)^2\cdot 2
  + 1^2\cdot 5 + 2^2\cdot 4
= 4 + 3 + 0 + \tfrac12 + 5 + 16 = 28{,}5
$$

$$
\hat{\sigma}^2 = \frac{1}{16}\left( 28{,}5 - 17\cdot\left(\frac{9}{17}\right)^2 \right)
= \frac{1}{16}\left( 28{,}5 - \frac{81}{17} \right)
= \frac{1}{16}\cdot \frac{807}{34}
= \frac{807}{544}
$$

\textbf{Lösung:}

$$
\boxed{\;\hat{\sigma}^2 = \frac{807}{544} \approx 1{,}483\;}
$$

\end{beispiel}

% ---- e) --------------------------------------------------------
\subsubsection*{e) Empirische Standardabweichung}

\begin{beispiel}[Empirische Standardabweichung]

\textbf{Gegeben:} $\hat{\sigma}^2 = \tfrac{807}{544} \approx 1{,}483$.

\textbf{Gesucht:} Die empirische Standardabweichung $\hat{\sigma}$.

\textbf{Formel} (Abschnitt~\ref{sec:streumasse}):

$$
\hat{\sigma} = \sqrt{\hat{\sigma}^2}
$$

\textbf{Rechnung:}

$$
\hat{\sigma} = \sqrt{\frac{807}{544}} = \sqrt{1{,}483\ldots}
$$

\textbf{Lösung:}

$$
\boxed{\;\hat{\sigma} \approx 1{,}218\;}
$$

\end{beispiel}

% ================================================================
\subsection*{Aufgabe 3.1 – Statistische Kenngrößen von Signalen}
% ================================================================

Für $k = 0, \ldots, 15$ (also $N = 16$ Abtastwerte) liest man aus dem Stem-Plot ab:

\begin{center}
\renewcommand{\arraystretch}{1.2}
\begin{tabular}{c|cccccccccccccccc}
$k$    & 0 & 1 & 2 & 3 & 4 & 5 & 6 & 7 & 8 & 9 & 10 & 11 & 12 & 13 & 14 & 15 \\
\hline
$x[k]$ & $0$ & $1$ & $0$ & $0$ & $1$ & $1$ & $5$ & $1$ & $0$ & $1$
       & $1$ & $7$ & $0$ & $0$ & $0$ & $1$ \\
\end{tabular}
\end{center}

\begin{center}
\begin{tikzpicture}
\begin{axis}[
  width=13cm, height=6cm,
  xmin=-0.6, xmax=15.6, ymin=0, ymax=7.6,
  axis lines=middle,
  xlabel={$k$}, ylabel={$x[k]$},
  xtick={0,1,2,3,4,5,6,7,8,9,10,11,12,13,14,15},
  ytick={1,2,3,4,5,6,7},
  tick label style={font=\small},
]
\addplot+[ycomb, blue, thick, mark=*, mark options={blue}] coordinates {
  (0,0) (1,1) (2,0) (3,0) (4,1) (5,1) (6,5) (7,1) (8,0) (9,1)
  (10,1) (11,7) (12,0) (13,0) (14,0) (15,1)
};
\end{axis}
\end{tikzpicture}
\end{center}

Häufigkeiten: der Wert $0$ tritt $7$-mal, der Wert $1$ tritt $7$-mal, der Wert $5$
einmal und der Wert $7$ einmal auf ($7+7+1+1 = 16$).

% ---- a) --------------------------------------------------------
\subsubsection*{a) Empirischer Mittelwert}

\begin{beispiel}[Empirischer Mittelwert]

\textbf{Gegeben:} $N = 16$ Amplitudenwerte $x[k]$.

\textbf{Gesucht:} $\hat{\bar{x}}$.

\textbf{Formel} (Abschnitt~\ref{sec:lagemasse}):

$$
\hat{\bar{x}} = \frac{1}{N} \sum_{k=0}^{N-1} x[k]
$$

\textbf{Rechnung:}

$$
\sum_{k=0}^{15} x[k] = 7\cdot 1 + 1\cdot 5 + 1\cdot 7 = 7 + 5 + 7 = 19,
\qquad
\hat{\bar{x}} = \frac{19}{16}
$$

\textbf{Lösung:}

$$
\boxed{\;\hat{\bar{x}} = \frac{19}{16} = 1{,}1875\;}
$$

\end{beispiel}

% ---- b) --------------------------------------------------------
\subsubsection*{b) Median}

\begin{beispiel}[Median]

\textbf{Gegeben:} $N = 16$ Amplitudenwerte $x[k]$.

\textbf{Gesucht:} $\tilde{x}$.

\textbf{Formel} (Abschnitt~\ref{sec:lagemasse}): Bei geradem $N$ ist der Median das
arithmetische Mittel der beiden mittleren Werte:

$$
\tilde{x} = \frac{1}{2}\left( x_{\left(\frac{N}{2}\right)}^{\text{sort}}
                              + x_{\left(\frac{N}{2}+1\right)}^{\text{sort}} \right)
$$

\textbf{Rechnung:} Sortieren der $16$ Werte:

$$
0,0,0,0,0,0,0,\; \underbrace{1}_{8.},\, \underbrace{1}_{9.},\, 1,1,1,1,1,\; 5,\, 7
$$

Der $8.$ und $9.$ Wert sind beide $1$:

$$
\tilde{x} = \frac{1 + 1}{2} = 1
$$

\textbf{Lösung:}

$$
\boxed{\;\tilde{x} = 1\;}
$$

\end{beispiel}

% ---- c) --------------------------------------------------------
\subsubsection*{c) Empirische Varianz}

\begin{beispiel}[Empirische Varianz]

\textbf{Gegeben:} $N = 16$, $\hat{\bar{x}} = \tfrac{19}{16}$.

\textbf{Gesucht:} $\hat{\sigma}^2$.

\textbf{Formel} (Abschnitt~\ref{sec:streumasse}):

$$
\hat{\sigma}^2 = \frac{1}{N-1}\left( \sum_{k=0}^{N-1} x[k]^2 - N\,\hat{\bar{x}}^{\,2} \right)
$$

\textbf{Rechnung:}

$$
\sum_{k=0}^{15} x[k]^2 = 7\cdot 1^2 + 1\cdot 5^2 + 1\cdot 7^2 = 7 + 25 + 49 = 81
$$

$$
\hat{\sigma}^2 = \frac{1}{15}\left( 81 - 16\cdot\left(\frac{19}{16}\right)^2 \right)
= \frac{1}{15}\left( 81 - \frac{361}{16} \right)
= \frac{1}{15}\cdot \frac{935}{16}
= \frac{935}{240}
$$

\textbf{Lösung:}

$$
\boxed{\;\hat{\sigma}^2 = \frac{935}{240} \approx 3{,}896\;}
$$

\end{beispiel}

% ---- d) --------------------------------------------------------
\subsubsection*{d) Empirische Standardabweichung}

\begin{beispiel}[Empirische Standardabweichung]

\textbf{Gegeben:} $\hat{\sigma}^2 = \tfrac{935}{240} \approx 3{,}896$.

\textbf{Gesucht:} $\hat{\sigma}$.

\textbf{Formel} (Abschnitt~\ref{sec:streumasse}):

$$
\hat{\sigma} = \sqrt{\hat{\sigma}^2}
$$

\textbf{Rechnung:}

$$
\hat{\sigma} = \sqrt{\frac{935}{240}} = \sqrt{3{,}896\ldots}
$$

\textbf{Lösung:}

$$
\boxed{\;\hat{\sigma} \approx 1{,}974\;}
$$

\end{beispiel}

% ---- e) --------------------------------------------------------
\subsubsection*{e) Mittlere absolute Abweichung vom Median}

\begin{beispiel}[Mittlere absolute Abweichung $\tilde{\Delta}$]

\textbf{Gegeben:} $N = 16$, Median $\tilde{x} = 1$.

\textbf{Gesucht:} Die mittlere absolute Abweichung

$$
\tilde{\Delta} = \frac{1}{N} \sum_{k=0}^{N-1} \bigl| x[k] - \tilde{x} \bigr|.
$$

\textbf{Formel:} wie in der Aufgabenstellung angegeben (Streumaß relativ zum
Median, vgl.\ Abschnitt~\ref{sec:streumasse}).

\textbf{Rechnung:} Die Beträge $|x[k] - 1|$ nach Häufigkeiten:

$$
\sum_{k=0}^{15} |x[k] - 1|
= 7\cdot|0-1| + 7\cdot|1-1| + 1\cdot|5-1| + 1\cdot|7-1|
= 7\cdot 1 + 0 + 4 + 6 = 17
$$

$$
\tilde{\Delta} = \frac{17}{16}
$$

\textbf{Lösung:}

$$
\boxed{\;\tilde{\Delta} = \frac{17}{16} = 1{,}0625\;}
$$

\end{beispiel}

% ================================================================
\subsection*{Aufgabe 3.2 – Histogramm}
% ================================================================

Gegeben ist das Histogramm der Signalamplituden mit den absoluten Häufigkeiten
$f[x]$ für die Amplitudenwerte $x = -3, \ldots, 3$:

\begin{center}
\renewcommand{\arraystretch}{1.2}
\begin{tabular}{c|ccccccc|c}
$x$    & $-3$ & $-2$ & $-1$ & $0$ & $1$ & $2$ & $3$ & $\sum$ \\
\hline
$f[x]$ & $1$  & $5$  & $13$ & $14$ & $9$ & $7$ & $1$ & $50$ \\
\end{tabular}
\end{center}

\begin{center}
\begin{tikzpicture}
\begin{axis}[
  width=11cm, height=6.5cm,
  ybar, bar width=20pt,
  xmin=-3.8, xmax=3.8, ymin=0, ymax=15,
  axis lines=left,
  xlabel={$x$}, ylabel={$f[x]$ (Häufigkeit)},
  xtick={-3,-2,-1,0,1,2,3},
  ytick={0,2,4,6,8,10,12,14},
  tick label style={font=\small},
  nodes near coords, nodes near coords style={font=\footnotesize},
]
\addplot[fill=blue!30, draw=blue!70] coordinates {
  (-3,1) (-2,5) (-1,13) (0,14) (1,9) (2,7) (3,1)
};
\end{axis}
\end{tikzpicture}
\end{center}

Die Gesamtzahl der Abtastwerte ist $N = \sum_x f[x] = 1+5+13+14+9+7+1 = 50$.

% ---- a) --------------------------------------------------------
\subsubsection*{a) Empirischer Mittelwert aus dem Histogramm}

\begin{beispiel}[Mittelwert aus Histogramm-Häufigkeiten]

\textbf{Gegeben:} Häufigkeiten $f[x]$, $N = 50$.

\textbf{Gesucht:} Formel für $\hat{\bar{x}}$ aus den Histogrammwerten sowie der
Zahlenwert.

\textbf{Formel:} Fasst man in $\hat{\bar{x}} = \tfrac{1}{N}\sum_k x[k]$ alle
gleichen Amplituden zusammen, so tritt jeder Wert $x$ genau $f[x]$-mal auf
(Abschnitt~\ref{sec:lagemasse} zusammen mit dem Histogramm-Begriff aus
Abschnitt~\ref{sec:kenngroessen-zufaellige-signale}):

$$
\hat{\bar{x}} = \frac{1}{N} \sum_{x} x \cdot f[x],
\qquad N = \sum_{x} f[x]
$$

\textbf{Rechnung:}

$$
\sum_{x} x \cdot f[x]
= (-3)(1) + (-2)(5) + (-1)(13) + (0)(14) + (1)(9) + (2)(7) + (3)(1)
$$
$$
= -3 - 10 - 13 + 0 + 9 + 14 + 3 = 0
$$

$$
\hat{\bar{x}} = \frac{0}{50} = 0
$$

\textbf{Lösung:}

$$
\boxed{\;\hat{\bar{x}} = \frac{1}{N}\sum_x x\, f[x] = 0\;}
$$

Das Signal ist empirisch mittelwertfrei.

\end{beispiel}

% ---- b) --------------------------------------------------------
\subsubsection*{b) Empirische Varianz aus dem Histogramm}

\begin{beispiel}[Varianz aus Histogramm-Häufigkeiten]

\textbf{Gegeben:} $f[x]$, $N = 50$, $\hat{\bar{x}} = 0$.

\textbf{Gesucht:} Formel für $\hat{\sigma}^2$ aus den Histogrammwerten sowie der
Zahlenwert.

\textbf{Formel} (gewichtete Summe der quadratischen Abweichungen,
Abschnitt~\ref{sec:streumasse}):

$$
\hat{\sigma}^2 = \frac{1}{N-1} \sum_{x} \bigl(x - \hat{\bar{x}}\bigr)^2 f[x]
$$

\textbf{Rechnung:} Mit $\hat{\bar{x}} = 0$ vereinfacht sich die Summe zu
$\sum_x x^2 f[x]$:

$$
\sum_{x} x^2 f[x]
= 9\cdot 1 + 4\cdot 5 + 1\cdot 13 + 0\cdot 14 + 1\cdot 9 + 4\cdot 7 + 9\cdot 1
$$
$$
= 9 + 20 + 13 + 0 + 9 + 28 + 9 = 88
$$

$$
\hat{\sigma}^2 = \frac{88}{50 - 1} = \frac{88}{49}
$$

\textbf{Lösung:}

$$
\boxed{\;\hat{\sigma}^2 = \frac{1}{N-1}\sum_x (x-\hat{\bar{x}})^2 f[x]
= \frac{88}{49} \approx 1{,}796\;}
$$

\end{beispiel}

% ---- c) --------------------------------------------------------
\subsubsection*{c) Empirische Standardabweichung aus dem Histogramm}

\begin{beispiel}[Standardabweichung aus Histogramm-Häufigkeiten]

\textbf{Gegeben:} $\hat{\sigma}^2 = \tfrac{88}{49} \approx 1{,}796$.

\textbf{Gesucht:} Formel für $\hat{\sigma}$ sowie der Zahlenwert.

\textbf{Formel} (Abschnitt~\ref{sec:streumasse}):

$$
\hat{\sigma} = \sqrt{\hat{\sigma}^2}
= \sqrt{\frac{1}{N-1} \sum_{x} \bigl(x - \hat{\bar{x}}\bigr)^2 f[x]}
$$

\textbf{Rechnung:}

$$
\hat{\sigma} = \sqrt{\frac{88}{49}} = \frac{\sqrt{88}}{7}
$$

\textbf{Lösung:}

$$
\boxed{\;\hat{\sigma} = \frac{\sqrt{88}}{7} \approx 1{,}340\;}
$$

\end{beispiel}

% ================================================================
\subsection*{Aufgabe 3.3 – Korrelation}
% ================================================================

Gegeben ist das Signal $x[k]$ für $k = 0, \ldots, 14$ ($N = 15$) sowie die
Codefolge $y[k] = [\,1,\, 1,\, -1\,]$, d.\,h.\ $y[0] = 1$, $y[1] = 1$,
$y[2] = -1$ (sonst $0$).

\begin{center}
\renewcommand{\arraystretch}{1.2}
\begin{tabular}{c|ccccccccccccccc}
$k$    & 0 & 1 & 2 & 3 & 4 & 5 & 6 & 7 & 8 & 9 & 10 & 11 & 12 & 13 & 14 \\
\hline
$x[k]$ & $0$ & $-\tfrac12$ & $1$ & $0$ & $1$ & $\tfrac12$ & $-1$ & $1$ & $-\tfrac32$
       & $0$ & $\tfrac32$ & $0$ & $0$ & $1$ & $0$ \\
\end{tabular}
\end{center}

\begin{center}
\begin{tikzpicture}
\begin{axis}[
  width=13cm, height=5.5cm,
  xmin=-0.6, xmax=14.6, ymin=-1.9, ymax=1.9,
  axis lines=middle,
  xlabel={$k$}, ylabel={$x[k]$},
  xtick={0,1,2,3,4,5,6,7,8,9,10,11,12,13,14},
  ytick={-1.5,-1,-0.5,0.5,1,1.5},
  yticklabels={$-\tfrac32$,$-1$,$-\tfrac12$,$\tfrac12$,$1$,$\tfrac32$},
  tick label style={font=\small},
]
\addplot+[ycomb, blue, thick, mark=*, mark options={blue}] coordinates {
  (0,0) (1,-0.5) (2,1) (3,0) (4,1) (5,0.5) (6,-1) (7,1) (8,-1.5)
  (9,0) (10,1.5) (11,0) (12,0) (13,1) (14,0)
};
\end{axis}
\end{tikzpicture}
\end{center}

\begin{center}
\begin{tikzpicture}
\begin{axis}[
  width=6cm, height=4cm,
  xmin=-0.6, xmax=4.6, ymin=-1.4, ymax=1.4,
  axis lines=middle,
  xlabel={$k$}, ylabel={$y[k]$},
  xtick={0,1,2,3,4}, ytick={-1,1},
  tick label style={font=\small},
]
\addplot+[ycomb, blue, thick, mark=*, mark options={blue}] coordinates {
  (0,1) (1,1) (2,-1)
};
\end{axis}
\end{tikzpicture}
\end{center}

% ---- a) --------------------------------------------------------
\subsubsection*{a) Empirischer Mittelwert von $x[k]$}

\begin{beispiel}[Empirischer Mittelwert des Signals]

\textbf{Gegeben:} $N = 15$ Amplitudenwerte $x[k]$.

\textbf{Gesucht:} $\hat{\bar{x}}$.

\textbf{Formel} (Abschnitt~\ref{sec:lagemasse}):

$$
\hat{\bar{x}} = \frac{1}{N} \sum_{k=0}^{N-1} x[k]
$$

\textbf{Rechnung:}

$$
\sum_{k=0}^{14} x[k]
= -\tfrac12 + 1 + 1 + \tfrac12 - 1 + 1 - \tfrac32 + \tfrac32 + 1 = 3
$$

$$
\hat{\bar{x}} = \frac{3}{15} = \frac{1}{5}
$$

\textbf{Lösung:}

$$
\boxed{\;\hat{\bar{x}} = \frac{1}{5} = 0{,}2\;}
$$

\end{beispiel}

% ---- b) --------------------------------------------------------
\subsubsection*{b) Empirische Varianz von $x[k]$}

\begin{beispiel}[Empirische Varianz des Signals]

\textbf{Gegeben:} $N = 15$, $\hat{\bar{x}} = \tfrac15$.

\textbf{Gesucht:} $\hat{\sigma}^2$.

\textbf{Formel} (Abschnitt~\ref{sec:streumasse}):

$$
\hat{\sigma}^2 = \frac{1}{N-1}\left( \sum_{k=0}^{N-1} x[k]^2 - N\,\hat{\bar{x}}^{\,2} \right)
$$

\textbf{Rechnung:}

$$
\sum_{k=0}^{14} x[k]^2
= \tfrac14 + 1 + 1 + \tfrac14 + 1 + 1 + \tfrac94 + \tfrac94 + 1 = 10
$$

$$
\hat{\sigma}^2 = \frac{1}{14}\left( 10 - 15\cdot\left(\tfrac15\right)^2 \right)
= \frac{1}{14}\left( 10 - \tfrac{3}{5} \right)
= \frac{1}{14}\cdot \frac{47}{5}
= \frac{47}{70}
$$

\textbf{Lösung:}

$$
\boxed{\;\hat{\sigma}^2 = \frac{47}{70} \approx 0{,}671\;}
$$

\end{beispiel}

% ---- c) --------------------------------------------------------
\subsubsection*{c) Median von $x[k]$}

\begin{beispiel}[Median des Signals]

\textbf{Gegeben:} $N = 15$ Amplitudenwerte $x[k]$.

\textbf{Gesucht:} $\tilde{x}$.

\textbf{Formel} (Abschnitt~\ref{sec:lagemasse}): Bei ungeradem $N$ ist der Median
der $\tfrac{N+1}{2} = 8$-te Wert der sortierten Folge.

\textbf{Rechnung:} Sortieren der $15$ Werte:

$$
-\tfrac32,\, -1,\, -\tfrac12,\; 0,\,0,\,0,\,0,\;
\underbrace{0}_{8.\ \text{Wert}},\, 0,\; \tfrac12,\; 1,\,1,\,1,\,1,\; \tfrac32
$$

(Es gibt sechs Nullen an den Positionen $k = 0,3,9,11,12,14$.) Der $8.$ Wert ist
$0$.

\textbf{Lösung:}

$$
\boxed{\;\tilde{x} = 0\;}
$$

\end{beispiel}

% ---- d) --------------------------------------------------------
\subsubsection*{d) Empirische Kreuzkorrelation und Detektion der Codefolge}

\begin{beispiel}[Kreuzkorrelation zur Detektion der Codefolge]

\textbf{Gegeben:} Signal $x[k]$ ($k = 0, \ldots, 14$) und Codefolge
$y[k] = [\,1,1,-1\,]$.

\textbf{Gesucht:} Die empirische Kreuzkorrelation $\hat{\varphi}_{xy}[\varkappa]$
und daraus die Stelle $k$, an der die Codefolge $[\,1,1,-1\,]$ in $x[k]$ beginnt.

\textbf{Formel} (empirische Kreuzkorrelation, Abschnitt~\ref{sec:empirische-korrelation}):

$$
\hat{\varphi}_{xy}[\varkappa] = \sum_{k} x[k]\, y[k + \varkappa]
$$

\textbf{Rechnung:} Da $y[n]$ nur für $n \in \{0,1,2\}$ von null verschieden ist
($y[0]=y[1]=1$, $y[2]=-1$), tragen in der Summe nur die Indizes $k$ mit
$k+\varkappa \in \{0,1,2\}$ bei. Mit $k = -\varkappa,\, 1-\varkappa,\, 2-\varkappa$ folgt

$$
\hat{\varphi}_{xy}[\varkappa]
= x[-\varkappa]\cdot 1 + x[1-\varkappa]\cdot 1 + x[2-\varkappa]\cdot(-1)
= x[-\varkappa] + x[1-\varkappa] - x[2-\varkappa].
$$

Ein positiver Beitrag entsteht also dort, wo $x$ dem Muster $[\,1,1,-1\,]$ ähnelt.
Auswertung für $\varkappa = 0, -1, \ldots, -12$:

\begin{center}
\renewcommand{\arraystretch}{1.3}
\begin{tabular}{c|c|c}
$\varkappa$ & $x[-\varkappa] + x[1-\varkappa] - x[2-\varkappa]$ & $\hat{\varphi}_{xy}[\varkappa]$ \\
\hline
$0$    & $x[0]+x[1]-x[2] = 0-\tfrac12-1$        & $-\tfrac32$ \\
$-1$   & $x[1]+x[2]-x[3] = -\tfrac12+1-0$       & $\tfrac12$ \\
$-2$   & $x[2]+x[3]-x[4] = 1+0-1$               & $0$ \\
$-3$   & $x[3]+x[4]-x[5] = 0+1-\tfrac12$        & $\tfrac12$ \\
$-4$   & $x[4]+x[5]-x[6] = 1+\tfrac12-(-1)$     & $\mathbf{\tfrac52}$ \\
$-5$   & $x[5]+x[6]-x[7] = \tfrac12-1-1$        & $-\tfrac32$ \\
$-6$   & $x[6]+x[7]-x[8] = -1+1-(-\tfrac32)$    & $\tfrac32$ \\
$-7$   & $x[7]+x[8]-x[9] = 1-\tfrac32-0$        & $-\tfrac12$ \\
$-8$   & $x[8]+x[9]-x[10] = -\tfrac32+0-\tfrac32$ & $-3$ \\
$-9$   & $x[9]+x[10]-x[11] = 0+\tfrac32-0$      & $\tfrac32$ \\
$-10$  & $x[10]+x[11]-x[12] = \tfrac32+0-0$     & $\tfrac32$ \\
$-11$  & $x[11]+x[12]-x[13] = 0+0-1$            & $-1$ \\
$-12$  & $x[12]+x[13]-x[14] = 0+1-0$            & $1$ \\
\end{tabular}
\end{center}

\begin{center}
\begin{tikzpicture}
\begin{axis}[
  width=13cm, height=5cm,
  xmin=-12.6, xmax=0.6, ymin=-3.4, ymax=3.4,
  axis lines=middle,
  xlabel={$\varkappa$}, ylabel={$\hat{\varphi}_{xy}[\varkappa]$},
  xtick={-12,-11,-10,-9,-8,-7,-6,-5,-4,-3,-2,-1,0},
  ytick={-3,-2,-1,1,2,3},
  tick label style={font=\small},
]
\addplot+[ycomb, blue, thick, mark=*, mark options={blue}] coordinates {
  (0,-1.5) (-1,0.5) (-2,0) (-3,0.5) (-4,2.5) (-5,-1.5) (-6,1.5)
  (-7,-0.5) (-8,-3) (-9,1.5) (-10,1.5) (-11,-1) (-12,1)
};
\addplot[red, dashed, thick] coordinates {(-4,0) (-4,2.5)};
\node[red, anchor=south, font=\small] at (axis cs:-4,2.5) {Maximum};
\end{axis}
\end{tikzpicture}
\end{center}

Das ausgeprägte Maximum liegt bei $\varkappa = -4$ mit
$\hat{\varphi}_{xy}[-4] = \tfrac52$. Für $\varkappa = -4$ überlappt die
Codefolge $y$ gerade die Werte $x[4], x[5], x[6]$; die Codefolge beginnt also bei
$k = 4$. Tatsächlich passt das Muster $[\,1,1,-1\,]$ dort am besten:
$x[4] = 1$, $x[5] = \tfrac12$, $x[6] = -1$.

\textbf{Lösung:}

$$
\boxed{\;\max_{\varkappa} \hat{\varphi}_{xy}[\varkappa]
= \hat{\varphi}_{xy}[-4] = \frac{5}{2}
\quad\Longrightarrow\quad
\text{Codefolge beginnt bei } k = 4.\;}
$$

Dies ist das Prinzip eines Korrelationsdetektors (vgl.\
Abschnitt~\ref{sec:signalentdeckung-rauschen}): Die bekannte Codefolge wird über
das Signal geschoben, und das Maximum der Kreuzkorrelation markiert die Stelle
der besten Übereinstimmung.

\end{beispiel}

\section{Übung 4 – Abtastung und Quantisierung}

\subsection*{Grundprinzip}

Um ein wert- und zeitkontinuierliches Signal zu digitalisieren, wird es in zwei Schritten 
verarbeitet:

\begin{enumerate}
    \item \textbf{Abtastung:} Das Signal wird in gleichmäßigen Zeitabständen $T_A$ abgetastet.
          Ergebnis: wertkontinuierliches, zeitdiskretes Signal $x[n] = x(n \cdot T_A)$.
    \item \textbf{Quantisierung:} Die Abtastwerte werden einer von $N$ diskreten Stufen zugeordnet.
          Ergebnis: vollständig diskretes Signal $\hat{x}[n]$.
\end{enumerate}

\begin{figure}[H]
\centering
\begin{tikzpicture}
\begin{axis}[
    width=11cm, height=4.5cm,
    xlabel={$t$},
    ylabel={Amplitude},
    xmin=0, xmax=5.5,
    ymin=-1.3, ymax=1.3,
    axis x line=middle,
    axis y line=left,
    xtick={1,2,3,4,5},
    xticklabels={$T_A$, $2T_A$, $3T_A$, $4T_A$, $5T_A$},
    ytick={-1, 1},
    yticklabels={$-A$, $A$},
    tick align=outside,
    legend style={at={(0.99,0.97)}, anchor=north east, font=\small},
    every axis plot/.append style={line width=1.2pt},
]
\addplot[blue, smooth, domain=0:5.4, samples=200] {sin(deg(x*1.15))};
\addlegendentry{Analogsignal $x(t)$}
\addplot[red, ycomb, mark=*, mark size=2.8pt, mark options={fill=red}]
    coordinates {
        (1, {sin(deg(1.15))})
        (2, {sin(deg(2.30))})
        (3, {sin(deg(3.45))})
        (4, {sin(deg(4.60))})
        (5, {sin(deg(5.75))})
    };
\addlegendentry{Abtastwerte $x[n]$}
\end{axis}
\end{tikzpicture}
\caption{Abtastung: Aus dem kontinuierlichen Signal (blau) werden in Abständen $T_A$ Werte 
    entnommen (rot).}
\end{figure}

\subsection*{Beispiel: Abtastperiode und Stufenbreite}

Gegeben sei die Abtastfrequenz $f_A = 10\,\text{kHz}$, wobei $\text{Hz} = \text{s}^{-1}$.
Die Abtastperiode beträgt:

$$
T_A = \frac{1}{f_A} = \frac{1}{10\,000\,\text{s}^{-1}} = 0{,}1\,\text{ms}
$$

Für den Wertebereich $[-0{,}4,\; 0{,}4]$ (Gesamtbreite $x_{\max} - x_{\min} = 0{,}8$)
und eine Stufenbreite $\delta x = 0{,}2$ ergibt sich die Stufenanzahl:

$$
N = \frac{x_{\max} - x_{\min}}{\delta x} = \frac{0{,}8}{0{,}2} = 4,
\qquad
\delta x = \frac{0{,}8}{N} = \frac{0{,}8}{4} = 0{,}2
$$

Der Rekonstruktionswert der untersten Stufe ergibt sich als Mittelpunkt des ersten Intervalls:

$$
x_1 = \frac{-0{,}4 + (-0{,}2)}{2} = -0{,}3
$$

\subsection*{Midrise-Kennlinie}

Die \textbf{Midrise-Kennlinie} ist eine Treppenfunktion, bei der die steigende Flanke durch den
Ursprung $(0,\,0)$ verläuft. Beim \textbf{Midtread-Quantisierer} liegt im Ursprung hingegen
eine konstante Stufe – die Kurve verläuft dort flach.

\begin{figure}[H]
\centering
\begin{tikzpicture}
\begin{axis}[
    width=9cm, height=9cm,
    xmin=-0.50, xmax=0.50,
    ymin=-0.50, ymax=0.50,
    axis lines=center,
    xtick={-0.4,-0.2,0.2,0.4},
    xticklabels={$-0{,}4$, $-0{,}2$, $0{,}2$, $0{,}4$},
    ytick={-0.3,-0.1,0.1,0.3},
    yticklabels={$-0{,}3$, $-0{,}1$, $0{,}1$, $0{,}3$},
    x tick label style={anchor=north, yshift=-4pt, font=\small},
    y tick label style={anchor=east, xshift=-4pt, font=\small},
    grid=both,
    grid style={dashed, gray!30},
    tick align=outside,
    every axis plot/.append style={line width=1.4pt},
]
% Horizontale Stufen
\addplot[blue] coordinates {(-0.50,-0.3) (-0.20,-0.3)};
\addplot[blue] coordinates {(-0.20,-0.1) ( 0.00,-0.1)};
\addplot[blue] coordinates {( 0.00, 0.1) ( 0.20, 0.1)};
\addplot[blue] coordinates {( 0.20, 0.3) ( 0.50, 0.3)};
% Gestrichelte Übergänge
\addplot[blue, dashed, line width=0.8pt] coordinates {(-0.2,-0.3) (-0.2,-0.1)};
\addplot[blue, dashed, line width=0.8pt] coordinates {( 0.0,-0.1) ( 0.0, 0.1)};
\addplot[blue, dashed, line width=0.8pt] coordinates {( 0.2, 0.1) ( 0.2, 0.3)};
% Offene Kreise (nicht inbegriffen)
\addplot[only marks, mark=o, mark size=3.5pt,
         mark options={draw=blue, fill=white, line width=1.2pt}]
    coordinates {(-0.2,-0.3) (0.0,-0.1) (0.2,0.1)};
% Gefüllte Kreise (inbegriffen)
\addplot[only marks, mark=*, mark size=3.5pt, mark options={fill=blue}]
    coordinates {(-0.2,-0.1) (0.0,0.1) (0.2,0.3)};
\end{axis}
\end{tikzpicture}
\caption{Midrise-Kennlinie: $N = 4$ Stufen, $\delta x = 0{,}2$, Wertebereich $[-0{,}4,\,0{,}4]$.
Die Treppenfunktion steigt am Ursprung und erfüllt $Q(x) \neq 0$ für alle $x \neq 0$.}
\end{figure}

Für das Beispiel ($N = 4$, $\delta x = 0{,}2$, Midrise) ergeben sich folgende
Quantisierungsintervalle und Rekonstruktionswerte $\hat{x}$:

\begin{itemize}
    \item $x \in [{-0{,}4},\; {-0{,}2})$ \quad$\Rightarrow$\quad $\hat{x} = {-0{,}3}$
    \item $x \in [{-0{,}2},\; 0)$         \quad$\Rightarrow$\quad $\hat{x} = {-0{,}1}$
    \item $x \in [0,\; 0{,}2)$            \quad$\Rightarrow$\quad $\hat{x} = 0{,}1$
    \item $x \in [0{,}2,\; 0{,}4]$        \quad$\Rightarrow$\quad $\hat{x} = 0{,}3$
\end{itemize}

\begin{figure}[H]
\centering
\begin{tikzpicture}
\begin{axis}[
    width=11cm, height=5cm,
    xlabel={Index $n$},
    ylabel={Amplitude},
    xmin=-0.5, xmax=9.5,
    ymin=-0.48, ymax=0.52,
    axis lines=left,
    xtick={0,1,...,9},
    ytick={-0.3,-0.1,0.1,0.3},
    yticklabels={$-0{,}3$, $-0{,}1$, $0{,}1$, $0{,}3$},
    ymajorgrids=true,
    grid style={dashed, gray!40},
    legend style={at={(0.99,0.98)}, anchor=north east, font=\small},
    every axis plot/.append style={line width=1.2pt},
]
% Originale Abtastwerte
\addplot[red, ycomb, mark=o, mark size=3pt,
         mark options={draw=red, fill=white, line width=1.2pt}]
    coordinates {
        (0,  0.38) (1,  0.22) (2, -0.05) (3, -0.31) (4, -0.39)
        (5, -0.17) (6,  0.08) (7,  0.28) (8,  0.39) (9,  0.15)
    };
\addlegendentry{Abtastwerte $x[n]$}
% Quantisierte Werte (leicht versetzt zur besseren Lesbarkeit)
\addplot[blue, ycomb, mark=*, mark size=3pt, mark options={fill=blue},
         transform canvas={xshift=5pt}]
    coordinates {
        (0,  0.3) (1,  0.3) (2, -0.1) (3, -0.3) (4, -0.3)
        (5, -0.1) (6,  0.1) (7,  0.3) (8,  0.3) (9,  0.1)
    };
\addlegendentry{Quantisierte Werte $\hat{x}[n]$}
\end{axis}
\end{tikzpicture}
\caption{Quantisierung: Die Abtastwerte (rot) werden den nächstgelegenen Quantisierungsstufen
(blau, gestrichelte Linien) zugeordnet. Jede Stufe entspricht dem Rekonstruktionswert $\hat{x}$.}
\end{figure}

Nach der Quantisierung aller Abtastwerte entsteht aus dem ursprünglich kontinuierlichen Signal
ein vollständig diskretes, digitales Signal.




% -------------------------------------------------- %
%                    Tutorium 4                      %
% -------------------------------------------------- %

\subsection{Tutorium 4}

\subsubsection*{Motivation: Warum abtasten und quantisieren?}

Digitale Systeme können nur mit endlich vielen, diskreten Werten umgehen. Um ein physikalisches
Signal – das in der Natur wert- und zeitkontinuierlich vorliegt – digital verarbeiten zu können,
sind zwei Schritte nötig:

\begin{enumerate}
    \item \textbf{Abtastung} (zeitdiskret machen): Das Signal wird nur zu Zeitpunkten $n \cdot T_A$
          ausgewertet. Ergebnis: wertkontinuierliches, zeitdiskretes Signal $x[n]$.
    \item \textbf{Quantisierung} (wertdiskret machen): Die Abtastwerte werden auf $N$ diskrete
          Stufen abgebildet. Ergebnis: vollständig diskretes Signal $\hat{x}[n]$.
\end{enumerate}

\subsubsection*{Midrise vs.\ Midtread}

Beide Quantisierungskennlinien ordnen jedem Eingangswert $x$ einen diskreten Rekonstruktionswert
$\hat{x}$ zu, unterscheiden sich aber im Verhalten am Ursprung:

\begin{itemize}
    \item \textbf{Midrise:} Die Stufengrenze liegt genau bei $x = 0$. Es gibt keine eigene Stufe
          für den Nullpunkt – die Kennlinie steigt dort durch. Sinnvoll bei gleichmäßig verteilten
          Signalen mit gerader Stufenanzahl $N$.
    \item \textbf{Midtread:} Die Stufe liegt mittig über $x = 0$, d.\,h.\ $Q(0) = 0$. Vorteil:
          Schwache Signale nahe Null werden exakt auf $0$ quantisiert – kein Rauschen bei Stille.
          Sinnvoll bei Signalen, die häufig um die Null schwingen.
\end{itemize}

\subsubsection*{Rekonstruktion durch Interpolation}

Um das zeitdiskrete Signal $x[n]$ wieder in ein zeitkontinuierliches Signal $x(t)$ umzuwandeln,
werden gewichtete Basisfunktionen $g_k(t)$ übereinandergelegt:

$$
x(t) = \sum_{k=-\infty}^{\infty} x[k] \cdot g_k(t)
$$

Die Wahl der Basisfunktion $g_k(t)$ bestimmt Qualität und Glattheitsgrad der Rekonstruktion:

\begin{itemize}
    \item \textbf{0.\,Ordnung – Zero-Order Hold (ZOH):} $g_k(t)$ ist ein Rechteckimpuls. Der
          Abtastwert wird bis zum nächsten Zeitpunkt konstant gehalten. Ergebnis: zeitkontinuierlich,
          aber wertdiskret (Treppenfunktion).
    \item \textbf{1.\,Ordnung – Lineare Interpolation:} $g_k(t)$ ist ein Dreieckimpuls. Benachbarte
          Abtastwerte werden linear verbunden. Ergebnis: zeitkontinuierlich und stetig,
          aber nicht differenzierbar (Knicke an den Abtastzeitpunkten).
    \item \textbf{Ideale Interpolation – $\operatorname{si}$-Funktion:} Mit
          $g_k(t) = \operatorname{si}\!\left(\tfrac{t - k T_A}{T_A}\right)$ wird bei Einhaltung
          des Nyquist-Kriteriums das Originalsignal exakt rekonstruiert.
\end{itemize}

\begin{figure}[H]
\centering
\begin{tikzpicture}
\begin{axis}[
    width=11cm, height=5cm,
    xmin=-0.3, xmax=4.3,
    ymin=-0.15, ymax=1.35,
    axis x line=middle,
    axis y line=left,
    xtick={0,1,2,3,4},
    xticklabels={$0$, $T_A$, $2T_A$, $3T_A$, $4T_A$},
    ytick={0,1},
    tick align=outside,
    legend style={at={(0.99,0.97)}, anchor=north east, font=\small},
    every axis plot/.append style={line width=1.2pt},
]
\addplot[blue, dashed, smooth, domain=0:4.2, samples=200] {0.5*sin(deg(x*1.8)) + 0.5};
\addlegendentry{Original $x(t)$}
\addplot[red, const plot, mark=*, mark size=2.5pt, mark options={fill=red}]
    coordinates {(0,0.50) (1,0.99) (2,0.28) (3,0.11) (4,0.90) (4.3,0.90)};
\addlegendentry{0.\,Ordnung (ZOH)}
\addplot[orange!90!black, sharp plot, mark=*, mark size=2.5pt,
         mark options={fill=orange!90!black}]
    coordinates {(0,0.50) (1,0.99) (2,0.28) (3,0.11) (4,0.90)};
\addlegendentry{1.\,Ordnung (linear)}
\end{axis}
\end{tikzpicture}
\caption{Interpolationsvergleich: Original (blau, gestrichelt), Zero-Order Hold (rot) und lineare
Interpolation (orange) anhand identischer Abtastwerte.}
\end{figure}

\subsubsection*{Die \texorpdfstring{$\operatorname{si}$}{si}-Funktion}

Die $\operatorname{si}$-Funktion (normierter Sinc) ist definiert als:

$$
\operatorname{si}(x) = \frac{\sin(\pi x)}{\pi x} \quad (x \neq 0),
\qquad \operatorname{si}(0) = 1
$$

Entscheidend ist die \textbf{Interpolationseigenschaft}: An ganzzahligen Argumenten gilt:

$$
\operatorname{si}(n) =
\begin{cases}
    1 & n = 0 \\
    0 & n \in \mathbb{Z},\; n \neq 0
\end{cases}
$$

Damit ist jede Basisfunktion $g_k(n \cdot T_A) = \operatorname{si}(n - k)$ genau an ihrem
eigenen Abtastzeitpunkt gleich $1$ und an allen anderen gleich $0$ – die Abtastwerte stören
sich gegenseitig nicht. Die vollständige Rekonstruktionsformel lautet:

$$
x(t) = \sum_{k=-\infty}^{\infty} x[k] \cdot \operatorname{si}\!\left(\frac{t - k T_A}{T_A}\right)
$$

\begin{figure}[H]
\centering
\begin{tikzpicture}
\begin{axis}[
    width=10cm, height=4cm,
    xmin=-3.5, xmax=3.5,
    ymin=-0.30, ymax=1.2,
    axis x line=middle,
    axis y line=left,
    xtick={-3,-2,-1,1,2,3},
    xticklabels={$-3$, $-2$, $-1$, $1$, $2$, $3$},
    ytick={1},
    yticklabels={$1$},
    tick align=outside,
    every axis plot/.append style={line width=1.2pt},
]
\addplot[blue, smooth, domain=-3.4:-0.03, samples=150] {sin(deg(pi*x))/(pi*x)};
\addplot[blue, smooth, domain= 0.03: 3.4, samples=150] {sin(deg(pi*x))/(pi*x)};
\addplot[only marks, mark=*, mark size=3pt, mark options={fill=blue}] coordinates {(0,1)};
\end{axis}
\end{tikzpicture}
\caption{$\operatorname{si}$-Funktion: Maximum bei $(0,\,1)$, Nullstellen bei allen
ganzzahligen $x \neq 0$.}
\end{figure}

\textbf{Beispiel Aufgabe}

Substituion üben:

Korrelation mit einem cos signal und einem rechtecksignal, wo substituiert werden muss.
Höllisch schwierig.

$$
\begin{array}{rl}
x_1(t) = & B \cdot rect(\frac{t+\frac{3}{2}T}{5T}) \\
x_2(t) = & \pi \cdot cos(\frac{5 \pi t}{T}) \\
=>  & \Phi_{x_1x_2}(\tau)
\end{array}
$$

Die Lösung ist:
$$
- \frac{2}{5} BT \cdot cos(\frac{5 \pi \tau}{T})
$$

Additionstheorem vollständig erlernen.

$$
sin(a-b) + sin(a+b) = 2 \cdot sin(a) \cdot cos(b)
$$

$$
\begin{array}{rl}
    sin(t+n \cdot 2 \pi) = & sin(t) \\
    sin(t) = & -sin(t)
\end{array}
$$

% ================================================================
%   Vorrechenaufgabe 4 – Ergänzung: Teil d) und e)
% ================================================================

\subsection*{Vorrechenaufgabe 4 – Ergänzung (Teil d und e)}

Die Teile a)–c) (Abtastung, Midrise-Kennlinie, Quantisierung) sowie f)–g)
(Interpolation) sind in den obigen Abschnitten bereits ausgearbeitet. Es fehlen
noch der Quantisierungsfehler (Teil~d) und die entstehende Datenmenge (Teil~e).

\subsubsection*{d) Quantisierungsfehler}

\begin{beispiel}[Quantisierungsfehler der abgetasteten und quantisierten Folge]

\textbf{Gegeben:} Die Abtastwerte $x[n]$ und die zugehörigen quantisierten Werte
$\hat{x}[n]$ aus Teil~c) (Midrise, $N = 4$, $\Delta x = 0{,}2$,
Rekonstruktionsniveaus $\{-0{,}3;\,-0{,}1;\,0{,}1;\,0{,}3\}$).

\textbf{Gesucht:} Der Quantisierungsfehler $e[n] = x[n] - \hat{x}[n]$ und der
Nachweis der Fehlerschranke.

\textbf{Formel} (Quantisierungsfehler und Fehlerschranke,
Abschnitt~\ref{sec:quantisierungsfehler}):

$$
e[n] = x[n] - \hat{x}[n], \qquad |e[n]| \leq \frac{\Delta x}{2}
$$

\textbf{Rechnung:}

Mit $\Delta x = 0{,}2$ beträgt die Fehlerschranke $\tfrac{\Delta x}{2} = 0{,}1$.
Differenzbildung für jeden Abtastwert:

\begin{center}
\begin{tabular}{c|cccccccccc}
\toprule
$n$ & 0 & 1 & 2 & 3 & 4 & 5 & 6 & 7 & 8 & 9 \\
\midrule
$x[n]$        & $0{,}38$ & $0{,}22$ & $-0{,}05$ & $-0{,}31$ & $-0{,}39$ & $-0{,}17$ & $0{,}08$ & $0{,}28$ & $0{,}39$ & $0{,}15$ \\
$\hat{x}[n]$  & $0{,}3$  & $0{,}3$  & $-0{,}1$  & $-0{,}3$  & $-0{,}3$  & $-0{,}1$  & $0{,}1$  & $0{,}3$  & $0{,}3$  & $0{,}1$ \\
$e[n]$        & $0{,}08$ & $-0{,}08$& $0{,}05$  & $-0{,}01$ & $-0{,}09$ & $-0{,}07$ & $-0{,}02$& $-0{,}02$& $0{,}09$ & $0{,}05$ \\
\bottomrule
\end{tabular}
\end{center}

\textbf{Lösung:}

$$
\boxed{\;|e[n]| \leq \frac{\Delta x}{2} = 0{,}1 \quad \text{für alle } n\;}
$$

Der Betrag des Quantisierungsfehlers überschreitet nie die halbe Stufenhöhe.

\begin{figure}[H]
\centering
\begin{tikzpicture}
\begin{axis}[
    width=11cm, height=4.5cm,
    xlabel={Index $n$},
    ylabel={$e[n]$},
    xmin=-0.5, xmax=9.5,
    ymin=-0.13, ymax=0.13,
    axis lines=left,
    xtick={0,1,...,9},
    ytick={-0.1,0,0.1},
    yticklabels={$-0{,}1$, $0$, $0{,}1$},
    ymajorgrids=true,
    grid style={dashed, gray!40},
    every axis plot/.append style={line width=1.2pt},
]
% Fehlerschranken
\addplot[red, dashed, line width=0.8pt, domain=-0.5:9.5] {0.1};
\addplot[red, dashed, line width=0.8pt, domain=-0.5:9.5] {-0.1};
\addplot[blue, ycomb, mark=*, mark size=2.8pt, mark options={fill=blue}]
    coordinates {
        (0, 0.08) (1, -0.08) (2, 0.05) (3, -0.01) (4, -0.09)
        (5, -0.07) (6, -0.02) (7, -0.02) (8, 0.09) (9, 0.05)
    };
\end{axis}
\end{tikzpicture}
\caption{Quantisierungsfehler $e[n] = x[n] - \hat{x}[n]$. Alle Werte liegen
innerhalb der Schranken $\pm\Delta x/2 = \pm 0{,}1$ (rot gestrichelt).}
\end{figure}

\end{beispiel}

\subsubsection*{e) Datenmenge}

\begin{beispiel}[Datenrate bei binärer Codierung der Quantisierungsniveaus]

\textbf{Gegeben:} Abtastfrequenz $f_A = 10\,\text{kHz}$, Quantisierung mit
$N = 4$ Stufen. Die Niveaus werden natürlichen binären Zahlen zugeordnet.

\textbf{Gesucht:} Die entstehende Datenrate (Datenmenge pro Zeit).

\textbf{Formel} (Wortbreite und Datenrate, Abschnitt~\ref{sec:codierung-binaer}):

$$
w = \log_2 N, \qquad R = f_A \cdot w
$$

\textbf{Rechnung:}

Für $N = 4$ Quantisierungsniveaus werden

$$
w = \log_2 4 = 2 \ \text{Bit}
$$

pro Abtastwert benötigt. Bei $f_A = 10\,\text{kHz} = 10\,000\ \text{Abtastwerte/s}$
folgt

$$
R = f_A \cdot w = 10\,000\,\text{s}^{-1} \cdot 2\,\text{Bit}
  = 20\,000\ \frac{\text{Bit}}{\text{s}}.
$$

\textbf{Lösung:}

$$
\boxed{\;w = 2\ \text{Bit}, \qquad R = 20\,000\ \frac{\text{Bit}}{\text{s}}
= 20\ \frac{\text{kbit}}{\text{s}}\;}
$$

Jeder Abtastwert benötigt $2$~Bit; pro Sekunde entstehen somit $20$~kbit.

\end{beispiel}

% ================================================================
%   Aufgabe 4.1 – Datenrate, Datenmenge (Audio-CD)
% ================================================================

\subsection*{Aufgabe 4.1 – Datenrate und Datenmenge (Audio-CD)}

Auf einer Audio-CD wird zweikanalig (Stereo) gespeichert. Jeder der beiden Kanäle
wird mit $f_A = 44{,}1\,\text{kHz}$ abgetastet und mit $w = 16$~Bit quantisiert.

\subsubsection*{a) Datenrate des Stereosignals}

\begin{beispiel}[Datenrate eines Stereosignals]

\textbf{Gegeben:} $K = 2$ Kanäle, $f_A = 44{,}1\,\text{kHz}$, $w = 16$~Bit.

\textbf{Gesucht:} Die Datenrate $R$ in kbit/s.

\textbf{Formel} (Datenrate, Abschnitt~\ref{sec:codierung-binaer}):

$$
R = K \cdot f_A \cdot w
$$

\textbf{Rechnung:}

$$
R = 2 \cdot 44\,100\,\text{s}^{-1} \cdot 16\,\text{Bit}
  = 1\,411\,200\ \frac{\text{Bit}}{\text{s}}.
$$

\textbf{Lösung:}

$$
\boxed{\;R = 1\,411\,200\ \frac{\text{Bit}}{\text{s}}
= 1411{,}2\ \frac{\text{kbit}}{\text{s}} \approx 1{,}41\ \frac{\text{Mbit}}{\text{s}}\;}
$$

\end{beispiel}

\subsubsection*{b) Maximal reproduzierbare Signalfrequenz}

\begin{beispiel}[Maximale Signalfrequenz nach dem Abtasttheorem]

\textbf{Gegeben:} $f_A = 44{,}1\,\text{kHz}$ (Quantisierung vernachlässigt).

\textbf{Gesucht:} Die maximale fehlerfrei reproduzierbare Signalfrequenz
$f_{\max}$.

\textbf{Formel} (Abtasttheorem, Abschnitt~\ref{sec:abtasttheorem}):

$$
f_A > 2\, f_{\max} \quad\Longrightarrow\quad f_{\max} = \frac{f_A}{2}
$$

\textbf{Rechnung:}

$$
f_{\max} = \frac{44{,}1\,\text{kHz}}{2} = 22{,}05\,\text{kHz}.
$$

\textbf{Lösung:}

$$
\boxed{\;f_{\max} = \frac{f_A}{2} = 22{,}05\,\text{kHz}\;}
$$

Dies deckt den gesamten für den Menschen hörbaren Bereich (bis ca.\ $20\,\text{kHz}$) ab.

\end{beispiel}

\subsubsection*{c) Datenmenge der Goldberg-Variationen}

\begin{beispiel}[Datenmenge zweier Aufnahmen unterschiedlicher Dauer]

\textbf{Gegeben:} Datenrate $R = 1\,411\,200\ \text{Bit/s}$ (aus Teil~a).
Dauer der Aufnahme von 1955: $t_1 = 38\,\text{min}\,25\,\text{s}$; von 1982:
$t_2 = 51\,\text{min}\,20\,\text{s}$.

\textbf{Gesucht:} Die zu speichernde Datenmenge $D$ beider Aufnahmen.

\textbf{Formel} (Datenmenge aus Datenrate und Dauer):

$$
D = R \cdot t
$$

\textbf{Rechnung:}

Umrechnung der Spielzeiten in Sekunden:

$$
t_1 = 38 \cdot 60 + 25 = 2305\,\text{s}, \qquad
t_2 = 51 \cdot 60 + 20 = 3080\,\text{s}.
$$

Aufnahme 1955:

$$
D_1 = 1\,411\,200\ \tfrac{\text{Bit}}{\text{s}} \cdot 2305\,\text{s}
    = 3\,252\,816\,000\ \text{Bit}
    = \frac{D_1}{8} = 406\,602\,000\ \text{Byte} \approx 406{,}6\ \text{MB}.
$$

Aufnahme 1982:

$$
D_2 = 1\,411\,200\ \tfrac{\text{Bit}}{\text{s}} \cdot 3080\,\text{s}
    = 4\,346\,496\,000\ \text{Bit}
    = 543\,312\,000\ \text{Byte} \approx 543{,}3\ \text{MB}.
$$

\textbf{Lösung:}

$$
\boxed{\;D_1 \approx 406{,}6\ \text{MB} \; (2305\,\text{s}), \qquad
D_2 \approx 543{,}3\ \text{MB} \; (3080\,\text{s})\;}
$$

Die langsamere Aufnahme von 1982 benötigt entsprechend ihrer längeren Spieldauer
mehr Speicher (Umrechnung mit $1\,\text{MB} = 10^6\,\text{Byte}$; mit
$1\,\text{MiB} = 2^{20}\,\text{Byte}$ ergäben sich $387{,}8$ bzw.\ $518{,}1$~MiB).

\end{beispiel}

% ================================================================
%   Aufgabe 4.2 – Quantisierungsfehler
% ================================================================

\subsection*{Aufgabe 4.2 – Quantisierungsfehler}

Ein zeitdiskretes Signal mit Wertebereich $[-1,\,1]$ (Spannweite
$x_{\max} - x_{\min} = 2$) soll so quantisiert werden, dass die Leistung des
Quantisierungsfehlers $P_e = 0{,}01$ beträgt.

\subsubsection*{a) Erforderliche Stufenhöhe}

\begin{beispiel}[Stufenhöhe aus vorgegebener Fehlerleistung]

\textbf{Gegeben:} Fehlerleistung $P_e = 0{,}01$.

\textbf{Gesucht:} Die maximal zulässige Stufenhöhe $\Delta x$.

\textbf{Formel} (Leistung des Quantisierungsfehlers,
Abschnitt~\ref{sec:quantisierungsfehler}):

$$
P_e = \frac{\Delta x^2}{12}
$$

\textbf{Rechnung:}

Nach $\Delta x$ auflösen:

$$
\Delta x = \sqrt{12\, P_e} = \sqrt{12 \cdot 0{,}01} = \sqrt{0{,}12}.
$$

\textbf{Lösung:}

$$
\boxed{\;\Delta x = \sqrt{0{,}12} \approx 0{,}346\;}
$$

Um $P_e = 0{,}01$ nicht zu überschreiten, darf die Stufenhöhe höchstens
$\Delta x \approx 0{,}346$ betragen.

\end{beispiel}

\subsubsection*{b) Mindestanzahl der Bits}

\begin{beispiel}[Wortbreite bei binärer Codierung]

\textbf{Gegeben:} Spannweite $x_{\max} - x_{\min} = 2$, maximale Stufenhöhe
$\Delta x \approx 0{,}346$ (aus Teil~a).

\textbf{Gesucht:} Die mindestens erforderliche Wortbreite $w$ (Bit pro
Abtastwert).

\textbf{Formel} (Stufenzahl und Wortbreite,
Abschnitte~\ref{sec:quantisierung} und~\ref{sec:codierung-binaer}):

$$
N = \frac{x_{\max} - x_{\min}}{\Delta x}, \qquad w = \left\lceil \log_2 N \right\rceil
$$

\textbf{Rechnung:}

Erforderliche Stufenzahl:

$$
N = \frac{2}{\sqrt{0{,}12}} = \frac{2}{0{,}346} \approx 5{,}77
\quad\Longrightarrow\quad N \geq 6.
$$

Da nur ganze Bit vergeben werden, ist auf die nächste Zweierpotenz aufzurunden:

$$
w = \left\lceil \log_2 6 \right\rceil = \left\lceil 2{,}585 \right\rceil = 3\,\text{Bit}
\qquad (2^3 = 8 \geq 6).
$$

Kontrolle: Mit $w = 3$~Bit ($8$~Stufen) wird
$\Delta x = \tfrac{2}{8} = 0{,}25$ und

$$
P_e = \frac{0{,}25^2}{12} \approx 0{,}0052 \leq 0{,}01. \quad\checkmark
$$

Mit nur $w = 2$~Bit ($4$~Stufen) wäre $\Delta x = 0{,}5$ und
$P_e = \tfrac{0{,}25}{12} \approx 0{,}0208 > 0{,}01$ -- unzureichend.

\textbf{Lösung:}

$$
\boxed{\;w = 3\ \text{Bit pro Abtastwert}\;}
$$

\end{beispiel}

% ================================================================
%   Aufgabe 4.3 – Datenmenge und Übertragungsdauer
% ================================================================

\subsection*{Aufgabe 4.3 – Datenmenge und Übertragungsdauer}

\subsubsection*{Datenmenge eines Bildes}

\begin{beispiel}[Datenmenge eines Vollbildes]

\textbf{Gegeben:} Bild mit $1920 \times 1080$ Bildpunkten, jeder Bildpunkt mit
$24$~Bit dargestellt.

\textbf{Gesucht:} Die Datenmenge $D$ des Bildes.

\textbf{Formel} (Datenmenge, Abschnitt~\ref{sec:codierung-binaer}):

$$
D = N_{\text{Pixel}} \cdot w = (n_x \cdot n_y) \cdot w
$$

\textbf{Rechnung:}

$$
D = 1920 \cdot 1080 \cdot 24\,\text{Bit}
  = 2\,073\,600 \cdot 24\,\text{Bit}
  = 49\,766\,400\ \text{Bit}.
$$

In Byte:

$$
D = \frac{49\,766\,400}{8}\ \text{Byte} = 6\,220\,800\ \text{Byte}
  \approx 6{,}22\ \text{MB}.
$$

\textbf{Lösung:}

$$
\boxed{\;D = 49\,766\,400\ \text{Bit} \approx 6{,}22\ \text{MB}\;}
$$

\end{beispiel}

\subsubsection*{a)–d) Übertragungsdauer über verschiedene Verbindungen}

\begin{beispiel}[Übertragungsdauer aus Datenmenge und Übertragungsrate]

\textbf{Gegeben:} Datenmenge $D = 49\,766\,400\ \text{Bit}$; Übertragungsraten
EDGE $R_a = 236{,}8\,\text{kbit/s}$, UMTS $R_b = 7{,}2\,\text{Mbit/s}$, LTE
$R_c = 150\,\text{Mbit/s}$, VDSL $R_d = 210\,\text{Mbit/s}$.

\textbf{Gesucht:} Die jeweilige Übertragungsdauer $t = D / R$.

\textbf{Formel} (Übertragungsdauer):

$$
t = \frac{D}{R}
$$

\textbf{Rechnung:}

$$
t_a = \frac{49\,766\,400\,\text{Bit}}{236\,800\ \text{Bit/s}} \approx 210{,}2\,\text{s}
    \approx 3\,\text{min}\,30\,\text{s} \quad (\text{EDGE})
$$

$$
t_b = \frac{49\,766\,400\,\text{Bit}}{7\,200\,000\ \text{Bit/s}} \approx 6{,}91\,\text{s}
    \quad (\text{UMTS})
$$

$$
t_c = \frac{49\,766\,400\,\text{Bit}}{150\,000\,000\ \text{Bit/s}} \approx 0{,}332\,\text{s}
    \quad (\text{LTE})
$$

$$
t_d = \frac{49\,766\,400\,\text{Bit}}{210\,000\,000\ \text{Bit/s}} \approx 0{,}237\,\text{s}
    \quad (\text{VDSL})
$$

\textbf{Lösung:}

$$
\boxed{\;t_a \approx 210{,}2\,\text{s}, \quad t_b \approx 6{,}91\,\text{s}, \quad
t_c \approx 0{,}332\,\text{s}, \quad t_d \approx 0{,}237\,\text{s}\;}
$$

Mit steigender Übertragungsrate sinkt die Übertragungsdauer entsprechend: von über
drei Minuten (EDGE) auf einen Bruchteil einer Sekunde (VDSL).

\end{beispiel}

% ================================================================
%   Aufgabe 4.4 – Adaptives Scheinwerferlicht
% ================================================================

\subsection*{Aufgabe 4.4 – Adaptives Scheinwerferlicht}

Die Scheinwerferausrichtung wird über den Lenkradeinschlagwinkel gesteuert. Das
Lenkrad kann maximal um $[-360^\circ,\, 360^\circ]$ eingeschlagen werden. Bei
Maximalausschlag dürfen die Scheinwerfer für Geschwindigkeiten unter
$30\,\tfrac{\text{km}}{\text{h}}$ um $20^\circ$, für Geschwindigkeiten über
$30\,\tfrac{\text{km}}{\text{h}}$ um $10^\circ$ verstellt werden. Die Verstellung
erfolgt in $5^\circ$-Schritten.

\subsubsection*{a) Benötigte Sensoren und Informationen}

\begin{beispiel}[Sensorik des adaptiven Scheinwerfers]

\textbf{Gegeben:} Steuerung der Scheinwerfer nach Lenkradwinkel und
geschwindigkeitsabhängiger Fallunterscheidung.

\textbf{Gesucht:} Anzahl und Aufgabe der benötigten Sensoren.

\textbf{Lösung:}

Es werden \textbf{zwei} Sensoren benötigt:

\begin{itemize}
  \item \textbf{Lenkwinkelsensor:} misst den Lenkradeinschlagwinkel im Bereich
        $[-360^\circ,\, 360^\circ]$. Er liefert die Richtungsinformation für die
        Scheinwerferverstellung.
  \item \textbf{Geschwindigkeitssensor:} misst die Fahrgeschwindigkeit. Benötigt
        wird nur die Unterscheidung, ob $v < 30\,\tfrac{\text{km}}{\text{h}}$
        (Fall~1, $\pm 20^\circ$) oder $v > 30\,\tfrac{\text{km}}{\text{h}}$
        (Fall~2, $\pm 10^\circ$) vorliegt.
\end{itemize}

$$
\boxed{\;\text{2 Sensoren: Lenkwinkel (Richtung) und Geschwindigkeit (Fallwahl)}\;}
$$

\end{beispiel}

\subsubsection*{b) Quantisierungskennlinien und Bitbedarf}

\begin{beispiel}[Kennlinien und Wortbreite der Sensoren]

\textbf{Gegeben:} Lenkwinkelbereich $[-360^\circ, 360^\circ]$ (Spannweite
$720^\circ$); Scheinwerfer in $5^\circ$-Schritten, maximal $\pm 20^\circ$
(Fall~1) bzw.\ $\pm 10^\circ$ (Fall~2).

\textbf{Gesucht:} Die Quantisierungskennlinien sowie die Wortbreite $w$ für die
Codierung.

\textbf{Formel} (Stufenzahl und Wortbreite,
Abschnitte~\ref{sec:quantisierungskennlinie} und~\ref{sec:codierung-binaer}):

$$
w = \left\lceil \log_2 N \right\rceil
$$

\textbf{Rechnung:}

\emph{Fall 1} ($v < 30\,\tfrac{\text{km}}{\text{h}}$): Scheinwerferpositionen
$\{-20, -15, -10, -5, 0, 5, 10, 15, 20\}^\circ$, also $N_1 = 9$ Stufen. Der
gesamte Lenkbereich $720^\circ$ wird auf $9$~Stufen abgebildet (Midtread, damit
Geradeausfahrt $0^\circ \to 0^\circ$), Stufenbreite $720^\circ/9 = 80^\circ$.

\emph{Fall 2} ($v > 30\,\tfrac{\text{km}}{\text{h}}$): Scheinwerferpositionen
$\{-10, -5, 0, 5, 10\}^\circ$, also $N_2 = 5$ Stufen, Stufenbreite
$720^\circ/5 = 144^\circ$.

Der Lenkwinkelsensor muss den anspruchsvolleren Fall~1 auflösen:

$$
w_{\text{Lenk}} = \left\lceil \log_2 9 \right\rceil = 4\,\text{Bit}
\qquad (2^4 = 16 \geq 9).
$$

Der Geschwindigkeitssensor muss nur zwei Zustände unterscheiden:

$$
w_{\text{Geschw}} = \left\lceil \log_2 2 \right\rceil = 1\,\text{Bit}.
$$

\textbf{Lösung:}

$$
\boxed{\;w_{\text{Lenk}} = 4\ \text{Bit}, \quad w_{\text{Geschw}} = 1\ \text{Bit},
\quad w_{\text{ges}} = 5\ \text{Bit}\;}
$$

\begin{figure}[H]
\centering
\begin{tikzpicture}
\begin{axis}[
    width=8cm, height=6.5cm,
    xlabel={Lenkwinkel [Grad]},
    ylabel={Scheinwerfer [Grad]},
    title={Fall 1: $v < 30\,$km/h ($N=9$)},
    title style={font=\small},
    xmin=-360, xmax=360, ymin=-24, ymax=24,
    axis lines=center,
    xtick={-360,-180,180,360},
    ytick={-20,-10,10,20},
    tick label style={font=\footnotesize},
    every axis plot/.append style={line width=1.2pt},
]
\addplot[blue] coordinates {(-360,-20) (-280,-20)};
\addplot[blue] coordinates {(-280,-15) (-200,-15)};
\addplot[blue] coordinates {(-200,-10) (-120,-10)};
\addplot[blue] coordinates {(-120,-5) (-40,-5)};
\addplot[blue] coordinates {(-40,0) (40,0)};
\addplot[blue] coordinates {(40,5) (120,5)};
\addplot[blue] coordinates {(120,10) (200,10)};
\addplot[blue] coordinates {(200,15) (280,15)};
\addplot[blue] coordinates {(280,20) (360,20)};
\end{axis}
\end{tikzpicture}
\hfill
\begin{tikzpicture}
\begin{axis}[
    width=8cm, height=6.5cm,
    xlabel={Lenkwinkel [Grad]},
    ylabel={Scheinwerfer [Grad]},
    title={Fall 2: $v > 30\,$km/h ($N=5$)},
    title style={font=\small},
    xmin=-360, xmax=360, ymin=-14, ymax=14,
    axis lines=center,
    xtick={-360,-180,180,360},
    ytick={-10,-5,5,10},
    tick label style={font=\footnotesize},
    every axis plot/.append style={line width=1.2pt},
]
\addplot[red] coordinates {(-360,-10) (-216,-10)};
\addplot[red] coordinates {(-216,-5) (-72,-5)};
\addplot[red] coordinates {(-72,0) (72,0)};
\addplot[red] coordinates {(72,5) (216,5)};
\addplot[red] coordinates {(216,10) (360,10)};
\end{axis}
\end{tikzpicture}
\caption{Quantisierungskennlinien (Midtread) des Lenkwinkelsensors: links Fall~1
($9$~Stufen, $\pm 20^\circ$), rechts Fall~2 ($5$~Stufen, $\pm 10^\circ$).}
\end{figure}

\end{beispiel}

\subsubsection*{c) Mögliches Problem und Abhilfe}

\begin{beispiel}[Unterabtastung und Flackern]

\textbf{Gegeben:} Zeitdiskrete Abtastung des schnell veränderlichen
Lenkwinkelverlaufs.

\textbf{Gesucht:} Ein Realisierungsproblem und geeignete Gegenmaßnahmen.

\textbf{Lösung:}

Wird die Abtastrate zu niedrig gewählt (z.\,B.\ $1\,\text{Hz}$), so verletzt sie
für schnelle Lenkbewegungen das Abtasttheorem
(Abschnitt~\ref{sec:abtasttheorem}). Es kommt zu \textbf{Unterabtastung
(Aliasing)}: Die Scheinwerfer folgen dem Lenkeinschlag nur verzögert oder falsch,
weil zwischen zwei Abtastzeitpunkten große Winkeländerungen unbemerkt bleiben.
Zusätzlich kann der quantisierte Wert an einer Stufengrenze zwischen zwei Niveaus
\textbf{hin- und herspringen (Flackern)}.

$$
\boxed{\;\text{Abhilfe: höhere Abtastrate (Abtasttheorem) + Hysterese gegen Flackern}\;}
$$

Eine höhere Abtastfrequenz stellt sicher, dass die Scheinwerfer der Lenkbewegung
zeitgerecht folgen; eine Hysterese an den Stufengrenzen unterdrückt das Flackern.

\end{beispiel}

\subsubsection*{d) Ausrichtung bei 1 Hz Abtastung mit Beschleunigung}

\begin{beispiel}[Abtastung des Lenkverlaufs und Scheinwerferausrichtung]

\textbf{Gegeben:} Lenkwinkelverlauf gemäß Abbildung, Abtastrate $f_A = 1\,\text{Hz}$
(Abtastung bei $t = 0, 1, \ldots, 12\,\text{s}$). Das Fahrzeug fährt bis
$t = 6\,\text{s}$ mit $25\,\tfrac{\text{km}}{\text{h}}$ (Fall~1, $\pm 20^\circ$)
und beschleunigt danach auf $35\,\tfrac{\text{km}}{\text{h}}$ (Fall~2,
$\pm 10^\circ$).

\textbf{Gesucht:} Die quantisierte Scheinwerferausrichtung zu jedem
Abtastzeitpunkt.

\textbf{Formel} (lineare Zuordnung Lenkwinkel $\varphi$ $\to$ Scheinwerferwinkel,
anschließend Quantisierung in $5^\circ$-Schritten,
Abschnitt~\ref{sec:quantisierungskennlinie}):

$$
\text{Fall 1: } \hat{s} = \operatorname{round}_{5^\circ}\!\left(\frac{\varphi}{360^\circ}\cdot 20^\circ\right),
\qquad
\text{Fall 2: } \hat{s} = \operatorname{round}_{5^\circ}\!\left(\frac{\varphi}{360^\circ}\cdot 10^\circ\right)
$$

\textbf{Rechnung:}

Zunächst werden die Lenkwinkel $\varphi(t)$ zu den ganzzahligen Sekunden aus dem
Diagramm abgelesen (Werte gerundet). Für $t < 6\,\text{s}$ gilt Fall~1
(Skalierung $\varphi/18$), für $t \geq 6\,\text{s}$ Fall~2 (Skalierung
$\varphi/36$); anschließend wird auf das nächste Vielfache von $5^\circ$
gerundet und auf $\pm 20^\circ$ bzw.\ $\pm 10^\circ$ begrenzt:

\begin{center}
\begin{tabular}{c|c|c|c|c|c}
\toprule
$t$/s & $\varphi$/Grad & $v$/(km/h) & Fall & ideal/Grad & $\hat{s}$/Grad \\
\midrule
0  & $0$    & 25 & 1 & $0{,}0$   & $0$   \\
1  & $-80$  & 25 & 1 & $-4{,}4$  & $-5$  \\
2  & $+100$ & 25 & 1 & $+5{,}6$  & $+5$  \\
3  & $+190$ & 25 & 1 & $+10{,}6$ & $+10$ \\
4  & $-350$ & 25 & 1 & $-19{,}4$ & $-20$ \\
5  & $+30$  & 25 & 1 & $+1{,}7$  & $0$   \\
\midrule
6  & $+150$ & 35 & 2 & $+4{,}2$  & $+5$  \\
7  & $-40$  & 35 & 2 & $-1{,}1$  & $0$   \\
8  & $+120$ & 35 & 2 & $+3{,}3$  & $+5$  \\
9  & $+150$ & 35 & 2 & $+4{,}2$  & $+5$  \\
10 & $-110$ & 35 & 2 & $-3{,}1$  & $-5$  \\
11 & $-40$  & 35 & 2 & $-1{,}1$  & $0$   \\
12 & $+20$  & 35 & 2 & $+0{,}6$  & $0$   \\
\bottomrule
\end{tabular}
\end{center}

\textbf{Lösung:}

$$
\boxed{\;
\begin{array}{l}
\hat{s}(0\ldots5) = (0,\, -5,\, +5,\, +10,\, -20,\, 0)^\circ \quad (\text{Fall 1}) \\[4pt]
\hat{s}(6\ldots12) = (+5,\, 0,\, +5,\, +5,\, -5,\, 0,\, 0)^\circ \quad (\text{Fall 2})
\end{array}\;}
$$

Ab $t = 6\,\text{s}$ (Beschleunigung über $30\,\tfrac{\text{km}}{\text{h}}$) wird
der zulässige Verstellbereich auf $\pm 10^\circ$ halbiert; identische
Lenkwinkel führen dann zu kleineren Scheinwerferauslenkungen. Wegen der niedrigen
Abtastrate von $1\,\text{Hz}$ werden die schnellen Ausschläge des Lenkverlaufs nur
grob erfasst (vgl.\ Teil~c).

\end{beispiel}

\section{Übung 5 – Zustandsraummodell eines Masse-Feder-Systems}

Betrachtet wird ein gedämpfter Oszillator (Masse-Feder-System) mit der Masse
$M$, der Federkonstante $k_F$ und der Dämpfungskonstante $b$. Die Auslenkung
$x(t)$ und die Geschwindigkeit $v(t) = \dot{x}(t)$ bilden den Zustandsvektor
$\vec{z}(t) = (x(t),\, v(t))^{\mathsf{T}}$. Mit den Hilfsgrößen $\omega_0^2 =
\tfrac{k_F}{m}$ und $2\gamma = \tfrac{b}{m}$ lautet die Bewegungsgleichung des
freien Oszillators

$$
\frac{\mathrm{d}v(t)}{\mathrm{d}t} = -\omega_0^2\, x(t) - 2\gamma\, v(t).
\qquad (3)
$$

Wird das System zusätzlich über einen Exzenter mit der Kraft
$F(t) = \hat{F}\cos(\omega t)$ angeregt, ergibt sich

$$
\frac{\mathrm{d}v(t)}{\mathrm{d}t} = \frac{\hat{F}}{m}\cos(\omega t)
- \omega_0^2\, x(t) - 2\gamma\, v(t).
\qquad (4)
$$

% ================================================================
\subsection*{Vorüberlegung – Warum überhaupt ein Zustandsraummodell?}
% ================================================================

\begin{gedankenexperiment}[Die Grundidee]

Die Bewegungsgleichung des Oszillators ist eine \emph{Differentialgleichung
zweiter Ordnung} in $x(t)$:

$$
m\,\ddot{x}(t) = -k_F\, x(t) - b\,\dot{x}(t).
$$

Mit einer DGL zweiter Ordnung kann man schlecht „rechnen``: Sie enthält eine
zweite Ableitung, ist nicht direkt in Matrixform darstellbar und nicht ohne
Weiteres numerisch (Schritt für Schritt) lösbar.

Die zentrale Idee des Zustandsraummodells ist daher: \textbf{Wir ersetzen eine
DGL der Ordnung $n$ durch ein System aus $n$ gekoppelten DGLen erster Ordnung.}
DGLen erster Ordnung haben die Form „Ableitung $=$ Funktion der aktuellen
Werte`` und lassen sich kompakt als

$$
\frac{\mathrm{d}\vec{z}(t)}{\mathrm{d}t} = \mathbf{A}\,\vec{z}(t) + \mathbf{B}\,u(t)
$$

schreiben. Der \emph{Zustandsvektor} $\vec{z}(t)$ fasst dabei genau die
Information zusammen, die man zu einem Zeitpunkt $t$ kennen muss, um die gesamte
Zukunft des Systems zu berechnen.

\end{gedankenexperiment}

% ================================================================
\subsection*{Das universelle Rezept: Von der DGL zum Zustandsraummodell}
% ================================================================

Das folgende Vorgehen funktioniert für \emph{jede} lineare DGL mit konstanten
Koeffizienten und ist das Werkzeug für alle Aufgaben dieses Typs. Ziel ist immer
die Standardform (vgl.\ Abschnitt~\ref{sec:zustandsraum})

$$
\frac{\mathrm{d}\vec{z}(t)}{\mathrm{d}t} = \mathbf{A}\,\vec{z}(t) + \mathbf{B}\,u(t),
\qquad
\vec{y}(t) = \mathbf{C}\,\vec{z}(t) + \mathbf{D}\,u(t).
$$

\begin{bemerkung}[Kochrezept in 6 Schritten]

\textbf{Schritt 1 – Ordnung $n$ bestimmen.}
Suche die höchste vorkommende Ableitung der gesuchten Größe. Ihre Ordnung $n$
ist die Anzahl der benötigten Zustandsvariablen. \emph{Warum:} Der Zustand muss
die Funktion und alle Ableitungen bis zur $(n-1)$-ten enthalten, denn genau
diese Werte braucht man als „Startbedingungen``, um die Lösung eindeutig
fortzuschreiben.

\textbf{Schritt 2 – Zustandsvariablen wählen.}
Setze $z_1 = $ (die Größe selbst), $z_2 = \dot{z}_1$, $z_3 = \ddot{z}_1$, \ldots,
$z_n = z_1^{(n-1)}$. \emph{Warum:} Diese „Ableitungskette`` ist die natürliche
Wahl (Phasenraum) und führt direkt auf die Schritte 3 und 4.

\textbf{Schritt 3 – Die trivialen Zustandsgleichungen.}
Aus der Definition folgt unmittelbar $\dot{z}_1 = z_2$, $\dot{z}_2 = z_3$,
\ldots, $\dot{z}_{n-1} = z_n$. \emph{Warum:} Das sind keine neuen Annahmen,
sondern nur die Umbenennung „die Ableitung von $z_i$ heißt $z_{i+1}$``.

\textbf{Schritt 4 – Die „echte`` Zustandsgleichung.}
Löse die ursprüngliche DGL nach der höchsten Ableitung $z_1^{(n)} = \dot{z}_n$
auf und ersetze rechts alle Ableitungen durch die Zustandsvariablen $z_i$ (und
ggf.\ das Eingangssignal $u$). \emph{Warum:} Damit ist auch die letzte fehlende
Ableitung $\dot{z}_n$ allein durch aktuelle Werte ausgedrückt -- das System ist
geschlossen.

\textbf{Schritt 5 – Matrixform: $\mathbf{A}$ und $\mathbf{B}$ ablesen.}
Schreibe die $n$ Gleichungen $\dot{z}_i = \ldots$ untereinander und sortiere
nach den $z_j$ und nach $u$. Die Koeffizienten vor den $z_j$ bilden zeilenweise
die Systemmatrix $\mathbf{A}$, die Koeffizienten vor $u$ die Eingangsmatrix
$\mathbf{B}$. \emph{Warum:} Das ist nur das Umschreiben eines linearen
Gleichungssystems in Matrix-Vektor-Schreibweise.

\textbf{Schritt 6 – Ausgangsgleichung: $\mathbf{C}$ und $\mathbf{D}$.}
Frage: \emph{Welche Größe wird gemessen/beobachtet?} Drücke diese durch die
Zustände (und ggf.\ $u$) aus. Die Koeffizienten bilden $\mathbf{C}$ (vor $z_j$)
und $\mathbf{D}$ (vor $u$). \emph{Warum:} Der Zustand $\vec{z}$ ist intern; das
Modell muss noch sagen, was davon nach außen sichtbar ist.

\textbf{Kontrolle:} Dimensionen müssen passen --
$\mathbf{A}\!:\,n\times n$, $\mathbf{B}\!:\,n\times L$,
$\mathbf{C}\!:\,M\times n$, $\mathbf{D}\!:\,M\times L$
($n$ Zustände, $L$ Eingänge, $M$ Ausgänge).

\end{bemerkung}

% ================================================================
\subsection*{Teil a) – Zustandsraummodell des freien Oszillators}
% ================================================================

\begin{beispiel}[Zustandsraummodell aus DGL (3)]

\textbf{Gegeben:} Die DGL (3) des freien gedämpften Oszillators,
$\dot{v}(t) = -\omega_0^2\, x(t) - 2\gamma\, v(t)$, mit $v(t) = \dot{x}(t)$.
Vorgegebener Zustandsvektor $\vec{z}(t) = (x(t),\, v(t))^{\mathsf{T}}$. Gemessen
wird die Auslenkung $x(t)$ (Skala in Bild~1).

\textbf{Gesucht:} Das vollständige Zustandsraummodell, d.\,h.\ die Matrizen
$\mathbf{A}$, $\mathbf{B}$, $\mathbf{C}$, $\mathbf{D}$ der Zustands- und
Ausgangsgleichung.

\textbf{Formel} (zeitinvariantes Zustandsraummodell,
Abschnitt~\ref{sec:zustandsraum}):

$$
\frac{\mathrm{d}\vec{z}(t)}{\mathrm{d}t} = \mathbf{A}\,\vec{z}(t) + \mathbf{B}\,u(t),
\qquad
y(t) = \mathbf{C}\,\vec{z}(t) + \mathbf{D}\,u(t).
$$

\textbf{Rechnung:} Wir gehen das Kochrezept Schritt für Schritt durch.

\medskip
\textit{Schritt 1 – Ordnung.}
Die zugrunde liegende Bewegungsgleichung $m\ddot{x} = -k_F x - b\dot{x}$ enthält
die zweite Ableitung $\ddot{x}$. Die Ordnung ist also $n = 2$
$\Rightarrow$ wir brauchen \textbf{zwei} Zustandsvariablen. Dies ist in der
Aufgabe bereits durch $\vec{z} = (x, v)^{\mathsf{T}}$ vorweggenommen.

\medskip
\textit{Schritt 2 – Zustandsvariablen.}
Gemäß der Ableitungskette wählen wir

$$
z_1(t) = x(t) \quad\text{(Auslenkung)}, \qquad
z_2(t) = \dot{x}(t) = v(t) \quad\text{(Geschwindigkeit)}.
$$

\emph{Woher:} Position und Geschwindigkeit sind genau die zwei Größen, die man
zu einem Zeitpunkt kennen muss, um die weitere Schwingung eindeutig zu bestimmen
(physikalischer Anfangszustand).

\medskip
\textit{Schritt 3 – Triviale Zustandsgleichung.}
Aus $z_2 = \dot{z}_1$ folgt unmittelbar

$$
\dot{z}_1(t) = \dot{x}(t) = v(t) = z_2(t).
$$

\emph{Warum:} Das ist keine Physik, sondern reine Definition -- die erste Zeile
des Modells „verkettet`` nur $x$ und $v$.

\medskip
\textit{Schritt 4 – Echte Zustandsgleichung.}
Die zweite Zeile ist die gegebene DGL (3), bereits nach der höchsten Ableitung
$\dot{v} = \ddot{x}$ aufgelöst:

$$
\dot{z}_2(t) = \dot{v}(t) = -\omega_0^2\, x(t) - 2\gamma\, v(t)
            = -\omega_0^2\, z_1(t) - 2\gamma\, z_2(t).
$$

\emph{Woher:} direkt aus der Bewegungsgleichung; rechts wurde nur
$x \to z_1$, $v \to z_2$ eingesetzt.

\medskip
\textit{Schritt 5 – Matrixform.}
Wir schreiben beide Zeilen untereinander und sortieren nach $z_1, z_2$:

$$
\begin{aligned}
\dot{z}_1 &= 0\cdot z_1 + 1\cdot z_2, \\
\dot{z}_2 &= -\omega_0^2\cdot z_1 - 2\gamma\cdot z_2.
\end{aligned}
$$

Die Koeffizienten in die Matrix übertragen (erste Zeile: $0,\,1$; zweite Zeile:
$-\omega_0^2,\,-2\gamma$):

$$
\frac{\mathrm{d}}{\mathrm{d}t}
\begin{pmatrix} z_1 \\ z_2 \end{pmatrix}
=
\begin{pmatrix} 0 & 1 \\ -\omega_0^2 & -2\gamma \end{pmatrix}
\begin{pmatrix} z_1 \\ z_2 \end{pmatrix}.
$$

Der freie Oszillator hat \emph{kein} Eingangssignal, daher $\mathbf{B} = \vec{0}$
und der Term $\mathbf{B}\,u(t)$ entfällt.

\medskip
\textit{Schritt 6 – Ausgangsgleichung.}
Gemessen wird laut Aufgabe nur die Auslenkung, also $y(t) = x(t) = z_1(t)$:

$$
y(t) = \underbrace{(1\;\; 0)}_{\mathbf{C}}
\begin{pmatrix} z_1 \\ z_2 \end{pmatrix} = z_1(t),
\qquad \mathbf{D} = 0.
$$

\emph{Warum:} Der Zustand enthält auch die (nicht abgelesene) Geschwindigkeit
$z_2$; $\mathbf{C}$ wählt aus, dass davon nur $z_1 = x$ sichtbar ist.

\textbf{Lösung:}

$$
\boxed{
\frac{\mathrm{d}\vec{z}(t)}{\mathrm{d}t}
= \underbrace{\begin{pmatrix} 0 & 1 \\ -\omega_0^2 & -2\gamma \end{pmatrix}}_{\mathbf{A}}
\,\vec{z}(t),
\qquad
y(t) = \underbrace{(1\;\; 0)}_{\mathbf{C}}\,\vec{z}(t)
}
$$

\textit{Kontrolle:} $\mathbf{A}$ ist $2\times 2$, $\mathbf{C}$ ist $1\times 2$ --
passend zu $n=2$ Zuständen und $M=1$ Ausgang. \checkmark

Das System ist \emph{autonom} (homogen): Ohne äußere Anregung wird die Dynamik
allein durch $\mathbf{A}$ und den Anfangszustand $\vec{z}_0$ bestimmt. Die
abklingende Schwingung in Bild~1 ist die Eigenbewegung aus einer impulsartigen
Anfangsauslenkung.

\end{beispiel}

% ================================================================
\subsection*{Teil b) – Zustandsraummodell des angeregten Oszillators}
% ================================================================

\begin{beispiel}[Zustandsraummodell aus DGL (4)]

\textbf{Gegeben:} Die DGL (4) des fremderregten Oszillators,
$\dot{v}(t) = \tfrac{\hat{F}}{m}\cos(\omega t) - \omega_0^2 x(t) - 2\gamma v(t)$,
mit der Anregungskraft $F(t) = \hat{F}\cos(\omega t)$ als Eingangssignal.
Zustandsvektor $\vec{z}(t) = (x(t),\, v(t))^{\mathsf{T}}$.

\textbf{Gesucht:} Das vollständige Zustandsraummodell mit $\mathbf{A}$,
$\mathbf{B}$, $\mathbf{C}$, $\mathbf{D}$.

\textbf{Formel} (zeitinvariantes Zustandsraummodell,
Abschnitt~\ref{sec:zustandsraum}):

$$
\frac{\mathrm{d}\vec{z}(t)}{\mathrm{d}t} = \mathbf{A}\,\vec{z}(t) + \mathbf{B}\,u(t),
\qquad
y(t) = \mathbf{C}\,\vec{z}(t) + \mathbf{D}\,u(t).
$$

\textbf{Rechnung:} Dasselbe Rezept; neu ist nur, dass jetzt ein Eingangssignal
existiert.

\medskip
\textit{Schritt 1+2 – Ordnung und Zustandsvariablen.}
Unverändert: Die DGL ist weiterhin zweiter Ordnung in $x$, also $n = 2$ mit

$$
z_1(t) = x(t), \qquad z_2(t) = v(t).
$$

\medskip
\textit{Eingangssignal identifizieren.}
Die Aufgabe sagt ausdrücklich: „Die zusätzliche Kraft stellt das Eingangssignal
dar.`` Also setzen wir

$$
u(t) = F(t) = \hat{F}\cos(\omega t).
$$

\emph{Warum wichtig:} In Schritt 5 müssen wir die Terme nach Zuständen $z_j$
\emph{und} Eingang $u$ trennen -- daher muss $u$ vorher eindeutig benannt sein.

\medskip
\textit{Schritt 3 – Triviale Zustandsgleichung.}
Wie in Teil a):

$$
\dot{z}_1(t) = v(t) = z_2(t).
$$

\medskip
\textit{Schritt 4 – Echte Zustandsgleichung.}
DGL (4) nach $\dot{v}$ aufgelöst und der Anregungsterm über
$\tfrac{\hat{F}}{m}\cos(\omega t) = \tfrac{1}{m}F(t) = \tfrac{1}{m}u(t)$
geschrieben:

$$
\dot{z}_2(t) = \dot{v}(t)
= -\omega_0^2\, z_1(t) - 2\gamma\, z_2(t) + \frac{1}{m}\,u(t).
$$

\emph{Woher kommt der $\tfrac{1}{m}$-Term:} Aus dem zweiten Newtonschen Gesetz
$F_{\text{ges}} = m\,\dot{v}$ folgt für die \emph{äußere} Kraft der Beitrag
$\dot{v} = \ldots + \tfrac{1}{m}F(t)$ -- eine Kraft wirkt über $a = F/m$ auf die
Geschwindigkeitsänderung.

\medskip
\textit{Schritt 5 – Matrixform.}
Beide Zeilen, sortiert nach $z_1, z_2$ und $u$:

$$
\begin{aligned}
\dot{z}_1 &= 0\cdot z_1 + 1\cdot z_2 + 0\cdot u, \\
\dot{z}_2 &= -\omega_0^2\cdot z_1 - 2\gamma\cdot z_2 + \tfrac{1}{m}\cdot u.
\end{aligned}
$$

Koeffizienten vor $z_1, z_2$ $\to \mathbf{A}$; Koeffizienten vor $u$ $\to
\mathbf{B}$:

$$
\frac{\mathrm{d}}{\mathrm{d}t}
\begin{pmatrix} z_1 \\ z_2 \end{pmatrix}
=
\begin{pmatrix} 0 & 1 \\ -\omega_0^2 & -2\gamma \end{pmatrix}
\begin{pmatrix} z_1 \\ z_2 \end{pmatrix}
+
\begin{pmatrix} 0 \\ \tfrac{1}{m} \end{pmatrix}
u(t).
$$

\emph{Beachte:} Die erste Komponente von $\mathbf{B}$ ist $0$, weil die Kraft
\emph{nicht direkt} auf die Position $z_1$ wirkt, sondern erst über die
Beschleunigung auf $z_2$.

\medskip
\textit{Schritt 6 – Ausgangsgleichung.}
Gemessen wird weiterhin die Auslenkung, $y(t) = x(t) = z_1(t)$:

$$
y(t) = (1\;\; 0)\,\vec{z}(t), \qquad \mathbf{D} = 0.
$$

\textbf{Lösung:}

$$
\boxed{
\frac{\mathrm{d}\vec{z}(t)}{\mathrm{d}t}
= \underbrace{\begin{pmatrix} 0 & 1 \\ -\omega_0^2 & -2\gamma \end{pmatrix}}_{\mathbf{A}}
\,\vec{z}(t)
+ \underbrace{\begin{pmatrix} 0 \\ \tfrac{1}{m} \end{pmatrix}}_{\mathbf{B}}
F(t),
\qquad
y(t) = \underbrace{(1\;\; 0)}_{\mathbf{C}}\,\vec{z}(t)
}
$$

\textit{Kontrolle:} $\mathbf{A}\!:\,2\times2$, $\mathbf{B}\!:\,2\times1$,
$\mathbf{C}\!:\,1\times2$ -- passend zu $n=2$, $L=1$, $M=1$. \checkmark

\textbf{Vergleich mit Teil a):} Die Systemmatrix $\mathbf{A}$ ist
\emph{identisch} -- die innere Dynamik (Eigenfrequenz $\omega_0$, Dämpfung
$\gamma$) ändert sich durch die Anregung nicht. Hinzugekommen ist allein die
Eingangsmatrix $\mathbf{B}$, die beschreibt, \emph{wie} das Eingangssignal in das
System eingreift. Die erzwungene Schwingung in Bild~2 ist daher die Summe aus
abklingender Eigenschwingung (homogene Lösung, $\mathbf{A}$) und der stationären
Antwort auf die periodische Anregung (partikuläre Lösung, $\mathbf{B}\,u$).

\textbf{Bemerkung:} Wählt man statt der Kraft das normierte Signal
$u(t) = \cos(\omega t)$ als Eingang, so wandert die Amplitude in die
Eingangsmatrix: $\mathbf{B} = (0,\; \tfrac{\hat{F}}{m})^{\mathsf{T}}$. Beide
Darstellungen beschreiben dasselbe System; sie unterscheiden sich nur in der
Frage, ob der Faktor $\hat{F}$ zum Signal $u$ oder zur Matrix $\mathbf{B}$
gezählt wird.

\end{beispiel}

\section{Übung 7 – Zufallsereignisse und statistische Abhängigkeit}

Zwei bzw.\ sechs ideale Würfel werden geworfen. Alle Teilaufgaben folgen dem
Laplace-Modell (vgl.\ Abschnitt~\ref{sec:axiome-wahrscheinlichkeit}).

% ----------------------------------------------------------------
\subsection*{Teil A – Zwei Würfel}
% ----------------------------------------------------------------

Der Ergebnisraum ist das kartesische Produkt
$\Omega = \Omega_w \times \Omega_w$ mit $|\Omega| = 6^2 = 36$
(vgl.\ Abschnitt~\ref{sec:kartesische-produktmengen}).
Jedes Elementarereignis $(\xi_i,\xi_j)$ hat die Wahrscheinlichkeit
$P(\xi_i,\xi_j) = \tfrac{1}{36}$.

% ---- A1 --------------------------------------------------------
\subsubsection*{A1 – Genau ein Würfel zeigt eine 5}

\begin{beispiel}[Genau eine 5 bei zwei Würfeln]

\textbf{Gegeben:} $n = 2$ unabhängige ideale Würfel, $p = P(\text{5}) = \tfrac{1}{6}$.

\textbf{Gesucht:} $P(\text{genau eine 5})$

\textbf{Formel} (Binomialformel, Abschnitt~\ref{sec:kombinatorik}):

$$
P = \binom{n}{s}\,p^{s}\,(1-p)^{n-s}
\qquad n=2,\; s=1,\; p=\tfrac{1}{6}
$$

\textbf{Rechnung:}

$$
P = \binom{2}{1}\cdot\frac{1}{6}\cdot\frac{5}{6}
  = 2 \cdot \frac{5}{36}
  = \frac{10}{36}
$$

\textbf{Lösung:}

$$
\boxed{P(\text{genau eine 5}) = \frac{5}{18} \approx 0{,}278}
$$

\end{beispiel}

% ---- A2 --------------------------------------------------------
\subsubsection*{A2 – Mindestens ein Würfel zeigt eine 5}

\begin{beispiel}[Mindestens eine 5 bei zwei Würfeln]

\textbf{Gegeben:} wie A1.

\textbf{Gesucht:} $P(\text{mindestens eine 5})$

\textbf{Formel} (Komplementregel + statistische Unabhängigkeit,
Abschnitt~\ref{sec:unabhaengige-versuche}):

$$
P(\text{mindestens eine 5}) = 1 - P(\text{keine 5})
= 1 - \left(\frac{5}{6}\right)^{2}
$$

\textbf{Rechnung:}

$$
P(\text{keine 5}) = \frac{5}{6}\cdot\frac{5}{6} = \frac{25}{36}
\qquad\Rightarrow\qquad
P = 1 - \frac{25}{36} = \frac{11}{36}
$$

\textbf{Lösung:}

$$
\boxed{P(\text{mindestens eine 5}) = \frac{11}{36} \approx 0{,}306}
$$

\end{beispiel}

% ---- A3 --------------------------------------------------------
\subsubsection*{A3 – Beide Würfel zeigen die gleiche Zahl}

\begin{beispiel}[Beide Würfel zeigen die gleiche Zahl]

\textbf{Gegeben:} $|\Omega| = 36$.

\textbf{Gesucht:} $P(W_1 = W_2)$

\textbf{Formel} (Laplace-Experiment, Abschnitt~\ref{sec:axiome-wahrscheinlichkeit}):

$$
P(A) = \frac{|A|}{|\Omega|}
$$

\textbf{Rechnung:}

Günstige Ereignisse: $A = \{(1,1),(2,2),(3,3),(4,4),(5,5),(6,6)\}$, also $|A| = 6$.

$$
P = \frac{6}{36}
$$

\textbf{Lösung:}

$$
\boxed{P(W_1 = W_2) = \frac{1}{6} \approx 0{,}167}
$$

\end{beispiel}

% ---- A4 --------------------------------------------------------
\subsubsection*{A4 – Beide Würfel zeigen unterschiedliche Zahlen}

\begin{beispiel}[Beide Würfel zeigen unterschiedliche Zahlen]

\textbf{Gegeben:} Ergebnis aus A3.

\textbf{Gesucht:} $P(W_1 \neq W_2)$

\textbf{Formel} (Komplementregel):

$$
P(W_1 \neq W_2) = 1 - P(W_1 = W_2)
$$

\textbf{Rechnung:}

$$
P = 1 - \frac{1}{6} = \frac{5}{6}
$$

\textbf{Lösung:}

$$
\boxed{P(W_1 \neq W_2) = \frac{5}{6} \approx 0{,}833}
$$

\end{beispiel}

% ----------------------------------------------------------------
\subsection*{Teil B – Sechs Würfel}
% ----------------------------------------------------------------

Der Ergebnisraum erweitert sich auf
$\Omega = \Omega_w^{\times 6}$ mit $|\Omega| = 6^6 = 46\,656$.

% ---- B1 --------------------------------------------------------
\subsubsection*{B1 – Genau ein Würfel zeigt eine 5}

\begin{beispiel}[Genau eine 5 bei sechs Würfeln]

\textbf{Gegeben:} $n = 6$ unabhängige ideale Würfel, $p = \tfrac{1}{6}$.

\textbf{Gesucht:} $P(\text{genau eine 5})$

\textbf{Formel} (Binomialformel, Abschnitt~\ref{sec:kombinatorik}):

$$
P = \binom{6}{1}\cdot\left(\frac{1}{6}\right)^{1}\cdot\left(\frac{5}{6}\right)^{5}
$$

\textbf{Rechnung:}

$$
P = 6 \cdot \frac{1}{6} \cdot \frac{5^5}{6^5}
  = 1 \cdot \frac{3125}{7776}
$$

\textbf{Lösung:}

$$
\boxed{P(\text{genau eine 5}) = \frac{3125}{7776} \approx 0{,}402}
$$

\end{beispiel}

% ---- B2 --------------------------------------------------------
\subsubsection*{B2 – Mindestens ein Würfel zeigt eine 5}

\begin{beispiel}[Mindestens eine 5 bei sechs Würfeln]

\textbf{Gegeben:} $n = 6$, $p = \tfrac{1}{6}$.

\textbf{Gesucht:} $P(\text{mindestens eine 5})$

\textbf{Formel} (Komplementregel, Abschnitt~\ref{sec:unabhaengige-versuche}):

$$
P(\text{mindestens eine 5}) = 1 - \left(\frac{5}{6}\right)^{6}
$$

\textbf{Rechnung:}

$$
\left(\frac{5}{6}\right)^{6} = \frac{15\,625}{46\,656}
\qquad\Rightarrow\qquad
P = 1 - \frac{15\,625}{46\,656} = \frac{31\,031}{46\,656}
$$

\textbf{Lösung:}

$$
\boxed{P(\text{mindestens eine 5}) = \frac{31\,031}{46\,656} \approx 0{,}665}
$$

\end{beispiel}

% ---- B3 --------------------------------------------------------
\subsubsection*{B3 – Alle sechs Würfel zeigen die gleiche Zahl}

\begin{beispiel}[Alle sechs Würfel zeigen die gleiche Zahl]

\textbf{Gegeben:} $|\Omega| = 6^6 = 46\,656$.

\textbf{Gesucht:} $P(W_1 = W_2 = \cdots = W_6)$

\textbf{Formel} (Laplace-Experiment, Abschnitt~\ref{sec:axiome-wahrscheinlichkeit}):

$$
P(A) = \frac{|A|}{|\Omega|}
$$

\textbf{Rechnung:}

Günstige Ereignisse: $\{(1,\ldots,1),(2,\ldots,2),\ldots,(6,\ldots,6)\}$, also $|A| = 6$.

$$
P = \frac{6}{6^6} = \frac{6}{46\,656} = \frac{1}{7\,776}
$$

\textbf{Lösung:}

$$
\boxed{P(W_1 = \cdots = W_6) = \frac{1}{7\,776} \approx 1{,}29 \cdot 10^{-4}}
$$

\end{beispiel}

% ---- B4 --------------------------------------------------------
\subsubsection*{B4 – Alle sechs Würfel zeigen unterschiedliche Zahlen}

\begin{beispiel}[Alle sechs Würfel zeigen unterschiedliche Zahlen]

\textbf{Gegeben:} $|\Omega| = 6^6 = 46\,656$.

\textbf{Gesucht:} $P(\text{alle Würfel verschieden})$

\textbf{Formel} (Permutationen ohne Wiederholung, Abschnitt~\ref{sec:kombinatorik}):

$$
|A| = n! \qquad \text{(alle Anordnungen von } \{1,2,3,4,5,6\}\text{)}
$$

$$
P(A) = \frac{n!}{|\Omega|}
$$

\textbf{Rechnung:}

$$
|A| = 6! = 720
\qquad\Rightarrow\qquad
P = \frac{720}{46\,656} = \frac{5}{324}
$$

\textbf{Lösung:}

$$
\boxed{P(\text{alle verschieden}) = \frac{5}{324} \approx 0{,}0154}
$$

\end{beispiel}

% ----------------------------------------------------------------
\subsection*{Teil C – Bildpunkte eines Schwarz-Weiß-Bildes}
% ----------------------------------------------------------------

\subsubsection*{C1 – Elementarereignisse eines einzelnen Bildpunkts}

\begin{beispiel}[Ereignismenge eines Bildpunkts]

\textbf{Gegeben:} Ein digitalisiertes Schwarz-Weiß-Bild. Jeder Bildpunkt kann genau
einen von zwei Zuständen annehmen: schwarz~(S) oder weiß~(W). Beide Zustände treten
mit gleicher Wahrscheinlichkeit auf.

\textbf{Gesucht:} Menge der Elementarereignisse $\Omega_P$ für einen einzelnen Bildpunkt.

\textbf{Formel} (Definition Elementarereignis, Abschnitt~\ref{sec:mengen-ereignisse}):

$$
\Omega_P = \{\xi_1, \xi_2, \ldots\}
$$

\textbf{Rechnung:}

Da schwarz und weiß die einzigen möglichen Zustände sind und sich gegenseitig ausschließen:

$$
\Omega_P = \{S,\, W\}
$$

Da es sich um ein Laplace-Experiment handelt (Abschnitt~\ref{sec:axiome-wahrscheinlichkeit}):

$$
P(S) = P(W) = \frac{1}{|\Omega_P|} = \frac{1}{2}
$$

\textbf{Lösung:}

$$
\boxed{\Omega_P = \{S,\, W\}, \qquad P(S) = P(W) = \frac{1}{2}}
$$

\end{beispiel}

% ---- C2 --------------------------------------------------------
\subsubsection*{C2 – Verbundwahrscheinlichkeiten zweier unabhängiger Bildpunkte}

\begin{beispiel}[Gemeinsame Betrachtung zweier Bildpunkte]

\textbf{Gegeben:} Zwei benachbarte Bildpunkte mit
$P(S) = P(W) = \tfrac{1}{2}$; die Bildpunkte sind statistisch unabhängig.

\textbf{Gesucht:} Wahrscheinlichkeiten aller Elementarereignisse des kombinierten
Experiments.

\textbf{Formel} (Kartesisches Produkt + Verbundwahrscheinlichkeit bei Unabhängigkeit,
Abschnitte~\ref{sec:kartesische-produktmengen} und~\ref{sec:verbundwahrscheinlichkeiten}):

$$
\Omega = \Omega_P \times \Omega_P, \qquad
P(A_i, B_j) = P(A_i)\cdot P(B_j)
$$

\textbf{Rechnung:}

$$
\Omega = \{(S,S),\,(S,W),\,(W,S),\,(W,W)\}, \qquad |\Omega| = 4
$$

\begin{center}
\begin{tabular}{c|cc}
  & $B = S$ & $B = W$ \\
  \hline
  $A = S$ & $P(S,S) = \tfrac{1}{2}\cdot\tfrac{1}{2} = \tfrac{1}{4}$
           & $P(S,W) = \tfrac{1}{2}\cdot\tfrac{1}{2} = \tfrac{1}{4}$ \\[4pt]
  $A = W$ & $P(W,S) = \tfrac{1}{2}\cdot\tfrac{1}{2} = \tfrac{1}{4}$
           & $P(W,W) = \tfrac{1}{2}\cdot\tfrac{1}{2} = \tfrac{1}{4}$ \\
\end{tabular}
\end{center}

Probe: $\displaystyle\sum_{i,j} P(A_i,B_j) = 4\cdot\frac{1}{4} = 1$ \checkmark

\textbf{Lösung:}

$$
\boxed{P(S,S) = P(S,W) = P(W,S) = P(W,W) = \frac{1}{4}}
$$

Alle vier Kombinationen sind gleichwahrscheinlich, da die Bildpunkte statistisch
unabhängig sind und jeder Zustand mit $\tfrac{1}{2}$ auftritt.

\end{beispiel}

% ---- C3–C5 Neue Verbundwahrscheinlichkeiten -------------------

\medskip
Für die folgenden Teilaufgaben C3–C5 gelten die Verbundwahrscheinlichkeiten:

\begin{center}
\begin{tabular}{c|cc|c}
  & $B = S$ & $B = W$ & $P(A_i)$ \\
  \hline
  $A = S$ & $0{,}3$ & $0{,}2$ & \\
  $A = W$ & $0{,}2$ & $0{,}3$ & \\
  \hline
  $P(B_j)$ & & & $1$ \\
\end{tabular}
\end{center}

wobei $A$ den linken und $B$ den rechten Bildpunkt bezeichnet.

\subsubsection*{C3 – Wahrscheinlichkeit: linker Bildpunkt ist schwarz}

\begin{beispiel}[Randwahrscheinlichkeit des linken Bildpunkts]

\textbf{Gegeben:} $P(S,S) = 0{,}3$,\; $P(W,S) = 0{,}2$,\; $P(S,W) = 0{,}2$,\; $P(W,W) = 0{,}3$.

\textbf{Gesucht:} $P(A = S)$

\textbf{Formel} (Randwahrscheinlichkeit, Abschnitt~\ref{sec:verbundwahrscheinlichkeiten}):

$$
P(A_i) = \sum_{j} P(A_i, B_j)
$$

\textbf{Rechnung:}

$$
P(A = S) = P(S,S) + P(S,W) = 0{,}3 + 0{,}2 = 0{,}5
$$

\textbf{Lösung:}

$$
\boxed{P(A = S) = 0{,}5}
$$

\end{beispiel}

\subsubsection*{C4 – Wahrscheinlichkeit: rechter Bildpunkt ist weiß}

\begin{beispiel}[Randwahrscheinlichkeit des rechten Bildpunkts]

\textbf{Gegeben:} wie C3.

\textbf{Gesucht:} $P(B = W)$

\textbf{Formel} (Randwahrscheinlichkeit, Abschnitt~\ref{sec:verbundwahrscheinlichkeiten}):

$$
P(B_j) = \sum_{i} P(A_i, B_j)
$$

\textbf{Rechnung:}

$$
P(B = W) = P(S,W) + P(W,W) = 0{,}2 + 0{,}3 = 0{,}5
$$

\textbf{Lösung:}

$$
\boxed{P(B = W) = 0{,}5}
$$

\end{beispiel}

\subsubsection*{C5 – Statistische Unabhängigkeit der beiden Bildpunkte}

\begin{beispiel}[Unabhängigkeitstest]

\textbf{Gegeben:} Verbundwahrscheinlichkeiten wie C3; Randwahrscheinlichkeiten aus C3/C4.

\textbf{Gesucht:} Sind die beiden Bildpunkte statistisch unabhängig?

\textbf{Formel} (Bedingung für statistische Unabhängigkeit,
Abschnitt~\ref{sec:verbundwahrscheinlichkeiten}):

$$
P(A_i, B_j) \stackrel{!}{=} P(A_i)\cdot P(B_j) \quad \forall\, i, j
$$

\textbf{Rechnung:}

Zunächst alle Randwahrscheinlichkeiten:

$$
P(A=S) = 0{,}3 + 0{,}2 = 0{,}5, \qquad P(A=W) = 0{,}2 + 0{,}3 = 0{,}5
$$
$$
P(B=S) = 0{,}3 + 0{,}2 = 0{,}5, \qquad P(B=W) = 0{,}2 + 0{,}3 = 0{,}5
$$

Gegenbeispiel: Es genügt ein einziger Fall, bei dem die Bedingung verletzt ist:

$$
P(A=S)\cdot P(B=S) = 0{,}5 \cdot 0{,}5 = 0{,}25 \;\neq\; P(S,S) = 0{,}3
$$

\textbf{Lösung:}

$$
\boxed{\text{Die Bildpunkte sind statistisch \textbf{abhängig}.}}
$$

Obwohl beide Randwahrscheinlichkeiten identisch zu Teil~C2 sind ($P = 0{,}5$), weichen
die Verbundwahrscheinlichkeiten vom Produkt der Randwahrscheinlichkeiten ab. Benachbarte
Bildpunkte in einem realen Bild sind typischerweise statistisch abhängig, da ähnliche
Farben in räumlicher Nähe häufiger gemeinsam auftreten.

\end{beispiel}

\section{Übung 8 – Bedingte Wahrscheinlichkeiten und Satz von Bayes}

In dieser Übung werden bedingte Wahrscheinlichkeiten beim Ziehen ohne Zurücklegen
(Teil~A) sowie die Binomialverteilung beim Ziehen mit Zurücklegen (Teil~B)
betrachtet.

% ----------------------------------------------------------------
\subsection*{Teil A – Urne mit farbigen Kugeln (ohne Zurücklegen)}
% ----------------------------------------------------------------

Eine Urne enthält $2$ rote, $5$ schwarze und $3$ weiße Kugeln, insgesamt also
$N = 10$ Kugeln. Es wird zweimal \emph{ohne Zurücklegen} gezogen.

% ---- A1 --------------------------------------------------------
\subsubsection*{A1 – Bedingte Wahrscheinlichkeiten für die 2.\ Ziehung}

\begin{beispiel}[Bedingte Wahrscheinlichkeit nach einer schwarzen Kugel]

\textbf{Gegeben:} Urne mit $2$ roten, $5$ schwarzen, $3$ weißen Kugeln
($N = 10$). Die erste gezogene Kugel ist schwarz und wird \emph{nicht}
zurückgelegt.

\textbf{Gesucht:} Die bedingten Wahrscheinlichkeiten
$P(\text{rot} \mid S_1)$, $P(\text{schwarz} \mid S_1)$ und
$P(\text{weiß} \mid S_1)$ für die zweite Ziehung, wobei $S_1$ das Ereignis
„erste Kugel schwarz`` bezeichnet.

\textbf{Formel} (bedingte Wahrscheinlichkeit als Laplace-Experiment auf der
reduzierten Urne, Abschnitt~\ref{sec:bedingte-wahrscheinlichkeit}):

$$
P(\text{Farbe} \mid S_1) = \frac{\text{Anzahl Kugeln dieser Farbe}}{\text{Restzahl Kugeln}}
$$

\textbf{Rechnung:}

Nach dem Entnehmen einer schwarzen Kugel verbleiben $N' = 9$ Kugeln, davon
$2$ rote, $4$ schwarze und $3$ weiße. Somit gilt:

$$
P(\text{rot} \mid S_1) = \frac{2}{9}, \qquad
P(\text{schwarz} \mid S_1) = \frac{4}{9}, \qquad
P(\text{weiß} \mid S_1) = \frac{3}{9} = \frac{1}{3}
$$

Probe: $\dfrac{2}{9} + \dfrac{4}{9} + \dfrac{3}{9} = \dfrac{9}{9} = 1$ \checkmark

\textbf{Lösung:}

$$
\boxed{P(\text{rot} \mid S_1) = \frac{2}{9} \approx 0{,}222, \quad
P(\text{schwarz} \mid S_1) = \frac{4}{9} \approx 0{,}444, \quad
P(\text{weiß} \mid S_1) = \frac{1}{3} \approx 0{,}333}
$$

\end{beispiel}

% ---- A2 --------------------------------------------------------
\subsubsection*{A2 – Wahrscheinlichkeit, dass die 2.\ Kugel schwarz ist}

\begin{beispiel}[Satz von der totalen Wahrscheinlichkeit]

\textbf{Gegeben:} Urne wie oben ($2$ rot, $5$ schwarz, $3$ weiß, $N = 10$),
zweimal ohne Zurücklegen.

\textbf{Gesucht:} $P(S_2)$ -- die Wahrscheinlichkeit, dass die \emph{zweite}
gezogene Kugel schwarz ist (ohne Kenntnis der ersten Kugel).

\textbf{Formel} (Satz von der totalen Wahrscheinlichkeit,
Abschnitt~\ref{sec:totale-wahrscheinlichkeit}):

$$
P(S_2) = \sum_{i} P(S_2 \mid F_i)\, P(F_i)
$$

wobei $F_i \in \{\text{rot}, \text{schwarz}, \text{weiß}\}$ die Farbe der ersten
Kugel ist.

\textbf{Rechnung:}

Die a-priori-Wahrscheinlichkeiten der ersten Ziehung sind
$P(R_1) = \tfrac{2}{10}$, $P(S_1) = \tfrac{5}{10}$, $P(W_1) = \tfrac{3}{10}$.
Die bedingten Wahrscheinlichkeiten für eine schwarze zweite Kugel:

$$
P(S_2 \mid R_1) = \frac{5}{9}, \qquad
P(S_2 \mid S_1) = \frac{4}{9}, \qquad
P(S_2 \mid W_1) = \frac{5}{9}
$$

Einsetzen:

$$
P(S_2) = \frac{5}{9}\cdot\frac{2}{10} + \frac{4}{9}\cdot\frac{5}{10}
       + \frac{5}{9}\cdot\frac{3}{10}
       = \frac{10}{90} + \frac{20}{90} + \frac{15}{90}
       = \frac{45}{90}
$$

\textbf{Lösung:}

$$
\boxed{P(S_2) = \frac{1}{2} = 0{,}5}
$$

Die zweite Kugel ist mit derselben Wahrscheinlichkeit schwarz wie die erste
($P(S_1) = \tfrac{5}{10} = \tfrac{1}{2}$). Dieses Symmetrieergebnis gilt
allgemein: Ohne Information über die erste Ziehung ist die Wahrscheinlichkeit
für jede Position dieselbe.

\end{beispiel}

% ---- A3 --------------------------------------------------------
\subsubsection*{A3 – Rückschluss auf die erste Kugel (Satz von Bayes)}

\begin{beispiel}[Satz von Bayes]

\textbf{Gegeben:} Aus der Urne ($2$ rot, $5$ schwarz, $3$ weiß) wird eine Kugel
\emph{verdeckt} gezogen und beiseitegelegt. Anschließend wird eine schwarze
Kugel gezogen (Ereignis $S_2$).

\textbf{Gesucht:} $P(S_1 \mid S_2)$ -- die Wahrscheinlichkeit, dass die
zuerst gezogene (verdeckte) Kugel ebenfalls schwarz war.

\textbf{Formel} (Satz von Bayes, Abschnitt~\ref{sec:satz-von-bayes}):

$$
P(S_1 \mid S_2) = \frac{P(S_2 \mid S_1)\, P(S_1)}{P(S_2)}
$$

\textbf{Rechnung:}

Die benötigten Größen sind aus A1 und A2 bekannt:

$$
P(S_2 \mid S_1) = \frac{4}{9}, \qquad
P(S_1) = \frac{5}{10} = \frac{1}{2}, \qquad
P(S_2) = \frac{1}{2}
$$

Einsetzen:

$$
P(S_1 \mid S_2) = \frac{\dfrac{4}{9}\cdot\dfrac{1}{2}}{\dfrac{1}{2}}
               = \frac{4}{9}
$$

\textbf{Lösung:}

$$
\boxed{P(S_1 \mid S_2) = \frac{4}{9} \approx 0{,}444}
$$

Da $P(S_2) = P(S_1)$ ist, vereinfacht sich der Bayes-Ausdruck zu
$P(S_1 \mid S_2) = P(S_2 \mid S_1) = \tfrac{4}{9}$. Die Information, dass die
zweite Kugel schwarz ist, senkt die Wahrscheinlichkeit für eine schwarze erste
Kugel von $\tfrac{1}{2}$ auf $\tfrac{4}{9}$, da nun eine schwarze Kugel
„verbraucht`` ist.

\end{beispiel}

% ----------------------------------------------------------------
\subsection*{Teil B – Zweite Urne (mit Zurücklegen, Binomialverteilung)}
% ----------------------------------------------------------------

Eine zweite Urne enthält grüne und blaue Kugeln. Die Wahrscheinlichkeit, eine
grüne Kugel zu ziehen, beträgt $P_g$. Es wird $n = 16$ Mal eine Kugel gezogen
und jeweils wieder \emph{zurückgelegt}. Die einzelnen Ziehungen sind damit
statistisch unabhängig und identisch verteilt (Bernoulli-Kette).

% ---- B1 --------------------------------------------------------
\subsubsection*{B1 – Genau $k$-mal grün}

\begin{beispiel}[Binomialverteilung]

\textbf{Gegeben:} $n = 16$ unabhängige Ziehungen mit Zurücklegen,
Erfolgswahrscheinlichkeit $P_g$ für „grün`` pro Ziehung.

\textbf{Gesucht:} $P(X = k)$ -- die Wahrscheinlichkeit, dass genau $k$-mal eine
grüne Kugel gezogen wurde.

\textbf{Formel} (Binomialverteilung, Abschnitt~\ref{sec:binomialverteilung}):

$$
P(X = k) = \binom{n}{k}\, P_g^{\,k}\,(1 - P_g)^{\,n-k}
$$

\textbf{Rechnung:}

Mit $n = 16$ ergibt sich unmittelbar:

$$
P(X = k) = \binom{16}{k}\, P_g^{\,k}\,(1 - P_g)^{\,16-k},
\qquad k \in \{0, 1, \ldots, 16\}
$$

\textbf{Lösung:}

$$
\boxed{P(X = k) = \binom{16}{k}\, P_g^{\,k}\,(1 - P_g)^{\,16-k}}
$$

\end{beispiel}

% ---- B2 --------------------------------------------------------
\subsubsection*{B2 – Höchstens eine grüne Kugel bei $P_g = 0{,}4$}

\begin{beispiel}[Kumulierte Binomialwahrscheinlichkeit]

\textbf{Gegeben:} $n = 16$, $P_g = 0{,}4$, also $1 - P_g = 0{,}6$.

\textbf{Gesucht:} $P(X \le 1)$ -- die Wahrscheinlichkeit, dass höchstens eine
grüne Kugel gezogen wurde.

\textbf{Formel} (Summe der Einzelwahrscheinlichkeiten,
Abschnitt~\ref{sec:binomialverteilung}):

$$
P(X \le 1) = P(X = 0) + P(X = 1)
$$

\textbf{Rechnung:}

$$
P(X = 0) = \binom{16}{0}\,(0{,}4)^{0}\,(0{,}6)^{16}
         = (0{,}6)^{16} \approx 0{,}000282
$$

$$
P(X = 1) = \binom{16}{1}\,(0{,}4)^{1}\,(0{,}6)^{15}
         = 16 \cdot 0{,}4 \cdot (0{,}6)^{15} \approx 0{,}003009
$$

Summe:

$$
P(X \le 1) \approx 0{,}000282 + 0{,}003009 = 0{,}003291
$$

\textbf{Lösung:}

$$
\boxed{P(X \le 1) \approx 0{,}00329 \approx 0{,}33\,\%}
$$

Bei einer recht hohen Grün-Wahrscheinlichkeit von $40\,\%$ pro Ziehung ist es
über $16$ Versuche äußerst unwahrscheinlich, höchstens eine grüne Kugel zu
ziehen. Dies zeigt sich auch am \emph{Erwartungswert}
(Abschnitt~\ref{sec:erwartungswert}): Für eine binomialverteilte Zufallsvariable
$X = \sum_{i=1}^{n} X_i$ als Summe von $n$ unabhängigen Bernoulli-Variablen mit
$\mathrm{E}\{X_i\} = P_g$ folgt aus der Linearität des Erwartungswerts

$$
\mu_X = \mathrm{E}\{X\} = \sum_{k=0}^{n} k\, p_X(k) = n\,P_g
$$

und damit konkret

$$
\mu_X = n\,P_g = 16 \cdot 0{,}4 = 6{,}4 \text{ grüne Kugeln}.
$$

Der Erwartungswert liegt also weit über $1$, weshalb $P(X \le 1)$ verschwindend
klein ausfällt.

\end{beispiel}

% ================================================================
\subsection*{Aufgaben 8.1 – 8.8}
% ================================================================

Die folgenden Aufgaben vertiefen den Satz von Bayes, den Satz von der totalen
Wahrscheinlichkeit sowie die Binomial- und die hypergeometrische Verteilung.

% ----------------------------------------------------------------
\subsubsection*{8.1 – Fahren ohne gültigen Fahrschein}

\begin{beispiel}[Satz von Bayes -- Schwarzfahrer]

\textbf{Gegeben:} In Bochum wird zu $7\,\%$ schwarzgefahren, also
$P(S) = 0{,}07$ (Schwarzfahrer) und $P(E) = 0{,}93$ (ehrlicher Fahrgast).
Von den Schwarzfahrern haben $80\,\%$ gar keine Fahrkarte, die übrigen
$20\,\%$ besitzen eine gefälschte bzw.\ illegal besorgte Fahrkarte. Von den
ehrlichen Fahrgästen haben $5\,\%$ ihren Fahrschein vergessen. Betrachtet wird
das Ereignis $K$: „der kontrollierte Fahrgast kann keinen Fahrschein
vorzeigen``. Ein Schwarzfahrer mit gefälschtem Schein kann diesen vorzeigen,
sodass

$$
P(K \mid S) = 0{,}80, \qquad P(K \mid E) = 0{,}05 .
$$

\textbf{Gesucht:} a) $P(S \mid K)$ und b) $P(E \mid K)$.

\textbf{Formel} (Satz von Bayes mit totaler Wahrscheinlichkeit im Nenner,
Abschnitte~\ref{sec:satz-von-bayes} und~\ref{sec:totale-wahrscheinlichkeit}):

$$
P(S \mid K) = \frac{P(K \mid S)\, P(S)}{P(K \mid S)\, P(S) + P(K \mid E)\, P(E)}
$$

\textbf{Rechnung:}

Zunächst die totale Wahrscheinlichkeit für „kein Fahrschein vorzeigbar``:

$$
P(K) = 0{,}80 \cdot 0{,}07 + 0{,}05 \cdot 0{,}93
     = 0{,}056 + 0{,}0465 = 0{,}1025
$$

a) Rückschluss auf den Schwarzfahrer:

$$
P(S \mid K) = \frac{0{,}056}{0{,}1025} = \frac{112}{205} \approx 0{,}5463
$$

b) Rückschluss auf den ehrlichen Fahrgast:

$$
P(E \mid K) = \frac{0{,}0465}{0{,}1025} = \frac{93}{205} \approx 0{,}4537
$$

Probe: $\tfrac{112}{205} + \tfrac{93}{205} = \tfrac{205}{205} = 1$ \checkmark

\textbf{Lösung:}

$$
\boxed{P(S \mid K) = \frac{112}{205} \approx 0{,}5463, \qquad
P(E \mid K) = \frac{93}{205} \approx 0{,}4537}
$$

Wer keinen Fahrschein vorzeigen kann, ist mit knapp $55\,\%$ ein Schwarzfahrer,
obwohl der Anteil der Schwarzfahrer insgesamt nur $7\,\%$ beträgt.

\end{beispiel}

% ----------------------------------------------------------------
\subsubsection*{8.2 – Bestehen einer Klausur}

\begin{beispiel}[Totale Wahrscheinlichkeit und Bayes -- Klausur]

\textbf{Gegeben:} $V$ = „Student ist vorbereitet`` mit $P(V) = 0{,}25$,
$P(\bar V) = 0{,}75$. $B$ = „Prüfung bestanden`` mit den bedingten
Wahrscheinlichkeiten

$$
P(B \mid V) = 0{,}98, \qquad P(B \mid \bar V) = 0{,}01 .
$$

\textbf{Gesucht:} a) $P(B)$, b) $P(\bar B \cap V)$, c) $P(V \mid \bar B)$.

\textbf{Formel} (Satz von der totalen Wahrscheinlichkeit
\ref{sec:totale-wahrscheinlichkeit}, Multiplikationssatz und Bayes
\ref{sec:satz-von-bayes}):

$$
P(B) = P(B \mid V)\,P(V) + P(B \mid \bar V)\,P(\bar V), \qquad
P(V \mid \bar B) = \frac{P(\bar B \cap V)}{P(\bar B)}
$$

\textbf{Rechnung:}

a)
$$
P(B) = 0{,}98 \cdot 0{,}25 + 0{,}01 \cdot 0{,}75 = 0{,}245 + 0{,}0075 = 0{,}2525
$$

b) Ein vorbereiteter Student fällt mit $P(\bar B \mid V) = 0{,}02$ durch:

$$
P(\bar B \cap V) = P(\bar B \mid V)\, P(V) = 0{,}02 \cdot 0{,}25 = 0{,}005
$$

c) Mit $P(\bar B) = 1 - P(B) = 0{,}7475$:

$$
P(V \mid \bar B) = \frac{0{,}005}{0{,}7475} = \frac{2}{299} \approx 0{,}00669
$$

\textbf{Lösung:}

$$
\boxed{P(B) = 0{,}2525, \qquad P(\bar B \cap V) = 0{,}005, \qquad
P(V \mid \bar B) = \frac{2}{299} \approx 0{,}00669}
$$

Wer durchfällt, war nur mit rund $0{,}67\,\%$ Wahrscheinlichkeit vorbereitet --
Durchfallen spricht sehr stark für mangelnde Vorbereitung.

\end{beispiel}

% ----------------------------------------------------------------
\subsubsection*{8.3 – Bitfehlerwahrscheinlichkeiten}

\begin{beispiel}[Binomialverteilung -- Übertragungsfehler]

\textbf{Gegeben:} Ein Übertragungskanal überträgt jedes Bit unabhängig mit der
Bitfehlerwahrscheinlichkeit $P_e$. Ein Datenpaket besteht aus $n$ Bits. Die
Anzahl $X$ der fehlerhaften Bits ist damit binomialverteilt (Bernoulli-Kette).

\textbf{Gesucht:} Wahrscheinlichkeit für a) keinen Fehler, b) genau einen
Fehler, c) genau $k$ Fehler, d) höchstens $3$ Fehler.

\textbf{Formel} (Binomialverteilung, Abschnitt~\ref{sec:binomialverteilung}):

$$
P(X = k) = \binom{n}{k}\, P_e^{\,k}\,(1 - P_e)^{\,n-k}
$$

\textbf{Rechnung:}

a) Kein Fehler ($k = 0$):

$$
P(X = 0) = \binom{n}{0}\, P_e^{0}\,(1 - P_e)^{n} = (1 - P_e)^{n}
$$

b) Genau ein Fehler ($k = 1$):

$$
P(X = 1) = \binom{n}{1}\, P_e^{1}\,(1 - P_e)^{n-1} = n\,P_e\,(1 - P_e)^{n-1}
$$

c) Genau $k$ Fehler:

$$
P(X = k) = \binom{n}{k}\, P_e^{\,k}\,(1 - P_e)^{\,n-k}
$$

d) Höchstens $3$ Fehler ($X \le 3$):

$$
P(X \le 3) = \sum_{k=0}^{3} \binom{n}{k}\, P_e^{\,k}\,(1 - P_e)^{\,n-k}
$$

\textbf{Lösung:}

$$
\boxed{%
\begin{aligned}
P(X = 0) &= (1 - P_e)^{n}, \\
P(X = 1) &= n\,P_e\,(1 - P_e)^{n-1}, \\
P(X = k) &= \binom{n}{k}\, P_e^{\,k}\,(1 - P_e)^{\,n-k}, \\
P(X \le 3) &= \sum_{k=0}^{3} \binom{n}{k}\, P_e^{\,k}\,(1 - P_e)^{\,n-k} .
\end{aligned}}
$$

\end{beispiel}

% ----------------------------------------------------------------
\subsubsection*{8.4 – Restposten ICs}

\begin{beispiel}[Hypergeometrische Verteilung -- defekte ICs]

\textbf{Gegeben:} Ein Posten enthält $N = 10$ ICs, davon $D = 4$ defekt und
$6$ intakt. Es werden \emph{ohne Zurücklegen} $n = 4$ ICs zufällig ausgewählt.
Die Anzahl $X$ der defekten unter den gezogenen ICs ist hypergeometrisch
verteilt.

\textbf{Gesucht:} $P(X = 4)$, $P(X = 1)$, $P(X = 2)$.

\textbf{Formel} (Laplace-Wahrscheinlichkeit über Kombinationen,
Abschnitt~\ref{sec:kombinatorik}):

$$
P(X = m) = \frac{\dbinom{D}{m}\dbinom{N-D}{\,n-m\,}}{\dbinom{N}{n}}
$$

\textbf{Rechnung:}

Anzahl aller Auswahlen: $\dbinom{10}{4} = 210$.

a) Alle vier defekt ($m = 4$):

$$
P(X = 4) = \frac{\dbinom{4}{4}\dbinom{6}{0}}{\dbinom{10}{4}}
         = \frac{1 \cdot 1}{210} = \frac{1}{210} \approx 0{,}00476
$$

b) Genau eins defekt ($m = 1$):

$$
P(X = 1) = \frac{\dbinom{4}{1}\dbinom{6}{3}}{\dbinom{10}{4}}
         = \frac{4 \cdot 20}{210} = \frac{80}{210} = \frac{8}{21} \approx 0{,}3810
$$

c) Genau zwei defekt ($m = 2$):

$$
P(X = 2) = \frac{\dbinom{4}{2}\dbinom{6}{2}}{\dbinom{10}{4}}
         = \frac{6 \cdot 15}{210} = \frac{90}{210} = \frac{3}{7} \approx 0{,}4286
$$

\textbf{Lösung:}

$$
\boxed{P(X = 4) = \frac{1}{210} \approx 0{,}00476, \quad
P(X = 1) = \frac{8}{21} \approx 0{,}3810, \quad
P(X = 2) = \frac{3}{7} \approx 0{,}4286}
$$

\end{beispiel}

% ----------------------------------------------------------------
\subsubsection*{8.5 – Radarstation}

\begin{beispiel}[Satz von Bayes -- Radarstation]

\textbf{Gegeben:} $N$ = „störungsüberlagertes Nutzsignal`` mit $P(N) = 0{,}90$,
$T$ = „reine Störung`` mit $P(T) = 0{,}10$. Die Anzeige $A$ = „Nutzsignal
angezeigt`` tritt ein mit

$$
P(A \mid N) = 0{,}98, \qquad P(A \mid T) = 0{,}10 .
$$

\textbf{Gesucht:} $P(N \mid A)$ -- Wahrscheinlichkeit, dass ein als Nutzsignal
angezeigtes Signal wirklich ein Nutzsignal ist.

\textbf{Formel} (Satz von Bayes, Abschnitt~\ref{sec:satz-von-bayes}):

$$
P(N \mid A) = \frac{P(A \mid N)\,P(N)}{P(A \mid N)\,P(N) + P(A \mid T)\,P(T)}
$$

\textbf{Rechnung:}

$$
P(A) = 0{,}98 \cdot 0{,}90 + 0{,}10 \cdot 0{,}10 = 0{,}882 + 0{,}010 = 0{,}892
$$

$$
P(N \mid A) = \frac{0{,}882}{0{,}892} = \frac{441}{446} \approx 0{,}9888
$$

\textbf{Lösung:}

$$
\boxed{P(N \mid A) = \frac{441}{446} \approx 0{,}9888}
$$

Eine „Nutzsignal``-Anzeige ist mit rund $98{,}9\,\%$ zuverlässig, da echte
Nutzsignale sowohl häufig als auch fast sicher korrekt angezeigt werden.

\end{beispiel}

% ----------------------------------------------------------------
\subsubsection*{8.6 – Kontrolle am Flughafen}

\begin{beispiel}[Satz von Bayes -- Sicherheitskontrolle]

\textbf{Gegeben:} $V$ = „verbotene Gegenstände im Handgepäck`` mit
$P(V) = 0{,}20$, $P(\bar V) = 0{,}80$. Das Scan-Gerät löst Alarm $A$ aus mit

$$
P(A \mid V) = 0{,}98, \qquad P(A \mid \bar V) = 0{,}01 \quad (\text{Fehlalarm}).
$$

\textbf{Gesucht:} a) $P(A)$, b) $P(V \mid A)$, c) Anzahl der Fehlalarme bei
$1\,000\,000$ Fluggästen, d) $P(\bar V \mid A)$ (ohne Ergebnis aus b).

\textbf{Formel} (totale Wahrscheinlichkeit \ref{sec:totale-wahrscheinlichkeit}
und Bayes \ref{sec:satz-von-bayes}):

$$
P(A) = P(A \mid V)\,P(V) + P(A \mid \bar V)\,P(\bar V), \qquad
P(V \mid A) = \frac{P(A \mid V)\,P(V)}{P(A)}
$$

\textbf{Rechnung:}

a)
$$
P(A) = 0{,}98 \cdot 0{,}20 + 0{,}01 \cdot 0{,}80 = 0{,}196 + 0{,}008 = 0{,}204
$$

b)
$$
P(V \mid A) = \frac{0{,}196}{0{,}204} = \frac{49}{51} \approx 0{,}9608
$$

c) Ein Fehlalarm ist das Ereignis $A \cap \bar V$ mit

$$
P(A \cap \bar V) = P(A \mid \bar V)\,P(\bar V) = 0{,}01 \cdot 0{,}80 = 0{,}008 .
$$

Bei $1\,000\,000$ Fluggästen also

$$
0{,}008 \cdot 10^{6} = 8\,000 \text{ Fehlalarme}.
$$

d) Direkt über Bayes (unabhängig von b):

$$
P(\bar V \mid A) = \frac{P(A \mid \bar V)\,P(\bar V)}{P(A)}
                 = \frac{0{,}008}{0{,}204} = \frac{2}{51} \approx 0{,}0392
$$

Probe: $\tfrac{49}{51} + \tfrac{2}{51} = 1$ \checkmark

\textbf{Lösung:}

$$
\boxed{P(A) = 0{,}204, \quad P(V \mid A) = \frac{49}{51} \approx 0{,}9608, \quad
8\,000 \text{ Fehlalarme}, \quad P(\bar V \mid A) = \frac{2}{51} \approx 0{,}0392}
$$

\end{beispiel}

% ----------------------------------------------------------------
\subsubsection*{8.7 – Diabetes-Test}

\begin{beispiel}[Satz von Bayes -- Diabetes-Test]

\textbf{Gegeben:} $K$ = „an Diabetes erkrankt`` mit $P(K) = 0{,}072$,
$P(\bar K) = 0{,}928$. Der Test $T$ (positiv) erfüllt

$$
P(T \mid K) = 0{,}95, \qquad P(\bar K \cap T) = 0{,}10 .
$$

\textbf{Gesucht:} $P(K \mid T)$ -- Wahrscheinlichkeit, tatsächlich krank zu
sein, wenn der Test positiv ausfällt.

\textbf{Formel} (Satz von Bayes, Abschnitt~\ref{sec:satz-von-bayes}):

$$
P(K \mid T) = \frac{P(T \cap K)}{P(T)}, \qquad
P(T) = P(T \mid K)\,P(K) + P(\bar K \cap T)
$$

\textbf{Rechnung:}

Wahrscheinlichkeit für richtig-positiv:

$$
P(T \cap K) = P(T \mid K)\,P(K) = 0{,}95 \cdot 0{,}072 = 0{,}0684
$$

Totale Wahrscheinlichkeit eines positiven Tests (der falsch-positive Anteil ist
mit $0{,}10$ direkt gegeben):

$$
P(T) = 0{,}0684 + 0{,}10 = 0{,}1684
$$

$$
P(K \mid T) = \frac{0{,}0684}{0{,}1684} = \frac{171}{421} \approx 0{,}4062
$$

\textbf{Lösung:}

$$
\boxed{P(K \mid T) = \frac{171}{421} \approx 0{,}4062}
$$

Trotz eines positiven Tests ist man nur mit rund $40{,}6\,\%$ tatsächlich krank,
da die große gesunde Bevölkerungsgruppe viele falsch-positive Ergebnisse liefert
(Basisraten-Effekt).

\end{beispiel}

% ----------------------------------------------------------------
\subsubsection*{8.8 – Relative Häufigkeiten}

\begin{beispiel}[Relative Häufigkeiten im Schwarz-Weiß-Bild]

\textbf{Gegeben:} Ein Schwarz-Weiß-Bild (Bild~1) besteht ausschließlich aus
schwarzen und weißen Bildpunkten. Aus früheren Untersuchungen ist die
Wahrscheinlichkeit für einen schwarzen Bildpunkt mit $P(S) = 0{,}3$ bekannt.

\emph{Annahme (aus dem gerenderten Bild abgelesen):} Das Punktraster ist
quadratisch mit $20$ Spalten und $20$ Zeilen. Ein exaktes pixelgenaues
Auszählen ist aus der Vorlage nicht zuverlässig möglich; die im Folgenden mit
\glqq{}angenommen\grqq{} gekennzeichneten Zahlen sind daher plausible,
konsistent gewählte Werte, mit denen das methodische Vorgehen demonstriert
wird.

\textbf{Gesucht:} a) Anzahl aller Bildpunkte, b) relative Häufigkeiten
schwarzer und weißer Bildpunkte, c) Prüfung der Binomial-Näherung für die
Anzahl schwarzer Punkte pro Zeile, d) akkumulierte relative Häufigkeit über die
Bildzeilen.

\textbf{Formel} (relative Häufigkeit \ref{sec:wahrscheinlichkeit-einfuehrung};
Erwartungswert und Varianz der Binomialverteilung
\ref{sec:binomialverteilung},~\ref{sec:erwartungswert}):

$$
h(A) = \frac{n_A}{N}, \qquad \mu = n\,p, \qquad \sigma^2 = n\,p\,(1-p)
$$

\textbf{Rechnung:}

\emph{a) Anzahl aller Bildpunkte.} Bei $20$ Spalten und $20$ Zeilen:

$$
N = 20 \cdot 20 = 400 \text{ Bildpunkte (Annahme quadratisches Raster).}
$$

\emph{b) Relative Häufigkeiten.} Aus $P(S) = 0{,}3$ folgt für die erwartete
Anzahl schwarzer Punkte $n_S \approx P(S)\cdot N = 0{,}3 \cdot 400 = 120$ und
somit

$$
h(S) = \frac{120}{400} = 0{,}3, \qquad
h(W) = \frac{280}{400} = 0{,}7 = 1 - h(S).
$$

\emph{c) Binomial-Näherung pro Zeile.} Eine Bildzeile besteht aus $n = 20$
Punkten; deutet man jeden Punkt als unabhängiges Bernoulli-Experiment mit
$p = P(S) = 0{,}3$, so ist die Anzahl schwarzer Punkte pro Zeile
binomialverteilt mit den theoretischen Kennwerten

$$
\mu = n\,p = 20 \cdot 0{,}3 = 6, \qquad
\sigma^2 = n\,p\,(1-p) = 20 \cdot 0{,}3 \cdot 0{,}7 = 4{,}2 .
$$

Für die $20$ Zeilen werden die (angenommenen) schwarzen Punktzahlen

$$
x_i = 8,\,4,\,9,\,3,\,7,\,5,\,10,\,2,\,6,\,7,\,5,\,6,\,8,\,4,\,6,\,5,\,7,\,3,\,6,\,9
$$

herangezogen (Summe $\sum x_i = 120$). Damit ergibt sich der empirische
Mittelwert

$$
\bar{x} = \frac{1}{20}\sum_{i=1}^{20} x_i = \frac{120}{20} = 6
$$

und die empirische Varianz

$$
s^2 = \frac{1}{20}\sum_{i=1}^{20} (x_i - \bar{x})^2 = \frac{90}{20} = 4{,}5 .
$$

Der Vergleich $\bar{x} = 6 = \mu$ und $s^2 = 4{,}5 \approx 4{,}2 = \sigma^2$ zeigt
eine sehr gute Übereinstimmung. Die Anzahl schwarzer Bildpunkte pro Zeile lässt
sich damit gut durch eine Binomialverteilung $\mathcal{B}(n = 20,\,p = 0{,}3)$
annähern.

\emph{d) Akkumulierte relative Häufigkeit.} Werden die schwarzen Bildpunkte von
oben nach unten aufsummiert und durch die bis dahin betrachtete Punktzahl
($\text{Zeile} \cdot 20$) geteilt, so ergibt sich die kumulierte relative
Häufigkeit $h_m = \tfrac{1}{20\,m}\sum_{i=1}^{m} x_i$. Diese pendelt sich mit
wachsender Zeilenzahl auf den Wert $P(S) = 0{,}3$ ein (empirisches Gesetz der
großen Zahlen):

\begin{center}
\begin{tikzpicture}
\begin{axis}[
  width=0.9\linewidth, height=6cm,
  xlabel={Bildzeile $m$}, ylabel={akkum.\ rel.\ Häufigkeit $h_m$},
  xmin=0, xmax=21, ymin=0.25, ymax=0.42,
  xtick={2,4,6,8,10,12,14,16,18,20},
  ytick={0.28,0.30,0.32,0.34,0.36,0.38,0.40},
  grid=both, grid style={gray!25},
  legend pos=north east, legend cell align=left,
]
% angenommene kumulierte relative Häufigkeit
\addplot[thick, mark=*, mark size=1.2pt, color=accent] coordinates {
  (1,0.400) (2,0.300) (3,0.350) (4,0.300) (5,0.310)
  (6,0.300) (7,0.329) (8,0.300) (9,0.300) (10,0.305)
  (11,0.300) (12,0.300) (13,0.308) (14,0.300) (15,0.300)
  (16,0.297) (17,0.300) (18,0.292) (19,0.292) (20,0.300)
};
\addlegendentry{$h_m$ (angenommen)}
\addplot[dashed, thick, color=satzred-fr, domain=0:21] {0.3};
\addlegendentry{$P(S) = 0{,}3$}
\end{axis}
\end{tikzpicture}
\end{center}

\textbf{Lösung:}

$$
\boxed{N = 400, \quad h(S) = 0{,}3, \quad h(W) = 0{,}7, \quad
\bar{x} = 6 \approx \mu, \quad s^2 = 4{,}5 \approx \sigma^2 = 4{,}2}
$$

Die Anzahl schwarzer Bildpunkte pro Zeile ist gut binomialverteilt, und die
akkumulierte relative Häufigkeit konvergiert gegen $P(S) = 0{,}3$.

\end{beispiel}

\section{Übung 10 – Funktion und Transformation von Zufallsvariablen}

In dieser Übung wird die Vorrechenaufgabe~10 gelöst. Es geht darum, aus der
bekannten Dichte einer Zufallsvariablen $X$ die Dichte einer transformierten
Zufallsvariablen $Y = g(X)$ zu bestimmen. Teil~A behandelt Funktionen einer
exponentialverteilten Zufallsvariablen (eine quadratische und eine lineare
Transformation), Teil~B die Standardisierung einer normalverteilten
Zufallsvariablen. Grundlage ist der Transformationssatz für streng monotone
Funktionen aus Abschnitt~\ref{sec:monotone-transformation}.

% ----------------------------------------------------------------
\subsection*{Teil A – Funktion einer exponentialverteilten Zufallsvariablen}
% ----------------------------------------------------------------

Die Zufallsvariable $X$ ist exponentialverteilt
(Abschnitt~\ref{sec:exponentialverteilung}) mit Parameter $\lambda > 0$:

$$
p_X(x) = \begin{cases} \lambda\, e^{-\lambda x} & , \quad x \geq 0 \\[2pt]
0 & , \quad \text{sonst.} \end{cases}
$$

Der Träger von $X$ ist also $[0, \infty)$. Diese Einschränkung ist für die
Wahl der Umkehrfunktion entscheidend.

% ---- A1 --------------------------------------------------------
\subsubsection*{A1 – Quadratische Transformation $Y = X^2$}

\begin{beispiel}[Funktion einer Zufallsvariablen -- $Y = X^2$]

\textbf{Gegeben:} Exponentialverteilte Zufallsvariable $X$ mit Dichte
$p_X(x) = \lambda e^{-\lambda x}$ für $x \geq 0$ und Parameter $\lambda > 0$.
Transformation $Y = g(X) = X^2$.

\textbf{Gesucht:} Die Wahrscheinlichkeitsdichte $p_Y(y)$ der transformierten
Zufallsvariablen $Y$.

\textbf{Formel} (Transformationssatz für streng monotone Funktionen,
Abschnitt~\ref{sec:monotone-transformation}):

$$
p_Y(y) = \frac{p_X\!\left(g^{-1}(y)\right)}{\left|\, g'\!\left(g^{-1}(y)\right)\right|}
$$

\textbf{Rechnung:}

Zwar ist $g(x) = x^2$ auf ganz $\mathbb{R}$ \emph{nicht} streng monoton, auf dem
Träger $x \geq 0$ von $X$ jedoch streng monoton steigend. Es kommt daher nur die
\emph{positive} Wurzel als Umkehrfunktion in Frage:

$$
x = g^{-1}(y) = \sqrt{y}, \qquad y \geq 0.
$$

Die Ableitung von $g(x) = x^2$ ist $g'(x) = 2x$, an der Stelle $x = \sqrt{y}$
also

$$
\left|\, g'\!\left(\sqrt{y}\right)\right| = 2\sqrt{y}.
$$

Einsetzen in den Transformationssatz liefert für $y \geq 0$:

$$
p_Y(y) = \frac{p_X\!\left(\sqrt{y}\right)}{2\sqrt{y}}
       = \frac{\lambda\, e^{-\lambda \sqrt{y}}}{2\sqrt{y}}.
$$

Für $y < 0$ existiert kein reelles Urbild, also $p_Y(y) = 0$.

\textbf{Lösung:}

$$
\boxed{\;p_Y(y) = \begin{cases}
\dfrac{\lambda}{2\sqrt{y}}\, e^{-\lambda \sqrt{y}} & , \quad y > 0 \\[8pt]
0 & , \quad \text{sonst.}
\end{cases}\;}
$$

An der Stelle $y \to 0^+$ besitzt die Dichte wegen des Faktors
$1/(2\sqrt{y})$ eine (integrierbare) Polstelle -- die Wahrscheinlichkeitsmasse
wird durch das Quadrieren nahe null stark „gestaucht``. Zur Kontrolle gilt
weiterhin $\int_0^\infty p_Y(y)\,\mathrm{d}y = 1$ (Substitution
$y = x^2$ führt auf das Integral der Ausgangsdichte).

\end{beispiel}

% ---- A2 --------------------------------------------------------
\subsubsection*{A2 – Lineare Transformation $Y = aX + b$}

\begin{beispiel}[Funktion einer Zufallsvariablen -- $Y = aX + b$]

\textbf{Gegeben:} Exponentialverteilte Zufallsvariable $X$ mit Dichte
$p_X(x) = \lambda e^{-\lambda x}$ für $x \geq 0$. Lineare Transformation
$Y = g(X) = aX + b$ mit Konstanten $a \neq 0$ und $b \in \mathbb{R}$.

\textbf{Gesucht:} Die Wahrscheinlichkeitsdichte $p_Y(y)$ der transformierten
Zufallsvariablen $Y$.

\textbf{Formel} (Transformationssatz für streng monotone Funktionen,
Abschnitt~\ref{sec:monotone-transformation}):

$$
p_Y(y) = \frac{p_X\!\left(g^{-1}(y)\right)}{\left|\, g'\!\left(g^{-1}(y)\right)\right|}
$$

\textbf{Rechnung:}

Die lineare Funktion $g(x) = ax + b$ ist wegen $a \neq 0$ streng monoton und
eindeutig umkehrbar:

$$
x = g^{-1}(y) = \frac{y - b}{a}, \qquad g'(x) = a, \qquad |g'(x)| = |a|.
$$

Einsetzen ergibt zunächst formal

$$
p_Y(y) = \frac{1}{|a|}\, p_X\!\left(\frac{y - b}{a}\right).
$$

Die Ausgangsdichte $p_X$ ist nur für ein nichtnegatives Argument von null
verschieden, also für $\dfrac{y - b}{a} \geq 0$. Für den typischen Fall
$a > 0$ bedeutet das $y \geq b$, und man erhält

$$
p_Y(y) = \frac{\lambda}{a}\,
         \exp\!\left(-\lambda\,\frac{y - b}{a}\right), \qquad y \geq b.
$$

\textbf{Lösung} (Fall $a > 0$):

$$
\boxed{\;p_Y(y) = \begin{cases}
\dfrac{\lambda}{a}\, \exp\!\left(-\dfrac{\lambda}{a}\,(y - b)\right)
  & , \quad y \geq b \\[8pt]
0 & , \quad y < b.
\end{cases}\;}
$$

$Y$ ist damit wieder \emph{exponentialverteilt}, allerdings um $b$ nach rechts
verschoben und mit dem neuen Parameter $\lambda' = \lambda/a$. Der Faktor $a$
skaliert die Zufallsvariable und verkleinert entsprechend den Ratenparameter
(im Fall $a < 0$ kehrt sich der Träger zu $y \leq b$ um; die kompakte
Bedingung lautet allgemein $\tfrac{y-b}{a} \geq 0$, und der Vorfaktor ist
$\lambda/|a|$).

\end{beispiel}

% ----------------------------------------------------------------
\subsection*{Teil B – Standardisierung einer normalverteilten Zufallsvariablen}
% ----------------------------------------------------------------

% ---- B1 --------------------------------------------------------
\subsubsection*{B1 – Transformation auf die Standardnormalverteilung}

\begin{beispiel}[Standardisierung $X \sim \mathcal{N}(\mu_X, \sigma_X^2) \to Y \sim \mathcal{N}(0,1)$]

\textbf{Gegeben:} Normalverteilte Zufallsvariable $X$
(Abschnitt~\ref{sec:normalverteilung}) mit Mittelwert $\mu_X$ und Varianz
$\sigma_X^2 > 0$:

$$
p_X(x) = \frac{1}{\sqrt{2\pi}\,\sigma_X}\,
         \exp\!\left(-\frac{(x - \mu_X)^2}{2\sigma_X^2}\right).
$$

\textbf{Gesucht:} Eine lineare Transformation $Y = g(X)$, so dass $Y$
standardnormalverteilt ist ($\mu_Y = 0$, $\sigma_Y^2 = 1$), sowie der Nachweis
über die transformierte Dichte $p_Y(y)$.

\textbf{Formel} (Standardisierung sowie Transformationssatz,
Abschnitt~\ref{sec:monotone-transformation}):

$$
Y = \frac{X - \mu_X}{\sigma_X}, \qquad
p_Y(y) = \frac{p_X\!\left(g^{-1}(y)\right)}{\left|\, g'\!\left(g^{-1}(y)\right)\right|}
$$

\textbf{Rechnung:}

\emph{Ansatz.} Die Standardisierung subtrahiert den Mittelwert und normiert auf
die Standardabweichung:

$$
Y = g(X) = \frac{X - \mu_X}{\sigma_X}.
$$

Dass diese Wahl die geforderten Momente erzeugt, folgt direkt aus der
Linearität von Erwartungswert (Abschnitt~\ref{sec:erwartungswert}) und Varianz:

$$
\mu_Y = \mathrm{E}\{Y\} = \frac{\mathrm{E}\{X\} - \mu_X}{\sigma_X} = 0,
\qquad
\sigma_Y^2 = \mathrm{var}\{Y\} = \frac{\mathrm{var}\{X\}}{\sigma_X^2}
           = \frac{\sigma_X^2}{\sigma_X^2} = 1.
$$

\emph{Dichte.} Die Funktion $g(x) = (x - \mu_X)/\sigma_X$ ist streng monoton
steigend mit Umkehrfunktion und Ableitung

$$
x = g^{-1}(y) = \sigma_X\, y + \mu_X, \qquad
g'(x) = \frac{1}{\sigma_X}, \qquad |g'(x)| = \frac{1}{\sigma_X}.
$$

Einsetzen in den Transformationssatz:

$$
p_Y(y) = \frac{p_X\!\left(\sigma_X\, y + \mu_X\right)}{1/\sigma_X}
       = \sigma_X\, p_X\!\left(\sigma_X\, y + \mu_X\right).
$$

Nun wird $x = \sigma_X y + \mu_X$ in die Normaldichte eingesetzt; im Exponenten
gilt $(x - \mu_X)^2 = (\sigma_X y)^2 = \sigma_X^2\, y^2$:

$$
p_Y(y) = \sigma_X \cdot \frac{1}{\sqrt{2\pi}\,\sigma_X}\,
         \exp\!\left(-\frac{\sigma_X^2\, y^2}{2\sigma_X^2}\right)
       = \frac{1}{\sqrt{2\pi}}\, \exp\!\left(-\frac{y^2}{2}\right).
$$

\textbf{Lösung:}

$$
\boxed{\;Y = \frac{X - \mu_X}{\sigma_X}
\quad\Longrightarrow\quad
p_Y(y) = \frac{1}{\sqrt{2\pi}}\, e^{-y^2/2}
= \mathcal{N}(y;\, 0,\, 1)\;}
$$

$Y$ ist also standardnormalverteilt. Die Standardabweichung $\sigma_X$ aus dem
Betrag der Ableitung kürzt sich exakt gegen den Vorfaktor der Ausgangsdichte,
und die Verschiebung um $\mu_X$ verschwindet im Exponenten. Diese
Standardisierung ist die Grundlage für die Nutzung tabellierter Werte der
Standardnormalverteilung $\Phi(y)$ für beliebige normalverteilte Größen.

\end{beispiel}

% ----------------------------------------------------------------
\subsection*{Aufgabe 10.1 – Transformation auf eine Standard-Gleichverteilung}
% ----------------------------------------------------------------

\begin{beispiel}[Gleichverteilung auf einem Intervall zur Standard-Gleichverteilung]

\textbf{Gegeben:} Eine auf dem Intervall $[a, b]$ (mit $a < b$) gleichverteilte
Zufallsvariable $X$ (Abschnitt~\ref{sec:gleichverteilung}):

$$
p_X(x) = \begin{cases} \dfrac{1}{b - a} & , \quad a \leq x \leq b \\[6pt]
0 & , \quad \text{sonst.} \end{cases}
$$

\textbf{Gesucht:} Eine Transformation $Y = g(X)$, so dass $Y$ auf dem Intervall
$\left[-\tfrac{1}{2}, \tfrac{1}{2}\right]$ gleichverteilt ist, sowie der
Nachweis über die transformierte Dichte $p_Y(y)$.

\textbf{Formel} (lineare Transformation $Y = cX + d$ und Transformationssatz,
Abschnitt~\ref{sec:monotone-transformation}):

$$
p_Y(y) = \frac{p_X\!\left(g^{-1}(y)\right)}{\left|\, g'\!\left(g^{-1}(y)\right)\right|}
$$

\textbf{Rechnung:}

\emph{Ansatz.} Da beide Intervalle die Endpunkte einer Gleichverteilung
festlegen, genügt eine \emph{lineare} und streng monoton steigende
Transformation $Y = cX + d$, welche die Intervallgrenzen aufeinander abbildet:

$$
g(a) = -\tfrac{1}{2}, \qquad g(b) = +\tfrac{1}{2}.
$$

Die Steigung ergibt sich aus dem Verhältnis der Intervalllängen:

$$
c = \frac{\tfrac{1}{2} - \left(-\tfrac{1}{2}\right)}{b - a} = \frac{1}{b - a}.
$$

Der Versatz $d$ folgt aus $g(a) = -\tfrac{1}{2}$:

$$
d = -\tfrac{1}{2} - \frac{a}{b - a} = -\frac{a + b}{2\,(b - a)}.
$$

Zusammengefasst lässt sich die Transformation kompakt als Zentrierung um den
Mittelwert $\mu_X = \tfrac{a+b}{2}$ und Normierung auf die Intervalllänge
schreiben:

$$
Y = g(X) = \frac{X - \dfrac{a + b}{2}}{b - a}.
$$

\emph{Kontrolle der Grenzen:}

$$
g(a) = \frac{a - \tfrac{a+b}{2}}{b - a} = \frac{-\tfrac{b-a}{2}}{b-a} = -\tfrac{1}{2},
\qquad
g(b) = \frac{b - \tfrac{a+b}{2}}{b - a} = \frac{\tfrac{b-a}{2}}{b-a} = +\tfrac{1}{2}.
\quad\checkmark
$$

\emph{Dichte.} Die Funktion $g$ ist streng monoton steigend mit

$$
x = g^{-1}(y) = (b - a)\, y + \frac{a + b}{2}, \qquad
g'(x) = \frac{1}{b - a}, \qquad |g'(x)| = \frac{1}{b - a}.
$$

Einsetzen in den Transformationssatz für $y \in \left[-\tfrac12, \tfrac12\right]$
(dort ist $g^{-1}(y) \in [a, b]$, also $p_X = \tfrac{1}{b-a}$):

$$
p_Y(y) = \frac{p_X\!\left(g^{-1}(y)\right)}{|g'|}
       = \frac{\dfrac{1}{b - a}}{\dfrac{1}{b - a}}
       = 1.
$$

\textbf{Lösung:}

$$
\boxed{\;Y = \frac{X - \dfrac{a + b}{2}}{b - a}
\quad\Longrightarrow\quad
p_Y(y) = \begin{cases} 1 & , \quad -\tfrac{1}{2} \leq y \leq \tfrac{1}{2} \\[4pt]
0 & , \quad \text{sonst.} \end{cases}\;}
$$

$Y$ ist also auf $\left[-\tfrac{1}{2}, \tfrac{1}{2}\right]$ gleichverteilt; die
Dichte hat dort den Wert $1$, da das Zielintervall die Länge $1$ besitzt. Die
Transformation entspricht einer Standardisierung der Gleichverteilung: Sie
verschiebt den Mittelwert nach $0$ und skaliert die Spannweite auf $1$
(analog zur Standardisierung der Normalverteilung in Teil~B, dort jedoch mit
$\sigma_X$ statt der Intervalllänge $b - a$).

\end{beispiel}


\chapter{Tutorien}

\section{Tutorium 7 – Linearität, Zeitinvarianz, BIBO-Stabilität}

% ================================================================
\subsection*{Teil A – Linearität und Zeitinvarianz}
% ================================================================

Für jedes System wird geprüft:
\begin{itemize}
  \item \textbf{Linearität} (Abschnitt~\ref{sec:zeitinvarianz-linearitaet}):
  $\mathcal{S}[a\,x_1(t) + b\,x_2(t)] \stackrel{!}{=} a\,y_1(t) + b\,y_2(t)$
  \item \textbf{Zeitinvarianz} (Abschnitt~\ref{sec:zeitinvarianz-linearitaet}):
  $\mathcal{S}[x(t - \tau_0)] \stackrel{!}{=} y(t - \tau_0)$
\end{itemize}

% ----------------------------------------------------------------
\subsubsection*{A1 – $y(t) = x(t) + a$, \quad $a \in \mathbb{R}$, $a \neq 0$}
% ----------------------------------------------------------------

\begin{beispiel}[Linearität und Zeitinvarianz: $y(t) = x(t) + a$]

\textbf{Gegeben:} $\mathcal{S}[x(t)] = x(t) + a$,\quad $a \in \mathbb{R}$, $a \neq 0$.

\textbf{Gesucht:} Ist das System linear? Ist es zeitinvariant?

\textbf{Formel} (Abschnitt~\ref{sec:zeitinvarianz-linearitaet}): Superpositionsprinzip
und Verschiebungsinvarianz.

\textbf{Rechnung – Linearität:}

Sei $y_1(t) = x_1(t) + a$ und $y_2(t) = x_2(t) + a$. Dann:

$$
\mathcal{S}[a_1 x_1(t) + a_2 x_2(t)]
= a_1 x_1(t) + a_2 x_2(t) + a
$$

$$
a_1 y_1(t) + a_2 y_2(t)
= a_1(x_1(t)+a) + a_2(x_2(t)+a)
= a_1 x_1(t) + a_2 x_2(t) + (a_1+a_2)\,a
$$

Da $(a_1+a_2)\,a \neq a$ im Allgemeinen: $\mathcal{S}[a_1 x_1 + a_2 x_2] \neq a_1 y_1 + a_2 y_2$.

\textbf{Rechnung – Zeitinvarianz:}

$$
\mathcal{S}[x(t-\tau_0)] = x(t-\tau_0) + a
\qquad
y(t-\tau_0) = x(t-\tau_0) + a
$$

Beide Seiten sind identisch.

\textbf{Lösung:}

$$
\boxed{\text{nicht linear},\quad \text{zeitinvariant (NL-Z)}}
$$

Der konstante Offset $a$ verletzt das Superpositionsprinzip (Homogenitätsbedingung: $\mathcal{S}[0] = a \neq 0$).

\end{beispiel}

% ----------------------------------------------------------------
\subsubsection*{A2 – $y(t) = t\,x(t)$}
% ----------------------------------------------------------------

\begin{beispiel}[Linearität und Zeitinvarianz: $y(t) = t\,x(t)$]

\textbf{Gegeben:} $\mathcal{S}[x(t)] = t\,x(t)$.

\textbf{Gesucht:} Ist das System linear? Ist es zeitinvariant?

\textbf{Formel} (Abschnitt~\ref{sec:zeitinvarianz-linearitaet}).

\textbf{Rechnung – Linearität:}

$$
\mathcal{S}[a_1 x_1(t) + a_2 x_2(t)]
= t\bigl(a_1 x_1(t) + a_2 x_2(t)\bigr)
= a_1\,t\,x_1(t) + a_2\,t\,x_2(t)
= a_1 y_1(t) + a_2 y_2(t)
$$

Das Superpositionsprinzip ist erfüllt.

\textbf{Rechnung – Zeitinvarianz:}

$$
\mathcal{S}[x(t-\tau_0)] = t\,x(t-\tau_0)
$$

$$
y(t-\tau_0) = (t-\tau_0)\,x(t-\tau_0)
$$

Da $t \neq t - \tau_0$ für $\tau_0 \neq 0$: $\mathcal{S}[x(t-\tau_0)] \neq y(t-\tau_0)$.

\textbf{Lösung:}

$$
\boxed{\text{linear},\quad \text{nicht zeitinvariant (L-NZ)}}
$$

Der zeitabhängige Koeffizient $t$ bewirkt, dass eine Verschiebung des Eingangs
nicht zur gleichen Verschiebung des Ausgangs führt. Das System ist \textbf{zeitvariant}.

\end{beispiel}

% ----------------------------------------------------------------
\subsubsection*{A3 – $y(t) = \dfrac{\mathrm{d}}{\mathrm{d}t}\,x(t)$}
% ----------------------------------------------------------------

\begin{beispiel}[Linearität und Zeitinvarianz: $y(t) = \dot{x}(t)$]

\textbf{Gegeben:} $\mathcal{S}[x(t)] = \dfrac{\mathrm{d}}{\mathrm{d}t}\,x(t)$.

\textbf{Gesucht:} Ist das System linear? Ist es zeitinvariant?

\textbf{Formel} (Abschnitt~\ref{sec:zeitinvarianz-linearitaet}).

\textbf{Rechnung – Linearität:}

Die Ableitung ist ein linearer Operator (Linearität der Differentiation):

$$
\mathcal{S}[a_1 x_1(t) + a_2 x_2(t)]
= \frac{\mathrm{d}}{\mathrm{d}t}\bigl(a_1 x_1(t) + a_2 x_2(t)\bigr)
= a_1\,\dot{x}_1(t) + a_2\,\dot{x}_2(t)
= a_1 y_1(t) + a_2 y_2(t)
$$

\textbf{Rechnung – Zeitinvarianz:}

$$
\mathcal{S}[x(t-\tau_0)]
= \frac{\mathrm{d}}{\mathrm{d}t}\,x(t-\tau_0)
= \dot{x}(t-\tau_0)
= y(t-\tau_0)
$$

(Kettenregel: $\frac{\mathrm{d}}{\mathrm{d}t}x(t-\tau_0) = \dot{x}(t-\tau_0)\cdot 1$)

\textbf{Lösung:}

$$
\boxed{\text{linear},\quad \text{zeitinvariant} \;\Rightarrow\; \textbf{LTI-System}}
$$

Der Differentiationsoperator ist linear und verschiebt das Signal ohne
zusätzliche Zeitabhängigkeit — das System ist ein LTI-System.

\end{beispiel}

% ================================================================
\subsection*{Teil B – BIBO-Stabilität}
% ================================================================

Kriterium (Abschnitt~\ref{sec:faltung-eigenschaften}): Ein zeitkontinuierliches LTI-System
mit Impulsantwort $h(t)$ ist BIBO-stabil genau dann, wenn

$$
\int_{-\infty}^{\infty} |h(\xi)|\,\mathrm{d}\xi < \infty\,.
$$

% ----------------------------------------------------------------
\subsubsection*{B1 – $h(t) = 2\,\varepsilon(t-2T)\,\sin\!\left(\dfrac{\pi(t-2T)}{2T}\right)$}
% ----------------------------------------------------------------

\begin{beispiel}[BIBO-Stabilität: ungedämpfte Sinusantwort]

\textbf{Gegeben:}

$$
h(t) = 2\,\varepsilon(t-2T)\,\sin\!\left(\frac{\pi(t-2T)}{2T}\right)
$$

\textbf{Gesucht:} Ist das System BIBO-stabil?

\textbf{Formel} (Abschnitt~\ref{sec:faltung-eigenschaften}):
$\displaystyle\int_{-\infty}^{\infty} |h(\xi)|\,\mathrm{d}\xi \stackrel{?}{<} \infty$

\textbf{Rechnung:}

Substitution $\tau = t - 2T$ (der Sprung verschiebt den Träger auf $\tau \geq 0$):

$$
\int_{-\infty}^{\infty} |h(t)|\,\mathrm{d}t
= \int_{0}^{\infty} 2\left|\sin\!\left(\frac{\pi\tau}{2T}\right)\right|\mathrm{d}\tau
$$

Der Integrand $2\left|\sin\!\left(\tfrac{\pi\tau}{2T}\right)\right|$ ist periodisch mit Periode
$2T$ und hat konstante Amplitude $2$ — er klingt \textbf{nicht} ab. Das Integral über $[0,\infty)$
wächst ohne Schranke:

$$
\int_{0}^{2nT} 2\left|\sin\!\left(\frac{\pi\tau}{2T}\right)\right|\mathrm{d}\tau
= n \cdot \frac{8T}{\pi} \;\xrightarrow{n\to\infty}\; \infty
$$

\textbf{Lösung:}

$$
\boxed{\int_{-\infty}^{\infty}|h(t)|\,\mathrm{d}t = \infty \;\Rightarrow\; \textbf{nicht BIBO-stabil}}
$$

Eine ungedämpfte Schwingung in der Impulsantwort führt stets zur Instabilität, da
der Energieinhalt unbegrenzt wächst.

\end{beispiel}

% ----------------------------------------------------------------
\subsubsection*{B2 – $h(t) = 4A\!\left[\operatorname{rect}\!\left(\dfrac{t}{T}\right) + e^{-(t-T)/T}\,\varepsilon(t-T)\right]$}
% ----------------------------------------------------------------

\begin{beispiel}[BIBO-Stabilität: Rechteck + abklingende Exponentialfunktion]

\textbf{Gegeben:}

$$
h(t) = 4A\!\left[\operatorname{rect}\!\left(\frac{t}{T}\right)
       + e^{-(t-T)/T}\,\varepsilon(t-T)\right],
\quad A > 0,\; T > 0
$$

\textbf{Gesucht:} Ist das System BIBO-stabil?

\textbf{Formel} (Abschnitt~\ref{sec:faltung-eigenschaften}):
$\displaystyle\int_{-\infty}^{\infty}|h(\xi)|\,\mathrm{d}\xi \stackrel{?}{<} \infty$

\textbf{Rechnung:}

Die beiden Terme haben disjunkte Träger und können separat integriert werden.

\textbf{Term 1 –} $\operatorname{rect}(t/T)$: Träger $|t| \leq T/2$, Amplitude $4A$:

$$
\int_{-\infty}^{\infty} 4A\left|\operatorname{rect}\!\left(\frac{t}{T}\right)\right|\mathrm{d}t
= 4A \int_{-T/2}^{T/2}\mathrm{d}t = 4AT
$$

\textbf{Term 2 –} $e^{-(t-T)/T}\,\varepsilon(t-T)$: Träger $t \geq T$, Substitution $u = (t-T)/T$:

$$
\int_{-\infty}^{\infty} 4A\,e^{-(t-T)/T}\,\varepsilon(t-T)\,\mathrm{d}t
= 4A\int_{T}^{\infty} e^{-(t-T)/T}\,\mathrm{d}t
= 4AT\int_{0}^{\infty} e^{-u}\,\mathrm{d}u
= 4AT
$$

\textbf{Gesamt:}

$$
\int_{-\infty}^{\infty}|h(t)|\,\mathrm{d}t = 4AT + 4AT = 8AT < \infty
$$

\textbf{Lösung:}

$$
\boxed{\int_{-\infty}^{\infty}|h(t)|\,\mathrm{d}t = 8AT < \infty \;\Rightarrow\; \textbf{BIBO-stabil}}
$$

Das Rechteck ist zeitlich begrenzt (endliche Energie), die Exponentialfunktion klingt
für $t \geq T$ monoton ab — beide Terme liefern endliche Beiträge.

\end{beispiel}


% ...

% ── Zusätzliches Wissen ──────────────────────────────────────
% \include{knowledge/Thema1}
% ...

% ── Formelsammlung ───────────────────────────────────────────
% ============================================================
%  Formelsammlung  –  Systemtheorie 1
%  Querformat, 3-spaltig, thematisch gruppiert
% ============================================================
\clearpage
\newgeometry{landscape, top=1.1cm, bottom=1.1cm, left=1.1cm, right=1.1cm}
\pagestyle{empty}
\addcontentsline{toc}{chapter}{Formelsammlung}

% ── kompaktes Layout ─────────────────────────────────────────
\setlength{\parindent}{0pt}
\setlength{\abovedisplayskip}{2pt}
\setlength{\belowdisplayskip}{2pt}
\setlength{\abovedisplayshortskip}{1pt}
\setlength{\belowdisplayshortskip}{1pt}
\setlength{\columnsep}{1.4em}
\setlength{\columnseprule}{0.4pt}

% Themen-Kopf: Titel + kleiner Skript-Verweis + Trennlinie
\newcommand{\ft}[2]{%
  \par\smallskip\penalty-200
  {\bfseries\footnotesize #1}\nobreak\hfill{\tiny\hyperref[#2]{\S\,\ref*{#2}}}\par\nobreak
  \vspace{-0.35em}{\color{accent}\hrule height 0.4pt}\nobreak\vspace{0.2em}\nobreak}

\begin{center}
  {\color{accent}\rule{\linewidth}{1pt}}\\[0.2em]
  {\Large\bfseries Formelsammlung}\quad\small Systemtheorie\,1 -- Signale und Systeme\,I\\[0.1em]
  {\color{accent}\rule{\linewidth}{1pt}}
\end{center}
\vspace{0.4em}

\footnotesize
\begin{multicols}{3}

% ─────────────────────────────────────────────────────────────
\ft{Elementare analoge Signale}{sec:elementare-analoge-signale}
Gleichsignal $x(t)=A$.\quad Sinus:
\[ x(t)=A\sin(\omega t+\varphi)=A\cos\!\big(\omega t+\varphi-\tfrac{\pi}{2}\big) \]
\[ \omega=2\pi f=\tfrac{2\pi}{T},\quad f=\tfrac{1}{T} \]
Sprung $\varepsilon(t)=\begin{cases}0&t<0\\1&t\ge0\end{cases}$\;
Rechteck $\operatorname{rect}\!\big(\tfrac{t}{T}\big)$, Fläche $=T$.\\
Exponential $x(t)=A\,e^{-t/\tau}\varepsilon(t)$, Steigung $|_{0^+}=-\tfrac{A}{\tau}$.\\
$\operatorname{si}\!\big(\tfrac{t}{T}\big)=\dfrac{\sin(\pi t/T)}{\pi t/T}$.\quad
Dirac (Siebeigenschaft):
\[ \int_{-\infty}^{\infty}\! f(\xi)\,\delta(\xi-a)\,\mathrm d\xi=f(a),\ \int\delta\,\mathrm d\xi=1 \]

% ─────────────────────────────────────────────────────────────
\ft{Elementare diskrete Signale}{sec:elementare-diskrete-signale}
$\delta[k]=\begin{cases}1&k=0\\0&\text{sonst}\end{cases}$\quad
$\varepsilon[k]=\begin{cases}1&k\ge0\\0&k<0\end{cases}$\\
Rechteck $x[k]=\varepsilon[k]-\varepsilon[k-M]$, Energie $E=MA^2$.\\
Exp.-Folge $x[k]=a^k\varepsilon[k]$.\quad Sinus $x[k]=A\sin\!\big(2\pi k\tfrac{f}{f_A}\big)$.\\
Geometrische Reihe:
\[ \sum_{k=0}^{M-1}\! q^k=\frac{1-q^M}{1-q},\quad \sum_{k=0}^{\infty} q^k=\frac{1}{1-q}\ (|q|<1) \]

% ─────────────────────────────────────────────────────────────
\ft{Komplexe Signale \& Zeiger}{sec:komplexwertige-signale}
\[ \underline{x}(t)=x_R+jx_I=|\underline{x}|\,e^{j\phi},\quad j=\sqrt{-1} \]
\[ x_R=\operatorname{Re}\{\underline x\}=\tfrac12(\underline x+\underline x^{*}),\
   x_I=\operatorname{Im}\{\underline x\}=\tfrac{1}{2j}(\underline x-\underline x^{*}) \]
\[ |\underline x|=\sqrt{x_R^2+x_I^2},\quad \phi=\arg\{\underline x\} \]
Komplexes Exp.: $\underline x(t)=A\,e^{j(\omega t+\varphi)}=\underline A\,e^{j\omega t}$, $\underline A=A e^{j\varphi}$.

% ─────────────────────────────────────────────────────────────
\ft{Symmetrie}{sec:symmetrie-signale}
gerade $x(t)=x(-t)$;\quad ungerade $x(t)=-x(-t)$.\\
konj.\ gerade $\underline x(t)=\underline x^{*}(-t)$;\ konj.\ ungerade $\underline x(t)=-\underline x^{*}(-t)$.\\
Zerlegung $x=x_{\mathrm g}+x_{\mathrm u}$ mit
\[ x_{\mathrm g}=\tfrac12\big(x(t)+x(-t)\big),\ x_{\mathrm u}=\tfrac12\big(x(t)-x(-t)\big) \]

% ─────────────────────────────────────────────────────────────
\ft{Mittelwert, Energie, Leistung, Norm}{sec:norm-energie-leistung}
$\bar x=\dfrac{1}{t_2-t_1}\!\int_{t_1}^{t_2}\! x\,\mathrm dt=\dfrac{1}{k_2-k_1+1}\!\displaystyle\sum_{k=k_1}^{k_2}\! x[k]$;\\
periodisch $\bar x=\tfrac1T\!\int_{-T/2}^{T/2}\! x\,\mathrm dt=\tfrac1N\!\sum_{k=0}^{N-1}\! x[k]$.\\
$L_p$-Norm $\|\underline x\|_p=\big(\int|\underline x|^p\mathrm dt\big)^{1/p}$.\\
Energie $E=\int|\underline x(t)|^2\mathrm dt=\sum|\underline x[k]|^2$ ($E<\infty$: Energiesignal).\\
Leistung $P=\lim_{T\to\infty}\tfrac1T\!\int_{-T/2}^{T/2}\!|\underline x|^2\mathrm dt$;\;
periodisch $P=\tfrac1T\!\int_T\!|\underline x|^2\mathrm dt$\\
($0<P<\infty$: Leistungssignal).

% ─────────────────────────────────────────────────────────────
\ft{Inneres Produkt \& Korrelation}{chap:inneres-produkt}
$\langle \underline x,\underline y\rangle=\int \underline x^{*}(t)\underline y(t)\,\mathrm dt=\sum \underline x^{*}[k]\underline y[k]$; \ orthogonal $\Leftrightarrow \langle\cdot\rangle=0$.\\
Kreuzkorrelation (KKF):
\[ \varphi_{xy}(\tau)=\!\int_{-\infty}^{\infty}\!\underline x^{*}(t)\,\underline y(t+\tau)\,\mathrm dt \]
\[ \varphi_{xy}[\varkappa]=\sum_{k}\underline x^{*}[k]\,\underline y[k+\varkappa] \]
Autokorrelation (AKF): $\varphi_{xx}$ (KKF mit $y=x$),
$\varphi_{xx}(0)=E$, $\varphi_{xx}(\tau)=\varphi_{xx}^{*}(-\tau)$.\\
Sinus: $\tilde\varphi_{xx}(\tau)=\tfrac{A^2}{2}\cos(\omega\tau)$.

% ─────────────────────────────────────────────────────────────
\ft{Empirische Kenngrößen}{sec:streumasse}
$\hat{\bar x}=\dfrac1N\displaystyle\sum_{k=0}^{N-1}\! x[k]$;\quad
$\hat\sigma^2=\dfrac{1}{N-1}\displaystyle\sum_{k=0}^{N-1}\!\big(x[k]-\hat{\bar x}\big)^2$;\\
$\hat\sigma=\sqrt{\hat\sigma^2}$;\quad
empir.\ KKF $\hat\varphi_{xy}[\varkappa]=\sum_{k=0}^{N} x[k]\,y[k+\varkappa]$.

% ─────────────────────────────────────────────────────────────
\ft{Fourier-Reihe \& DFT}{sec:fourierreihe}
reell: $x(t)=\bar x+\sum_{n\ge1}X_a[n]\cos(n\omega_0t)+X_b[n]\sin(n\omega_0t)$.\\
komplex: $x(t)=\sum_{n=-\infty}^{\infty}\! X[n]e^{jn\omega_0t}$, $\omega_0=\tfrac{2\pi}{T_0}$,
\[ X[n]=\frac{1}{T_0}\!\int_{t_1}^{t_1+T_0}\! x(t)\,e^{-jn\omega_0t}\,\mathrm dt \]
$X_a=\tfrac{2}{T_0}\!\int x\cos$, $X_b=\tfrac{2}{T_0}\!\int x\sin$.\\
DFT:
\[ x[k]=\sum_{m=0}^{M-1}\! X[m]\,e^{j\frac{2\pi}{M}mk} \]
\[ X[m]=\tfrac1M\sum_{k=0}^{M-1}\! x[k]\,e^{-j\frac{2\pi}{M}mk} \]
Bandbreite $B=f_{\max}-f_{\min}$; \ $f_m=\tfrac{m}{M}\tfrac{1}{T_A}$.

% ─────────────────────────────────────────────────────────────
\ft{Abtastung \& Quantisierung}{sec:abtasttheorem}
$f_A=\tfrac{\text{Abtastwerte}}{\text{Zeit}}$, $T_A=\tfrac{1}{f_A}$.\\
Abtasttheorem $f_A>2f_g$; \ Aliasing/Spiegelfreq.\ $\tilde f=f_A-f$.\\
Ideale Rekonstruktion:
\[ \tilde x(t)=\sum_k x[k]\,\operatorname{si}\!\big(\tfrac{t-kT_A}{T_A}\big) \]
Halteglied $\tilde x(t)=x(kT_A)$, $kT_A\le t<(k{+}1)T_A$.\\
Quant.-Fehler $e=x-x_q$, $|e|\le\tfrac{\Delta x}{2}$, Leistung $P_e=\dfrac{\Delta x^2}{12}$.\\
Dezibel: $10\lg\!\tfrac{P_2}{P_1}$ [dB] $=20\lg\!\tfrac{A_2}{A_1}$ [dB].

% ─────────────────────────────────────────────────────────────
\ft{Systeme \& Zustandsraum}{sec:zustandsraum}
System $\mathbb S=\{(x(t),y(t))\mid t_0\le t<t_1\}$.\\
kontinuierlich:
\[ \tfrac{\mathrm d\vec z}{\mathrm dt}=\mathbf A\vec z+\mathbf B\vec x,\quad \vec y=\mathbf C\vec z+\mathbf D\vec x \]
diskret:
\[ \vec z[k{+}1]=\mathbf A\vec z[k]+\mathbf B\vec x[k] \]
\[ \vec y[k]=\mathbf C\vec z[k]+\mathbf D\vec x[k] \]
\emph{LTI}: zeitinvariant $y(t-\tau_0)=\mathcal S[x(t-\tau_0)]$ und linear
$a y_1+b y_2=\mathcal S[a x_1+b x_2]$.\\
kausal: $y(t_0)$ nur von $t\le t_0$;\ BIBO: $|x|<\infty\Rightarrow|y|<\infty$.

% ─────────────────────────────────────────────────────────────
\ft{LTI-Systeme \& Faltung}{sec:faltungsintegral}
Impulsantwort $h=\mathcal S[\delta]$. Faltung:
\[ y(t)=\!\int_{-\infty}^{\infty}\! x(\xi)\,h(t-\xi)\,\mathrm d\xi=x*h \]
\[ y[k]=\sum_{l} x[l]\,h[k-l]=x[k]*h[k] \]
$x*\delta=x$;\ $x*h=h*x$;\ assoziativ;\ distributiv.\\
kausal $h(t)=0\ \forall t<0$;\ BIBO $\int|h|\mathrm dt<\infty$ ($\sum|h[l]|<\infty$).\\
Frequenzgang $\underline H(j\omega)=\int h(\xi)e^{-j\omega\xi}\mathrm d\xi$ (Eigenfkt.\ $e^{j\omega t}$).\\
Wiener--Lee: $\varphi_{yy}=\varphi_{hh}*\varphi_{xx}$.

% ─────────────────────────────────────────────────────────────
\ft{Parametrische Systeme (FIR/IIR)}{sec:parametrische-lsi}
allg.\ LSI: $y[k]=\sum_{\nu=0}^{N} x[k{-}\nu]b_\nu+\sum_{\mu=1}^{M} y[k{-}\mu]a_\mu$.\\
FIR: $y[k]=\sum_{\nu=0}^{N} x[k{-}\nu]b_\nu$.\\
IIR: $y[k]=x[k]b_0+\sum_{\mu=1}^{M} y[k{-}\mu]a_\mu$.

% ─────────────────────────────────────────────────────────────
\ft{Wahrscheinlichkeit -- Grundlagen}{sec:axiome-wahrscheinlichkeit}
Axiome: $P(A)\ge0$;\ $P(\Omega)=1$;\ $P(A{+}B)=P(A)+P(B)$ (disjunkt).\\
Laplace: $P(\xi)=\tfrac{1}{|\Omega|}$.\quad
unabh.\ Versuche: $P((\xi_i,\xi_j))=P(\xi_i)P(\xi_j)$.\\
Rand aus Verbund: $P(A_i)=\sum_j P(A_i,B_j)$;\ unabh.\ $P(A,B)=P(A)P(B)$.

% ─────────────────────────────────────────────────────────────
\ft{Kombinatorik}{sec:kombinatorik}
\renewcommand{\arraystretch}{1.9}
\begin{tabular}{@{}p{1.5cm}cc@{}}
  \toprule
   & \scriptsize o.\ Zurückl. & \scriptsize m.\ Zurückl. \\
  \midrule
  \scriptsize geordnet   & $\dfrac{n!}{(n-k)!}$ & $n^k$ \\
  \scriptsize ungeordnet & $\dbinom{n}{k}$      & $\dbinom{n{-}1{+}k}{k}$ \\
  \bottomrule
\end{tabular}
\renewcommand{\arraystretch}{1}\\[0.2em]
Permutationen (alle $n$): $n!$.

% ─────────────────────────────────────────────────────────────
\ft{Bedingte W.\ \& Satz von Bayes}{sec:satz-von-bayes}
\[ P(A_i|B_j)=\frac{P(A_i,B_j)}{P(B_j)} \]
totale W.: $P(A_k)=\sum_{i} P(A_k|B_i)\,P(B_i)$.
\[ P(B_j|A_k)=\frac{P(A_k|B_j)P(B_j)}{\sum_{i} P(A_k|B_i)P(B_i)} \]

% ─────────────────────────────────────────────────────────────
\ft{Zufallsvariablen \& Momente}{sec:momente}
diskret: $\mathrm E\{X\}=\sum_i x_i p_X(x_i)$, $\mathrm E\{X^2\}=\sum_i x_i^2 p_X(x_i)$.\\
kontinuierlich: $\mathrm E\{g(X)\}=\int g(x)p_X(x)\,\mathrm dx$.
\[ \operatorname{var}\{X\}=\mathrm E\{(X-\mu)^2\}=\mathrm E\{X^2\}-\mathrm E\{X\}^2=\sigma^2 \]
$\sigma=\sqrt{\operatorname{var}\{X\}}$;\\
$\mathrm E\{aX{+}bY\}=a\,\mathrm E\{X\}+b\,\mathrm E\{Y\}$.

% ─────────────────────────────────────────────────────────────
\ft{Diskrete Verteilungen}{sec:binomialverteilung}
\emph{Binomial:} $P(k)=\dbinom{N}{k}P^k(1-P)^{N-k}$,
\[ \mathrm E\{X\}=NP,\qquad \operatorname{var}\{X\}=NP(1-P) \]
\emph{Poisson:} $p(k)=e^{-\lambda}\dfrac{\lambda^k}{k!}$,\;
$\mathrm E\{X\}=\lambda$, $\mathrm E\{X^2\}=\lambda^2+\lambda$, $\operatorname{var}=\lambda$.

% ─────────────────────────────────────────────────────────────
\ft{Stetige Verteilungen}{sec:normalverteilung}
{\scriptsize
\renewcommand{\arraystretch}{2.1}
\begin{tabular}{@{}p{1.15cm}ccc@{}}
  \toprule
  Vert. & $p_X(x)$ & $\mathrm E\{X\}$ & $\operatorname{var}\{X\}$ \\
  \midrule
  Gleich $[a,b]$ & $\tfrac{1}{b-a}$ & $\tfrac{a+b}{2}$ & $\tfrac{(b-a)^2}{12}$ \\
  Normal & $\tfrac{1}{\sqrt{2\pi\sigma^2}}e^{-\frac{(x-\mu)^2}{2\sigma^2}}$ & $\mu$ & $\sigma^2$ \\
  Exp.\ ($x{\ge}0$) & $\lambda e^{-\lambda x}$ & $\tfrac{1}{\lambda}$ & $\tfrac{1}{\lambda^2}$ \\
  Laplace & $\tfrac{1}{2\sigma}e^{-\frac{|x-\mu|}{\sigma}}$ & $\mu$ & $2\sigma^2$ \\
  \bottomrule
\end{tabular}
\renewcommand{\arraystretch}{1}}

% ─────────────────────────────────────────────────────────────
\ft{Verteilungsfkt.\ \& Transformation}{sec:transformation-zv}
$F_X(x)=P(X\le x)=\int_{-\infty}^{x} p_X(u)\,\mathrm du$,\ $p_X=\tfrac{\mathrm d}{\mathrm dx}F_X$.\\
monoton $Y=g(X)$: $p_Y(y)=\dfrac{p_X(g^{-1}(y))}{|g'(g^{-1}(y))|}$.\\
nichtmonoton: $p_Y(y)=\sum_i \dfrac{p_X(x_i)}{|g'(x_i)|}$.

% ─────────────────────────────────────────────────────────────
\ft{Zwei Zufallsvariablen}{sec:kovarianz-korrelation}
$F_{XY}=P(X\le x,Y\le y)$,\ $p_{XY}=\dfrac{\partial^2 F_{XY}}{\partial x\,\partial y}$.\\
Rand $p_X=\int p_{XY}\,\mathrm dy$;\ unabhängig $p_{XY}=p_X p_Y$.\\
Kovarianz $C=\mathrm E\{(X-\mu_X)(Y-\mu_Y)\}$;\\
$r_{XY}=\dfrac{C}{\sigma_X\sigma_Y}$, $|r_{XY}|\le1$;\ orthogonal $\mathrm E\{XY\}=0$.\\
Summe unabh.\ ZV: $p_Z=p_X*p_Y$.

% ─────────────────────────────────────────────────────────────
\ft{Informationstheorie}{sec:entropie-redundanz}
Entscheidungsgehalt $H_0=\log_2 N$.\quad
Info-Gehalt $I(a_i)=-\log_2 P(a_i)$.
\[ H=-\sum_{i} P(a_i)\log_2 P(a_i)\ \le H_0 \]
rel.\ Entropie $\tfrac{H}{H_0}$;\ Redundanz $r_q=1-\tfrac{H}{H_0}$.\\
Kanalkapazität $C=\lim_{\tau\to\infty}\tfrac{\log_2 N(\tau)}{\tau}$.\\
Coderedundanz $r_c=\tfrac{\mathrm E\{L\}-H}{H}$;\
KL $D=\sum_i P\log_2\tfrac{P}{\widehat P}$.

\end{multicols}

\restoregeometry


\end{document}
